\documentclass[a4paper,10pt,twoside,openany]{book}

\usepackage[utf8]{inputenc}
\usepackage[T1]{fontenc}
\usepackage[english]{babel}
\usepackage[margin=20mm,inner=19mm,outer=19mm,marginparwidth=12mm,marginparsep=4mm]{geometry}
\usepackage{array}
\usepackage{booktabs}
\usepackage{longtable}
\usepackage{graphicx}
\usepackage{capt-of}
\usepackage{needspace}
\usepackage{makeidx}
\usepackage{hyperref}
\usepackage[pdftex,outerbars,grey]{changebar}

\newcommand{\mr}{\textsc{MR}}
\newcommand{\mrmac}{\textsc{MRMAC}}
\newcommand{\cmd}[1]{\textsc{#1}\index{#1@\textsc{#1}}\allowbreak}
\newcommand{\syntax}[1]{\emph{#1}}
\newenvironment{mrmacextension}{\begin{changebar}}{\end{changebar}}
\newcommand{\legacy}[1]{\footnote{MEMAC compatibility: \textsc{#1} is not available in MRMAC.}}
\newcommand{\tablehead}[2]{\toprule #1 & #2 \\\midrule\endhead}
\newcommand{\tabletail}{\bottomrule}
\newcommand{\commandrow}[2]{\cmd{#1} & #2 \\}
\newcommand{\primitiveoverview}[4]{{\scriptsize\textsc{#1}} & #3 & #4 & \hyperref[prim:#2]{p.~\pageref*{prim:#2}} \\}
\newcolumntype{L}[1]{>{\raggedright\arraybackslash}p{#1}}
\newcommand{\debuggerscreenshot}[3]{%
\Needspace{0.56\textheight}%
\begin{center}
\begin{minipage}{\textwidth}
\begin{changebar}
\centering
\includegraphics[width=\linewidth,height=0.50\textheight,keepaspectratio]{../pngsjpegs/#1}
\captionof{figure}{#3}\label{#2}
\end{changebar}
\end{minipage}
\end{center}
}

\hypersetup{colorlinks=true,linkcolor=black,urlcolor=black,pdftitle={MRMAC Reference Manual},pdfauthor={MR Project}}
\makeindex
\setlength{\emergencystretch}{7em}
\renewcommand{\arraystretch}{1.16}

\title{\mrmac{} Reference Manual\\\large Multi-Edit Revisited Macro Language}
\author{Dr.\ Michael 'iDoc' Raus}
\date{14 July 2026}

\begin{document}
\begin{titlepage}
\begin{center}
\vspace*{\fill}
{\LARGE \mrmac{} Reference Manual\par}
\vspace{0.25\baselineskip}
{\large Multi-Edit Revisited Macro Language\par}
\vspace{3\baselineskip}
{\large Dr.\ Michael 'iDoc' Raus\par}
\vspace{2\baselineskip}
{\large 14 July 2026\par}
\end{center}

\vfill
\begin{minipage}{0.82\textwidth}
\small
This manual describes the executable MRMAC language in the current MR source
tree. The original Multi-Edit reference guide supplies historical context, not
the authority for MRMAC semantics.

\begin{changebar}
An outer-margin bar marks an MRMAC extension or a documented semantic change
relative to MEMAC. A footnote identifies a close historical command that is not
implemented. The bars are deliberately mechanical: they remain useful when a
reader prints only selected pages.
\end{changebar}
\end{minipage}
\vspace*{1.5cm}
\end{titlepage}

\tableofcontents
\clearpage



\part{Language Fundamentals}

\chapter{Purpose and Source Form}

MRMAC source is ordinary text compiled before it runs. A source file may name itself, define macro entries, and define closures for scheduled work. Begin here when writing a new macro: every source-form primitive is listed in the local map and then documented immediately below.

\section{Primitive Map}
The map is a local index for this chapter. The final column is a generated page
reference and hyperlink to the full entry, where the source form, parameters,
semantics and a complete example macro appear.

\begin{longtable}{@{}L{0.38\textwidth}L{0.14\textwidth}L{0.27\textwidth}L{0.10\textwidth}@{}}
\toprule
Primitive & Role & Purpose & Reference \\
\midrule
\endhead
\toprule
Primitive & Role & Purpose & Reference \\
\midrule
\endhead
\primitiveoverview{\$MACRO}{keyword-dollar-macro}{keyword}{Source structure.}
\primitiveoverview{\$MACRO\_FILE}{keyword-dollar-macro-file}{keyword}{Source structure.}
\primitiveoverview{\$CLOSURE}{keyword-dollar-closure}{keyword}{Source structure.}
\primitiveoverview{DEF\_TICK}{keyword-def-tick}{keyword}{Source structure.}
\primitiveoverview{END\_MACRO}{keyword-end-macro}{keyword}{Macro terminator.}
\primitiveoverview{END\_CLOSURE}{keyword-end-closure}{keyword}{Closure terminator.}
\primitiveoverview{FROM}{keyword-from}{header keyword}{Macro context.}
\primitiveoverview{TO}{keyword-to}{header keyword}{Key binding.}
\primitiveoverview{TRANS}{keyword-trans}{header keyword}{TRANS attribute.}
\primitiveoverview{DUMP}{keyword-dump}{header keyword}{DUMP attribute.}
\primitiveoverview{PERM}{keyword-perm}{header keyword}{PERM attribute.}
\bottomrule
\end{longtable}

\section{Reference Entries}

\phantomsection
\label{prim:keyword-dollar-macro}
\subsection*{\$MACRO}
\index{\$MACRO@\texttt{\$MACRO}}
\begin{verbatim}
$MACRO macroName FROM EDIT;
    statement;
END_MACRO;
\end{verbatim}
Parameters: macroName names the entry; FROM EDIT selects the editor context.

Starts a named executable macro entry.

\paragraph{Example macro.}
\begin{verbatim}
$MACRO REFERENCE_MACRO FROM EDIT;
    MAKE_MESSAGE('reference');
END_MACRO;
\end{verbatim}

\phantomsection
\label{prim:keyword-dollar-macro-file}
\subsection*{\$MACRO\_FILE}
\index{\$MACRO\_FILE@\texttt{\$MACRO\_FILE}}
\begin{verbatim}
$MACRO_FILE fileName;
\end{verbatim}
Parameters: fileName is the single source-file identity.

Declares the identity of one macro source file.

\paragraph{Example macro.}
\begin{verbatim}
$MACRO_FILE REFERENCE_FILE;
$MACRO REFERENCE_MACRO FROM EDIT;
END_MACRO;
\end{verbatim}

\begin{mrmacextension}
\phantomsection
\label{prim:keyword-dollar-closure}
\subsection*{\$CLOSURE}
\index{\$CLOSURE@\texttt{\$CLOSURE}}
\begin{verbatim}
$CLOSURE closureName;
DEF_TICK(milliseconds);
    statement;
END_CLOSURE;
\end{verbatim}
Parameters: closureName names the scheduled unit; milliseconds is a positive tick interval.

Starts a named closure with local execution state.

\paragraph{Example macro.}
\begin{verbatim}
$CLOSURE REFERENCE_CLOSURE;
DEF_TICK(1000);
    MAKE_MESSAGE('tick');
END_CLOSURE;
\end{verbatim}
\end{mrmacextension}

\begin{mrmacextension}
\phantomsection
\label{prim:keyword-def-tick}
\subsection*{DEF\_TICK}
\index{DEF\_TICK@\texttt{DEF\_TICK}}
\begin{verbatim}
DEF_TICK(milliseconds);
\end{verbatim}
Parameters: milliseconds is a positive integer interval.

Declares the tick interval of the enclosing closure.

\paragraph{Example macro.}
\begin{verbatim}
$CLOSURE REFERENCE_TICK;
DEF_TICK(1000);
END_CLOSURE;
\end{verbatim}
\end{mrmacextension}

\phantomsection
\label{prim:keyword-end-macro}
\subsection*{END\_MACRO}
\index{END\_MACRO!keyword@\textsc{keyword}}
\begin{verbatim}
END_MACRO;
\end{verbatim}
Parameters: None.

Terminates the current macro entry. It is required after the entry body.

\paragraph{Example macro.}
\begin{verbatim}
$MACRO REFERENCE_END_MACRO FROM EDIT;
END_MACRO;
\end{verbatim}

\begin{mrmacextension}
\phantomsection
\label{prim:keyword-end-closure}
\subsection*{END\_CLOSURE}
\index{END\_CLOSURE!keyword@\textsc{keyword}}
\begin{verbatim}
END_CLOSURE;
\end{verbatim}
Parameters: None.

Terminates the current closure entry after its tick declaration and body.

\paragraph{Example macro.}
\begin{verbatim}
$CLOSURE REFERENCE_END_CLOSURE;
DEF_TICK(1000);
END_CLOSURE;
\end{verbatim}
\end{mrmacextension}

\phantomsection
\label{prim:keyword-from}
\subsection*{FROM}
\index{FROM!header keyword@\textsc{header keyword}}
\begin{verbatim}
$MACRO macroName FROM mode;
\end{verbatim}
Parameters: macroName names the macro; mode selects its context.

Selects the declared execution context. EDIT is the normal editor context.

\paragraph{Example macro.}
\begin{verbatim}
$MACRO REFERENCE_FROM FROM EDIT;
END_MACRO;
\end{verbatim}

\phantomsection
\label{prim:keyword-to}
\subsection*{TO}
\index{TO!header keyword@\textsc{header keyword}}
\begin{verbatim}
$MACRO macroName TO <keySpec> FROM EDIT;
\end{verbatim}
Parameters: macroName names the macro; keySpec is a supported key specification.

Associates an entry with a key specification when its source is loaded.

\paragraph{Example macro.}
\begin{verbatim}
$MACRO REFERENCE_TO TO <F1> FROM EDIT;
END_MACRO;
\end{verbatim}

\phantomsection
\label{prim:keyword-trans}
\subsection*{TRANS}
\index{TRANS!header keyword@\textsc{header keyword}}
\begin{verbatim}
$MACRO macroName TRANS FROM EDIT;
\end{verbatim}
Parameters: macroName names the macro.

Sets the compiler-supported TRANS attribute on the macro header.

\paragraph{Example macro.}
\begin{verbatim}
$MACRO REFERENCE_TRANS TRANS FROM EDIT;
END_MACRO;
\end{verbatim}

\phantomsection
\label{prim:keyword-dump}
\subsection*{DUMP}
\index{DUMP!header keyword@\textsc{header keyword}}
\begin{verbatim}
$MACRO macroName DUMP FROM EDIT;
\end{verbatim}
Parameters: macroName names the macro.

Sets the compiler-supported DUMP attribute on the macro header.

\paragraph{Example macro.}
\begin{verbatim}
$MACRO REFERENCE_DUMP DUMP FROM EDIT;
END_MACRO;
\end{verbatim}

\phantomsection
\label{prim:keyword-perm}
\subsection*{PERM}
\index{PERM!header keyword@\textsc{header keyword}}
\begin{verbatim}
$MACRO macroName PERM FROM EDIT;
\end{verbatim}
Parameters: macroName names the macro.

Sets the compiler-supported PERM attribute on the macro header.

\paragraph{Example macro.}
\begin{verbatim}
$MACRO REFERENCE_PERM PERM FROM EDIT;
END_MACRO;
\end{verbatim}


\chapter{Values, Variables and Collections}

A macro works with typed values. Declare local values before using them; declarations belong to the current execution, not to a source file. Integers, real numbers, strings and characters are scalar values. Hashes associate keys with values, and typed arrays contain one declared element type. Built-in constants and variables provide runtime values.

\section{Primitive Map}
The map is a local index for this chapter. The final column is a generated page
reference and hyperlink to the full entry, where the source form, parameters,
semantics and a complete example macro appear.

\begin{longtable}{@{}L{0.38\textwidth}L{0.14\textwidth}L{0.27\textwidth}L{0.10\textwidth}@{}}
\toprule
Primitive & Role & Purpose & Reference \\
\midrule
\endhead
\toprule
Primitive & Role & Purpose & Reference \\
\midrule
\endhead
\primitiveoverview{DEF\_INT}{keyword-def-int}{keyword}{Typed declaration.}
\primitiveoverview{DEF\_REAL}{keyword-def-real}{keyword}{Typed declaration.}
\primitiveoverview{DEF\_STR}{keyword-def-str}{keyword}{Typed declaration.}
\primitiveoverview{DEF\_CHAR}{keyword-def-char}{keyword}{Typed declaration.}
\primitiveoverview{DEF\_HASH}{keyword-def-hash}{keyword}{Typed declaration.}
\primitiveoverview{assignment}{form-assignment}{source form}{Assignment.}
\primitiveoverview{indexing}{form-indexing}{source form}{Collection access.}
\primitiveoverview{TRUE}{constant-true}{constant}{Symbolic constant.}
\primitiveoverview{FALSE}{constant-false}{constant}{Symbolic constant.}
\primitiveoverview{BLACK}{constant-black}{constant}{Symbolic constant.}
\primitiveoverview{BLUE}{constant-blue}{constant}{Symbolic constant.}
\primitiveoverview{GREEN}{constant-green}{constant}{Symbolic constant.}
\primitiveoverview{CYAN}{constant-cyan}{constant}{Symbolic constant.}
\primitiveoverview{RED}{constant-red}{constant}{Symbolic constant.}
\primitiveoverview{MAGENTA}{constant-magenta}{constant}{Symbolic constant.}
\primitiveoverview{BROWN}{constant-brown}{constant}{Symbolic constant.}
\primitiveoverview{LIGHTGRAY}{constant-lightgray}{constant}{Symbolic constant.}
\primitiveoverview{DARKGRAY}{constant-darkgray}{constant}{Symbolic constant.}
\primitiveoverview{LIGHTBLUE}{constant-lightblue}{constant}{Symbolic constant.}
\primitiveoverview{LIGHTGREEN}{constant-lightgreen}{constant}{Symbolic constant.}
\primitiveoverview{LIGHTCYAN}{constant-lightcyan}{constant}{Symbolic constant.}
\primitiveoverview{LIGHTRED}{constant-lightred}{constant}{Symbolic constant.}
\primitiveoverview{LIGHTMAGENTA}{constant-lightmagenta}{constant}{Symbolic constant.}
\primitiveoverview{YELLOW}{constant-yellow}{constant}{Symbolic constant.}
\primitiveoverview{WHITE}{constant-white}{constant}{Symbolic constant.}
\primitiveoverview{EDIT}{constant-edit}{constant}{Symbolic constant.}
\primitiveoverview{DOS\_SHELL}{constant-dos-shell}{constant}{Symbolic constant.}
\primitiveoverview{BACK\_SPACE}{constant-back-space}{constant}{Symbolic constant.}
\primitiveoverview{BLOCK\_BEGIN}{constant-block-begin}{constant}{Symbolic constant.}
\primitiveoverview{BLOCK\_END}{constant-block-end}{constant}{Symbolic constant.}
\primitiveoverview{BLOCK\_OFF}{constant-block-off}{constant}{Symbolic constant.}
\primitiveoverview{COL\_BLOCK\_BEGIN}{constant-col-block-begin}{constant}{Symbolic constant.}
\primitiveoverview{COPY\_BLOCK}{constant-copy-block}{constant}{Symbolic constant.}
\primitiveoverview{CR}{constant-cr}{constant}{Symbolic constant.}
\primitiveoverview{DELETE\_BLOCK}{constant-delete-block}{constant}{Symbolic constant.}
\primitiveoverview{DEL\_CHAR}{constant-del-char}{constant}{Symbolic constant.}
\primitiveoverview{DEL\_LINE}{constant-del-line}{constant}{Symbolic constant.}
\primitiveoverview{DOWN}{constant-down}{constant}{Symbolic constant.}
\primitiveoverview{EOF}{constant-eof}{constant}{Symbolic constant.}
\primitiveoverview{EOL}{constant-eol}{constant}{Symbolic constant.}
\primitiveoverview{FIRST\_WORD}{constant-first-word}{constant}{Symbolic constant.}
\primitiveoverview{GOTO\_MARK}{constant-goto-mark}{constant}{Symbolic constant.}
\primitiveoverview{HOME}{constant-home}{constant}{Symbolic constant.}
\primitiveoverview{INDENT}{constant-indent}{constant}{Symbolic constant.}
\primitiveoverview{KEY\_RECORD}{constant-key-record}{constant}{Symbolic constant.}
\primitiveoverview{LAST\_PAGE\_BREAK}{constant-last-page-break}{constant}{Symbolic constant.}
\primitiveoverview{LEFT}{constant-left}{constant}{Symbolic constant.}
\primitiveoverview{MARK\_POS}{constant-mark-pos}{constant}{Symbolic constant.}
\primitiveoverview{MOVE\_BLOCK}{constant-move-block}{constant}{Symbolic constant.}
\primitiveoverview{NEXT\_PAGE\_BREAK}{constant-next-page-break}{constant}{Symbolic constant.}
\primitiveoverview{PAGE\_DOWN}{constant-page-down}{constant}{Symbolic constant.}
\primitiveoverview{PAGE\_UP}{constant-page-up}{constant}{Symbolic constant.}
\primitiveoverview{RIGHT}{constant-right}{constant}{Symbolic constant.}
\primitiveoverview{SAVE\_FILE}{constant-save-file}{constant}{Symbolic constant.}
\primitiveoverview{STR\_BLOCK\_BEGIN}{constant-str-block-begin}{constant}{Symbolic constant.}
\primitiveoverview{TAB\_LEFT}{constant-tab-left}{constant}{Symbolic constant.}
\primitiveoverview{TAB\_RIGHT}{constant-tab-right}{constant}{Symbolic constant.}
\primitiveoverview{TOF}{constant-tof}{constant}{Symbolic constant.}
\primitiveoverview{UNDENT}{constant-undent}{constant}{Symbolic constant.}
\primitiveoverview{UNDO}{constant-undo}{constant}{Symbolic constant.}
\primitiveoverview{UP}{constant-up}{constant}{Symbolic constant.}
\primitiveoverview{WORD\_LEFT}{constant-word-left}{constant}{Symbolic constant.}
\primitiveoverview{WORD\_RIGHT}{constant-word-right}{constant}{Symbolic constant.}
\primitiveoverview{ALL}{constant-all}{constant}{Symbolic constant.}
\primitiveoverview{RETURN\_INT}{variable-return-int}{variable}{Runtime variable.}
\primitiveoverview{ERROR\_LEVEL}{variable-error-level}{variable}{Runtime variable.}
\primitiveoverview{FIRST\_RUN}{variable-first-run}{variable}{Runtime variable.}
\primitiveoverview{IGNORE\_CASE}{variable-ignore-case}{variable}{Runtime variable.}
\primitiveoverview{REG\_EXP\_STAT}{variable-reg-exp-stat}{variable}{Runtime variable.}
\primitiveoverview{FIRST\_SAVE}{variable-first-save}{variable}{Runtime variable.}
\primitiveoverview{BUFFER\_ID}{variable-buffer-id}{variable}{Runtime variable.}
\primitiveoverview{TMP\_FILE}{variable-tmp-file}{variable}{Runtime variable.}
\primitiveoverview{FILE\_CHANGED}{variable-file-changed}{variable}{Runtime variable.}
\primitiveoverview{LAST\_FILE\_ATTR}{variable-last-file-attr}{variable}{Runtime variable.}
\primitiveoverview{LAST\_FILE\_SIZE}{variable-last-file-size}{variable}{Runtime variable.}
\primitiveoverview{LAST\_FILE\_TIME}{variable-last-file-time}{variable}{Runtime variable.}
\primitiveoverview{CUR\_FILE\_ATTR}{variable-cur-file-attr}{variable}{Runtime variable.}
\primitiveoverview{CUR\_FILE\_SIZE}{variable-cur-file-size}{variable}{Runtime variable.}
\primitiveoverview{READ\_ONLY}{variable-read-only}{variable}{Runtime variable.}
\primitiveoverview{DOC\_MODE}{variable-doc-mode}{variable}{Runtime variable.}
\primitiveoverview{PRINT\_MARGIN}{variable-print-margin}{variable}{Runtime variable.}
\primitiveoverview{SHADOW\_CHAR}{variable-shadow-char}{variable}{Runtime variable.}
\primitiveoverview{REFRESH}{variable-refresh}{variable}{Runtime variable.}
\primitiveoverview{MESSAGES}{variable-messages}{variable}{Runtime variable.}
\primitiveoverview{MOUSE}{variable-mouse}{variable}{Runtime variable.}
\primitiveoverview{LOGO\_SCREEN}{variable-logo-screen}{variable}{Runtime variable.}
\primitiveoverview{EXPLOSIONS}{variable-explosions}{variable}{Runtime variable.}
\primitiveoverview{TRUNCATE\_SPACES}{variable-truncate-spaces}{variable}{Runtime variable.}
\primitiveoverview{BACKUPS}{variable-backups}{variable}{Runtime variable.}
\primitiveoverview{AUTOSAVE}{variable-autosave}{variable}{Runtime variable.}
\primitiveoverview{UNDO\_STAT}{variable-undo-stat}{variable}{Runtime variable.}
\primitiveoverview{FORMAT\_STAT}{variable-format-stat}{variable}{Runtime variable.}
\primitiveoverview{WRAP\_STAT}{variable-wrap-stat}{variable}{Runtime variable.}
\primitiveoverview{MEM\_ALLOC}{variable-mem-alloc}{variable}{Runtime variable.}
\primitiveoverview{LEFT\_MARGIN}{variable-left-margin}{variable}{Runtime variable.}
\primitiveoverview{RIGHT\_MARGIN}{variable-right-margin}{variable}{Runtime variable.}
\primitiveoverview{FORMAT\_RULER}{variable-format-ruler}{variable}{Runtime variable.}
\primitiveoverview{INDENT\_STYLE}{variable-indent-style}{variable}{Runtime variable.}
\primitiveoverview{INS\_CURSOR}{variable-ins-cursor}{variable}{Runtime variable.}
\primitiveoverview{OVR\_CURSOR}{variable-ovr-cursor}{variable}{Runtime variable.}
\primitiveoverview{CTRL\_HELP}{variable-ctrl-help}{variable}{Runtime variable.}
\primitiveoverview{MOUSE\_H\_SENSE}{variable-mouse-h-sense}{variable}{Runtime variable.}
\primitiveoverview{MOUSE\_V\_SENSE}{variable-mouse-v-sense}{variable}{Runtime variable.}
\primitiveoverview{WINDOW\_ATTR}{variable-window-attr}{variable}{Runtime variable.}
\primitiveoverview{TEXT\_COLOR}{variable-text-color}{variable}{Runtime variable.}
\primitiveoverview{CHANGE\_COLOR}{variable-change-color}{variable}{Runtime variable.}
\primitiveoverview{BACK\_COLOR}{variable-back-color}{variable}{Runtime variable.}
\primitiveoverview{MENU\_COLOR}{variable-menu-color}{variable}{Runtime variable.}
\primitiveoverview{STAT\_COLOR}{variable-stat-color}{variable}{Runtime variable.}
\primitiveoverview{ERROR\_COLOR}{variable-error-color}{variable}{Runtime variable.}
\primitiveoverview{SHADOW\_COLOR}{variable-shadow-color}{variable}{Runtime variable.}
\primitiveoverview{STATUS\_ROW}{variable-status-row}{variable}{Runtime variable.}
\primitiveoverview{MESSAGE\_ROW}{variable-message-row}{variable}{Runtime variable.}
\primitiveoverview{MAX\_WINDOW\_ROW}{variable-max-window-row}{variable}{Runtime variable.}
\primitiveoverview{MIN\_WINDOW\_ROW}{variable-min-window-row}{variable}{Runtime variable.}
\primitiveoverview{NAME\_LINE}{variable-name-line}{variable}{Runtime variable.}
\primitiveoverview{PARAM\_COUNT}{variable-param-count}{variable}{Runtime variable.}
\primitiveoverview{CPU}{variable-cpu}{variable}{Runtime variable.}
\primitiveoverview{C\_COL}{variable-c-col}{variable}{Runtime variable.}
\primitiveoverview{C\_LINE}{variable-c-line}{variable}{Runtime variable.}
\primitiveoverview{C\_ROW}{variable-c-row}{variable}{Runtime variable.}
\primitiveoverview{C\_PAGE}{variable-c-page}{variable}{Runtime variable.}
\primitiveoverview{PG\_LINE}{variable-pg-line}{variable}{Runtime variable.}
\primitiveoverview{AT\_EOF}{variable-at-eof}{variable}{Runtime variable.}
\primitiveoverview{AT\_EOL}{variable-at-eol}{variable}{Runtime variable.}
\primitiveoverview{BLOCK\_STAT}{variable-block-stat}{variable}{Runtime variable.}
\primitiveoverview{BLOCK\_LINE1}{variable-block-line1}{variable}{Runtime variable.}
\primitiveoverview{BLOCK\_LINE2}{variable-block-line2}{variable}{Runtime variable.}
\primitiveoverview{BLOCK\_COL1}{variable-block-col1}{variable}{Runtime variable.}
\primitiveoverview{BLOCK\_COL2}{variable-block-col2}{variable}{Runtime variable.}
\primitiveoverview{MARKING}{variable-marking}{variable}{Runtime variable.}
\primitiveoverview{TAB\_EXPAND}{variable-tab-expand}{variable}{Runtime variable.}
\primitiveoverview{INSERT\_MODE}{variable-insert-mode}{variable}{Runtime variable.}
\primitiveoverview{INDENT\_LEVEL}{variable-indent-level}{variable}{Runtime variable.}
\primitiveoverview{CUR\_WINDOW}{variable-cur-window}{variable}{Runtime variable.}
\primitiveoverview{LINK\_STAT}{variable-link-stat}{variable}{Runtime variable.}
\primitiveoverview{WIN\_X1}{variable-win-x1}{variable}{Runtime variable.}
\primitiveoverview{WIN\_Y1}{variable-win-y1}{variable}{Runtime variable.}
\primitiveoverview{WIN\_X2}{variable-win-x2}{variable}{Runtime variable.}
\primitiveoverview{WIN\_Y2}{variable-win-y2}{variable}{Runtime variable.}
\primitiveoverview{WINDOW\_COUNT}{variable-window-count}{variable}{Runtime variable.}
\primitiveoverview{VIRTUAL\_DESKTOPS}{variable-virtual-desktops}{variable}{Runtime variable.}
\primitiveoverview{CYCLIC\_VIRTUAL\_DESKTOPS}{variable-cyclic-virtual-desktops}{variable}{Runtime variable.}
\primitiveoverview{KEY1}{variable-key1}{variable}{Runtime variable.}
\primitiveoverview{KEY2}{variable-key2}{variable}{Runtime variable.}
\primitiveoverview{RETURN\_STR}{variable-return-str}{variable}{Runtime variable.}
\primitiveoverview{MPARM\_STR}{variable-mparm-str}{variable}{Runtime variable.}
\primitiveoverview{DATE}{variable-date}{variable}{Runtime variable.}
\primitiveoverview{TIME}{variable-time}{variable}{Runtime variable.}
\primitiveoverview{FIRST\_MACRO}{variable-first-macro}{variable}{Runtime variable.}
\primitiveoverview{NEXT\_MACRO}{variable-next-macro}{variable}{Runtime variable.}
\primitiveoverview{LAST\_FILE\_NAME}{variable-last-file-name}{variable}{Runtime variable.}
\primitiveoverview{TMP\_FILE\_NAME}{variable-tmp-file-name}{variable}{Runtime variable.}
\primitiveoverview{FILE\_NAME}{variable-file-name}{variable}{Runtime variable.}
\primitiveoverview{COMSPEC}{variable-comspec}{variable}{Runtime variable.}
\primitiveoverview{TEMP\_PATH}{variable-temp-path}{variable}{Runtime variable.}
\primitiveoverview{FOUND\_STR}{variable-found-str}{variable}{Runtime variable.}
\primitiveoverview{SEARCH\_FILE}{variable-search-file}{variable}{Runtime variable.}
\primitiveoverview{MR\_PATH}{variable-mr-path}{variable}{Runtime variable.}
\primitiveoverview{OS\_VERSION}{variable-os-version}{variable}{Runtime variable.}
\primitiveoverview{GET\_LINE}{variable-get-line}{variable}{Runtime variable.}
\primitiveoverview{FORMAT\_LINE}{variable-format-line}{variable}{Runtime variable.}
\primitiveoverview{DEFAULT\_FORMAT}{variable-default-format}{variable}{Runtime variable.}
\primitiveoverview{PAGE\_STR}{variable-page-str}{variable}{Runtime variable.}
\primitiveoverview{WORD\_DELIMITS}{variable-word-delimits}{variable}{Runtime variable.}
\primitiveoverview{FOUND\_X}{variable-found-x}{variable}{Runtime variable.}
\primitiveoverview{FOUND\_Y}{variable-found-y}{variable}{Runtime variable.}
\primitiveoverview{CUR\_CHAR}{variable-cur-char}{variable}{Runtime variable.}
\bottomrule
\end{longtable}

\section{Reference Entries}

\phantomsection
\label{prim:keyword-def-int}
\subsection*{DEF\_INT}
\index{DEF\_INT@\texttt{DEF\_INT}}
\begin{verbatim}
DEF_INT(name);
DEF_INT(name[]);
\end{verbatim}
Parameters: name is a local integer scalar or typed-array name.

Declares an integer variable or array.

\paragraph{Example macro.}
\begin{verbatim}
$MACRO REFERENCE_DEF_INT FROM EDIT;
    DEF_INT(Value);
    Value := 1;
END_MACRO;
\end{verbatim}

\phantomsection
\label{prim:keyword-def-real}
\subsection*{DEF\_REAL}
\index{DEF\_REAL@\texttt{DEF\_REAL}}
\begin{verbatim}
DEF_REAL(name);
DEF_REAL(name[]);
\end{verbatim}
Parameters: name is a local real scalar or typed-array name.

Declares a real variable or array.

\paragraph{Example macro.}
\begin{verbatim}
$MACRO REFERENCE_DEF_REAL FROM EDIT;
    DEF_REAL(Value);
    Value := 1.0;
END_MACRO;
\end{verbatim}

\phantomsection
\label{prim:keyword-def-str}
\subsection*{DEF\_STR}
\index{DEF\_STR@\texttt{DEF\_STR}}
\begin{verbatim}
DEF_STR(name);
DEF_STR(name[]);
\end{verbatim}
Parameters: name is a local string scalar or typed-array name.

Declares a string variable or array.

\paragraph{Example macro.}
\begin{verbatim}
$MACRO REFERENCE_DEF_STR FROM EDIT;
    DEF_STR(Value);
    Value := 'reference';
END_MACRO;
\end{verbatim}

\phantomsection
\label{prim:keyword-def-char}
\subsection*{DEF\_CHAR}
\index{DEF\_CHAR@\texttt{DEF\_CHAR}}
\begin{verbatim}
DEF_CHAR(name);
DEF_CHAR(name[]);
\end{verbatim}
Parameters: name is a local character scalar or typed-array name.

Declares a character variable or array.

\paragraph{Example macro.}
\begin{verbatim}
$MACRO REFERENCE_DEF_CHAR FROM EDIT;
    DEF_CHAR(Value);
    Value := 'r';
END_MACRO;
\end{verbatim}

\begin{mrmacextension}
\phantomsection
\label{prim:keyword-def-hash}
\subsection*{DEF\_HASH}
\index{DEF\_HASH@\texttt{DEF\_HASH}}
\begin{verbatim}
DEF_HASH(name);
\end{verbatim}
Parameters: name is a local hash name.

Declares a local hash.

\paragraph{Example macro.}
\begin{verbatim}
$MACRO REFERENCE_DEF_HASH FROM EDIT;
    DEF_HASH(State);
    State['name'] := 'reference';
END_MACRO;
\end{verbatim}
\end{mrmacextension}

\begin{mrmacextension}
\phantomsection
\label{prim:form-assignment}
\subsection*{Assignment}
\index{assignment!source form@\textsc{source form}}
\begin{verbatim}
target := expression;
\end{verbatim}
Parameters: target is a declared variable, array element or writable hash element; expression supplies a compatible value.

Assigns the evaluated right-hand expression to the left-hand target.

\paragraph{Example macro.}
\begin{verbatim}
$MACRO REFERENCE_ASSIGNMENT FROM EDIT;
    DEF_INT(Value);
    Value := 1;
END_MACRO;
\end{verbatim}
\end{mrmacextension}

\begin{mrmacextension}
\phantomsection
\label{prim:form-indexing}
\subsection*{Array and Hash Indexing}
\index{indexing!source form@\textsc{source form}}
\begin{verbatim}
collection[indexOrKey]
\end{verbatim}
Parameters: collection is an array or hash; indexOrKey is a positive integer array index or string-like hash key.

Addresses an element of a typed array or hash. Array reads outside the current bounds are runtime errors.

\paragraph{Example macro.}
\begin{verbatim}
$MACRO REFERENCE_INDEXING FROM EDIT;
    DEF_INT(Values[]);
    Values[1] := 42;
END_MACRO;
\end{verbatim}
\end{mrmacextension}

\phantomsection
\label{prim:constant-true}
\subsection*{TRUE --- constant}
\index{TRUE!constant@\textsc{constant}}
\begin{verbatim}
TRUE
\end{verbatim}
Parameters: None. Result: integer constant.

Evaluates to integer 1, the true result used by conditional expressions.

\paragraph{Example macro.}
\begin{verbatim}
$MACRO REFERENCE_CONSTANT_TRUE FROM EDIT;
    DEF_INT(Value);
    Value := TRUE;
END_MACRO;
\end{verbatim}

\phantomsection
\label{prim:constant-false}
\subsection*{FALSE --- constant}
\index{FALSE!constant@\textsc{constant}}
\begin{verbatim}
FALSE
\end{verbatim}
Parameters: None. Result: integer constant.

Evaluates to integer 0, the false result used by conditional expressions.

\paragraph{Example macro.}
\begin{verbatim}
$MACRO REFERENCE_CONSTANT_FALSE FROM EDIT;
    DEF_INT(Value);
    Value := FALSE;
END_MACRO;
\end{verbatim}

\phantomsection
\subsection*{Colour Constants}
\label{prim:color-constants}
\label{prim:constant-black}
\label{prim:constant-blue}
\label{prim:constant-green}
\label{prim:constant-cyan}
\label{prim:constant-red}
\label{prim:constant-magenta}
\label{prim:constant-brown}
\label{prim:constant-lightgray}
\label{prim:constant-darkgray}
\label{prim:constant-lightblue}
\label{prim:constant-lightgreen}
\label{prim:constant-lightcyan}
\label{prim:constant-lightred}
\label{prim:constant-lightmagenta}
\label{prim:constant-yellow}
\label{prim:constant-white}
\index{BLACK!constant@\textsc{constant}}
\index{BLUE!constant@\textsc{constant}}
\index{GREEN!constant@\textsc{constant}}
\index{CYAN!constant@\textsc{constant}}
\index{RED!constant@\textsc{constant}}
\index{MAGENTA!constant@\textsc{constant}}
\index{BROWN!constant@\textsc{constant}}
\index{LIGHTGRAY!constant@\textsc{constant}}
\index{DARKGRAY!constant@\textsc{constant}}
\index{LIGHTBLUE!constant@\textsc{constant}}
\index{LIGHTGREEN!constant@\textsc{constant}}
\index{LIGHTCYAN!constant@\textsc{constant}}
\index{LIGHTRED!constant@\textsc{constant}}
\index{LIGHTMAGENTA!constant@\textsc{constant}}
\index{YELLOW!constant@\textsc{constant}}
\index{WHITE!constant@\textsc{constant}}
\begin{verbatim}
BLACK          0       LIGHTGRAY      7
BLUE           1       DARKGRAY       8
GREEN          2       LIGHTBLUE      9
CYAN           3       LIGHTGREEN    10
RED            4       LIGHTCYAN     11
MAGENTA        5       LIGHTRED      12
BROWN          6       LIGHTMAGENTA  13
YELLOW        14       WHITE         15
\end{verbatim}
Parameters: None. Result: integer constant.

These constants denote the fixed values of the 16-colour palette. Use them only where a primitive documents a palette-colour value.

\paragraph{Example macro.}
\begin{verbatim}
$MACRO REFERENCE_COLOUR_CONSTANTS FROM EDIT;
    DEF_INT(Colour);
    Colour := LIGHTGREEN;
END_MACRO;
\end{verbatim}

\phantomsection
\label{prim:constant-edit}
\subsection*{EDIT --- constant}
\index{EDIT!constant@\textsc{constant}}
\begin{verbatim}
EDIT
\end{verbatim}
Parameters: None. Result: integer constant.

Provides compiler-defined macro mode value 0. DOS\_SHELL is historical source syntax; MRMAC provides no DOS shell runtime.

\paragraph{Example macro.}
\begin{verbatim}
$MACRO REFERENCE_CONSTANT_EDIT FROM EDIT;
    DEF_INT(Value);
    Value := EDIT;
END_MACRO;
\end{verbatim}

\phantomsection
\label{prim:constant-dos-shell}
\subsection*{DOS\_SHELL --- constant}
\index{DOS\_SHELL!constant@\textsc{constant}}
\begin{verbatim}
DOS_SHELL
\end{verbatim}
Parameters: None. Result: integer constant.

Provides compiler-defined macro mode value 1. DOS\_SHELL is historical source syntax; MRMAC provides no DOS shell runtime.

\paragraph{Example macro.}
\begin{verbatim}
$MACRO REFERENCE_CONSTANT_DOS_SHELL FROM EDIT;
    DEF_INT(Value);
    Value := DOS_SHELL;
END_MACRO;
\end{verbatim}

\phantomsection
\label{prim:constant-back-space}
\subsection*{BACK\_SPACE --- constant}
\index{BACK\_SPACE!constant@\textsc{constant}}
\begin{verbatim}
BACK_SPACE
\end{verbatim}
Parameters: None. Result: integer constant.

Provides legacy integer command identifier 0x7001. Use it only where a procedure explicitly accepts a command identifier.

\paragraph{Example macro.}
\begin{verbatim}
$MACRO REFERENCE_CONSTANT_BACK_SPACE FROM EDIT;
    DEF_INT(Value);
    Value := BACK_SPACE;
END_MACRO;
\end{verbatim}

\phantomsection
\label{prim:constant-block-begin}
\subsection*{BLOCK\_BEGIN --- constant}
\index{BLOCK\_BEGIN!constant@\textsc{constant}}
\begin{verbatim}
BLOCK_BEGIN
\end{verbatim}
Parameters: None. Result: integer constant.

Provides legacy integer command identifier 0x7002. Use it only where a procedure explicitly accepts a command identifier.

\paragraph{Example macro.}
\begin{verbatim}
$MACRO REFERENCE_CONSTANT_BLOCK_BEGIN FROM EDIT;
    DEF_INT(Value);
    Value := BLOCK_BEGIN;
END_MACRO;
\end{verbatim}

\phantomsection
\label{prim:constant-block-end}
\subsection*{BLOCK\_END --- constant}
\index{BLOCK\_END!constant@\textsc{constant}}
\begin{verbatim}
BLOCK_END
\end{verbatim}
Parameters: None. Result: integer constant.

Provides legacy integer command identifier 0x7003. Use it only where a procedure explicitly accepts a command identifier.

\paragraph{Example macro.}
\begin{verbatim}
$MACRO REFERENCE_CONSTANT_BLOCK_END FROM EDIT;
    DEF_INT(Value);
    Value := BLOCK_END;
END_MACRO;
\end{verbatim}

\phantomsection
\label{prim:constant-block-off}
\subsection*{BLOCK\_OFF --- constant}
\index{BLOCK\_OFF!constant@\textsc{constant}}
\begin{verbatim}
BLOCK_OFF
\end{verbatim}
Parameters: None. Result: integer constant.

Provides legacy integer command identifier 0x7004. Use it only where a procedure explicitly accepts a command identifier.

\paragraph{Example macro.}
\begin{verbatim}
$MACRO REFERENCE_CONSTANT_BLOCK_OFF FROM EDIT;
    DEF_INT(Value);
    Value := BLOCK_OFF;
END_MACRO;
\end{verbatim}

\phantomsection
\label{prim:constant-col-block-begin}
\subsection*{COL\_BLOCK\_BEGIN --- constant}
\index{COL\_BLOCK\_BEGIN!constant@\textsc{constant}}
\begin{verbatim}
COL_BLOCK_BEGIN
\end{verbatim}
Parameters: None. Result: integer constant.

Provides legacy integer command identifier 0x7005. Use it only where a procedure explicitly accepts a command identifier.

\paragraph{Example macro.}
\begin{verbatim}
$MACRO REFERENCE_CONSTANT_COL_BLOCK_BEGIN FROM EDIT;
    DEF_INT(Value);
    Value := COL_BLOCK_BEGIN;
END_MACRO;
\end{verbatim}

\phantomsection
\label{prim:constant-copy-block}
\subsection*{COPY\_BLOCK --- constant}
\index{COPY\_BLOCK!constant@\textsc{constant}}
\begin{verbatim}
COPY_BLOCK
\end{verbatim}
Parameters: None. Result: integer constant.

Provides legacy integer command identifier 0x7006. Use it only where a procedure explicitly accepts a command identifier.

\paragraph{Example macro.}
\begin{verbatim}
$MACRO REFERENCE_CONSTANT_COPY_BLOCK FROM EDIT;
    DEF_INT(Value);
    Value := COPY_BLOCK;
END_MACRO;
\end{verbatim}

\phantomsection
\label{prim:constant-cr}
\subsection*{CR --- constant}
\index{CR!constant@\textsc{constant}}
\begin{verbatim}
CR
\end{verbatim}
Parameters: None. Result: integer constant.

Provides legacy integer command identifier 0x7007. Use it only where a procedure explicitly accepts a command identifier.

\paragraph{Example macro.}
\begin{verbatim}
$MACRO REFERENCE_CONSTANT_CR FROM EDIT;
    DEF_INT(Value);
    Value := CR;
END_MACRO;
\end{verbatim}

\phantomsection
\label{prim:constant-delete-block}
\subsection*{DELETE\_BLOCK --- constant}
\index{DELETE\_BLOCK!constant@\textsc{constant}}
\begin{verbatim}
DELETE_BLOCK
\end{verbatim}
Parameters: None. Result: integer constant.

Provides legacy integer command identifier 0x7008. Use it only where a procedure explicitly accepts a command identifier.

\paragraph{Example macro.}
\begin{verbatim}
$MACRO REFERENCE_CONSTANT_DELETE_BLOCK FROM EDIT;
    DEF_INT(Value);
    Value := DELETE_BLOCK;
END_MACRO;
\end{verbatim}

\phantomsection
\label{prim:constant-del-char}
\subsection*{DEL\_CHAR --- constant}
\index{DEL\_CHAR!constant@\textsc{constant}}
\begin{verbatim}
DEL_CHAR
\end{verbatim}
Parameters: None. Result: integer constant.

Provides legacy integer command identifier 0x7009. Use it only where a procedure explicitly accepts a command identifier.

\paragraph{Example macro.}
\begin{verbatim}
$MACRO REFERENCE_CONSTANT_DEL_CHAR FROM EDIT;
    DEF_INT(Value);
    Value := DEL_CHAR;
END_MACRO;
\end{verbatim}

\phantomsection
\label{prim:constant-del-line}
\subsection*{DEL\_LINE --- constant}
\index{DEL\_LINE!constant@\textsc{constant}}
\begin{verbatim}
DEL_LINE
\end{verbatim}
Parameters: None. Result: integer constant.

Provides legacy integer command identifier 0x700A. Use it only where a procedure explicitly accepts a command identifier.

\paragraph{Example macro.}
\begin{verbatim}
$MACRO REFERENCE_CONSTANT_DEL_LINE FROM EDIT;
    DEF_INT(Value);
    Value := DEL_LINE;
END_MACRO;
\end{verbatim}

\phantomsection
\label{prim:constant-down}
\subsection*{DOWN --- constant}
\index{DOWN!constant@\textsc{constant}}
\begin{verbatim}
DOWN
\end{verbatim}
Parameters: None. Result: integer constant.

Provides legacy integer command identifier 0x700B. Use it only where a procedure explicitly accepts a command identifier.

\paragraph{Example macro.}
\begin{verbatim}
$MACRO REFERENCE_CONSTANT_DOWN FROM EDIT;
    DEF_INT(Value);
    Value := DOWN;
END_MACRO;
\end{verbatim}

\phantomsection
\label{prim:constant-eof}
\subsection*{EOF --- constant}
\index{EOF!constant@\textsc{constant}}
\begin{verbatim}
EOF
\end{verbatim}
Parameters: None. Result: integer constant.

Provides legacy integer command identifier 0x700C. Use it only where a procedure explicitly accepts a command identifier.

\paragraph{Example macro.}
\begin{verbatim}
$MACRO REFERENCE_CONSTANT_EOF FROM EDIT;
    DEF_INT(Value);
    Value := EOF;
END_MACRO;
\end{verbatim}

\phantomsection
\label{prim:constant-eol}
\subsection*{EOL --- constant}
\index{EOL!constant@\textsc{constant}}
\begin{verbatim}
EOL
\end{verbatim}
Parameters: None. Result: integer constant.

Provides legacy integer command identifier 0x700D. Use it only where a procedure explicitly accepts a command identifier.

\paragraph{Example macro.}
\begin{verbatim}
$MACRO REFERENCE_CONSTANT_EOL FROM EDIT;
    DEF_INT(Value);
    Value := EOL;
END_MACRO;
\end{verbatim}

\phantomsection
\label{prim:constant-first-word}
\subsection*{FIRST\_WORD --- constant}
\index{FIRST\_WORD!constant@\textsc{constant}}
\begin{verbatim}
FIRST_WORD
\end{verbatim}
Parameters: None. Result: integer constant.

Provides legacy integer command identifier 0x700E. Use it only where a procedure explicitly accepts a command identifier.

\paragraph{Example macro.}
\begin{verbatim}
$MACRO REFERENCE_CONSTANT_FIRST_WORD FROM EDIT;
    DEF_INT(Value);
    Value := FIRST_WORD;
END_MACRO;
\end{verbatim}

\phantomsection
\label{prim:constant-goto-mark}
\subsection*{GOTO\_MARK --- constant}
\index{GOTO\_MARK!constant@\textsc{constant}}
\begin{verbatim}
GOTO_MARK
\end{verbatim}
Parameters: None. Result: integer constant.

Provides legacy integer command identifier 0x700F. Use it only where a procedure explicitly accepts a command identifier.

\paragraph{Example macro.}
\begin{verbatim}
$MACRO REFERENCE_CONSTANT_GOTO_MARK FROM EDIT;
    DEF_INT(Value);
    Value := GOTO_MARK;
END_MACRO;
\end{verbatim}

\phantomsection
\label{prim:constant-home}
\subsection*{HOME --- constant}
\index{HOME!constant@\textsc{constant}}
\begin{verbatim}
HOME
\end{verbatim}
Parameters: None. Result: integer constant.

Provides legacy integer command identifier 0x7010. Use it only where a procedure explicitly accepts a command identifier.

\paragraph{Example macro.}
\begin{verbatim}
$MACRO REFERENCE_CONSTANT_HOME FROM EDIT;
    DEF_INT(Value);
    Value := HOME;
END_MACRO;
\end{verbatim}

\phantomsection
\label{prim:constant-indent}
\subsection*{INDENT --- constant}
\index{INDENT!constant@\textsc{constant}}
\begin{verbatim}
INDENT
\end{verbatim}
Parameters: None. Result: integer constant.

Provides legacy integer command identifier 0x7011. Use it only where a procedure explicitly accepts a command identifier.

\paragraph{Example macro.}
\begin{verbatim}
$MACRO REFERENCE_CONSTANT_INDENT FROM EDIT;
    DEF_INT(Value);
    Value := INDENT;
END_MACRO;
\end{verbatim}

\phantomsection
\label{prim:constant-key-record}
\subsection*{KEY\_RECORD --- constant}
\index{KEY\_RECORD!constant@\textsc{constant}}
\begin{verbatim}
KEY_RECORD
\end{verbatim}
Parameters: None. Result: integer constant.

Provides legacy integer command identifier 0x7012. Use it only where a procedure explicitly accepts a command identifier.

\paragraph{Example macro.}
\begin{verbatim}
$MACRO REFERENCE_CONSTANT_KEY_RECORD FROM EDIT;
    DEF_INT(Value);
    Value := KEY_RECORD;
END_MACRO;
\end{verbatim}

\phantomsection
\label{prim:constant-last-page-break}
\subsection*{LAST\_PAGE\_BREAK --- constant}
\index{LAST\_PAGE\_BREAK!constant@\textsc{constant}}
\begin{verbatim}
LAST_PAGE_BREAK
\end{verbatim}
Parameters: None. Result: integer constant.

Provides legacy integer command identifier 0x7013. Use it only where a procedure explicitly accepts a command identifier.

\paragraph{Example macro.}
\begin{verbatim}
$MACRO REFERENCE_CONSTANT_LAST_PAGE_BREAK FROM EDIT;
    DEF_INT(Value);
    Value := LAST_PAGE_BREAK;
END_MACRO;
\end{verbatim}

\phantomsection
\label{prim:constant-left}
\subsection*{LEFT --- constant}
\index{LEFT!constant@\textsc{constant}}
\begin{verbatim}
LEFT
\end{verbatim}
Parameters: None. Result: integer constant.

Provides legacy integer command identifier 0x7014. Use it only where a procedure explicitly accepts a command identifier.

\paragraph{Example macro.}
\begin{verbatim}
$MACRO REFERENCE_CONSTANT_LEFT FROM EDIT;
    DEF_INT(Value);
    Value := LEFT;
END_MACRO;
\end{verbatim}

\phantomsection
\label{prim:constant-mark-pos}
\subsection*{MARK\_POS --- constant}
\index{MARK\_POS!constant@\textsc{constant}}
\begin{verbatim}
MARK_POS
\end{verbatim}
Parameters: None. Result: integer constant.

Provides legacy integer command identifier 0x7015. Use it only where a procedure explicitly accepts a command identifier.

\paragraph{Example macro.}
\begin{verbatim}
$MACRO REFERENCE_CONSTANT_MARK_POS FROM EDIT;
    DEF_INT(Value);
    Value := MARK_POS;
END_MACRO;
\end{verbatim}

\phantomsection
\label{prim:constant-move-block}
\subsection*{MOVE\_BLOCK --- constant}
\index{MOVE\_BLOCK!constant@\textsc{constant}}
\begin{verbatim}
MOVE_BLOCK
\end{verbatim}
Parameters: None. Result: integer constant.

Provides legacy integer command identifier 0x7016. Use it only where a procedure explicitly accepts a command identifier.

\paragraph{Example macro.}
\begin{verbatim}
$MACRO REFERENCE_CONSTANT_MOVE_BLOCK FROM EDIT;
    DEF_INT(Value);
    Value := MOVE_BLOCK;
END_MACRO;
\end{verbatim}

\phantomsection
\label{prim:constant-next-page-break}
\subsection*{NEXT\_PAGE\_BREAK --- constant}
\index{NEXT\_PAGE\_BREAK!constant@\textsc{constant}}
\begin{verbatim}
NEXT_PAGE_BREAK
\end{verbatim}
Parameters: None. Result: integer constant.

Provides legacy integer command identifier 0x7017. Use it only where a procedure explicitly accepts a command identifier.

\paragraph{Example macro.}
\begin{verbatim}
$MACRO REFERENCE_CONSTANT_NEXT_PAGE_BREAK FROM EDIT;
    DEF_INT(Value);
    Value := NEXT_PAGE_BREAK;
END_MACRO;
\end{verbatim}

\phantomsection
\label{prim:constant-page-down}
\subsection*{PAGE\_DOWN --- constant}
\index{PAGE\_DOWN!constant@\textsc{constant}}
\begin{verbatim}
PAGE_DOWN
\end{verbatim}
Parameters: None. Result: integer constant.

Provides legacy integer command identifier 0x7018. Use it only where a procedure explicitly accepts a command identifier.

\paragraph{Example macro.}
\begin{verbatim}
$MACRO REFERENCE_CONSTANT_PAGE_DOWN FROM EDIT;
    DEF_INT(Value);
    Value := PAGE_DOWN;
END_MACRO;
\end{verbatim}

\phantomsection
\label{prim:constant-page-up}
\subsection*{PAGE\_UP --- constant}
\index{PAGE\_UP!constant@\textsc{constant}}
\begin{verbatim}
PAGE_UP
\end{verbatim}
Parameters: None. Result: integer constant.

Provides legacy integer command identifier 0x7019. Use it only where a procedure explicitly accepts a command identifier.

\paragraph{Example macro.}
\begin{verbatim}
$MACRO REFERENCE_CONSTANT_PAGE_UP FROM EDIT;
    DEF_INT(Value);
    Value := PAGE_UP;
END_MACRO;
\end{verbatim}

\phantomsection
\label{prim:constant-right}
\subsection*{RIGHT --- constant}
\index{RIGHT!constant@\textsc{constant}}
\begin{verbatim}
RIGHT
\end{verbatim}
Parameters: None. Result: integer constant.

Provides legacy integer command identifier 0x701A. Use it only where a procedure explicitly accepts a command identifier.

\paragraph{Example macro.}
\begin{verbatim}
$MACRO REFERENCE_CONSTANT_RIGHT FROM EDIT;
    DEF_INT(Value);
    Value := RIGHT;
END_MACRO;
\end{verbatim}

\phantomsection
\label{prim:constant-save-file}
\subsection*{SAVE\_FILE --- constant}
\index{SAVE\_FILE!constant@\textsc{constant}}
\begin{verbatim}
SAVE_FILE
\end{verbatim}
Parameters: None. Result: integer constant.

Provides legacy integer command identifier 0x701B. Use it only where a procedure explicitly accepts a command identifier.

\paragraph{Example macro.}
\begin{verbatim}
$MACRO REFERENCE_CONSTANT_SAVE_FILE FROM EDIT;
    DEF_INT(Value);
    Value := SAVE_FILE;
END_MACRO;
\end{verbatim}

\phantomsection
\label{prim:constant-str-block-begin}
\subsection*{STR\_BLOCK\_BEGIN --- constant}
\index{STR\_BLOCK\_BEGIN!constant@\textsc{constant}}
\begin{verbatim}
STR_BLOCK_BEGIN
\end{verbatim}
Parameters: None. Result: integer constant.

Provides legacy integer command identifier 0x701C. Use it only where a procedure explicitly accepts a command identifier.

\paragraph{Example macro.}
\begin{verbatim}
$MACRO REFERENCE_CONSTANT_STR_BLOCK_BEGIN FROM EDIT;
    DEF_INT(Value);
    Value := STR_BLOCK_BEGIN;
END_MACRO;
\end{verbatim}

\phantomsection
\label{prim:constant-tab-left}
\subsection*{TAB\_LEFT --- constant}
\index{TAB\_LEFT!constant@\textsc{constant}}
\begin{verbatim}
TAB_LEFT
\end{verbatim}
Parameters: None. Result: integer constant.

Provides legacy integer command identifier 0x701D. Use it only where a procedure explicitly accepts a command identifier.

\paragraph{Example macro.}
\begin{verbatim}
$MACRO REFERENCE_CONSTANT_TAB_LEFT FROM EDIT;
    DEF_INT(Value);
    Value := TAB_LEFT;
END_MACRO;
\end{verbatim}

\phantomsection
\label{prim:constant-tab-right}
\subsection*{TAB\_RIGHT --- constant}
\index{TAB\_RIGHT!constant@\textsc{constant}}
\begin{verbatim}
TAB_RIGHT
\end{verbatim}
Parameters: None. Result: integer constant.

Provides legacy integer command identifier 0x701E. Use it only where a procedure explicitly accepts a command identifier.

\paragraph{Example macro.}
\begin{verbatim}
$MACRO REFERENCE_CONSTANT_TAB_RIGHT FROM EDIT;
    DEF_INT(Value);
    Value := TAB_RIGHT;
END_MACRO;
\end{verbatim}

\phantomsection
\label{prim:constant-tof}
\subsection*{TOF --- constant}
\index{TOF!constant@\textsc{constant}}
\begin{verbatim}
TOF
\end{verbatim}
Parameters: None. Result: integer constant.

Provides legacy integer command identifier 0x701F. Use it only where a procedure explicitly accepts a command identifier.

\paragraph{Example macro.}
\begin{verbatim}
$MACRO REFERENCE_CONSTANT_TOF FROM EDIT;
    DEF_INT(Value);
    Value := TOF;
END_MACRO;
\end{verbatim}

\phantomsection
\label{prim:constant-undent}
\subsection*{UNDENT --- constant}
\index{UNDENT!constant@\textsc{constant}}
\begin{verbatim}
UNDENT
\end{verbatim}
Parameters: None. Result: integer constant.

Provides legacy integer command identifier 0x7020. Use it only where a procedure explicitly accepts a command identifier.

\paragraph{Example macro.}
\begin{verbatim}
$MACRO REFERENCE_CONSTANT_UNDENT FROM EDIT;
    DEF_INT(Value);
    Value := UNDENT;
END_MACRO;
\end{verbatim}

\phantomsection
\label{prim:constant-undo}
\subsection*{UNDO --- constant}
\index{UNDO!constant@\textsc{constant}}
\begin{verbatim}
UNDO
\end{verbatim}
Parameters: None. Result: integer constant.

Provides legacy integer command identifier 0x7021. Use it only where a procedure explicitly accepts a command identifier.

\paragraph{Example macro.}
\begin{verbatim}
$MACRO REFERENCE_CONSTANT_UNDO FROM EDIT;
    DEF_INT(Value);
    Value := UNDO;
END_MACRO;
\end{verbatim}

\phantomsection
\label{prim:constant-up}
\subsection*{UP --- constant}
\index{UP!constant@\textsc{constant}}
\begin{verbatim}
UP
\end{verbatim}
Parameters: None. Result: integer constant.

Provides legacy integer command identifier 0x7022. Use it only where a procedure explicitly accepts a command identifier.

\paragraph{Example macro.}
\begin{verbatim}
$MACRO REFERENCE_CONSTANT_UP FROM EDIT;
    DEF_INT(Value);
    Value := UP;
END_MACRO;
\end{verbatim}

\phantomsection
\label{prim:constant-word-left}
\subsection*{WORD\_LEFT --- constant}
\index{WORD\_LEFT!constant@\textsc{constant}}
\begin{verbatim}
WORD_LEFT
\end{verbatim}
Parameters: None. Result: integer constant.

Provides legacy integer command identifier 0x7023. Use it only where a procedure explicitly accepts a command identifier.

\paragraph{Example macro.}
\begin{verbatim}
$MACRO REFERENCE_CONSTANT_WORD_LEFT FROM EDIT;
    DEF_INT(Value);
    Value := WORD_LEFT;
END_MACRO;
\end{verbatim}

\phantomsection
\label{prim:constant-word-right}
\subsection*{WORD\_RIGHT --- constant}
\index{WORD\_RIGHT!constant@\textsc{constant}}
\begin{verbatim}
WORD_RIGHT
\end{verbatim}
Parameters: None. Result: integer constant.

Provides legacy integer command identifier 0x7024. Use it only where a procedure explicitly accepts a command identifier.

\paragraph{Example macro.}
\begin{verbatim}
$MACRO REFERENCE_CONSTANT_WORD_RIGHT FROM EDIT;
    DEF_INT(Value);
    Value := WORD_RIGHT;
END_MACRO;
\end{verbatim}

\phantomsection
\label{prim:constant-all}
\subsection*{ALL --- constant}
\index{ALL!constant@\textsc{constant}}
\begin{verbatim}
ALL
\end{verbatim}
Parameters: None. Result: integer constant.

Provides compiler-defined macro mode value 255. DOS\_SHELL is historical source syntax; MRMAC provides no DOS shell runtime.

\paragraph{Example macro.}
\begin{verbatim}
$MACRO REFERENCE_CONSTANT_ALL FROM EDIT;
    DEF_INT(Value);
    Value := ALL;
END_MACRO;
\end{verbatim}

\phantomsection
\label{prim:variable-return-int}
\subsection*{RETURN\_INT --- variable}
\index{RETURN\_INT!variable@\textsc{variable}}
\begin{verbatim}
RETURN_INT
\end{verbatim}
Parameters: None. Result: integer value.

Provides the integer return slot of the current macro. Assign a value before RET to return an integer result.

\paragraph{Example macro.}
\begin{verbatim}
$MACRO REFERENCE_VARIABLE_RETURN_INT FROM EDIT;
    DEF_INT(Value);
    Value := RETURN_INT;
END_MACRO;
\end{verbatim}

\phantomsection
\label{prim:variable-error-level}
\subsection*{ERROR\_LEVEL --- variable}
\index{ERROR\_LEVEL!variable@\textsc{variable}}
\begin{verbatim}
ERROR_LEVEL
\end{verbatim}
Parameters: None. Result: integer value.

Provides the current macro error level. The VM permits assignment to set the level for the active execution.

\paragraph{Example macro.}
\begin{verbatim}
$MACRO REFERENCE_VARIABLE_ERROR_LEVEL FROM EDIT;
    DEF_INT(Value);
    Value := ERROR_LEVEL;
END_MACRO;
\end{verbatim}

\phantomsection
\label{prim:variable-first-run}
\subsection*{FIRST\_RUN --- variable}
\index{FIRST\_RUN!variable@\textsc{variable}}
\begin{verbatim}
FIRST_RUN
\end{verbatim}
Parameters: None. Result: integer value.

Returns 1 only on the first run of the current macro stack entry; otherwise 0.

\paragraph{Example macro.}
\begin{verbatim}
$MACRO REFERENCE_VARIABLE_FIRST_RUN FROM EDIT;
    DEF_INT(Value);
    Value := FIRST_RUN;
END_MACRO;
\end{verbatim}

\phantomsection
\label{prim:variable-ignore-case}
\subsection*{IGNORE\_CASE --- variable}
\index{IGNORE\_CASE!variable@\textsc{variable}}
\begin{verbatim}
IGNORE_CASE
\end{verbatim}
Parameters: None. Result: integer value.

Provides the active search case-sensitivity flag as 1 or 0. The VM permits assignment.

\paragraph{Example macro.}
\begin{verbatim}
$MACRO REFERENCE_VARIABLE_IGNORE_CASE FROM EDIT;
    DEF_INT(Value);
    Value := IGNORE_CASE;
END_MACRO;
\end{verbatim}

\phantomsection
\label{prim:variable-reg-exp-stat}
\subsection*{REG\_EXP\_STAT --- variable}
\index{REG\_EXP\_STAT!variable@\textsc{variable}}
\begin{verbatim}
REG_EXP_STAT
\end{verbatim}
Parameters: None. Result: integer value.

Provides the active regular-expression status flag as 1 or 0. The VM permits assignment.

\paragraph{Example macro.}
\begin{verbatim}
$MACRO REFERENCE_VARIABLE_REG_EXP_STAT FROM EDIT;
    DEF_INT(Value);
    Value := REG_EXP_STAT;
END_MACRO;
\end{verbatim}

\phantomsection
\label{prim:variable-first-save}
\subsection*{FIRST\_SAVE --- variable}
\index{FIRST\_SAVE!variable@\textsc{variable}}
\begin{verbatim}
FIRST_SAVE
\end{verbatim}
Parameters: None. Result: integer value.

Provides current or last-file metadata named first save.

\paragraph{Example macro.}
\begin{verbatim}
$MACRO REFERENCE_VARIABLE_FIRST_SAVE FROM EDIT;
    DEF_INT(Value);
    Value := FIRST_SAVE;
END_MACRO;
\end{verbatim}

\phantomsection
\label{prim:variable-buffer-id}
\subsection*{BUFFER\_ID --- variable}
\index{BUFFER\_ID!variable@\textsc{variable}}
\begin{verbatim}
BUFFER_ID
\end{verbatim}
Parameters: None. Result: integer value.

Provides current or last-file metadata named buffer id.

\paragraph{Example macro.}
\begin{verbatim}
$MACRO REFERENCE_VARIABLE_BUFFER_ID FROM EDIT;
    DEF_INT(Value);
    Value := BUFFER_ID;
END_MACRO;
\end{verbatim}

\phantomsection
\label{prim:variable-tmp-file}
\subsection*{TMP\_FILE --- variable}
\index{TMP\_FILE!variable@\textsc{variable}}
\begin{verbatim}
TMP_FILE
\end{verbatim}
Parameters: None. Result: integer value.

Provides current or last-file metadata named tmp file.

\paragraph{Example macro.}
\begin{verbatim}
$MACRO REFERENCE_VARIABLE_TMP_FILE FROM EDIT;
    DEF_INT(Value);
    Value := TMP_FILE;
END_MACRO;
\end{verbatim}

\phantomsection
\label{prim:variable-file-changed}
\subsection*{FILE\_CHANGED --- variable}
\index{FILE\_CHANGED!variable@\textsc{variable}}
\begin{verbatim}
FILE_CHANGED
\end{verbatim}
Parameters: None. Result: integer value.

Provides whether the active file has unsaved changes as 1 or 0. The VM permits assignment when an active file context exists.

\paragraph{Example macro.}
\begin{verbatim}
$MACRO REFERENCE_VARIABLE_FILE_CHANGED FROM EDIT;
    DEF_INT(Value);
    Value := FILE_CHANGED;
END_MACRO;
\end{verbatim}

\phantomsection
\label{prim:variable-last-file-attr}
\subsection*{LAST\_FILE\_ATTR --- variable}
\index{LAST\_FILE\_ATTR!variable@\textsc{variable}}
\begin{verbatim}
LAST_FILE_ATTR
\end{verbatim}
Parameters: None. Result: integer value.

Provides current or last-file metadata named last file attr.

\paragraph{Example macro.}
\begin{verbatim}
$MACRO REFERENCE_VARIABLE_LAST_FILE_ATTR FROM EDIT;
    DEF_INT(Value);
    Value := LAST_FILE_ATTR;
END_MACRO;
\end{verbatim}

\phantomsection
\label{prim:variable-last-file-size}
\subsection*{LAST\_FILE\_SIZE --- variable}
\index{LAST\_FILE\_SIZE!variable@\textsc{variable}}
\begin{verbatim}
LAST_FILE_SIZE
\end{verbatim}
Parameters: None. Result: integer value.

Provides current or last-file metadata named last file size.

\paragraph{Example macro.}
\begin{verbatim}
$MACRO REFERENCE_VARIABLE_LAST_FILE_SIZE FROM EDIT;
    DEF_INT(Value);
    Value := LAST_FILE_SIZE;
END_MACRO;
\end{verbatim}

\phantomsection
\label{prim:variable-last-file-time}
\subsection*{LAST\_FILE\_TIME --- variable}
\index{LAST\_FILE\_TIME!variable@\textsc{variable}}
\begin{verbatim}
LAST_FILE_TIME
\end{verbatim}
Parameters: None. Result: integer value.

Provides current or last-file metadata named last file time.

\paragraph{Example macro.}
\begin{verbatim}
$MACRO REFERENCE_VARIABLE_LAST_FILE_TIME FROM EDIT;
    DEF_INT(Value);
    Value := LAST_FILE_TIME;
END_MACRO;
\end{verbatim}

\phantomsection
\label{prim:variable-cur-file-attr}
\subsection*{CUR\_FILE\_ATTR --- variable}
\index{CUR\_FILE\_ATTR!variable@\textsc{variable}}
\begin{verbatim}
CUR_FILE_ATTR
\end{verbatim}
Parameters: None. Result: integer value.

Provides current or last-file metadata named cur file attr.

\paragraph{Example macro.}
\begin{verbatim}
$MACRO REFERENCE_VARIABLE_CUR_FILE_ATTR FROM EDIT;
    DEF_INT(Value);
    Value := CUR_FILE_ATTR;
END_MACRO;
\end{verbatim}

\phantomsection
\label{prim:variable-cur-file-size}
\subsection*{CUR\_FILE\_SIZE --- variable}
\index{CUR\_FILE\_SIZE!variable@\textsc{variable}}
\begin{verbatim}
CUR_FILE_SIZE
\end{verbatim}
Parameters: None. Result: integer value.

Provides current or last-file metadata named cur file size.

\paragraph{Example macro.}
\begin{verbatim}
$MACRO REFERENCE_VARIABLE_CUR_FILE_SIZE FROM EDIT;
    DEF_INT(Value);
    Value := CUR_FILE_SIZE;
END_MACRO;
\end{verbatim}

\phantomsection
\label{prim:variable-read-only}
\subsection*{READ\_ONLY --- variable}
\index{READ\_ONLY!variable@\textsc{variable}}
\begin{verbatim}
READ_ONLY
\end{verbatim}
Parameters: None. Result: integer value.

Provides current or last-file metadata named read only.

\paragraph{Example macro.}
\begin{verbatim}
$MACRO REFERENCE_VARIABLE_READ_ONLY FROM EDIT;
    DEF_INT(Value);
    Value := READ_ONLY;
END_MACRO;
\end{verbatim}

\phantomsection
\label{prim:variable-doc-mode}
\subsection*{DOC\_MODE --- variable}
\index{DOC\_MODE!variable@\textsc{variable}}
\begin{verbatim}
DOC_MODE
\end{verbatim}
Parameters: None. Result: integer value.

Provides the current runtime value named doc mode.

\paragraph{Example macro.}
\begin{verbatim}
$MACRO REFERENCE_VARIABLE_DOC_MODE FROM EDIT;
    DEF_INT(Value);
    Value := DOC_MODE;
END_MACRO;
\end{verbatim}

\phantomsection
\label{prim:variable-print-margin}
\subsection*{PRINT\_MARGIN --- variable}
\index{PRINT\_MARGIN!variable@\textsc{variable}}
\begin{verbatim}
PRINT_MARGIN
\end{verbatim}
Parameters: None. Result: integer value.

Provides the current runtime value named print margin.

\paragraph{Example macro.}
\begin{verbatim}
$MACRO REFERENCE_VARIABLE_PRINT_MARGIN FROM EDIT;
    DEF_INT(Value);
    Value := PRINT_MARGIN;
END_MACRO;
\end{verbatim}

\phantomsection
\label{prim:variable-shadow-char}
\subsection*{SHADOW\_CHAR --- variable}
\index{SHADOW\_CHAR!variable@\textsc{variable}}
\begin{verbatim}
SHADOW_CHAR
\end{verbatim}
Parameters: None. Result: integer value.

Provides the current runtime value named shadow char.

\paragraph{Example macro.}
\begin{verbatim}
$MACRO REFERENCE_VARIABLE_SHADOW_CHAR FROM EDIT;
    DEF_INT(Value);
    Value := SHADOW_CHAR;
END_MACRO;
\end{verbatim}

\phantomsection
\label{prim:variable-refresh}
\subsection*{REFRESH --- variable}
\index{REFRESH!variable@\textsc{variable}}
\begin{verbatim}
REFRESH
\end{verbatim}
Parameters: None. Result: integer value.

Provides the current runtime value named refresh.

\paragraph{Example macro.}
\begin{verbatim}
$MACRO REFERENCE_VARIABLE_REFRESH FROM EDIT;
    DEF_INT(Value);
    Value := REFRESH;
END_MACRO;
\end{verbatim}

\phantomsection
\label{prim:variable-messages}
\subsection*{MESSAGES --- variable}
\index{MESSAGES!variable@\textsc{variable}}
\begin{verbatim}
MESSAGES
\end{verbatim}
Parameters: None. Result: integer value.

Provides the current runtime value named messages.

\paragraph{Example macro.}
\begin{verbatim}
$MACRO REFERENCE_VARIABLE_MESSAGES FROM EDIT;
    DEF_INT(Value);
    Value := MESSAGES;
END_MACRO;
\end{verbatim}

\phantomsection
\label{prim:variable-mouse}
\subsection*{MOUSE --- variable}
\index{MOUSE!variable@\textsc{variable}}
\begin{verbatim}
MOUSE
\end{verbatim}
Parameters: None. Result: integer value.

Provides the current runtime value named mouse.

\paragraph{Example macro.}
\begin{verbatim}
$MACRO REFERENCE_VARIABLE_MOUSE FROM EDIT;
    DEF_INT(Value);
    Value := MOUSE;
END_MACRO;
\end{verbatim}

\phantomsection
\label{prim:variable-logo-screen}
\subsection*{LOGO\_SCREEN --- variable}
\index{LOGO\_SCREEN!variable@\textsc{variable}}
\begin{verbatim}
LOGO_SCREEN
\end{verbatim}
Parameters: None. Result: integer value.

Provides the current runtime value named logo screen.

\paragraph{Example macro.}
\begin{verbatim}
$MACRO REFERENCE_VARIABLE_LOGO_SCREEN FROM EDIT;
    DEF_INT(Value);
    Value := LOGO_SCREEN;
END_MACRO;
\end{verbatim}

\phantomsection
\label{prim:variable-explosions}
\subsection*{EXPLOSIONS --- variable}
\index{EXPLOSIONS!variable@\textsc{variable}}
\begin{verbatim}
EXPLOSIONS
\end{verbatim}
Parameters: None. Result: integer value.

Provides the current runtime value named explosions.

\paragraph{Example macro.}
\begin{verbatim}
$MACRO REFERENCE_VARIABLE_EXPLOSIONS FROM EDIT;
    DEF_INT(Value);
    Value := EXPLOSIONS;
END_MACRO;
\end{verbatim}

\phantomsection
\label{prim:variable-truncate-spaces}
\subsection*{TRUNCATE\_SPACES --- variable}
\index{TRUNCATE\_SPACES!variable@\textsc{variable}}
\begin{verbatim}
TRUNCATE_SPACES
\end{verbatim}
Parameters: None. Result: integer value.

Provides the active editor-format value named truncate spaces; writable settings use the configured runtime boundary.

\paragraph{Example macro.}
\begin{verbatim}
$MACRO REFERENCE_VARIABLE_TRUNCATE_SPACES FROM EDIT;
    DEF_INT(Value);
    Value := TRUNCATE_SPACES;
END_MACRO;
\end{verbatim}

\phantomsection
\label{prim:variable-backups}
\subsection*{BACKUPS --- variable}
\index{BACKUPS!variable@\textsc{variable}}
\begin{verbatim}
BACKUPS
\end{verbatim}
Parameters: None. Result: integer value.

Provides the active editor-format value named backups; writable settings use the configured runtime boundary.

\paragraph{Example macro.}
\begin{verbatim}
$MACRO REFERENCE_VARIABLE_BACKUPS FROM EDIT;
    DEF_INT(Value);
    Value := BACKUPS;
END_MACRO;
\end{verbatim}

\phantomsection
\label{prim:variable-autosave}
\subsection*{AUTOSAVE --- variable}
\index{AUTOSAVE!variable@\textsc{variable}}
\begin{verbatim}
AUTOSAVE
\end{verbatim}
Parameters: None. Result: integer value.

Provides the active editor-format value named autosave; writable settings use the configured runtime boundary.

\paragraph{Example macro.}
\begin{verbatim}
$MACRO REFERENCE_VARIABLE_AUTOSAVE FROM EDIT;
    DEF_INT(Value);
    Value := AUTOSAVE;
END_MACRO;
\end{verbatim}

\phantomsection
\label{prim:variable-undo-stat}
\subsection*{UNDO\_STAT --- variable}
\index{UNDO\_STAT!variable@\textsc{variable}}
\begin{verbatim}
UNDO_STAT
\end{verbatim}
Parameters: None. Result: integer value.

Provides the current runtime value named undo stat.

\paragraph{Example macro.}
\begin{verbatim}
$MACRO REFERENCE_VARIABLE_UNDO_STAT FROM EDIT;
    DEF_INT(Value);
    Value := UNDO_STAT;
END_MACRO;
\end{verbatim}

\phantomsection
\label{prim:variable-format-stat}
\subsection*{FORMAT\_STAT --- variable}
\index{FORMAT\_STAT!variable@\textsc{variable}}
\begin{verbatim}
FORMAT_STAT
\end{verbatim}
Parameters: None. Result: integer value.

Provides the current runtime value named format stat.

\paragraph{Example macro.}
\begin{verbatim}
$MACRO REFERENCE_VARIABLE_FORMAT_STAT FROM EDIT;
    DEF_INT(Value);
    Value := FORMAT_STAT;
END_MACRO;
\end{verbatim}

\phantomsection
\label{prim:variable-wrap-stat}
\subsection*{WRAP\_STAT --- variable}
\index{WRAP\_STAT!variable@\textsc{variable}}
\begin{verbatim}
WRAP_STAT
\end{verbatim}
Parameters: None. Result: integer value.

Provides the active editor-format value named wrap stat; writable settings use the configured runtime boundary.

\paragraph{Example macro.}
\begin{verbatim}
$MACRO REFERENCE_VARIABLE_WRAP_STAT FROM EDIT;
    DEF_INT(Value);
    Value := WRAP_STAT;
END_MACRO;
\end{verbatim}

\phantomsection
\label{prim:variable-mem-alloc}
\subsection*{MEM\_ALLOC --- variable}
\index{MEM\_ALLOC!variable@\textsc{variable}}
\begin{verbatim}
MEM_ALLOC
\end{verbatim}
Parameters: None. Result: integer value.

Provides the current runtime value named mem alloc.

\paragraph{Example macro.}
\begin{verbatim}
$MACRO REFERENCE_VARIABLE_MEM_ALLOC FROM EDIT;
    DEF_INT(Value);
    Value := MEM_ALLOC;
END_MACRO;
\end{verbatim}

\phantomsection
\label{prim:variable-left-margin}
\subsection*{LEFT\_MARGIN --- variable}
\index{LEFT\_MARGIN!variable@\textsc{variable}}
\begin{verbatim}
LEFT_MARGIN
\end{verbatim}
Parameters: None. Result: integer value.

Provides the active editor-format value named left margin; writable settings use the configured runtime boundary.

\paragraph{Example macro.}
\begin{verbatim}
$MACRO REFERENCE_VARIABLE_LEFT_MARGIN FROM EDIT;
    DEF_INT(Value);
    Value := LEFT_MARGIN;
END_MACRO;
\end{verbatim}

\phantomsection
\label{prim:variable-right-margin}
\subsection*{RIGHT\_MARGIN --- variable}
\index{RIGHT\_MARGIN!variable@\textsc{variable}}
\begin{verbatim}
RIGHT_MARGIN
\end{verbatim}
Parameters: None. Result: integer value.

Provides the active editor-format value named right margin; writable settings use the configured runtime boundary.

\paragraph{Example macro.}
\begin{verbatim}
$MACRO REFERENCE_VARIABLE_RIGHT_MARGIN FROM EDIT;
    DEF_INT(Value);
    Value := RIGHT_MARGIN;
END_MACRO;
\end{verbatim}

\phantomsection
\label{prim:variable-format-ruler}
\subsection*{FORMAT\_RULER --- variable}
\index{FORMAT\_RULER!variable@\textsc{variable}}
\begin{verbatim}
FORMAT_RULER
\end{verbatim}
Parameters: None. Result: integer value.

Provides the active editor-format value named format ruler; writable settings use the configured runtime boundary.

\paragraph{Example macro.}
\begin{verbatim}
$MACRO REFERENCE_VARIABLE_FORMAT_RULER FROM EDIT;
    DEF_INT(Value);
    Value := FORMAT_RULER;
END_MACRO;
\end{verbatim}

\phantomsection
\label{prim:variable-indent-style}
\subsection*{INDENT\_STYLE --- variable}
\index{INDENT\_STYLE!variable@\textsc{variable}}
\begin{verbatim}
INDENT_STYLE
\end{verbatim}
Parameters: None. Result: integer value.

Provides the active editor-format value named indent style; writable settings use the configured runtime boundary.

\paragraph{Example macro.}
\begin{verbatim}
$MACRO REFERENCE_VARIABLE_INDENT_STYLE FROM EDIT;
    DEF_INT(Value);
    Value := INDENT_STYLE;
END_MACRO;
\end{verbatim}

\phantomsection
\label{prim:variable-ins-cursor}
\subsection*{INS\_CURSOR --- variable}
\index{INS\_CURSOR!variable@\textsc{variable}}
\begin{verbatim}
INS_CURSOR
\end{verbatim}
Parameters: None. Result: integer value.

Provides the current runtime value named ins cursor.

\paragraph{Example macro.}
\begin{verbatim}
$MACRO REFERENCE_VARIABLE_INS_CURSOR FROM EDIT;
    DEF_INT(Value);
    Value := INS_CURSOR;
END_MACRO;
\end{verbatim}

\phantomsection
\label{prim:variable-ovr-cursor}
\subsection*{OVR\_CURSOR --- variable}
\index{OVR\_CURSOR!variable@\textsc{variable}}
\begin{verbatim}
OVR_CURSOR
\end{verbatim}
Parameters: None. Result: integer value.

Provides the current runtime value named ovr cursor.

\paragraph{Example macro.}
\begin{verbatim}
$MACRO REFERENCE_VARIABLE_OVR_CURSOR FROM EDIT;
    DEF_INT(Value);
    Value := OVR_CURSOR;
END_MACRO;
\end{verbatim}

\phantomsection
\label{prim:variable-ctrl-help}
\subsection*{CTRL\_HELP --- variable}
\index{CTRL\_HELP!variable@\textsc{variable}}
\begin{verbatim}
CTRL_HELP
\end{verbatim}
Parameters: None. Result: integer value.

Provides the current runtime value named ctrl help.

\paragraph{Example macro.}
\begin{verbatim}
$MACRO REFERENCE_VARIABLE_CTRL_HELP FROM EDIT;
    DEF_INT(Value);
    Value := CTRL_HELP;
END_MACRO;
\end{verbatim}

\phantomsection
\label{prim:variable-mouse-h-sense}
\subsection*{MOUSE\_H\_SENSE --- variable}
\index{MOUSE\_H\_SENSE!variable@\textsc{variable}}
\begin{verbatim}
MOUSE_H_SENSE
\end{verbatim}
Parameters: None. Result: integer value.

Provides the current runtime value named mouse h sense.

\paragraph{Example macro.}
\begin{verbatim}
$MACRO REFERENCE_VARIABLE_MOUSE_H_SENSE FROM EDIT;
    DEF_INT(Value);
    Value := MOUSE_H_SENSE;
END_MACRO;
\end{verbatim}

\phantomsection
\label{prim:variable-mouse-v-sense}
\subsection*{MOUSE\_V\_SENSE --- variable}
\index{MOUSE\_V\_SENSE!variable@\textsc{variable}}
\begin{verbatim}
MOUSE_V_SENSE
\end{verbatim}
Parameters: None. Result: integer value.

Provides the current runtime value named mouse v sense.

\paragraph{Example macro.}
\begin{verbatim}
$MACRO REFERENCE_VARIABLE_MOUSE_V_SENSE FROM EDIT;
    DEF_INT(Value);
    Value := MOUSE_V_SENSE;
END_MACRO;
\end{verbatim}

\phantomsection
\label{prim:variable-window-attr}
\subsection*{WINDOW\_ATTR --- variable}
\index{WINDOW\_ATTR!variable@\textsc{variable}}
\begin{verbatim}
WINDOW_ATTR
\end{verbatim}
Parameters: None. Result: integer value.

Provides the current configured window attr as an integer colour or display attribute.

\paragraph{Example macro.}
\begin{verbatim}
$MACRO REFERENCE_VARIABLE_WINDOW_ATTR FROM EDIT;
    DEF_INT(Value);
    Value := WINDOW_ATTR;
END_MACRO;
\end{verbatim}

\phantomsection
\label{prim:variable-text-color}
\subsection*{TEXT\_COLOR --- variable}
\index{TEXT\_COLOR!variable@\textsc{variable}}
\begin{verbatim}
TEXT_COLOR
\end{verbatim}
Parameters: None. Result: integer value.

Provides the current configured text color as an integer colour or display attribute.

\paragraph{Example macro.}
\begin{verbatim}
$MACRO REFERENCE_VARIABLE_TEXT_COLOR FROM EDIT;
    DEF_INT(Value);
    Value := TEXT_COLOR;
END_MACRO;
\end{verbatim}

\phantomsection
\label{prim:variable-change-color}
\subsection*{CHANGE\_COLOR --- variable}
\index{CHANGE\_COLOR!variable@\textsc{variable}}
\begin{verbatim}
CHANGE_COLOR
\end{verbatim}
Parameters: None. Result: integer value.

Provides the current configured change color as an integer colour or display attribute.

\paragraph{Example macro.}
\begin{verbatim}
$MACRO REFERENCE_VARIABLE_CHANGE_COLOR FROM EDIT;
    DEF_INT(Value);
    Value := CHANGE_COLOR;
END_MACRO;
\end{verbatim}

\phantomsection
\label{prim:variable-back-color}
\subsection*{BACK\_COLOR --- variable}
\index{BACK\_COLOR!variable@\textsc{variable}}
\begin{verbatim}
BACK_COLOR
\end{verbatim}
Parameters: None. Result: integer value.

Provides the current configured back color as an integer colour or display attribute.

\paragraph{Example macro.}
\begin{verbatim}
$MACRO REFERENCE_VARIABLE_BACK_COLOR FROM EDIT;
    DEF_INT(Value);
    Value := BACK_COLOR;
END_MACRO;
\end{verbatim}

\phantomsection
\label{prim:variable-menu-color}
\subsection*{MENU\_COLOR --- variable}
\index{MENU\_COLOR!variable@\textsc{variable}}
\begin{verbatim}
MENU_COLOR
\end{verbatim}
Parameters: None. Result: integer value.

Provides the current configured menu color as an integer colour or display attribute.

\paragraph{Example macro.}
\begin{verbatim}
$MACRO REFERENCE_VARIABLE_MENU_COLOR FROM EDIT;
    DEF_INT(Value);
    Value := MENU_COLOR;
END_MACRO;
\end{verbatim}

\phantomsection
\label{prim:variable-stat-color}
\subsection*{STAT\_COLOR --- variable}
\index{STAT\_COLOR!variable@\textsc{variable}}
\begin{verbatim}
STAT_COLOR
\end{verbatim}
Parameters: None. Result: integer value.

Provides the current configured stat color as an integer colour or display attribute.

\paragraph{Example macro.}
\begin{verbatim}
$MACRO REFERENCE_VARIABLE_STAT_COLOR FROM EDIT;
    DEF_INT(Value);
    Value := STAT_COLOR;
END_MACRO;
\end{verbatim}

\phantomsection
\label{prim:variable-error-color}
\subsection*{ERROR\_COLOR --- variable}
\index{ERROR\_COLOR!variable@\textsc{variable}}
\begin{verbatim}
ERROR_COLOR
\end{verbatim}
Parameters: None. Result: integer value.

Provides the current configured error color as an integer colour or display attribute.

\paragraph{Example macro.}
\begin{verbatim}
$MACRO REFERENCE_VARIABLE_ERROR_COLOR FROM EDIT;
    DEF_INT(Value);
    Value := ERROR_COLOR;
END_MACRO;
\end{verbatim}

\phantomsection
\label{prim:variable-shadow-color}
\subsection*{SHADOW\_COLOR --- variable}
\index{SHADOW\_COLOR!variable@\textsc{variable}}
\begin{verbatim}
SHADOW_COLOR
\end{verbatim}
Parameters: None. Result: integer value.

Provides the current configured shadow color as an integer colour or display attribute.

\paragraph{Example macro.}
\begin{verbatim}
$MACRO REFERENCE_VARIABLE_SHADOW_COLOR FROM EDIT;
    DEF_INT(Value);
    Value := SHADOW_COLOR;
END_MACRO;
\end{verbatim}

\phantomsection
\label{prim:variable-status-row}
\subsection*{STATUS\_ROW --- variable}
\index{STATUS\_ROW!variable@\textsc{variable}}
\begin{verbatim}
STATUS_ROW
\end{verbatim}
Parameters: None. Result: integer value.

Provides the current runtime value named status row.

\paragraph{Example macro.}
\begin{verbatim}
$MACRO REFERENCE_VARIABLE_STATUS_ROW FROM EDIT;
    DEF_INT(Value);
    Value := STATUS_ROW;
END_MACRO;
\end{verbatim}

\phantomsection
\label{prim:variable-message-row}
\subsection*{MESSAGE\_ROW --- variable}
\index{MESSAGE\_ROW!variable@\textsc{variable}}
\begin{verbatim}
MESSAGE_ROW
\end{verbatim}
Parameters: None. Result: integer value.

Provides the current runtime value named message row.

\paragraph{Example macro.}
\begin{verbatim}
$MACRO REFERENCE_VARIABLE_MESSAGE_ROW FROM EDIT;
    DEF_INT(Value);
    Value := MESSAGE_ROW;
END_MACRO;
\end{verbatim}

\phantomsection
\label{prim:variable-max-window-row}
\subsection*{MAX\_WINDOW\_ROW --- variable}
\index{MAX\_WINDOW\_ROW!variable@\textsc{variable}}
\begin{verbatim}
MAX_WINDOW_ROW
\end{verbatim}
Parameters: None. Result: integer value.

Provides the current runtime value named max window row.

\paragraph{Example macro.}
\begin{verbatim}
$MACRO REFERENCE_VARIABLE_MAX_WINDOW_ROW FROM EDIT;
    DEF_INT(Value);
    Value := MAX_WINDOW_ROW;
END_MACRO;
\end{verbatim}

\phantomsection
\label{prim:variable-min-window-row}
\subsection*{MIN\_WINDOW\_ROW --- variable}
\index{MIN\_WINDOW\_ROW!variable@\textsc{variable}}
\begin{verbatim}
MIN_WINDOW_ROW
\end{verbatim}
Parameters: None. Result: integer value.

Provides the current runtime value named min window row.

\paragraph{Example macro.}
\begin{verbatim}
$MACRO REFERENCE_VARIABLE_MIN_WINDOW_ROW FROM EDIT;
    DEF_INT(Value);
    Value := MIN_WINDOW_ROW;
END_MACRO;
\end{verbatim}

\phantomsection
\label{prim:variable-name-line}
\subsection*{NAME\_LINE --- variable}
\index{NAME\_LINE!variable@\textsc{variable}}
\begin{verbatim}
NAME_LINE
\end{verbatim}
Parameters: None. Result: integer value.

Provides the current runtime value named name line.

\paragraph{Example macro.}
\begin{verbatim}
$MACRO REFERENCE_VARIABLE_NAME_LINE FROM EDIT;
    DEF_INT(Value);
    Value := NAME_LINE;
END_MACRO;
\end{verbatim}

\phantomsection
\label{prim:variable-param-count}
\subsection*{PARAM\_COUNT --- variable}
\index{PARAM\_COUNT!variable@\textsc{variable}}
\begin{verbatim}
PARAM_COUNT
\end{verbatim}
Parameters: None. Result: integer value.

Provides the number of process parameters visible to the runtime.

\paragraph{Example macro.}
\begin{verbatim}
$MACRO REFERENCE_VARIABLE_PARAM_COUNT FROM EDIT;
    DEF_INT(Value);
    Value := PARAM_COUNT;
END_MACRO;
\end{verbatim}

\phantomsection
\label{prim:variable-cpu}
\subsection*{CPU --- variable}
\index{CPU!variable@\textsc{variable}}
\begin{verbatim}
CPU
\end{verbatim}
Parameters: None. Result: integer value.

Provides the current runtime value named cpu.

\paragraph{Example macro.}
\begin{verbatim}
$MACRO REFERENCE_VARIABLE_CPU FROM EDIT;
    DEF_INT(Value);
    Value := CPU;
END_MACRO;
\end{verbatim}

\phantomsection
\label{prim:variable-c-col}
\subsection*{C\_COL --- variable}
\index{C\_COL!variable@\textsc{variable}}
\begin{verbatim}
C_COL
\end{verbatim}
Parameters: None. Result: integer value.

Provides the current editor cursor column.

\paragraph{Example macro.}
\begin{verbatim}
$MACRO REFERENCE_VARIABLE_C_COL FROM EDIT;
    DEF_INT(Value);
    Value := C_COL;
END_MACRO;
\end{verbatim}

\phantomsection
\label{prim:variable-c-line}
\subsection*{C\_LINE --- variable}
\index{C\_LINE!variable@\textsc{variable}}
\begin{verbatim}
C_LINE
\end{verbatim}
Parameters: None. Result: integer value.

Provides the current editor cursor line number.

\paragraph{Example macro.}
\begin{verbatim}
$MACRO REFERENCE_VARIABLE_C_LINE FROM EDIT;
    DEF_INT(Value);
    Value := C_LINE;
END_MACRO;
\end{verbatim}

\phantomsection
\label{prim:variable-c-row}
\subsection*{C\_ROW --- variable}
\index{C\_ROW!variable@\textsc{variable}}
\begin{verbatim}
C_ROW
\end{verbatim}
Parameters: None. Result: integer value.

Provides the current editor row within the visible view.

\paragraph{Example macro.}
\begin{verbatim}
$MACRO REFERENCE_VARIABLE_C_ROW FROM EDIT;
    DEF_INT(Value);
    Value := C_ROW;
END_MACRO;
\end{verbatim}

\phantomsection
\label{prim:variable-c-page}
\subsection*{C\_PAGE --- variable}
\index{C\_PAGE!variable@\textsc{variable}}
\begin{verbatim}
C_PAGE
\end{verbatim}
Parameters: None. Result: integer value.

Provides the current editor page number.

\paragraph{Example macro.}
\begin{verbatim}
$MACRO REFERENCE_VARIABLE_C_PAGE FROM EDIT;
    DEF_INT(Value);
    Value := C_PAGE;
END_MACRO;
\end{verbatim}

\phantomsection
\label{prim:variable-pg-line}
\subsection*{PG\_LINE --- variable}
\index{PG\_LINE!variable@\textsc{variable}}
\begin{verbatim}
PG_LINE
\end{verbatim}
Parameters: None. Result: integer value.

Provides the current line number within the current page.

\paragraph{Example macro.}
\begin{verbatim}
$MACRO REFERENCE_VARIABLE_PG_LINE FROM EDIT;
    DEF_INT(Value);
    Value := PG_LINE;
END_MACRO;
\end{verbatim}

\phantomsection
\label{prim:variable-at-eof}
\subsection*{AT\_EOF --- variable}
\index{AT\_EOF!variable@\textsc{variable}}
\begin{verbatim}
AT_EOF
\end{verbatim}
Parameters: None. Result: integer value.

Returns 1 when the cursor is at end of file; otherwise 0.

\paragraph{Example macro.}
\begin{verbatim}
$MACRO REFERENCE_VARIABLE_AT_EOF FROM EDIT;
    DEF_INT(Value);
    Value := AT_EOF;
END_MACRO;
\end{verbatim}

\phantomsection
\label{prim:variable-at-eol}
\subsection*{AT\_EOL --- variable}
\index{AT\_EOL!variable@\textsc{variable}}
\begin{verbatim}
AT_EOL
\end{verbatim}
Parameters: None. Result: integer value.

Returns 1 when the cursor is at end of line; otherwise 0.

\paragraph{Example macro.}
\begin{verbatim}
$MACRO REFERENCE_VARIABLE_AT_EOL FROM EDIT;
    DEF_INT(Value);
    Value := AT_EOL;
END_MACRO;
\end{verbatim}

\phantomsection
\label{prim:variable-block-stat}
\subsection*{BLOCK\_STAT --- variable}
\index{BLOCK\_STAT!variable@\textsc{variable}}
\begin{verbatim}
BLOCK_STAT
\end{verbatim}
Parameters: None. Result: integer value.

Provides current selection state named block stat.

\paragraph{Example macro.}
\begin{verbatim}
$MACRO REFERENCE_VARIABLE_BLOCK_STAT FROM EDIT;
    DEF_INT(Value);
    Value := BLOCK_STAT;
END_MACRO;
\end{verbatim}

\phantomsection
\label{prim:variable-block-line1}
\subsection*{BLOCK\_LINE1 --- variable}
\index{BLOCK\_LINE1!variable@\textsc{variable}}
\begin{verbatim}
BLOCK_LINE1
\end{verbatim}
Parameters: None. Result: integer value.

Provides current selection state named block line1.

\paragraph{Example macro.}
\begin{verbatim}
$MACRO REFERENCE_VARIABLE_BLOCK_LINE1 FROM EDIT;
    DEF_INT(Value);
    Value := BLOCK_LINE1;
END_MACRO;
\end{verbatim}

\phantomsection
\label{prim:variable-block-line2}
\subsection*{BLOCK\_LINE2 --- variable}
\index{BLOCK\_LINE2!variable@\textsc{variable}}
\begin{verbatim}
BLOCK_LINE2
\end{verbatim}
Parameters: None. Result: integer value.

Provides current selection state named block line2.

\paragraph{Example macro.}
\begin{verbatim}
$MACRO REFERENCE_VARIABLE_BLOCK_LINE2 FROM EDIT;
    DEF_INT(Value);
    Value := BLOCK_LINE2;
END_MACRO;
\end{verbatim}

\phantomsection
\label{prim:variable-block-col1}
\subsection*{BLOCK\_COL1 --- variable}
\index{BLOCK\_COL1!variable@\textsc{variable}}
\begin{verbatim}
BLOCK_COL1
\end{verbatim}
Parameters: None. Result: integer value.

Provides current selection state named block col1.

\paragraph{Example macro.}
\begin{verbatim}
$MACRO REFERENCE_VARIABLE_BLOCK_COL1 FROM EDIT;
    DEF_INT(Value);
    Value := BLOCK_COL1;
END_MACRO;
\end{verbatim}

\phantomsection
\label{prim:variable-block-col2}
\subsection*{BLOCK\_COL2 --- variable}
\index{BLOCK\_COL2!variable@\textsc{variable}}
\begin{verbatim}
BLOCK_COL2
\end{verbatim}
Parameters: None. Result: integer value.

Provides current selection state named block col2.

\paragraph{Example macro.}
\begin{verbatim}
$MACRO REFERENCE_VARIABLE_BLOCK_COL2 FROM EDIT;
    DEF_INT(Value);
    Value := BLOCK_COL2;
END_MACRO;
\end{verbatim}

\phantomsection
\label{prim:variable-marking}
\subsection*{MARKING --- variable}
\index{MARKING!variable@\textsc{variable}}
\begin{verbatim}
MARKING
\end{verbatim}
Parameters: None. Result: integer value.

Provides current selection state named marking.

\paragraph{Example macro.}
\begin{verbatim}
$MACRO REFERENCE_VARIABLE_MARKING FROM EDIT;
    DEF_INT(Value);
    Value := MARKING;
END_MACRO;
\end{verbatim}

\phantomsection
\label{prim:variable-tab-expand}
\subsection*{TAB\_EXPAND --- variable}
\index{TAB\_EXPAND!variable@\textsc{variable}}
\begin{verbatim}
TAB_EXPAND
\end{verbatim}
Parameters: None. Result: integer value.

Provides the active editor-format value named tab expand; writable settings use the configured runtime boundary.

\paragraph{Example macro.}
\begin{verbatim}
$MACRO REFERENCE_VARIABLE_TAB_EXPAND FROM EDIT;
    DEF_INT(Value);
    Value := TAB_EXPAND;
END_MACRO;
\end{verbatim}

\phantomsection
\label{prim:variable-insert-mode}
\subsection*{INSERT\_MODE --- variable}
\index{INSERT\_MODE!variable@\textsc{variable}}
\begin{verbatim}
INSERT_MODE
\end{verbatim}
Parameters: None. Result: integer value.

Provides the active editor-format value named insert mode; writable settings use the configured runtime boundary.

\paragraph{Example macro.}
\begin{verbatim}
$MACRO REFERENCE_VARIABLE_INSERT_MODE FROM EDIT;
    DEF_INT(Value);
    Value := INSERT_MODE;
END_MACRO;
\end{verbatim}

\phantomsection
\label{prim:variable-indent-level}
\subsection*{INDENT\_LEVEL --- variable}
\index{INDENT\_LEVEL!variable@\textsc{variable}}
\begin{verbatim}
INDENT_LEVEL
\end{verbatim}
Parameters: None. Result: integer value.

Provides the active editor-format value named indent level; writable settings use the configured runtime boundary.

\paragraph{Example macro.}
\begin{verbatim}
$MACRO REFERENCE_VARIABLE_INDENT_LEVEL FROM EDIT;
    DEF_INT(Value);
    Value := INDENT_LEVEL;
END_MACRO;
\end{verbatim}

\phantomsection
\label{prim:variable-cur-window}
\subsection*{CUR\_WINDOW --- variable}
\index{CUR\_WINDOW!variable@\textsc{variable}}
\begin{verbatim}
CUR_WINDOW
\end{verbatim}
Parameters: None. Result: integer value.

Provides current editor-window state named cur window.

\paragraph{Example macro.}
\begin{verbatim}
$MACRO REFERENCE_VARIABLE_CUR_WINDOW FROM EDIT;
    DEF_INT(Value);
    Value := CUR_WINDOW;
END_MACRO;
\end{verbatim}

\phantomsection
\label{prim:variable-link-stat}
\subsection*{LINK\_STAT --- variable}
\index{LINK\_STAT!variable@\textsc{variable}}
\begin{verbatim}
LINK_STAT
\end{verbatim}
Parameters: None. Result: integer value.

Provides current editor-window state named link stat.

\paragraph{Example macro.}
\begin{verbatim}
$MACRO REFERENCE_VARIABLE_LINK_STAT FROM EDIT;
    DEF_INT(Value);
    Value := LINK_STAT;
END_MACRO;
\end{verbatim}

\phantomsection
\label{prim:variable-win-x1}
\subsection*{WIN\_X1 --- variable}
\index{WIN\_X1!variable@\textsc{variable}}
\begin{verbatim}
WIN_X1
\end{verbatim}
Parameters: None. Result: integer value.

Provides current editor-window state named win x1.

\paragraph{Example macro.}
\begin{verbatim}
$MACRO REFERENCE_VARIABLE_WIN_X1 FROM EDIT;
    DEF_INT(Value);
    Value := WIN_X1;
END_MACRO;
\end{verbatim}

\phantomsection
\label{prim:variable-win-y1}
\subsection*{WIN\_Y1 --- variable}
\index{WIN\_Y1!variable@\textsc{variable}}
\begin{verbatim}
WIN_Y1
\end{verbatim}
Parameters: None. Result: integer value.

Provides current editor-window state named win y1.

\paragraph{Example macro.}
\begin{verbatim}
$MACRO REFERENCE_VARIABLE_WIN_Y1 FROM EDIT;
    DEF_INT(Value);
    Value := WIN_Y1;
END_MACRO;
\end{verbatim}

\phantomsection
\label{prim:variable-win-x2}
\subsection*{WIN\_X2 --- variable}
\index{WIN\_X2!variable@\textsc{variable}}
\begin{verbatim}
WIN_X2
\end{verbatim}
Parameters: None. Result: integer value.

Provides current editor-window state named win x2.

\paragraph{Example macro.}
\begin{verbatim}
$MACRO REFERENCE_VARIABLE_WIN_X2 FROM EDIT;
    DEF_INT(Value);
    Value := WIN_X2;
END_MACRO;
\end{verbatim}

\phantomsection
\label{prim:variable-win-y2}
\subsection*{WIN\_Y2 --- variable}
\index{WIN\_Y2!variable@\textsc{variable}}
\begin{verbatim}
WIN_Y2
\end{verbatim}
Parameters: None. Result: integer value.

Provides current editor-window state named win y2.

\paragraph{Example macro.}
\begin{verbatim}
$MACRO REFERENCE_VARIABLE_WIN_Y2 FROM EDIT;
    DEF_INT(Value);
    Value := WIN_Y2;
END_MACRO;
\end{verbatim}

\phantomsection
\label{prim:variable-window-count}
\subsection*{WINDOW\_COUNT --- variable}
\index{WINDOW\_COUNT!variable@\textsc{variable}}
\begin{verbatim}
WINDOW_COUNT
\end{verbatim}
Parameters: None. Result: integer value.

Provides current editor-window state named window count.

\paragraph{Example macro.}
\begin{verbatim}
$MACRO REFERENCE_VARIABLE_WINDOW_COUNT FROM EDIT;
    DEF_INT(Value);
    Value := WINDOW_COUNT;
END_MACRO;
\end{verbatim}

\phantomsection
\label{prim:variable-virtual-desktops}
\subsection*{VIRTUAL\_DESKTOPS --- variable}
\index{VIRTUAL\_DESKTOPS!variable@\textsc{variable}}
\begin{verbatim}
VIRTUAL_DESKTOPS
\end{verbatim}
Parameters: None. Result: integer value.

Provides the current runtime value named virtual desktops.

\paragraph{Example macro.}
\begin{verbatim}
$MACRO REFERENCE_VARIABLE_VIRTUAL_DESKTOPS FROM EDIT;
    DEF_INT(Value);
    Value := VIRTUAL_DESKTOPS;
END_MACRO;
\end{verbatim}

\phantomsection
\label{prim:variable-cyclic-virtual-desktops}
\subsection*{CYCLIC\_VIRTUAL\_DESKTOPS --- variable}
\index{CYCLIC\_VIRTUAL\_DESKTOPS!variable@\textsc{variable}}
\begin{verbatim}
CYCLIC_VIRTUAL_DESKTOPS
\end{verbatim}
Parameters: None. Result: integer value.

Provides the current runtime value named cyclic virtual desktops.

\paragraph{Example macro.}
\begin{verbatim}
$MACRO REFERENCE_VARIABLE_CYCLIC_VIRTUAL_DESKTOPS FROM EDIT;
    DEF_INT(Value);
    Value := CYCLIC_VIRTUAL_DESKTOPS;
END_MACRO;
\end{verbatim}

\phantomsection
\label{prim:variable-key1}
\subsection*{KEY1 --- variable}
\index{KEY1!variable@\textsc{variable}}
\begin{verbatim}
KEY1
\end{verbatim}
Parameters: None. Result: integer value.

Provides the current runtime value named key1.

\paragraph{Example macro.}
\begin{verbatim}
$MACRO REFERENCE_VARIABLE_KEY1 FROM EDIT;
    DEF_INT(Value);
    Value := KEY1;
END_MACRO;
\end{verbatim}

\phantomsection
\label{prim:variable-key2}
\subsection*{KEY2 --- variable}
\index{KEY2!variable@\textsc{variable}}
\begin{verbatim}
KEY2
\end{verbatim}
Parameters: None. Result: integer value.

Provides the current runtime value named key2.

\paragraph{Example macro.}
\begin{verbatim}
$MACRO REFERENCE_VARIABLE_KEY2 FROM EDIT;
    DEF_INT(Value);
    Value := KEY2;
END_MACRO;
\end{verbatim}

\phantomsection
\label{prim:variable-return-str}
\subsection*{RETURN\_STR --- variable}
\index{RETURN\_STR!variable@\textsc{variable}}
\begin{verbatim}
RETURN_STR
\end{verbatim}
Parameters: None. Result: string value.

Provides the string return slot of the current macro. Assign a value before RET to return a string result.

\paragraph{Example macro.}
\begin{verbatim}
$MACRO REFERENCE_VARIABLE_RETURN_STR FROM EDIT;
    DEF_STR(Value);
    Value := RETURN_STR;
END_MACRO;
\end{verbatim}

\phantomsection
\label{prim:variable-mparm-str}
\subsection*{MPARM\_STR --- variable}
\index{MPARM\_STR!variable@\textsc{variable}}
\begin{verbatim}
MPARM_STR
\end{verbatim}
Parameters: None. Result: string value.

Provides the current macro parameter string. The VM permits assignment for the active execution.

\paragraph{Example macro.}
\begin{verbatim}
$MACRO REFERENCE_VARIABLE_MPARM_STR FROM EDIT;
    DEF_STR(Value);
    Value := MPARM_STR;
END_MACRO;
\end{verbatim}

\phantomsection
\label{prim:variable-date}
\subsection*{DATE --- variable}
\index{DATE!variable@\textsc{variable}}
\begin{verbatim}
DATE
\end{verbatim}
Parameters: None. Result: string value.

Provides the current local date formatted by the runtime.

\paragraph{Example macro.}
\begin{verbatim}
$MACRO REFERENCE_VARIABLE_DATE FROM EDIT;
    DEF_STR(Value);
    Value := DATE;
END_MACRO;
\end{verbatim}

\phantomsection
\label{prim:variable-time}
\subsection*{TIME --- variable}
\index{TIME!variable@\textsc{variable}}
\begin{verbatim}
TIME
\end{verbatim}
Parameters: None. Result: string value.

Provides the current local time formatted by the runtime.

\paragraph{Example macro.}
\begin{verbatim}
$MACRO REFERENCE_VARIABLE_TIME FROM EDIT;
    DEF_STR(Value);
    Value := TIME;
END_MACRO;
\end{verbatim}

\phantomsection
\label{prim:variable-first-macro}
\subsection*{FIRST\_MACRO --- variable}
\index{FIRST\_MACRO!variable@\textsc{variable}}
\begin{verbatim}
FIRST_MACRO
\end{verbatim}
Parameters: None. Result: string value.

Starts loaded-macro enumeration and returns the first visible macro name, or an empty string.

\paragraph{Example macro.}
\begin{verbatim}
$MACRO REFERENCE_VARIABLE_FIRST_MACRO FROM EDIT;
    DEF_STR(Value);
    Value := FIRST_MACRO;
END_MACRO;
\end{verbatim}

\phantomsection
\label{prim:variable-next-macro}
\subsection*{NEXT\_MACRO --- variable}
\index{NEXT\_MACRO!variable@\textsc{variable}}
\begin{verbatim}
NEXT_MACRO
\end{verbatim}
Parameters: None. Result: string value.

Continues loaded-macro enumeration and returns the next visible macro name, or an empty string.

\paragraph{Example macro.}
\begin{verbatim}
$MACRO REFERENCE_VARIABLE_NEXT_MACRO FROM EDIT;
    DEF_STR(Value);
    Value := NEXT_MACRO;
END_MACRO;
\end{verbatim}

\phantomsection
\label{prim:variable-last-file-name}
\subsection*{LAST\_FILE\_NAME --- variable}
\index{LAST\_FILE\_NAME!variable@\textsc{variable}}
\begin{verbatim}
LAST_FILE_NAME
\end{verbatim}
Parameters: None. Result: string value.

Provides current or last-file metadata named last file name.

\paragraph{Example macro.}
\begin{verbatim}
$MACRO REFERENCE_VARIABLE_LAST_FILE_NAME FROM EDIT;
    DEF_STR(Value);
    Value := LAST_FILE_NAME;
END_MACRO;
\end{verbatim}

\phantomsection
\label{prim:variable-tmp-file-name}
\subsection*{TMP\_FILE\_NAME --- variable}
\index{TMP\_FILE\_NAME!variable@\textsc{variable}}
\begin{verbatim}
TMP_FILE_NAME
\end{verbatim}
Parameters: None. Result: string value.

Provides current or last-file metadata named tmp file name.

\paragraph{Example macro.}
\begin{verbatim}
$MACRO REFERENCE_VARIABLE_TMP_FILE_NAME FROM EDIT;
    DEF_STR(Value);
    Value := TMP_FILE_NAME;
END_MACRO;
\end{verbatim}

\phantomsection
\label{prim:variable-file-name}
\subsection*{FILE\_NAME --- variable}
\index{FILE\_NAME!variable@\textsc{variable}}
\begin{verbatim}
FILE_NAME
\end{verbatim}
Parameters: None. Result: string value.

Provides the current file name. The VM permits assignment when an active file context exists.

\paragraph{Example macro.}
\begin{verbatim}
$MACRO REFERENCE_VARIABLE_FILE_NAME FROM EDIT;
    DEF_STR(Value);
    Value := FILE_NAME;
END_MACRO;
\end{verbatim}

\phantomsection
\label{prim:variable-comspec}
\subsection*{COMSPEC --- variable}
\index{COMSPEC!variable@\textsc{variable}}
\begin{verbatim}
COMSPEC
\end{verbatim}
Parameters: None. Result: string value.

Provides the runtime environment value named comspec.

\paragraph{Example macro.}
\begin{verbatim}
$MACRO REFERENCE_VARIABLE_COMSPEC FROM EDIT;
    DEF_STR(Value);
    Value := COMSPEC;
END_MACRO;
\end{verbatim}

\phantomsection
\label{prim:variable-temp-path}
\subsection*{TEMP\_PATH --- variable}
\index{TEMP\_PATH!variable@\textsc{variable}}
\begin{verbatim}
TEMP_PATH
\end{verbatim}
Parameters: None. Result: string value.

Provides the runtime environment value named temp path.

\paragraph{Example macro.}
\begin{verbatim}
$MACRO REFERENCE_VARIABLE_TEMP_PATH FROM EDIT;
    DEF_STR(Value);
    Value := TEMP_PATH;
END_MACRO;
\end{verbatim}

\phantomsection
\label{prim:variable-found-str}
\subsection*{FOUND\_STR --- variable}
\index{FOUND\_STR!variable@\textsc{variable}}
\begin{verbatim}
FOUND_STR
\end{verbatim}
Parameters: None. Result: string value.

Provides the text of the current search match, or an empty string when no match exists.

\paragraph{Example macro.}
\begin{verbatim}
$MACRO REFERENCE_VARIABLE_FOUND_STR FROM EDIT;
    DEF_STR(Value);
    Value := FOUND_STR;
END_MACRO;
\end{verbatim}

\phantomsection
\label{prim:variable-search-file}
\subsection*{SEARCH\_FILE --- variable}
\index{SEARCH\_FILE!variable@\textsc{variable}}
\begin{verbatim}
SEARCH_FILE
\end{verbatim}
Parameters: None. Result: string value.

Provides the file name of the current search match, or an empty string when no match exists.

\paragraph{Example macro.}
\begin{verbatim}
$MACRO REFERENCE_VARIABLE_SEARCH_FILE FROM EDIT;
    DEF_STR(Value);
    Value := SEARCH_FILE;
END_MACRO;
\end{verbatim}

\phantomsection
\label{prim:variable-mr-path}
\subsection*{MR\_PATH --- variable}
\index{MR\_PATH!variable@\textsc{variable}}
\begin{verbatim}
MR_PATH
\end{verbatim}
Parameters: None. Result: string value.

Provides the runtime environment value named mr path.

\paragraph{Example macro.}
\begin{verbatim}
$MACRO REFERENCE_VARIABLE_MR_PATH FROM EDIT;
    DEF_STR(Value);
    Value := MR_PATH;
END_MACRO;
\end{verbatim}

\phantomsection
\label{prim:variable-os-version}
\subsection*{OS\_VERSION --- variable}
\index{OS\_VERSION!variable@\textsc{variable}}
\begin{verbatim}
OS_VERSION
\end{verbatim}
Parameters: None. Result: string value.

Provides the runtime environment value named os version.

\paragraph{Example macro.}
\begin{verbatim}
$MACRO REFERENCE_VARIABLE_OS_VERSION FROM EDIT;
    DEF_STR(Value);
    Value := OS_VERSION;
END_MACRO;
\end{verbatim}

\phantomsection
\label{prim:variable-get-line}
\subsection*{GET\_LINE --- variable}
\index{GET\_LINE!variable@\textsc{variable}}
\begin{verbatim}
GET_LINE
\end{verbatim}
Parameters: None. Result: string value.

Provides the complete current editor line as a string.

\paragraph{Example macro.}
\begin{verbatim}
$MACRO REFERENCE_VARIABLE_GET_LINE FROM EDIT;
    DEF_STR(Value);
    Value := GET_LINE;
END_MACRO;
\end{verbatim}

\phantomsection
\label{prim:variable-format-line}
\subsection*{FORMAT\_LINE --- variable}
\index{FORMAT\_LINE!variable@\textsc{variable}}
\begin{verbatim}
FORMAT_LINE
\end{verbatim}
Parameters: None. Result: string value.

Provides the active editor-format value named format line; writable settings use the configured runtime boundary.

\paragraph{Example macro.}
\begin{verbatim}
$MACRO REFERENCE_VARIABLE_FORMAT_LINE FROM EDIT;
    DEF_STR(Value);
    Value := FORMAT_LINE;
END_MACRO;
\end{verbatim}

\phantomsection
\label{prim:variable-default-format}
\subsection*{DEFAULT\_FORMAT --- variable}
\index{DEFAULT\_FORMAT!variable@\textsc{variable}}
\begin{verbatim}
DEFAULT_FORMAT
\end{verbatim}
Parameters: None. Result: string value.

Provides the active editor-format value named default format; writable settings use the configured runtime boundary.

\paragraph{Example macro.}
\begin{verbatim}
$MACRO REFERENCE_VARIABLE_DEFAULT_FORMAT FROM EDIT;
    DEF_STR(Value);
    Value := DEFAULT_FORMAT;
END_MACRO;
\end{verbatim}

\phantomsection
\label{prim:variable-page-str}
\subsection*{PAGE\_STR --- variable}
\index{PAGE\_STR!variable@\textsc{variable}}
\begin{verbatim}
PAGE_STR
\end{verbatim}
Parameters: None. Result: string value.

Provides the active editor-format value named page str; writable settings use the configured runtime boundary.

\paragraph{Example macro.}
\begin{verbatim}
$MACRO REFERENCE_VARIABLE_PAGE_STR FROM EDIT;
    DEF_STR(Value);
    Value := PAGE_STR;
END_MACRO;
\end{verbatim}

\phantomsection
\label{prim:variable-word-delimits}
\subsection*{WORD\_DELIMITS --- variable}
\index{WORD\_DELIMITS!variable@\textsc{variable}}
\begin{verbatim}
WORD_DELIMITS
\end{verbatim}
Parameters: None. Result: string value.

Provides the active editor-format value named word delimits; writable settings use the configured runtime boundary.

\paragraph{Example macro.}
\begin{verbatim}
$MACRO REFERENCE_VARIABLE_WORD_DELIMITS FROM EDIT;
    DEF_STR(Value);
    Value := WORD_DELIMITS;
END_MACRO;
\end{verbatim}

\phantomsection
\label{prim:variable-found-x}
\subsection*{FOUND\_X --- variable}
\index{FOUND\_X!variable@\textsc{variable}}
\begin{verbatim}
FOUND_X
\end{verbatim}
Parameters: None. Result: integer value.

Provides the column of the current search match, or 0 when no match exists.

\paragraph{Example macro.}
\begin{verbatim}
$MACRO REFERENCE_VARIABLE_FOUND_X FROM EDIT;
    DEF_INT(Value);
    Value := FOUND_X;
END_MACRO;
\end{verbatim}

\phantomsection
\label{prim:variable-found-y}
\subsection*{FOUND\_Y --- variable}
\index{FOUND\_Y!variable@\textsc{variable}}
\begin{verbatim}
FOUND_Y
\end{verbatim}
Parameters: None. Result: integer value.

Provides the line of the current search match, or 0 when no match exists.

\paragraph{Example macro.}
\begin{verbatim}
$MACRO REFERENCE_VARIABLE_FOUND_Y FROM EDIT;
    DEF_INT(Value);
    Value := FOUND_Y;
END_MACRO;
\end{verbatim}

\phantomsection
\label{prim:variable-cur-char}
\subsection*{CUR\_CHAR --- variable}
\index{CUR\_CHAR!variable@\textsc{variable}}
\begin{verbatim}
CUR_CHAR
\end{verbatim}
Parameters: None. Result: character value.

Provides the character at the current editor cursor position.

\paragraph{Example macro.}
\begin{verbatim}
$MACRO REFERENCE_VARIABLE_CUR_CHAR FROM EDIT;
    DEF_CHAR(Value);
    Value := CUR_CHAR;
END_MACRO;
\end{verbatim}


\chapter{Expressions and Conversion}

An expression computes a value from literals, variables, collection access, operators and intrinsics. Expressions are used in assignments, conditions and procedure arguments. The following entries document the intrinsic operations that create, inspect, convert and search values.

\section{Primitive Map}
The map is a local index for this chapter. The final column is a generated page
reference and hyperlink to the full entry, where the source form, parameters,
semantics and a complete example macro appear.

\begin{longtable}{@{}L{0.38\textwidth}L{0.14\textwidth}L{0.27\textwidth}L{0.10\textwidth}@{}}
\toprule
Primitive & Role & Purpose & Reference \\
\midrule
\endhead
\toprule
Primitive & Role & Purpose & Reference \\
\midrule
\endhead
\primitiveoverview{RVAL}{intrinsic-rval}{intrinsic}{Expression operation.}
\primitiveoverview{VAL}{intrinsic-val}{intrinsic}{Expression operation.}
\primitiveoverview{ASCII}{intrinsic-ascii}{intrinsic}{Expression operation.}
\primitiveoverview{BAR\_MENU}{intrinsic-bar-menu}{intrinsic}{Expression operation.}
\primitiveoverview{CAPS}{intrinsic-caps}{intrinsic}{Expression operation.}
\primitiveoverview{CHAR}{intrinsic-char}{intrinsic}{Expression operation.}
\primitiveoverview{CHECK\_KEY}{intrinsic-check-key}{intrinsic}{Expression operation.}
\primitiveoverview{COPY}{intrinsic-copy}{intrinsic}{Expression operation.}
\primitiveoverview{EXISTS}{intrinsic-exists}{intrinsic}{Expression operation.}
\primitiveoverview{HAS\_VALUE}{intrinsic-has-value}{intrinsic}{Expression operation.}
\primitiveoverview{INT\_R}{intrinsic-int-r}{intrinsic}{Expression operation.}
\primitiveoverview{KEYS}{intrinsic-keys}{intrinsic}{Expression operation.}
\primitiveoverview{LEN}{intrinsic-len}{intrinsic}{Expression operation.}
\primitiveoverview{LENGTH}{intrinsic-length}{intrinsic}{Expression operation.}
\primitiveoverview{PARSE\_INT}{intrinsic-parse-int}{intrinsic}{Expression operation.}
\primitiveoverview{PARSE\_STR}{intrinsic-parse-str}{intrinsic}{Expression operation.}
\primitiveoverview{POS}{intrinsic-pos}{intrinsic}{Expression operation.}
\primitiveoverview{REAL\_I}{intrinsic-real-i}{intrinsic}{Expression operation.}
\primitiveoverview{REMOVE\_SPACE}{intrinsic-remove-space}{intrinsic}{Expression operation.}
\primitiveoverview{RSTR}{intrinsic-rstr}{intrinsic}{Expression operation.}
\primitiveoverview{SCREEN\_LENGTH}{intrinsic-screen-length}{intrinsic}{Expression operation.}
\primitiveoverview{SCREEN\_WIDTH}{intrinsic-screen-width}{intrinsic}{Expression operation.}
\primitiveoverview{STR}{intrinsic-str}{intrinsic}{Expression operation.}
\primitiveoverview{STRING\_IN}{intrinsic-string-in}{intrinsic}{Expression operation.}
\primitiveoverview{STR\_DEL}{intrinsic-str-del}{intrinsic}{Expression operation.}
\primitiveoverview{STR\_INS}{intrinsic-str-ins}{intrinsic}{Expression operation.}
\primitiveoverview{UTF8}{intrinsic-utf8}{intrinsic}{Expression operation.}
\primitiveoverview{VALUES}{intrinsic-values}{intrinsic}{Expression operation.}
\primitiveoverview{V\_MENU}{intrinsic-v-menu}{intrinsic}{Expression operation.}
\primitiveoverview{WHEREX}{intrinsic-wherex}{intrinsic}{Expression operation.}
\primitiveoverview{WHEREY}{intrinsic-wherey}{intrinsic}{Expression operation.}
\primitiveoverview{XPOS}{intrinsic-xpos}{intrinsic}{Expression operation.}
\bottomrule
\end{longtable}

\section{Reference Entries}

\phantomsection
\label{prim:intrinsic-rval}
\subsection*{RVAL}
\index{RVAL@\texttt{RVAL}}
\begin{verbatim}
RVAL(target, text)
\end{verbatim}
Parameters:
\begin{description}
\item[target] Writable typed target variable.
\item[text] String-like text.
\end{description}

Parses text into the named real target and returns a status value.

Returns an integer.

\paragraph{Example macro.}
\begin{verbatim}
$MACRO REFERENCE_RVAL FROM EDIT;
    DEF_REAL(Target);
    DEF_INT(Result);
    Result := RVAL(Target, '42.0');
END_MACRO;
\end{verbatim}

\phantomsection
\label{prim:intrinsic-val}
\subsection*{VAL}
\index{VAL@\texttt{VAL}}
\begin{verbatim}
VAL(target, text)
\end{verbatim}
Parameters:
\begin{description}
\item[target] Writable typed target variable.
\item[text] String-like text.
\end{description}

Parses text into the named integer target and returns a status value.

Returns an integer.

\paragraph{Example macro.}
\begin{verbatim}
$MACRO REFERENCE_VAL FROM EDIT;
    DEF_INT(Target);
    DEF_INT(Result);
    Result := VAL(Target, '42');
END_MACRO;
\end{verbatim}

\phantomsection
\label{prim:intrinsic-ascii}
\subsection*{ASCII}
\index{ASCII@\texttt{ASCII}}
\begin{verbatim}
ASCII(text)
\end{verbatim}
Parameters:
\begin{description}
\item[text] String-like text.
\end{description}

Returns the numeric value of the first character.

Returns an integer.

\paragraph{Example macro.}
\begin{verbatim}
$MACRO REFERENCE_ASCII FROM EDIT;
    DEF_INT(Result);
    Result := ASCII('reference');
END_MACRO;
\end{verbatim}

\phantomsection
\label{prim:intrinsic-bar-menu}
\subsection*{BAR\_MENU}
\index{BAR\_MENU@\texttt{BAR\_MENU}}
\begin{verbatim}
BAR_MENU(column, row, width, items, title)
\end{verbatim}
Parameters:
\begin{description}
\item[column] Screen or editor column, as required by the primitive.
\item[row] Zero-based screen row.
\item[width] Control or region width.
\item[items] String-like encoded menu item list.
\item[title] User-visible title.
\end{description}

Runs a horizontal menu and returns the selected item index.

Returns an integer.

\paragraph{Example macro.}
\begin{verbatim}
$MACRO REFERENCE_BAR_MENU FROM EDIT;
    DEF_INT(Result);
    Result := BAR_MENU(0, 0, 0, 'reference', 'reference');
END_MACRO;
\end{verbatim}

\phantomsection
\label{prim:intrinsic-caps}
\subsection*{CAPS}
\index{CAPS@\texttt{CAPS}}
\begin{verbatim}
CAPS(text)
\end{verbatim}
Parameters:
\begin{description}
\item[text] String-like text.
\end{description}

Returns an upper-case copy of the text.

Returns a string.

\paragraph{Example macro.}
\begin{verbatim}
$MACRO REFERENCE_CAPS FROM EDIT;
    DEF_STR(Result);
    Result := CAPS('reference');
END_MACRO;
\end{verbatim}

\phantomsection
\label{prim:intrinsic-char}
\subsection*{CHAR}
\index{CHAR@\texttt{CHAR}}
\begin{verbatim}
CHAR(integerValue)
\end{verbatim}
Parameters:
\begin{description}
\item[integerValue] Integer source value for the conversion.
\end{description}

Converts an integer to a character value.

Returns a character.

\paragraph{Example macro.}
\begin{verbatim}
$MACRO REFERENCE_CHAR FROM EDIT;
    DEF_STR(Result);
    Result := CHAR(0);
END_MACRO;
\end{verbatim}

\phantomsection
\label{prim:intrinsic-check-key}
\subsection*{CHECK\_KEY}
\index{CHECK\_KEY@\texttt{CHECK\_KEY}}
\begin{verbatim}
CHECK_KEY
\end{verbatim}
Parameters: None.

Reports whether input is available.

Returns an integer.

\paragraph{Example macro.}
\begin{verbatim}
$MACRO REFERENCE_CHECK_KEY FROM EDIT;
    DEF_INT(Result);
    Result := CHECK_KEY;
END_MACRO;
\end{verbatim}

\phantomsection
\label{prim:intrinsic-copy}
\subsection*{COPY}
\index{COPY@\texttt{COPY}}
\begin{verbatim}
COPY(text, start, count)
\end{verbatim}
Parameters:
\begin{description}
\item[text] String-like text.
\item[start] Zero-based string start position.
\item[count] Number of characters to process.
\end{description}

Returns a substring using the specified start and count.

Returns a string.

\paragraph{Example macro.}
\begin{verbatim}
$MACRO REFERENCE_COPY FROM EDIT;
    DEF_STR(Result);
    Result := COPY('reference', 0, 0);
END_MACRO;
\end{verbatim}

\begin{mrmacextension}
\phantomsection
\label{prim:intrinsic-exists}
\subsection*{EXISTS}
\index{EXISTS@\texttt{EXISTS}}
\begin{verbatim}
EXISTS(hash, key)
\end{verbatim}
Parameters:
\begin{description}
\item[hash] Hash value.
\item[key] String-like hash key.
\end{description}

Tests whether a hash contains a key.

Returns an integer.

\paragraph{Example macro.}
\begin{verbatim}
$MACRO REFERENCE_EXISTS FROM EDIT;
    DEF_HASH(State);
    DEF_INT(Result);
    Result := EXISTS(State, 'reference');
END_MACRO;
\end{verbatim}
\end{mrmacextension}

\begin{mrmacextension}
\phantomsection
\label{prim:intrinsic-has-value}
\subsection*{HAS\_VALUE}
\index{HAS\_VALUE@\texttt{HAS\_VALUE}}
\begin{verbatim}
HAS_VALUE(hash, key)
\end{verbatim}
Parameters:
\begin{description}
\item[hash] Hash value.
\item[key] String-like hash key.
\end{description}

Tests whether a hash key has a non-default value.

Returns an integer.

\paragraph{Example macro.}
\begin{verbatim}
$MACRO REFERENCE_HAS_VALUE FROM EDIT;
    DEF_HASH(State);
    DEF_INT(Result);
    Result := HAS_VALUE(State, 'reference');
END_MACRO;
\end{verbatim}
\end{mrmacextension}

\phantomsection
\label{prim:intrinsic-int-r}
\subsection*{INT\_R}
\index{INT\_R@\texttt{INT\_R}}
\begin{verbatim}
INT_R(realValue)
\end{verbatim}
Parameters:
\begin{description}
\item[realValue] Real source value for the conversion or formatting operation.
\end{description}

Converts a real value to an integer.

Returns an integer.

\paragraph{Example macro.}
\begin{verbatim}
$MACRO REFERENCE_INT_R FROM EDIT;
    DEF_INT(Result);
    Result := INT_R(0.0);
END_MACRO;
\end{verbatim}

\begin{mrmacextension}
\phantomsection
\label{prim:intrinsic-keys}
\subsection*{KEYS}
\index{KEYS@\texttt{KEYS}}
\begin{verbatim}
KEYS(hash)
\end{verbatim}
Parameters:
\begin{description}
\item[hash] Hash value.
\end{description}

Returns the hash keys as a string array.

Returns a string array.

\paragraph{Example macro.}
\begin{verbatim}
$MACRO REFERENCE_KEYS FROM EDIT;
    DEF_HASH(State);
    DEF_INT(Result);
    Result := LEN(KEYS(State));
END_MACRO;
\end{verbatim}
\end{mrmacextension}

\begin{mrmacextension}
\phantomsection
\label{prim:intrinsic-len}
\subsection*{LEN}
\index{LEN@\texttt{LEN}}
\begin{verbatim}
LEN(value)
\end{verbatim}
Parameters:
\begin{description}
\item[value] String-like value or typed array whose length is requested.
\end{description}

Returns the number of characters or collection elements.

Returns an integer.

\paragraph{Example macro.}
\begin{verbatim}
$MACRO REFERENCE_LEN FROM EDIT;
    DEF_STR(Text);
    DEF_INT(Result);
    Text := 'reference';
    Result := LEN(Text);
END_MACRO;
\end{verbatim}
\end{mrmacextension}

\phantomsection
\label{prim:intrinsic-length}
\subsection*{LENGTH}
\index{LENGTH@\texttt{LENGTH}}
\begin{verbatim}
LENGTH(text)
\end{verbatim}
Parameters:
\begin{description}
\item[text] String-like text.
\end{description}

Returns the number of characters in the string-like value.

Returns an integer.

\paragraph{Example macro.}
\begin{verbatim}
$MACRO REFERENCE_LENGTH FROM EDIT;
    DEF_INT(Result);
    Result := LENGTH('reference');
END_MACRO;
\end{verbatim}

\phantomsection
\label{prim:intrinsic-parse-int}
\subsection*{PARSE\_INT}
\index{PARSE\_INT@\texttt{PARSE\_INT}}
\begin{verbatim}
PARSE_INT(text, delimiter)
\end{verbatim}
Parameters:
\begin{description}
\item[text] String-like text.
\item[delimiter] String-like delimiter separating parsed fields.
\end{description}

Returns the selected integer component from delimiter parsing.

Returns an integer.

\paragraph{Example macro.}
\begin{verbatim}
$MACRO REFERENCE_PARSE_INT FROM EDIT;
    DEF_INT(Result);
    Result := PARSE_INT('reference', 'reference');
END_MACRO;
\end{verbatim}

\phantomsection
\label{prim:intrinsic-parse-str}
\subsection*{PARSE\_STR}
\index{PARSE\_STR@\texttt{PARSE\_STR}}
\begin{verbatim}
PARSE_STR(text, delimiter)
\end{verbatim}
Parameters:
\begin{description}
\item[text] String-like text.
\item[delimiter] String-like delimiter separating parsed fields.
\end{description}

Returns the selected string component from delimiter parsing.

Returns a string.

\paragraph{Example macro.}
\begin{verbatim}
$MACRO REFERENCE_PARSE_STR FROM EDIT;
    DEF_STR(Result);
    Result := PARSE_STR('reference', 'reference');
END_MACRO;
\end{verbatim}

\phantomsection
\label{prim:intrinsic-pos}
\subsection*{POS}
\index{POS@\texttt{POS}}
\begin{verbatim}
POS(needle, text)
\end{verbatim}
Parameters:
\begin{description}
\item[needle] String-like text sought in the supplied text.
\item[text] String-like text.
\end{description}

Returns the first position of the needle in the text, or zero when absent.

Returns an integer.

\paragraph{Example macro.}
\begin{verbatim}
$MACRO REFERENCE_POS FROM EDIT;
    DEF_INT(Result);
    Result := POS('reference', 'reference');
END_MACRO;
\end{verbatim}

\phantomsection
\label{prim:intrinsic-real-i}
\subsection*{REAL\_I}
\index{REAL\_I@\texttt{REAL\_I}}
\begin{verbatim}
REAL_I(integerValue)
\end{verbatim}
Parameters:
\begin{description}
\item[integerValue] Integer source value for the conversion.
\end{description}

Converts an integer to a real value.

Returns a real value.

\paragraph{Example macro.}
\begin{verbatim}
$MACRO REFERENCE_REAL_I FROM EDIT;
    DEF_REAL(Result);
    Result := REAL_I(0);
END_MACRO;
\end{verbatim}

\phantomsection
\label{prim:intrinsic-remove-space}
\subsection*{REMOVE\_SPACE}
\index{REMOVE\_SPACE@\texttt{REMOVE\_SPACE}}
\begin{verbatim}
REMOVE_SPACE(text)
\end{verbatim}
Parameters:
\begin{description}
\item[text] String-like text.
\end{description}

Normalises spacing in a string.

Returns a string.

\paragraph{Example macro.}
\begin{verbatim}
$MACRO REFERENCE_REMOVE_SPACE FROM EDIT;
    DEF_STR(Result);
    Result := REMOVE_SPACE('reference');
END_MACRO;
\end{verbatim}

\phantomsection
\label{prim:intrinsic-rstr}
\subsection*{RSTR}
\index{RSTR@\texttt{RSTR}}
\begin{verbatim}
RSTR(realValue, significantDigits, decimalPlaces)
\end{verbatim}
Parameters:
\begin{description}
\item[realValue] Real source value for the conversion or formatting operation.
\item[significantDigits] Integer number of significant digits to emit.
\item[decimalPlaces] Integer number of decimal places to emit.
\end{description}

Formats a real value with the requested precision.

Returns a string.

\paragraph{Example macro.}
\begin{verbatim}
$MACRO REFERENCE_RSTR FROM EDIT;
    DEF_STR(Result);
    Result := RSTR(0.0, 0, 0);
END_MACRO;
\end{verbatim}

\phantomsection
\label{prim:intrinsic-screen-length}
\subsection*{SCREEN\_LENGTH}
\index{SCREEN\_LENGTH@\texttt{SCREEN\_LENGTH}}
\begin{verbatim}
SCREEN_LENGTH
\end{verbatim}
Parameters: None.

Returns the current screen row count.

Returns an integer.

\paragraph{Example macro.}
\begin{verbatim}
$MACRO REFERENCE_SCREEN_LENGTH FROM EDIT;
    DEF_INT(Result);
    Result := SCREEN_LENGTH;
END_MACRO;
\end{verbatim}

\phantomsection
\label{prim:intrinsic-screen-width}
\subsection*{SCREEN\_WIDTH}
\index{SCREEN\_WIDTH@\texttt{SCREEN\_WIDTH}}
\begin{verbatim}
SCREEN_WIDTH
\end{verbatim}
Parameters: None.

Returns the current screen column count.

Returns an integer.

\paragraph{Example macro.}
\begin{verbatim}
$MACRO REFERENCE_SCREEN_WIDTH FROM EDIT;
    DEF_INT(Result);
    Result := SCREEN_WIDTH;
END_MACRO;
\end{verbatim}

\phantomsection
\label{prim:intrinsic-str}
\subsection*{STR}
\index{STR@\texttt{STR}}
\begin{verbatim}
STR(integerValue)
\end{verbatim}
Parameters:
\begin{description}
\item[integerValue] Integer source value for the conversion.
\end{description}

Converts an integer to its decimal string representation.

Returns a string.

\paragraph{Example macro.}
\begin{verbatim}
$MACRO REFERENCE_STR FROM EDIT;
    DEF_STR(Result);
    Result := STR(0);
END_MACRO;
\end{verbatim}

\phantomsection
\label{prim:intrinsic-string-in}
\subsection*{STRING\_IN}
\index{STRING\_IN@\texttt{STRING\_IN}}
\begin{verbatim}
STRING_IN(column, row, width, initialText, prompt)
\end{verbatim}
Parameters:
\begin{description}
\item[column] Screen or editor column, as required by the primitive.
\item[row] Zero-based screen row.
\item[width] Control or region width.
\item[initialText] Initial text value.
\item[prompt] String or character positional argument.
\end{description}

Runs string input and returns the entered text.

Returns a string.

\paragraph{Example macro.}
\begin{verbatim}
$MACRO REFERENCE_STRING_IN FROM EDIT;
    DEF_STR(Result);
    Result := STRING_IN(0, 0, 0, 'reference', 'reference');
END_MACRO;
\end{verbatim}

\phantomsection
\label{prim:intrinsic-str-del}
\subsection*{STR\_DEL}
\index{STR\_DEL@\texttt{STR\_DEL}}
\begin{verbatim}
STR_DEL(text, start, count)
\end{verbatim}
Parameters:
\begin{description}
\item[text] String-like text.
\item[start] Zero-based string start position.
\item[count] Number of characters to process.
\end{description}

Returns text with count characters removed from start.

Returns a string.

\paragraph{Example macro.}
\begin{verbatim}
$MACRO REFERENCE_STR_DEL FROM EDIT;
    DEF_STR(Result);
    Result := STR_DEL('reference', 0, 0);
END_MACRO;
\end{verbatim}

\phantomsection
\label{prim:intrinsic-str-ins}
\subsection*{STR\_INS}
\index{STR\_INS@\texttt{STR\_INS}}
\begin{verbatim}
STR_INS(text, insertText, start)
\end{verbatim}
Parameters:
\begin{description}
\item[text] String-like text.
\item[insertText] String-like text inserted into the source text.
\item[start] Zero-based string start position.
\end{description}

Returns text with insertText inserted at start.

Returns a string.

\paragraph{Example macro.}
\begin{verbatim}
$MACRO REFERENCE_STR_INS FROM EDIT;
    DEF_STR(Result);
    Result := STR_INS('reference', 'reference', 0);
END_MACRO;
\end{verbatim}

\begin{mrmacextension}
\phantomsection
\label{prim:intrinsic-utf8}
\subsection*{UTF8}
\index{UTF8@\texttt{UTF8}}
\begin{verbatim}
UTF8(codePoint)
\end{verbatim}
Parameters:
\begin{description}
\item[codePoint] Integer Unicode code point.
\end{description}

Encodes one Unicode code point as a UTF-8 string.

Returns a string.

\paragraph{Example macro.}
\begin{verbatim}
$MACRO REFERENCE_UTF8 FROM EDIT;
    DEF_STR(Result);
    Result := UTF8(0);
END_MACRO;
\end{verbatim}
\end{mrmacextension}

\begin{mrmacextension}
\phantomsection
\label{prim:intrinsic-values}
\subsection*{VALUES}
\index{VALUES@\texttt{VALUES}}
\begin{verbatim}
VALUES(hash)
\end{verbatim}
Parameters:
\begin{description}
\item[hash] Hash value.
\end{description}

Returns the hash values as a string array.

Returns a string array.

\paragraph{Example macro.}
\begin{verbatim}
$MACRO REFERENCE_VALUES FROM EDIT;
    DEF_HASH(State);
    DEF_INT(Result);
    Result := LEN(VALUES(State));
END_MACRO;
\end{verbatim}
\end{mrmacextension}

\phantomsection
\label{prim:intrinsic-v-menu}
\subsection*{V\_MENU}
\index{V\_MENU@\texttt{V\_MENU}}
\begin{verbatim}
V_MENU(column, row, height, items, title)
\end{verbatim}
Parameters:
\begin{description}
\item[column] Screen or editor column, as required by the primitive.
\item[row] Zero-based screen row.
\item[height] Control or region height.
\item[items] String-like encoded menu item list.
\item[title] User-visible title.
\end{description}

Runs a vertical menu and returns the selected item index.

Returns an integer.

\paragraph{Example macro.}
\begin{verbatim}
$MACRO REFERENCE_V_MENU FROM EDIT;
    DEF_INT(Result);
    Result := V_MENU(0, 0, 0, 'reference', 'reference');
END_MACRO;
\end{verbatim}

\phantomsection
\label{prim:intrinsic-wherex}
\subsection*{WHEREX}
\index{WHEREX@\texttt{WHEREX}}
\begin{verbatim}
WHEREX
\end{verbatim}
Parameters: None.

Returns the current screen column.

Returns an integer.

\paragraph{Example macro.}
\begin{verbatim}
$MACRO REFERENCE_WHEREX FROM EDIT;
    DEF_INT(Result);
    Result := WHEREX;
END_MACRO;
\end{verbatim}

\phantomsection
\label{prim:intrinsic-wherey}
\subsection*{WHEREY}
\index{WHEREY@\texttt{WHEREY}}
\begin{verbatim}
WHEREY
\end{verbatim}
Parameters: None.

Returns the current screen row.

Returns an integer.

\paragraph{Example macro.}
\begin{verbatim}
$MACRO REFERENCE_WHEREY FROM EDIT;
    DEF_INT(Result);
    Result := WHEREY;
END_MACRO;
\end{verbatim}

\phantomsection
\label{prim:intrinsic-xpos}
\subsection*{XPOS}
\index{XPOS@\texttt{XPOS}}
\begin{verbatim}
XPOS(needle, text, start)
\end{verbatim}
Parameters:
\begin{description}
\item[needle] String-like text sought in the supplied text.
\item[text] String-like text.
\item[start] Zero-based string start position.
\end{description}

Searches from start and returns the matching position, or zero.

Returns an integer.

\paragraph{Example macro.}
\begin{verbatim}
$MACRO REFERENCE_XPOS FROM EDIT;
    DEF_INT(Result);
    Result := XPOS('reference', 'reference', 0);
END_MACRO;
\end{verbatim}


\chapter{Control Flow and Macro Composition}

Statements normally run in source order. Control-flow forms choose a branch, repeat a block, transfer to a label or return from a local call. Composition forms expose invocation parameters and run macros already loaded into the runtime catalogue.

\section{Primitive Map}
The map is a local index for this chapter. The final column is a generated page
reference and hyperlink to the full entry, where the source form, parameters,
semantics and a complete example macro appear.

\begin{longtable}{@{}L{0.38\textwidth}L{0.14\textwidth}L{0.27\textwidth}L{0.10\textwidth}@{}}
\toprule
Primitive & Role & Purpose & Reference \\
\midrule
\endhead
\toprule
Primitive & Role & Purpose & Reference \\
\midrule
\endhead
\primitiveoverview{IF}{keyword-if}{keyword}{Control flow.}
\primitiveoverview{WHILE}{keyword-while}{keyword}{Control flow.}
\primitiveoverview{GOTO}{keyword-goto}{keyword}{Control flow.}
\primitiveoverview{CALL}{keyword-call}{keyword}{Control flow.}
\primitiveoverview{RET}{keyword-ret}{keyword}{Control flow.}
\primitiveoverview{PARAM\_STR}{intrinsic-param-str}{intrinsic}{Control flow.}
\primitiveoverview{THEN}{keyword-then}{keyword}{True branch.}
\primitiveoverview{ELSE}{keyword-else}{keyword}{False branch.}
\primitiveoverview{END}{keyword-end}{keyword}{Block terminator.}
\primitiveoverview{DO}{keyword-do}{keyword}{Loop body.}
\primitiveoverview{label}{form-label}{source form}{Branch label.}
\bottomrule
\end{longtable}

\section{Reference Entries}

\phantomsection
\label{prim:keyword-if}
\subsection*{IF}
\index{IF@\texttt{IF}}
\begin{verbatim}
IF integerExpression THEN
    statement;
ELSE
    statement;
END;
\end{verbatim}
Parameters: integerExpression is non-zero for the true branch.

Selects one of two statement blocks.

\paragraph{Example macro.}
\begin{verbatim}
$MACRO REFERENCE_IF FROM EDIT;
    IF TRUE THEN
        MAKE_MESSAGE('true');
    END;
END_MACRO;
\end{verbatim}

\phantomsection
\label{prim:keyword-while}
\subsection*{WHILE}
\index{WHILE@\texttt{WHILE}}
\begin{verbatim}
WHILE integerExpression DO
    statement;
END;
\end{verbatim}
Parameters: integerExpression is evaluated before each iteration.

Repeats a statement block while its condition is non-zero.

\paragraph{Example macro.}
\begin{verbatim}
$MACRO REFERENCE_WHILE FROM EDIT;
    DEF_INT(Count);
    WHILE Count < 1 DO
        Count := Count + 1;
    END;
END_MACRO;
\end{verbatim}

\phantomsection
\label{prim:keyword-goto}
\subsection*{GOTO}
\index{GOTO@\texttt{GOTO}}
\begin{verbatim}
GOTO label;
\end{verbatim}
Parameters: label names a label in the same macro.

Branches to a local label.

\paragraph{Example macro.}
\begin{verbatim}
$MACRO REFERENCE_GOTO FROM EDIT;
    GOTO Done;
Done:
END_MACRO;
\end{verbatim}

\phantomsection
\label{prim:keyword-call}
\subsection*{CALL}
\index{CALL@\texttt{CALL}}
\begin{verbatim}
CALL label;
\end{verbatim}
Parameters: label names a label in the same macro.

Calls a local label and records a return point.

\paragraph{Example macro.}
\begin{verbatim}
$MACRO REFERENCE_CALL FROM EDIT;
    CALL Done;
Done:
    RET;
END_MACRO;
\end{verbatim}

\phantomsection
\label{prim:keyword-ret}
\subsection*{RET}
\index{RET@\texttt{RET}}
\begin{verbatim}
RET;
\end{verbatim}
Parameters: No parameters.

Returns from a local label call.

\paragraph{Example macro.}
\begin{verbatim}
$MACRO REFERENCE_RET FROM EDIT;
    CALL Done;
Done:
    RET;
END_MACRO;
\end{verbatim}

\phantomsection
\label{prim:intrinsic-param-str}
\subsection*{PARAM\_STR}
\index{PARAM\_STR@\texttt{PARAM\_STR}}
\begin{verbatim}
PARAM_STR(index)
\end{verbatim}
Parameters:
\begin{description}
\item[index] Zero-based parameter or selection index.
\end{description}

Returns one macro invocation parameter.

Returns a string.

\paragraph{Example macro.}
\begin{verbatim}
$MACRO REFERENCE_PARAM_STR FROM EDIT;
    DEF_STR(Result);
    Result := PARAM_STR(0);
END_MACRO;
\end{verbatim}

\phantomsection
\label{prim:keyword-then}
\subsection*{THEN}
\index{THEN!keyword@\textsc{keyword}}
\begin{verbatim}
IF expression THEN
    statement;
END;
\end{verbatim}
Parameters: expression is an integer truth value; statement is a macro statement.

Separates an IF condition from its true branch.

\paragraph{Example macro.}
\begin{verbatim}
$MACRO REFERENCE_THEN FROM EDIT;
    IF TRUE THEN
        MAKE_MESSAGE('true');
    END;
END_MACRO;
\end{verbatim}

\phantomsection
\label{prim:keyword-else}
\subsection*{ELSE}
\index{ELSE!keyword@\textsc{keyword}}
\begin{verbatim}
ELSE
    statement;
\end{verbatim}
Parameters: statement is a macro statement.

Introduces the optional false branch of the nearest IF block.

\paragraph{Example macro.}
\begin{verbatim}
$MACRO REFERENCE_ELSE FROM EDIT;
    IF FALSE THEN
        MAKE_MESSAGE('true');
    ELSE
        MAKE_MESSAGE('false');
    END;
END_MACRO;
\end{verbatim}

\phantomsection
\label{prim:keyword-end}
\subsection*{END}
\index{END!keyword@\textsc{keyword}}
\begin{verbatim}
END;
\end{verbatim}
Parameters: None.

Terminates the nearest IF or WHILE block.

\paragraph{Example macro.}
\begin{verbatim}
$MACRO REFERENCE_END FROM EDIT;
    IF TRUE THEN
        MAKE_MESSAGE('done');
    END;
END_MACRO;
\end{verbatim}

\phantomsection
\label{prim:keyword-do}
\subsection*{DO}
\index{DO!keyword@\textsc{keyword}}
\begin{verbatim}
WHILE expression DO
    statement;
END;
\end{verbatim}
Parameters: expression is an integer truth value; statement is a macro statement.

Separates a WHILE condition from its repeated body.

\paragraph{Example macro.}
\begin{verbatim}
$MACRO REFERENCE_DO FROM EDIT;
    WHILE FALSE DO
        MAKE_MESSAGE('not reached');
    END;
END_MACRO;
\end{verbatim}

\phantomsection
\label{prim:form-label}
\subsection*{Label}
\index{label!source form@\textsc{source form}}
\begin{verbatim}
labelName:
\end{verbatim}
Parameters: labelName is a unique local label identifier.

Defines the target used by GOTO and CALL in the same macro.

\paragraph{Example macro.}
\begin{verbatim}
$MACRO REFERENCE_LABEL_FORM FROM EDIT;
    GOTO REFERENCE_LABEL_DONE;
REFERENCE_LABEL_DONE:
END_MACRO;
\end{verbatim}


\part{Editor Automation}

\chapter{Text and String Operations}

Text primitives act through the current editor context. They obtain, insert and delete text and derive useful string or path values without exposing a document object. Each entry identifies the exact operation, its parameters and one complete macro.

\section{Primitive Map}
The map is a local index for this chapter. The final column is a generated page
reference and hyperlink to the full entry, where the source form, parameters,
semantics and a complete example macro appear.

\begin{longtable}{@{}L{0.38\textwidth}L{0.14\textwidth}L{0.27\textwidth}L{0.10\textwidth}@{}}
\toprule
Primitive & Role & Purpose & Reference \\
\midrule
\endhead
\toprule
Primitive & Role & Purpose & Reference \\
\midrule
\endhead
\primitiveoverview{GET\_EXTENSION}{intrinsic-get-extension}{intrinsic}{Text operation.}
\primitiveoverview{GET\_PATH}{intrinsic-get-path}{intrinsic}{Text operation.}
\primitiveoverview{GET\_WORD}{intrinsic-get-word}{intrinsic}{Text operation.}
\primitiveoverview{TRUNCATE\_EXTENSION}{intrinsic-truncate-extension}{intrinsic}{Text operation.}
\primitiveoverview{TRUNCATE\_PATH}{intrinsic-truncate-path}{intrinsic}{Text operation.}
\primitiveoverview{BACK\_SPACE}{procedure-back-space}{procedure}{Text operation.}
\primitiveoverview{CR}{procedure-cr}{procedure}{Text operation.}
\primitiveoverview{DEL\_CHAR}{procedure-del-char}{procedure}{Text operation.}
\primitiveoverview{DEL\_LINE}{procedure-del-line}{procedure}{Text operation.}
\primitiveoverview{DEL\_CHARS}{procedure-del-chars}{procedure}{Text operation.}
\primitiveoverview{PUT\_LINE}{procedure-put-line}{procedure}{Text operation.}
\primitiveoverview{TEXT}{procedure-text}{procedure}{Text operation.}
\bottomrule
\end{longtable}

\section{Reference Entries}

\phantomsection
\label{prim:intrinsic-get-extension}
\subsection*{GET\_EXTENSION}
\index{GET\_EXTENSION@\texttt{GET\_EXTENSION}}
\begin{verbatim}
GET_EXTENSION(path)
\end{verbatim}
Parameters:
\begin{description}
\item[path] A string-like file-system path.
\end{description}

Returns the filename extension.

Returns a string.

\paragraph{Example macro.}
\begin{verbatim}
$MACRO REFERENCE_GET_EXTENSION FROM EDIT;
    DEF_STR(Result);
    Result := GET_EXTENSION('reference');
END_MACRO;
\end{verbatim}

\phantomsection
\label{prim:intrinsic-get-path}
\subsection*{GET\_PATH}
\index{GET\_PATH@\texttt{GET\_PATH}}
\begin{verbatim}
GET_PATH(path)
\end{verbatim}
Parameters:
\begin{description}
\item[path] A string-like file-system path.
\end{description}

Returns the directory portion of a path.

Returns a string.

\paragraph{Example macro.}
\begin{verbatim}
$MACRO REFERENCE_GET_PATH FROM EDIT;
    DEF_STR(Result);
    Result := GET_PATH('reference');
END_MACRO;
\end{verbatim}

\phantomsection
\label{prim:intrinsic-get-word}
\subsection*{GET\_WORD}
\index{GET\_WORD@\texttt{GET\_WORD}}
\begin{verbatim}
GET_WORD(delimiters)
\end{verbatim}
Parameters:
\begin{description}
\item[delimiters] String-like set of word-delimiter characters.
\end{description}

Returns the current editor word using the supplied delimiters.

Returns a string.

\paragraph{Example macro.}
\begin{verbatim}
$MACRO REFERENCE_GET_WORD FROM EDIT;
    DEF_STR(Result);
    Result := GET_WORD('reference');
END_MACRO;
\end{verbatim}

\phantomsection
\label{prim:intrinsic-truncate-extension}
\subsection*{TRUNCATE\_EXTENSION}
\index{TRUNCATE\_EXTENSION@\texttt{TRUNCATE\_EXTENSION}}
\begin{verbatim}
TRUNCATE_EXTENSION(path)
\end{verbatim}
Parameters:
\begin{description}
\item[path] A string-like file-system path.
\end{description}

Removes the filename extension.

Returns a string.

\paragraph{Example macro.}
\begin{verbatim}
$MACRO REFERENCE_TRUNCATE_EXTENSION FROM EDIT;
    DEF_STR(Result);
    Result := TRUNCATE_EXTENSION('reference');
END_MACRO;
\end{verbatim}

\phantomsection
\label{prim:intrinsic-truncate-path}
\subsection*{TRUNCATE\_PATH}
\index{TRUNCATE\_PATH@\texttt{TRUNCATE\_PATH}}
\begin{verbatim}
TRUNCATE_PATH(path)
\end{verbatim}
Parameters:
\begin{description}
\item[path] A string-like file-system path.
\end{description}

Removes the directory portion.

Returns a string.

\paragraph{Example macro.}
\begin{verbatim}
$MACRO REFERENCE_TRUNCATE_PATH FROM EDIT;
    DEF_STR(Result);
    Result := TRUNCATE_PATH('reference');
END_MACRO;
\end{verbatim}

\phantomsection
\label{prim:procedure-back-space}
\subsection*{BACK\_SPACE}
\index{BACK\_SPACE@\texttt{BACK\_SPACE}}
\begin{verbatim}
BACK_SPACE;
\end{verbatim}
Parameters: None.

Performs the documented MR operation named by this primitive in the active execution context.

\paragraph{Example macro.}
\begin{verbatim}
$MACRO REFERENCE_BACK_SPACE FROM EDIT;
    BACK_SPACE;
END_MACRO;
\end{verbatim}

\phantomsection
\label{prim:procedure-cr}
\subsection*{CR}
\index{CR@\texttt{CR}}
\begin{verbatim}
CR;
\end{verbatim}
Parameters: None.

Performs the documented MR operation named by this primitive in the active execution context.

\paragraph{Example macro.}
\begin{verbatim}
$MACRO REFERENCE_CR FROM EDIT;
    CR;
END_MACRO;
\end{verbatim}

\phantomsection
\label{prim:procedure-del-char}
\subsection*{DEL\_CHAR}
\index{DEL\_CHAR@\texttt{DEL\_CHAR}}
\begin{verbatim}
DEL_CHAR;
\end{verbatim}
Parameters: None.

Performs the documented MR operation named by this primitive in the active execution context.

\paragraph{Example macro.}
\begin{verbatim}
$MACRO REFERENCE_DEL_CHAR FROM EDIT;
    DEL_CHAR;
END_MACRO;
\end{verbatim}

\phantomsection
\label{prim:procedure-del-line}
\subsection*{DEL\_LINE}
\index{DEL\_LINE@\texttt{DEL\_LINE}}
\begin{verbatim}
DEL_LINE;
\end{verbatim}
Parameters: None.

Performs the documented MR operation named by this primitive in the active execution context.

\paragraph{Example macro.}
\begin{verbatim}
$MACRO REFERENCE_DEL_LINE FROM EDIT;
    DEL_LINE;
END_MACRO;
\end{verbatim}

\phantomsection
\label{prim:procedure-del-chars}
\subsection*{DEL\_CHARS}
\index{DEL\_CHARS@\texttt{DEL\_CHARS}}
\begin{verbatim}
DEL_CHARS(count);
\end{verbatim}
Parameters:
\begin{description}
\item[count] Number of characters to process.
\end{description}

Deletes count characters at the current editor position.

\paragraph{Example macro.}
\begin{verbatim}
$MACRO REFERENCE_DEL_CHARS FROM EDIT;
    DEL_CHARS(0);
END_MACRO;
\end{verbatim}

\phantomsection
\label{prim:procedure-put-line}
\subsection*{PUT\_LINE}
\index{PUT\_LINE@\texttt{PUT\_LINE}}
\begin{verbatim}
PUT_LINE(text);
\end{verbatim}
Parameters:
\begin{description}
\item[text] String-like text.
\end{description}

Replaces or inserts the current editor line.

\paragraph{Example macro.}
\begin{verbatim}
$MACRO REFERENCE_PUT_LINE FROM EDIT;
    PUT_LINE('reference');
END_MACRO;
\end{verbatim}

\phantomsection
\label{prim:procedure-text}
\subsection*{TEXT}
\index{TEXT@\texttt{TEXT}}
\begin{verbatim}
TEXT(text);
\end{verbatim}
Parameters:
\begin{description}
\item[text] String-like text.
\end{description}

Inserts text at the current editor cursor.

\paragraph{Example macro.}
\begin{verbatim}
$MACRO REFERENCE_TEXT FROM EDIT;
    TEXT('reference');
END_MACRO;
\end{verbatim}


\chapter{Cursor, Tabs and Indentation}

Cursor primitives move within the active editor and view. Formatting primitives use the editor tab and indentation settings; they do not construct a private formatting model. Use the map to distinguish movement, marks, paging and indentation.

\section{Primitive Map}
The map is a local index for this chapter. The final column is a generated page
reference and hyperlink to the full entry, where the source form, parameters,
semantics and a complete example macro appear.

\begin{longtable}{@{}L{0.38\textwidth}L{0.14\textwidth}L{0.27\textwidth}L{0.10\textwidth}@{}}
\toprule
Primitive & Role & Purpose & Reference \\
\midrule
\endhead
\toprule
Primitive & Role & Purpose & Reference \\
\midrule
\endhead
\primitiveoverview{EXPAND\_TABS}{procedure-expand-tabs}{procedure}{Cursor or formatting.}
\primitiveoverview{TABS\_TO\_SPACES}{procedure-tabs-to-spaces}{procedure}{Cursor or formatting.}
\primitiveoverview{BACK\_HOME}{procedure-back-home}{procedure}{Cursor or formatting.}
\primitiveoverview{BACK\_WORD}{procedure-back-word}{procedure}{Cursor or formatting.}
\primitiveoverview{BOTTOM\_OF\_WINDOW}{procedure-bottom-of-window}{procedure}{Cursor or formatting.}
\primitiveoverview{CENTER\_LINE}{procedure-center-line}{procedure}{Cursor or formatting.}
\primitiveoverview{CENTER\_LINE\_ON\_SCREEN}{procedure-center-line-on-screen}{procedure}{Cursor or formatting.}
\primitiveoverview{DOWN}{procedure-down}{procedure}{Cursor or formatting.}
\primitiveoverview{EOF}{procedure-eof}{procedure}{Cursor or formatting.}
\primitiveoverview{EOL}{procedure-eol}{procedure}{Cursor or formatting.}
\primitiveoverview{FIRST\_WORD}{procedure-first-word}{procedure}{Cursor or formatting.}
\primitiveoverview{GET\_RANDOM\_MARK}{procedure-get-random-mark}{procedure}{Cursor or formatting.}
\primitiveoverview{GOTO\_MARK}{procedure-goto-mark}{procedure}{Cursor or formatting.}
\primitiveoverview{HOME}{procedure-home}{procedure}{Cursor or formatting.}
\primitiveoverview{INDENT}{procedure-indent}{procedure}{Cursor or formatting.}
\primitiveoverview{INDENT\_BLOCK}{procedure-indent-block}{procedure}{Cursor or formatting.}
\primitiveoverview{JUSTIFY\_PARAGRAPH}{procedure-justify-paragraph}{procedure}{Cursor or formatting.}
\primitiveoverview{LAST\_PAGE\_BREAK}{procedure-last-page-break}{procedure}{Cursor or formatting.}
\primitiveoverview{LEFT}{procedure-left}{procedure}{Cursor or formatting.}
\primitiveoverview{MARK\_POS}{procedure-mark-pos}{procedure}{Cursor or formatting.}
\primitiveoverview{MOVE\_VIEWPORT\_LEFT}{procedure-move-viewport-left}{procedure}{Cursor or formatting.}
\primitiveoverview{MOVE\_VIEWPORT\_RIGHT}{procedure-move-viewport-right}{procedure}{Cursor or formatting.}
\primitiveoverview{NEXT\_PAGE\_BREAK}{procedure-next-page-break}{procedure}{Cursor or formatting.}
\primitiveoverview{PAGE\_DOWN}{procedure-page-down}{procedure}{Cursor or formatting.}
\primitiveoverview{PAGE\_UP}{procedure-page-up}{procedure}{Cursor or formatting.}
\primitiveoverview{POP\_MARK}{procedure-pop-mark}{procedure}{Cursor or formatting.}
\primitiveoverview{REFORMAT\_DOCUMENT}{procedure-reformat-document}{procedure}{Cursor or formatting.}
\primitiveoverview{REFORMAT\_PARAGRAPH}{procedure-reformat-paragraph}{procedure}{Cursor or formatting.}
\primitiveoverview{REVERSE\_TAB}{procedure-reverse-tab}{procedure}{Cursor or formatting.}
\primitiveoverview{RIGHT}{procedure-right}{procedure}{Cursor or formatting.}
\primitiveoverview{SCROLL\_DOWN}{procedure-scroll-down}{procedure}{Cursor or formatting.}
\primitiveoverview{SCROLL\_UP}{procedure-scroll-up}{procedure}{Cursor or formatting.}
\primitiveoverview{SET\_INDENT\_LEVEL}{procedure-set-indent-level}{procedure}{Cursor or formatting.}
\primitiveoverview{SET\_LEFT\_MARGIN}{procedure-set-left-margin}{procedure}{Cursor or formatting.}
\primitiveoverview{SET\_RANDOM\_MARK}{procedure-set-random-mark}{procedure}{Cursor or formatting.}
\primitiveoverview{SET\_RIGHT\_MARGIN}{procedure-set-right-margin}{procedure}{Cursor or formatting.}
\primitiveoverview{TAB\_LEFT}{procedure-tab-left}{procedure}{Cursor or formatting.}
\primitiveoverview{TAB\_RIGHT}{procedure-tab-right}{procedure}{Cursor or formatting.}
\primitiveoverview{TOF}{procedure-tof}{procedure}{Cursor or formatting.}
\primitiveoverview{TOGGLE\_FORMAT\_RULER}{procedure-toggle-format-ruler}{procedure}{Cursor or formatting.}
\primitiveoverview{TOGGLE\_WORD\_WRAP}{procedure-toggle-word-wrap}{procedure}{Cursor or formatting.}
\primitiveoverview{TOP\_OF\_WINDOW}{procedure-top-of-window}{procedure}{Cursor or formatting.}
\primitiveoverview{UNDENT}{procedure-undent}{procedure}{Cursor or formatting.}
\primitiveoverview{UNDENT\_BLOCK}{procedure-undent-block}{procedure}{Cursor or formatting.}
\primitiveoverview{UP}{procedure-up}{procedure}{Cursor or formatting.}
\primitiveoverview{WORD\_LEFT}{procedure-word-left}{procedure}{Cursor or formatting.}
\primitiveoverview{WORD\_RIGHT}{procedure-word-right}{procedure}{Cursor or formatting.}
\primitiveoverview{WORD\_WRAP\_LINE}{procedure-word-wrap-line}{procedure}{Cursor or formatting.}
\primitiveoverview{GOTO\_COL}{procedure-goto-col}{procedure}{Cursor or formatting.}
\primitiveoverview{GOTO\_LINE}{procedure-goto-line}{procedure}{Cursor or formatting.}
\bottomrule
\end{longtable}

\section{Reference Entries}

\phantomsection
\label{prim:procedure-expand-tabs}
\subsection*{EXPAND\_TABS}
\index{EXPAND\_TABS@\texttt{EXPAND\_TABS}}
\begin{verbatim}
EXPAND_TABS(textVariable);
EXPAND_TABS(textVariable, tabWidthVariable);
\end{verbatim}
Parameters:
\begin{description}
\item[textVariable] Writable string variable.
\item[tabWidthVariable] Integer variable holding the tab width.
\end{description}

Expands tab characters in the supplied string variable in place.

\paragraph{Example macro.}
\begin{verbatim}
$MACRO REFERENCE_EXPAND_TABS FROM EDIT;
    DEF_STR(Text);
    DEF_INT(TabWidth);
    Text := 'a\tb';
    TabWidth := 4;
    EXPAND_TABS(Text, TabWidth);
END_MACRO;
\end{verbatim}

\phantomsection
\label{prim:procedure-tabs-to-spaces}
\subsection*{TABS\_TO\_SPACES}
\index{TABS\_TO\_SPACES@\texttt{TABS\_TO\_SPACES}}
\begin{verbatim}
TABS_TO_SPACES(textVariable);
\end{verbatim}
Parameters:
\begin{description}
\item[textVariable] Writable string variable.
\end{description}

Replaces tabs in the supplied string variable in place.

\paragraph{Example macro.}
\begin{verbatim}
$MACRO REFERENCE_TABS_TO_SPACES FROM EDIT;
    DEF_STR(Text);
    Text := 'a\tb';
    TABS_TO_SPACES(Text);
END_MACRO;
\end{verbatim}

\begin{mrmacextension}
\phantomsection
\label{prim:procedure-back-home}
\subsection*{BACK\_HOME}
\index{BACK\_HOME@\texttt{BACK\_HOME}}
\begin{verbatim}
BACK_HOME;
\end{verbatim}
Parameters: None.

Performs the named editor action through the active MR editor command route.

\paragraph{Example macro.}
\begin{verbatim}
$MACRO REFERENCE_BACK_HOME FROM EDIT;
    BACK_HOME;
END_MACRO;
\end{verbatim}
\end{mrmacextension}

\begin{mrmacextension}
\phantomsection
\label{prim:procedure-back-word}
\subsection*{BACK\_WORD}
\index{BACK\_WORD@\texttt{BACK\_WORD}}
\begin{verbatim}
BACK_WORD;
\end{verbatim}
Parameters: None.

Performs the named editor action through the active MR editor command route.

\paragraph{Example macro.}
\begin{verbatim}
$MACRO REFERENCE_BACK_WORD FROM EDIT;
    BACK_WORD;
END_MACRO;
\end{verbatim}
\end{mrmacextension}

\phantomsection
\label{prim:procedure-bottom-of-window}
\subsection*{BOTTOM\_OF\_WINDOW}
\index{BOTTOM\_OF\_WINDOW@\texttt{BOTTOM\_OF\_WINDOW}}
\begin{verbatim}
BOTTOM_OF_WINDOW;
\end{verbatim}
Parameters: None.

Performs the named editor action through the active MR editor command route.

\paragraph{Example macro.}
\begin{verbatim}
$MACRO REFERENCE_BOTTOM_OF_WINDOW FROM EDIT;
    BOTTOM_OF_WINDOW;
END_MACRO;
\end{verbatim}

\phantomsection
\label{prim:procedure-center-line}
\subsection*{CENTER\_LINE}
\index{CENTER\_LINE@\texttt{CENTER\_LINE}}
\begin{verbatim}
CENTER_LINE;
\end{verbatim}
Parameters: None.

Performs the named editor action through the active MR editor command route.

\paragraph{Example macro.}
\begin{verbatim}
$MACRO REFERENCE_CENTER_LINE FROM EDIT;
    CENTER_LINE;
END_MACRO;
\end{verbatim}

\phantomsection
\label{prim:procedure-center-line-on-screen}
\subsection*{CENTER\_LINE\_ON\_SCREEN}
\index{CENTER\_LINE\_ON\_SCREEN@\texttt{CENTER\_LINE\_ON\_SCREEN}}
\begin{verbatim}
CENTER_LINE_ON_SCREEN;
\end{verbatim}
Parameters: None.

Performs the named editor action through the active MR editor command route.

\paragraph{Example macro.}
\begin{verbatim}
$MACRO REFERENCE_CENTER_LINE_ON_SCREEN FROM EDIT;
    CENTER_LINE_ON_SCREEN;
END_MACRO;
\end{verbatim}

\phantomsection
\label{prim:procedure-down}
\subsection*{DOWN}
\index{DOWN@\texttt{DOWN}}
\begin{verbatim}
DOWN;
\end{verbatim}
Parameters: None.

Moves the editor cursor one position in the named direction.

\paragraph{Example macro.}
\begin{verbatim}
$MACRO REFERENCE_DOWN FROM EDIT;
    DOWN;
END_MACRO;
\end{verbatim}

\phantomsection
\label{prim:procedure-eof}
\subsection*{EOF}
\index{EOF@\texttt{EOF}}
\begin{verbatim}
EOF;
\end{verbatim}
Parameters: None.

Moves the editor cursor to the named line or file boundary.

\paragraph{Example macro.}
\begin{verbatim}
$MACRO REFERENCE_EOF FROM EDIT;
    EOF;
END_MACRO;
\end{verbatim}

\phantomsection
\label{prim:procedure-eol}
\subsection*{EOL}
\index{EOL@\texttt{EOL}}
\begin{verbatim}
EOL;
\end{verbatim}
Parameters: None.

Moves the editor cursor to the named line or file boundary.

\paragraph{Example macro.}
\begin{verbatim}
$MACRO REFERENCE_EOL FROM EDIT;
    EOL;
END_MACRO;
\end{verbatim}

\phantomsection
\label{prim:procedure-first-word}
\subsection*{FIRST\_WORD}
\index{FIRST\_WORD@\texttt{FIRST\_WORD}}
\begin{verbatim}
FIRST_WORD;
\end{verbatim}
Parameters: None.

Moves the editor cursor by word using the named direction or boundary.

\paragraph{Example macro.}
\begin{verbatim}
$MACRO REFERENCE_FIRST_WORD FROM EDIT;
    FIRST_WORD;
END_MACRO;
\end{verbatim}

\begin{mrmacextension}
\phantomsection
\label{prim:procedure-get-random-mark}
\subsection*{GET\_RANDOM\_MARK}
\index{GET\_RANDOM\_MARK@\texttt{GET\_RANDOM\_MARK}}
\begin{verbatim}
GET_RANDOM_MARK(markNumber);
\end{verbatim}
Parameters:
\begin{description}
\item[markNumber] Integer positional argument.
\end{description}

Performs the named editor action through the active MR editor command route.

\paragraph{Example macro.}
\begin{verbatim}
$MACRO REFERENCE_GET_RANDOM_MARK FROM EDIT;
    GET_RANDOM_MARK(0);
END_MACRO;
\end{verbatim}
\end{mrmacextension}

\phantomsection
\label{prim:procedure-goto-mark}
\subsection*{GOTO\_MARK}
\index{GOTO\_MARK@\texttt{GOTO\_MARK}}
\begin{verbatim}
GOTO_MARK;
\end{verbatim}
Parameters: None.

Performs the named operation on the editor mark stack.

\paragraph{Example macro.}
\begin{verbatim}
$MACRO REFERENCE_GOTO_MARK FROM EDIT;
    GOTO_MARK;
END_MACRO;
\end{verbatim}

\phantomsection
\label{prim:procedure-home}
\subsection*{HOME}
\index{HOME@\texttt{HOME}}
\begin{verbatim}
HOME;
\end{verbatim}
Parameters: None.

Moves the editor cursor to the named line or file boundary.

\paragraph{Example macro.}
\begin{verbatim}
$MACRO REFERENCE_HOME FROM EDIT;
    HOME;
END_MACRO;
\end{verbatim}

\phantomsection
\label{prim:procedure-indent}
\subsection*{INDENT}
\index{INDENT@\texttt{INDENT}}
\begin{verbatim}
INDENT;
\end{verbatim}
Parameters: None.

Performs the named tab, indentation or line-wrap action using active editor settings.

\paragraph{Example macro.}
\begin{verbatim}
$MACRO REFERENCE_INDENT FROM EDIT;
    INDENT;
END_MACRO;
\end{verbatim}

\begin{mrmacextension}
\phantomsection
\label{prim:procedure-indent-block}
\subsection*{INDENT\_BLOCK}
\index{INDENT\_BLOCK@\texttt{INDENT\_BLOCK}}
\begin{verbatim}
INDENT_BLOCK;
\end{verbatim}
Parameters: None.

Performs the named editor action through the active MR editor command route.

\paragraph{Example macro.}
\begin{verbatim}
$MACRO REFERENCE_INDENT_BLOCK FROM EDIT;
    INDENT_BLOCK;
END_MACRO;
\end{verbatim}
\end{mrmacextension}

\begin{mrmacextension}
\phantomsection
\label{prim:procedure-justify-paragraph}
\subsection*{JUSTIFY\_PARAGRAPH}
\index{JUSTIFY\_PARAGRAPH@\texttt{JUSTIFY\_PARAGRAPH}}
\begin{verbatim}
JUSTIFY_PARAGRAPH;
\end{verbatim}
Parameters: None.

Performs the named editor action through the active MR editor command route.

\paragraph{Example macro.}
\begin{verbatim}
$MACRO REFERENCE_JUSTIFY_PARAGRAPH FROM EDIT;
    JUSTIFY_PARAGRAPH;
END_MACRO;
\end{verbatim}
\end{mrmacextension}

\phantomsection
\label{prim:procedure-last-page-break}
\subsection*{LAST\_PAGE\_BREAK}
\index{LAST\_PAGE\_BREAK@\texttt{LAST\_PAGE\_BREAK}}
\begin{verbatim}
LAST_PAGE_BREAK;
\end{verbatim}
Parameters: None.

Moves the editor view by the named page operation.

\paragraph{Example macro.}
\begin{verbatim}
$MACRO REFERENCE_LAST_PAGE_BREAK FROM EDIT;
    LAST_PAGE_BREAK;
END_MACRO;
\end{verbatim}

\phantomsection
\label{prim:procedure-left}
\subsection*{LEFT}
\index{LEFT@\texttt{LEFT}}
\begin{verbatim}
LEFT;
\end{verbatim}
Parameters: None.

Moves the editor cursor one position in the named direction.

\paragraph{Example macro.}
\begin{verbatim}
$MACRO REFERENCE_LEFT FROM EDIT;
    LEFT;
END_MACRO;
\end{verbatim}

\phantomsection
\label{prim:procedure-mark-pos}
\subsection*{MARK\_POS}
\index{MARK\_POS@\texttt{MARK\_POS}}
\begin{verbatim}
MARK_POS;
\end{verbatim}
Parameters: None.

Performs the named operation on the editor mark stack.

\paragraph{Example macro.}
\begin{verbatim}
$MACRO REFERENCE_MARK_POS FROM EDIT;
    MARK_POS;
END_MACRO;
\end{verbatim}

\begin{mrmacextension}
\phantomsection
\label{prim:procedure-move-viewport-left}
\subsection*{MOVE\_VIEWPORT\_LEFT}
\index{MOVE\_VIEWPORT\_LEFT@\texttt{MOVE\_VIEWPORT\_LEFT}}
\begin{verbatim}
MOVE_VIEWPORT_LEFT;
\end{verbatim}
Parameters: None.

Performs the named MR window, screen or desktop operation in the active editor context.

\paragraph{Example macro.}
\begin{verbatim}
$MACRO REFERENCE_MOVE_VIEWPORT_LEFT FROM EDIT;
    MOVE_VIEWPORT_LEFT;
END_MACRO;
\end{verbatim}
\end{mrmacextension}

\begin{mrmacextension}
\phantomsection
\label{prim:procedure-move-viewport-right}
\subsection*{MOVE\_VIEWPORT\_RIGHT}
\index{MOVE\_VIEWPORT\_RIGHT@\texttt{MOVE\_VIEWPORT\_RIGHT}}
\begin{verbatim}
MOVE_VIEWPORT_RIGHT;
\end{verbatim}
Parameters: None.

Performs the named MR window, screen or desktop operation in the active editor context.

\paragraph{Example macro.}
\begin{verbatim}
$MACRO REFERENCE_MOVE_VIEWPORT_RIGHT FROM EDIT;
    MOVE_VIEWPORT_RIGHT;
END_MACRO;
\end{verbatim}
\end{mrmacextension}

\phantomsection
\label{prim:procedure-next-page-break}
\subsection*{NEXT\_PAGE\_BREAK}
\index{NEXT\_PAGE\_BREAK@\texttt{NEXT\_PAGE\_BREAK}}
\begin{verbatim}
NEXT_PAGE_BREAK;
\end{verbatim}
Parameters: None.

Moves the editor view by the named page operation.

\paragraph{Example macro.}
\begin{verbatim}
$MACRO REFERENCE_NEXT_PAGE_BREAK FROM EDIT;
    NEXT_PAGE_BREAK;
END_MACRO;
\end{verbatim}

\phantomsection
\label{prim:procedure-page-down}
\subsection*{PAGE\_DOWN}
\index{PAGE\_DOWN@\texttt{PAGE\_DOWN}}
\begin{verbatim}
PAGE_DOWN;
\end{verbatim}
Parameters: None.

Moves the editor view by the named page operation.

\paragraph{Example macro.}
\begin{verbatim}
$MACRO REFERENCE_PAGE_DOWN FROM EDIT;
    PAGE_DOWN;
END_MACRO;
\end{verbatim}

\phantomsection
\label{prim:procedure-page-up}
\subsection*{PAGE\_UP}
\index{PAGE\_UP@\texttt{PAGE\_UP}}
\begin{verbatim}
PAGE_UP;
\end{verbatim}
Parameters: None.

Moves the editor view by the named page operation.

\paragraph{Example macro.}
\begin{verbatim}
$MACRO REFERENCE_PAGE_UP FROM EDIT;
    PAGE_UP;
END_MACRO;
\end{verbatim}

\phantomsection
\label{prim:procedure-pop-mark}
\subsection*{POP\_MARK}
\index{POP\_MARK@\texttt{POP\_MARK}}
\begin{verbatim}
POP_MARK;
\end{verbatim}
Parameters: None.

Performs the named operation on the editor mark stack.

\paragraph{Example macro.}
\begin{verbatim}
$MACRO REFERENCE_POP_MARK FROM EDIT;
    POP_MARK;
END_MACRO;
\end{verbatim}

\begin{mrmacextension}
\phantomsection
\label{prim:procedure-reformat-document}
\subsection*{REFORMAT\_DOCUMENT}
\index{REFORMAT\_DOCUMENT@\texttt{REFORMAT\_DOCUMENT}}
\begin{verbatim}
REFORMAT_DOCUMENT;
\end{verbatim}
Parameters: None.

Performs the named editor action through the active MR editor command route.

\paragraph{Example macro.}
\begin{verbatim}
$MACRO REFERENCE_REFORMAT_DOCUMENT FROM EDIT;
    REFORMAT_DOCUMENT;
END_MACRO;
\end{verbatim}
\end{mrmacextension}

\begin{mrmacextension}
\phantomsection
\label{prim:procedure-reformat-paragraph}
\subsection*{REFORMAT\_PARAGRAPH}
\index{REFORMAT\_PARAGRAPH@\texttt{REFORMAT\_PARAGRAPH}}
\begin{verbatim}
REFORMAT_PARAGRAPH;
\end{verbatim}
Parameters: None.

Performs the named editor action through the active MR editor command route.

\paragraph{Example macro.}
\begin{verbatim}
$MACRO REFERENCE_REFORMAT_PARAGRAPH FROM EDIT;
    REFORMAT_PARAGRAPH;
END_MACRO;
\end{verbatim}
\end{mrmacextension}

\phantomsection
\label{prim:procedure-reverse-tab}
\subsection*{REVERSE\_TAB}
\index{REVERSE\_TAB@\texttt{REVERSE\_TAB}}
\begin{verbatim}
REVERSE_TAB;
\end{verbatim}
Parameters: None.

Performs the named tab, indentation or line-wrap action using active editor settings.

\paragraph{Example macro.}
\begin{verbatim}
$MACRO REFERENCE_REVERSE_TAB FROM EDIT;
    REVERSE_TAB;
END_MACRO;
\end{verbatim}

\phantomsection
\label{prim:procedure-right}
\subsection*{RIGHT}
\index{RIGHT@\texttt{RIGHT}}
\begin{verbatim}
RIGHT;
\end{verbatim}
Parameters: None.

Moves the editor cursor one position in the named direction.

\paragraph{Example macro.}
\begin{verbatim}
$MACRO REFERENCE_RIGHT FROM EDIT;
    RIGHT;
END_MACRO;
\end{verbatim}

\phantomsection
\label{prim:procedure-scroll-down}
\subsection*{SCROLL\_DOWN}
\index{SCROLL\_DOWN@\texttt{SCROLL\_DOWN}}
\begin{verbatim}
SCROLL_DOWN;
\end{verbatim}
Parameters: None.

Performs the named editor action through the active MR editor command route.

\paragraph{Example macro.}
\begin{verbatim}
$MACRO REFERENCE_SCROLL_DOWN FROM EDIT;
    SCROLL_DOWN;
END_MACRO;
\end{verbatim}

\phantomsection
\label{prim:procedure-scroll-up}
\subsection*{SCROLL\_UP}
\index{SCROLL\_UP@\texttt{SCROLL\_UP}}
\begin{verbatim}
SCROLL_UP;
\end{verbatim}
Parameters: None.

Performs the named editor action through the active MR editor command route.

\paragraph{Example macro.}
\begin{verbatim}
$MACRO REFERENCE_SCROLL_UP FROM EDIT;
    SCROLL_UP;
END_MACRO;
\end{verbatim}

\phantomsection
\label{prim:procedure-set-indent-level}
\subsection*{SET\_INDENT\_LEVEL}
\index{SET\_INDENT\_LEVEL@\texttt{SET\_INDENT\_LEVEL}}
\begin{verbatim}
SET_INDENT_LEVEL;
\end{verbatim}
Parameters: None.

Performs the named tab, indentation or line-wrap action using active editor settings.

\paragraph{Example macro.}
\begin{verbatim}
$MACRO REFERENCE_SET_INDENT_LEVEL FROM EDIT;
    SET_INDENT_LEVEL;
END_MACRO;
\end{verbatim}

\begin{mrmacextension}
\phantomsection
\label{prim:procedure-set-left-margin}
\subsection*{SET\_LEFT\_MARGIN}
\index{SET\_LEFT\_MARGIN@\texttt{SET\_LEFT\_MARGIN}}
\begin{verbatim}
SET_LEFT_MARGIN;
\end{verbatim}
Parameters: None.

Performs the named editor action through the active MR editor command route.

\end{mrmacextension}

\paragraph{Example macro.}
\begin{verbatim}
$MACRO REFERENCE_SET_LEFT_MARGIN FROM EDIT;
    SET_LEFT_MARGIN;
END_MACRO;
\end{verbatim}

\begin{mrmacextension}
\phantomsection
\label{prim:procedure-set-random-mark}
\subsection*{SET\_RANDOM\_MARK}
\index{SET\_RANDOM\_MARK@\texttt{SET\_RANDOM\_MARK}}
\begin{verbatim}
SET_RANDOM_MARK(markNumber);
\end{verbatim}
Parameters:
\begin{description}
\item[markNumber] Integer positional argument.
\end{description}

Performs the named editor action through the active MR editor command route.

\paragraph{Example macro.}
\begin{verbatim}
$MACRO REFERENCE_SET_RANDOM_MARK FROM EDIT;
    SET_RANDOM_MARK(0);
END_MACRO;
\end{verbatim}
\end{mrmacextension}

\begin{mrmacextension}
\phantomsection
\label{prim:procedure-set-right-margin}
\subsection*{SET\_RIGHT\_MARGIN}
\index{SET\_RIGHT\_MARGIN@\texttt{SET\_RIGHT\_MARGIN}}
\begin{verbatim}
SET_RIGHT_MARGIN;
\end{verbatim}
Parameters: None.

Performs the named editor action through the active MR editor command route.

\paragraph{Example macro.}
\begin{verbatim}
$MACRO REFERENCE_SET_RIGHT_MARGIN FROM EDIT;
    SET_RIGHT_MARGIN;
END_MACRO;
\end{verbatim}
\end{mrmacextension}

\phantomsection
\label{prim:procedure-tab-left}
\subsection*{TAB\_LEFT}
\index{TAB\_LEFT@\texttt{TAB\_LEFT}}
\begin{verbatim}
TAB_LEFT;
\end{verbatim}
Parameters: None.

Performs the named tab, indentation or line-wrap action using active editor settings.

\paragraph{Example macro.}
\begin{verbatim}
$MACRO REFERENCE_TAB_LEFT FROM EDIT;
    TAB_LEFT;
END_MACRO;
\end{verbatim}

\phantomsection
\label{prim:procedure-tab-right}
\subsection*{TAB\_RIGHT}
\index{TAB\_RIGHT@\texttt{TAB\_RIGHT}}
\begin{verbatim}
TAB_RIGHT;
\end{verbatim}
Parameters: None.

Performs the named tab, indentation or line-wrap action using active editor settings.

\paragraph{Example macro.}
\begin{verbatim}
$MACRO REFERENCE_TAB_RIGHT FROM EDIT;
    TAB_RIGHT;
END_MACRO;
\end{verbatim}

\phantomsection
\label{prim:procedure-tof}
\subsection*{TOF}
\index{TOF@\texttt{TOF}}
\begin{verbatim}
TOF;
\end{verbatim}
Parameters: None.

Moves the editor cursor to the named line or file boundary.

\paragraph{Example macro.}
\begin{verbatim}
$MACRO REFERENCE_TOF FROM EDIT;
    TOF;
END_MACRO;
\end{verbatim}

\begin{mrmacextension}
\phantomsection
\label{prim:procedure-toggle-format-ruler}
\subsection*{TOGGLE\_FORMAT\_RULER}
\index{TOGGLE\_FORMAT\_RULER@\texttt{TOGGLE\_FORMAT\_RULER}}
\begin{verbatim}
TOGGLE_FORMAT_RULER;
\end{verbatim}
Parameters: None.

Performs the named editor action through the active MR editor command route.

\paragraph{Example macro.}
\begin{verbatim}
$MACRO REFERENCE_TOGGLE_FORMAT_RULER FROM EDIT;
    TOGGLE_FORMAT_RULER;
END_MACRO;
\end{verbatim}
\end{mrmacextension}

\begin{mrmacextension}
\phantomsection
\label{prim:procedure-toggle-word-wrap}
\subsection*{TOGGLE\_WORD\_WRAP}
\index{TOGGLE\_WORD\_WRAP@\texttt{TOGGLE\_WORD\_WRAP}}
\begin{verbatim}
TOGGLE_WORD_WRAP;
\end{verbatim}
Parameters: None.

Performs the named editor action through the active MR editor command route.

\paragraph{Example macro.}
\begin{verbatim}
$MACRO REFERENCE_TOGGLE_WORD_WRAP FROM EDIT;
    TOGGLE_WORD_WRAP;
END_MACRO;
\end{verbatim}
\end{mrmacextension}

\phantomsection
\label{prim:procedure-top-of-window}
\subsection*{TOP\_OF\_WINDOW}
\index{TOP\_OF\_WINDOW@\texttt{TOP\_OF\_WINDOW}}
\begin{verbatim}
TOP_OF_WINDOW;
\end{verbatim}
Parameters: None.

Performs the named editor action through the active MR editor command route.

\paragraph{Example macro.}
\begin{verbatim}
$MACRO REFERENCE_TOP_OF_WINDOW FROM EDIT;
    TOP_OF_WINDOW;
END_MACRO;
\end{verbatim}

\phantomsection
\label{prim:procedure-undent}
\subsection*{UNDENT}
\index{UNDENT@\texttt{UNDENT}}
\begin{verbatim}
UNDENT;
\end{verbatim}
Parameters: None.

Performs the named tab, indentation or line-wrap action using active editor settings.

\paragraph{Example macro.}
\begin{verbatim}
$MACRO REFERENCE_UNDENT FROM EDIT;
    UNDENT;
END_MACRO;
\end{verbatim}

\begin{mrmacextension}
\phantomsection
\label{prim:procedure-undent-block}
\subsection*{UNDENT\_BLOCK}
\index{UNDENT\_BLOCK@\texttt{UNDENT\_BLOCK}}
\begin{verbatim}
UNDENT_BLOCK;
\end{verbatim}
Parameters: None.

Performs the documented MR operation named by this primitive in the active execution context.

\paragraph{Example macro.}
\begin{verbatim}
$MACRO REFERENCE_UNDENT_BLOCK FROM EDIT;
    UNDENT_BLOCK;
END_MACRO;
\end{verbatim}
\end{mrmacextension}

\phantomsection
\label{prim:procedure-up}
\subsection*{UP}
\index{UP@\texttt{UP}}
\begin{verbatim}
UP;
\end{verbatim}
Parameters: None.

Moves the editor cursor one position in the named direction.

\paragraph{Example macro.}
\begin{verbatim}
$MACRO REFERENCE_UP FROM EDIT;
    UP;
END_MACRO;
\end{verbatim}

\phantomsection
\label{prim:procedure-word-left}
\subsection*{WORD\_LEFT}
\index{WORD\_LEFT@\texttt{WORD\_LEFT}}
\begin{verbatim}
WORD_LEFT;
\end{verbatim}
Parameters: None.

Moves the editor cursor by word using the named direction or boundary.

\paragraph{Example macro.}
\begin{verbatim}
$MACRO REFERENCE_WORD_LEFT FROM EDIT;
    WORD_LEFT;
END_MACRO;
\end{verbatim}

\phantomsection
\label{prim:procedure-word-right}
\subsection*{WORD\_RIGHT}
\index{WORD\_RIGHT@\texttt{WORD\_RIGHT}}
\begin{verbatim}
WORD_RIGHT;
\end{verbatim}
Parameters: None.

Moves the editor cursor by word using the named direction or boundary.

\paragraph{Example macro.}
\begin{verbatim}
$MACRO REFERENCE_WORD_RIGHT FROM EDIT;
    WORD_RIGHT;
END_MACRO;
\end{verbatim}

\phantomsection
\label{prim:procedure-word-wrap-line}
\subsection*{WORD\_WRAP\_LINE}
\index{WORD\_WRAP\_LINE@\texttt{WORD\_WRAP\_LINE}}
\begin{verbatim}
WORD_WRAP_LINE;
\end{verbatim}
Parameters: None.

Performs the named tab, indentation or line-wrap action using active editor settings.

\paragraph{Example macro.}
\begin{verbatim}
$MACRO REFERENCE_WORD_WRAP_LINE FROM EDIT;
    WORD_WRAP_LINE;
END_MACRO;
\end{verbatim}

\phantomsection
\label{prim:procedure-goto-col}
\subsection*{GOTO\_COL}
\index{GOTO\_COL@\texttt{GOTO\_COL}}
\begin{verbatim}
GOTO_COL(column);
\end{verbatim}
Parameters:
\begin{description}
\item[column] Screen or editor column, as required by the primitive.
\end{description}

Moves the editor cursor to the specified column.

\paragraph{Example macro.}
\begin{verbatim}
$MACRO REFERENCE_GOTO_COL FROM EDIT;
    GOTO_COL(0);
END_MACRO;
\end{verbatim}

\phantomsection
\label{prim:procedure-goto-line}
\subsection*{GOTO\_LINE}
\index{GOTO\_LINE@\texttt{GOTO\_LINE}}
\begin{verbatim}
GOTO_LINE(line);
\end{verbatim}
Parameters:
\begin{description}
\item[line] One-based editor line number.
\end{description}

Moves the editor cursor to the specified line.

\paragraph{Example macro.}
\begin{verbatim}
$MACRO REFERENCE_GOTO_LINE FROM EDIT;
    GOTO_LINE(0);
END_MACRO;
\end{verbatim}


\chapter{Blocks, Clipboard and Undo}

A block is the current editor selection. Block primitives start, extend, inspect, copy, move, delete and persist that selection. Clipboard and undo operations use editor-owned state in the active execution context.

\section{Primitive Map}
The map is a local index for this chapter. The final column is a generated page
reference and hyperlink to the full entry, where the source form, parameters,
semantics and a complete example macro appear.

\begin{longtable}{@{}L{0.38\textwidth}L{0.14\textwidth}L{0.27\textwidth}L{0.10\textwidth}@{}}
\toprule
Primitive & Role & Purpose & Reference \\
\midrule
\endhead
\toprule
Primitive & Role & Purpose & Reference \\
\midrule
\endhead
\primitiveoverview{BLOCK\_TEXT}{intrinsic-block-text}{intrinsic}{Block operation.}
\primitiveoverview{APPEND\_BLOCK}{procedure-append-block}{procedure}{Block operation.}
\primitiveoverview{BLOCK\_BEGIN}{procedure-block-begin}{procedure}{Block operation.}
\primitiveoverview{BLOCK\_END}{procedure-block-end}{procedure}{Block operation.}
\primitiveoverview{BLOCK\_MARK\_LINES}{procedure-block-mark-lines}{procedure}{Block operation.}
\primitiveoverview{BLOCK\_OFF}{procedure-block-off}{procedure}{Block operation.}
\primitiveoverview{BLOCK\_STAT}{procedure-block-stat}{procedure}{Block operation.}
\primitiveoverview{BLOCK\_TOGGLE\_MARKING}{procedure-block-toggle-marking}{procedure}{Block operation.}
\primitiveoverview{BLOCK\_TOGGLE\_VISIBILITY}{procedure-block-toggle-visibility}{procedure}{Block operation.}
\primitiveoverview{COL\_BLOCK\_BEGIN}{procedure-col-block-begin}{procedure}{Block operation.}
\primitiveoverview{COPY\_BLOCK}{procedure-copy-block}{procedure}{Block operation.}
\primitiveoverview{COPY\_BLOCK\_TO\_CLIPBOARD}{procedure-copy-block-to-clipboard}{procedure}{Block operation.}
\primitiveoverview{CUT\_APPEND\_BLOCK}{procedure-cut-append-block}{procedure}{Block operation.}
\primitiveoverview{CUT\_BLOCK}{procedure-cut-block}{procedure}{Block operation.}
\primitiveoverview{DELETE\_BLOCK}{procedure-delete-block}{procedure}{Block operation.}
\primitiveoverview{DEL\_CHAR\_OR\_BLOCK}{procedure-del-char-or-block}{procedure}{Block operation.}
\primitiveoverview{DEL\_EOL}{procedure-del-eol}{procedure}{Block operation.}
\primitiveoverview{DEL\_WORD}{procedure-del-word}{procedure}{Block operation.}
\primitiveoverview{END\_OF\_BLOCK}{procedure-end-of-block}{procedure}{Block operation.}
\primitiveoverview{EXTEND\_BLOCK\_BY\_MOTION}{procedure-extend-block-by-motion}{procedure}{Block operation.}
\primitiveoverview{MARK\_WORD\_RIGHT}{procedure-mark-word-right}{procedure}{Block operation.}
\primitiveoverview{MOVE\_BLOCK}{procedure-move-block}{procedure}{Block operation.}
\primitiveoverview{PASTE\_BLOCK}{procedure-paste-block}{procedure}{Block operation.}
\primitiveoverview{PASTE\_FROM\_CLIPBOARD}{procedure-paste-from-clipboard}{procedure}{Block operation.}
\primitiveoverview{REDO}{procedure-redo}{procedure}{Block operation.}
\primitiveoverview{SORT\_COLUMN\_BLOCK\_TOGGLE}{procedure-sort-column-block-toggle}{procedure}{Block operation.}
\primitiveoverview{START\_OF\_BLOCK}{procedure-start-of-block}{procedure}{Block operation.}
\primitiveoverview{STR\_BLOCK\_BEGIN}{procedure-str-block-begin}{procedure}{Block operation.}
\primitiveoverview{UNDO}{procedure-undo}{procedure}{Block operation.}
\primitiveoverview{LOAD\_BLOCK}{procedure-load-block}{procedure}{Block operation.}
\primitiveoverview{SAVE\_BLOCK}{procedure-save-block}{procedure}{Block operation.}
\bottomrule
\end{longtable}

\section{Reference Entries}

\phantomsection
\label{prim:intrinsic-block-text}
\subsection*{BLOCK\_TEXT}
\index{BLOCK\_TEXT@\texttt{BLOCK\_TEXT}}
\begin{verbatim}
BLOCK_TEXT
\end{verbatim}
Parameters: None.

Returns the selected block text.

Returns a string.

\paragraph{Example macro.}
\begin{verbatim}
$MACRO REFERENCE_BLOCK_TEXT FROM EDIT;
    DEF_STR(Result);
    Result := BLOCK_TEXT;
END_MACRO;
\end{verbatim}

\phantomsection
\label{prim:procedure-append-block}
\subsection*{APPEND\_BLOCK}
\index{APPEND\_BLOCK@\texttt{APPEND\_BLOCK}}
\begin{verbatim}
APPEND_BLOCK;
\end{verbatim}
Parameters: None.

Performs the named editor action through the active MR editor command route.

\paragraph{Example macro.}
\begin{verbatim}
$MACRO REFERENCE_APPEND_BLOCK FROM EDIT;
    APPEND_BLOCK;
END_MACRO;
\end{verbatim}

\begin{mrmacextension}
This entry is an MRMAC editor-command extension.
\end{mrmacextension}

\phantomsection
\label{prim:procedure-block-begin}
\subsection*{BLOCK\_BEGIN}
\index{BLOCK\_BEGIN@\texttt{BLOCK\_BEGIN}}
\begin{verbatim}
BLOCK_BEGIN;
\end{verbatim}
Parameters: None.

Performs the named selection or block operation in the active editor context.

\paragraph{Example macro.}
\begin{verbatim}
$MACRO REFERENCE_BLOCK_BEGIN FROM EDIT;
    BLOCK_BEGIN;
END_MACRO;
\end{verbatim}

\phantomsection
\label{prim:procedure-block-end}
\subsection*{BLOCK\_END}
\index{BLOCK\_END@\texttt{BLOCK\_END}}
\begin{verbatim}
BLOCK_END;
\end{verbatim}
Parameters: None.

Performs the named selection or block operation in the active editor context.

\paragraph{Example macro.}
\begin{verbatim}
$MACRO REFERENCE_BLOCK_END FROM EDIT;
    BLOCK_END;
END_MACRO;
\end{verbatim}

\begin{mrmacextension}
\phantomsection
\label{prim:procedure-block-mark-lines}
\subsection*{BLOCK\_MARK\_LINES}
\index{BLOCK\_MARK\_LINES@\texttt{BLOCK\_MARK\_LINES}}
\begin{verbatim}
BLOCK_MARK_LINES;
\end{verbatim}
Parameters: None.

Performs the named selection or block operation in the active editor context.

\paragraph{Example macro.}
\begin{verbatim}
$MACRO REFERENCE_BLOCK_MARK_LINES FROM EDIT;
    BLOCK_MARK_LINES;
END_MACRO;
\end{verbatim}
\end{mrmacextension}

\phantomsection
\label{prim:procedure-block-off}
\subsection*{BLOCK\_OFF}
\index{BLOCK\_OFF@\texttt{BLOCK\_OFF}}
\begin{verbatim}
BLOCK_OFF;
\end{verbatim}
Parameters: None.

Performs the named selection or block operation in the active editor context.

\paragraph{Example macro.}
\begin{verbatim}
$MACRO REFERENCE_BLOCK_OFF FROM EDIT;
    BLOCK_OFF;
END_MACRO;
\end{verbatim}

\phantomsection
\label{prim:procedure-block-stat}
\subsection*{BLOCK\_STAT}
\index{BLOCK\_STAT@\texttt{BLOCK\_STAT}}
\begin{verbatim}
BLOCK_STAT;
\end{verbatim}
Parameters: None.

Performs the named selection or block operation in the active editor context.

\paragraph{Example macro.}
\begin{verbatim}
$MACRO REFERENCE_BLOCK_STAT FROM EDIT;
    BLOCK_STAT;
END_MACRO;
\end{verbatim}

\begin{mrmacextension}
\phantomsection
\label{prim:procedure-block-toggle-marking}
\subsection*{BLOCK\_TOGGLE\_MARKING}
\index{BLOCK\_TOGGLE\_MARKING@\texttt{BLOCK\_TOGGLE\_MARKING}}
\begin{verbatim}
BLOCK_TOGGLE_MARKING;
\end{verbatim}
Parameters: None.

Performs the named selection or block operation in the active editor context.

\paragraph{Example macro.}
\begin{verbatim}
$MACRO REFERENCE_BLOCK_TOGGLE_MARKING FROM EDIT;
    BLOCK_TOGGLE_MARKING;
END_MACRO;
\end{verbatim}
\end{mrmacextension}

\begin{mrmacextension}
\phantomsection
\label{prim:procedure-block-toggle-visibility}
\subsection*{BLOCK\_TOGGLE\_VISIBILITY}
\index{BLOCK\_TOGGLE\_VISIBILITY@\texttt{BLOCK\_TOGGLE\_VISIBILITY}}
\begin{verbatim}
BLOCK_TOGGLE_VISIBILITY;
\end{verbatim}
Parameters: None.

Performs the named selection or block operation in the active editor context.

\paragraph{Example macro.}
\begin{verbatim}
$MACRO REFERENCE_BLOCK_TOGGLE_VISIBILITY FROM EDIT;
    BLOCK_TOGGLE_VISIBILITY;
END_MACRO;
\end{verbatim}
\end{mrmacextension}

\phantomsection
\label{prim:procedure-col-block-begin}
\subsection*{COL\_BLOCK\_BEGIN}
\index{COL\_BLOCK\_BEGIN@\texttt{COL\_BLOCK\_BEGIN}}
\begin{verbatim}
COL_BLOCK_BEGIN;
\end{verbatim}
Parameters: None.

Performs the documented MR operation named by this primitive in the active execution context.

\paragraph{Example macro.}
\begin{verbatim}
$MACRO REFERENCE_COL_BLOCK_BEGIN FROM EDIT;
    COL_BLOCK_BEGIN;
END_MACRO;
\end{verbatim}

\phantomsection
\label{prim:procedure-copy-block}
\subsection*{COPY\_BLOCK}
\index{COPY\_BLOCK@\texttt{COPY\_BLOCK}}
\begin{verbatim}
COPY_BLOCK;
\end{verbatim}
Parameters: None.

Performs the named selection or block operation in the active editor context.

\paragraph{Example macro.}
\begin{verbatim}
$MACRO REFERENCE_COPY_BLOCK FROM EDIT;
    COPY_BLOCK;
END_MACRO;
\end{verbatim}

\begin{mrmacextension}
\phantomsection
\label{prim:procedure-copy-block-to-clipboard}
\subsection*{COPY\_BLOCK\_TO\_CLIPBOARD}
\index{COPY\_BLOCK\_TO\_CLIPBOARD@\texttt{COPY\_BLOCK\_TO\_CLIPBOARD}}
\begin{verbatim}
COPY_BLOCK_TO_CLIPBOARD;
\end{verbatim}
Parameters: None.

Performs the named selection or block operation in the active editor context.

\paragraph{Example macro.}
\begin{verbatim}
$MACRO REFERENCE_COPY_BLOCK_TO_CLIPBOARD FROM EDIT;
    COPY_BLOCK_TO_CLIPBOARD;
END_MACRO;
\end{verbatim}
\end{mrmacextension}

\begin{mrmacextension}
\phantomsection
\label{prim:procedure-cut-append-block}
\subsection*{CUT\_APPEND\_BLOCK}
\index{CUT\_APPEND\_BLOCK@\texttt{CUT\_APPEND\_BLOCK}}
\begin{verbatim}
CUT_APPEND_BLOCK;
\end{verbatim}
Parameters: None.

Performs the named selection or block operation in the active editor context.

\paragraph{Example macro.}
\begin{verbatim}
$MACRO REFERENCE_CUT_APPEND_BLOCK FROM EDIT;
    CUT_APPEND_BLOCK;
END_MACRO;
\end{verbatim}
\end{mrmacextension}

\begin{mrmacextension}
\phantomsection
\label{prim:procedure-cut-block}
\subsection*{CUT\_BLOCK}
\index{CUT\_BLOCK@\texttt{CUT\_BLOCK}}
\begin{verbatim}
CUT_BLOCK;
\end{verbatim}
Parameters: None.

Performs the named selection or block operation in the active editor context.

\paragraph{Example macro.}
\begin{verbatim}
$MACRO REFERENCE_CUT_BLOCK FROM EDIT;
    CUT_BLOCK;
END_MACRO;
\end{verbatim}
\end{mrmacextension}

\phantomsection
\label{prim:procedure-delete-block}
\subsection*{DELETE\_BLOCK}
\index{DELETE\_BLOCK@\texttt{DELETE\_BLOCK}}
\begin{verbatim}
DELETE_BLOCK;
\end{verbatim}
Parameters: None.

Performs the named selection or block operation in the active editor context.

\paragraph{Example macro.}
\begin{verbatim}
$MACRO REFERENCE_DELETE_BLOCK FROM EDIT;
    DELETE_BLOCK;
END_MACRO;
\end{verbatim}

\begin{mrmacextension}
\phantomsection
\label{prim:procedure-del-char-or-block}
\subsection*{DEL\_CHAR\_OR\_BLOCK}
\index{DEL\_CHAR\_OR\_BLOCK@\texttt{DEL\_CHAR\_OR\_BLOCK}}
\begin{verbatim}
DEL_CHAR_OR_BLOCK;
\end{verbatim}
Parameters: None.

Performs the named editor action through the active MR editor command route.

\paragraph{Example macro.}
\begin{verbatim}
$MACRO REFERENCE_DEL_CHAR_OR_BLOCK FROM EDIT;
    DEL_CHAR_OR_BLOCK;
END_MACRO;
\end{verbatim}
\end{mrmacextension}

\begin{mrmacextension}
\phantomsection
\label{prim:procedure-del-eol}
\subsection*{DEL\_EOL}
\index{DEL\_EOL@\texttt{DEL\_EOL}}
\begin{verbatim}
DEL_EOL;
\end{verbatim}
Parameters: None.

Performs the named editor action through the active MR editor command route.

\paragraph{Example macro.}
\begin{verbatim}
$MACRO REFERENCE_DEL_EOL FROM EDIT;
    DEL_EOL;
END_MACRO;
\end{verbatim}
\end{mrmacextension}

\begin{mrmacextension}
\phantomsection
\label{prim:procedure-del-word}
\subsection*{DEL\_WORD}
\index{DEL\_WORD@\texttt{DEL\_WORD}}
\begin{verbatim}
DEL_WORD;
\end{verbatim}
Parameters: None.

Performs the named editor action through the active MR editor command route.

\paragraph{Example macro.}
\begin{verbatim}
$MACRO REFERENCE_DEL_WORD FROM EDIT;
    DEL_WORD;
END_MACRO;
\end{verbatim}
\end{mrmacextension}

\begin{mrmacextension}
\phantomsection
\label{prim:procedure-end-of-block}
\subsection*{END\_OF\_BLOCK}
\index{END\_OF\_BLOCK@\texttt{END\_OF\_BLOCK}}
\begin{verbatim}
END_OF_BLOCK;
\end{verbatim}
Parameters: None.

Performs the named selection or block operation in the active editor context.

\paragraph{Example macro.}
\begin{verbatim}
$MACRO REFERENCE_END_OF_BLOCK FROM EDIT;
    END_OF_BLOCK;
END_MACRO;
\end{verbatim}
\end{mrmacextension}

\begin{mrmacextension}
\phantomsection
\label{prim:procedure-extend-block-by-motion}
\subsection*{EXTEND\_BLOCK\_BY\_MOTION}
\index{EXTEND\_BLOCK\_BY\_MOTION@\texttt{EXTEND\_BLOCK\_BY\_MOTION}}
\begin{verbatim}
EXTEND_BLOCK_BY_MOTION(motionName);
\end{verbatim}
Parameters:
\begin{description}
\item[motionName] String or character positional argument.
\end{description}

Performs the named selection or block operation in the active editor context.

\paragraph{Example macro.}
\begin{verbatim}
$MACRO REFERENCE_EXTEND_BLOCK_BY_MOTION FROM EDIT;
    EXTEND_BLOCK_BY_MOTION('reference');
END_MACRO;
\end{verbatim}
\end{mrmacextension}

\begin{mrmacextension}
\phantomsection
\label{prim:procedure-mark-word-right}
\subsection*{MARK\_WORD\_RIGHT}
\index{MARK\_WORD\_RIGHT@\texttt{MARK\_WORD\_RIGHT}}
\begin{verbatim}
MARK_WORD_RIGHT;
\end{verbatim}
Parameters: None.

Performs the named editor action through the active MR editor command route.

\paragraph{Example macro.}
\begin{verbatim}
$MACRO REFERENCE_MARK_WORD_RIGHT FROM EDIT;
    MARK_WORD_RIGHT;
END_MACRO;
\end{verbatim}
\end{mrmacextension}

\phantomsection
\label{prim:procedure-move-block}
\subsection*{MOVE\_BLOCK}
\index{MOVE\_BLOCK@\texttt{MOVE\_BLOCK}}
\begin{verbatim}
MOVE_BLOCK;
\end{verbatim}
Parameters: None.

Performs the named selection or block operation in the active editor context.

\paragraph{Example macro.}
\begin{verbatim}
$MACRO REFERENCE_MOVE_BLOCK FROM EDIT;
    MOVE_BLOCK;
END_MACRO;
\end{verbatim}

\begin{mrmacextension}
\phantomsection
\label{prim:procedure-paste-block}
\subsection*{PASTE\_BLOCK}
\index{PASTE\_BLOCK@\texttt{PASTE\_BLOCK}}
\begin{verbatim}
PASTE_BLOCK;
\end{verbatim}
Parameters: None.

Performs the named selection or block operation in the active editor context.

\paragraph{Example macro.}
\begin{verbatim}
$MACRO REFERENCE_PASTE_BLOCK FROM EDIT;
    PASTE_BLOCK;
END_MACRO;
\end{verbatim}
\end{mrmacextension}

\begin{mrmacextension}
\phantomsection
\label{prim:procedure-paste-from-clipboard}
\subsection*{PASTE\_FROM\_CLIPBOARD}
\index{PASTE\_FROM\_CLIPBOARD@\texttt{PASTE\_FROM\_CLIPBOARD}}
\begin{verbatim}
PASTE_FROM_CLIPBOARD;
\end{verbatim}
Parameters: None.

Performs the named selection or block operation in the active editor context.

\paragraph{Example macro.}
\begin{verbatim}
$MACRO REFERENCE_PASTE_FROM_CLIPBOARD FROM EDIT;
    PASTE_FROM_CLIPBOARD;
END_MACRO;
\end{verbatim}
\end{mrmacextension}

\phantomsection
\label{prim:procedure-redo}
\subsection*{REDO}
\index{REDO@\texttt{REDO}}
\begin{verbatim}
REDO;
\end{verbatim}
Parameters: None.

Performs the named editor action through the active MR editor command route.

\paragraph{Example macro.}
\begin{verbatim}
$MACRO REFERENCE_REDO FROM EDIT;
    REDO;
END_MACRO;
\end{verbatim}

\begin{mrmacextension}
\phantomsection
\label{prim:procedure-sort-column-block-toggle}
\subsection*{SORT\_COLUMN\_BLOCK\_TOGGLE}
\index{SORT\_COLUMN\_BLOCK\_TOGGLE@\texttt{SORT\_COLUMN\_BLOCK\_TOGGLE}}
\begin{verbatim}
SORT_COLUMN_BLOCK_TOGGLE;
\end{verbatim}
Parameters: None.

Performs the named editor action through the active MR editor command route.

\paragraph{Example macro.}
\begin{verbatim}
$MACRO REFERENCE_SORT_COLUMN_BLOCK_TOGGLE FROM EDIT;
    SORT_COLUMN_BLOCK_TOGGLE;
END_MACRO;
\end{verbatim}
\end{mrmacextension}

\begin{mrmacextension}
\phantomsection
\label{prim:procedure-start-of-block}
\subsection*{START\_OF\_BLOCK}
\index{START\_OF\_BLOCK@\texttt{START\_OF\_BLOCK}}
\begin{verbatim}
START_OF_BLOCK;
\end{verbatim}
Parameters: None.

Performs the named selection or block operation in the active editor context.

\paragraph{Example macro.}
\begin{verbatim}
$MACRO REFERENCE_START_OF_BLOCK FROM EDIT;
    START_OF_BLOCK;
END_MACRO;
\end{verbatim}
\end{mrmacextension}

\phantomsection
\label{prim:procedure-str-block-begin}
\subsection*{STR\_BLOCK\_BEGIN}
\index{STR\_BLOCK\_BEGIN@\texttt{STR\_BLOCK\_BEGIN}}
\begin{verbatim}
STR_BLOCK_BEGIN;
\end{verbatim}
Parameters: None.

Performs the documented MR operation named by this primitive in the active execution context.

\paragraph{Example macro.}
\begin{verbatim}
$MACRO REFERENCE_STR_BLOCK_BEGIN FROM EDIT;
    STR_BLOCK_BEGIN;
END_MACRO;
\end{verbatim}

\phantomsection
\label{prim:procedure-undo}
\subsection*{UNDO}
\index{UNDO@\texttt{UNDO}}
\begin{verbatim}
UNDO;
\end{verbatim}
Parameters: None.

Performs the named editor action through the active MR editor command route.

\paragraph{Example macro.}
\begin{verbatim}
$MACRO REFERENCE_UNDO FROM EDIT;
    UNDO;
END_MACRO;
\end{verbatim}

\begin{mrmacextension}
\phantomsection
\label{prim:procedure-load-block}
\subsection*{LOAD\_BLOCK}
\index{LOAD\_BLOCK@\texttt{LOAD\_BLOCK}}
\begin{verbatim}
LOAD_BLOCK(path);
\end{verbatim}
Parameters:
\begin{description}
\item[path] A string-like file-system path.
\end{description}

Loads a block from path.

\paragraph{Example macro.}
\begin{verbatim}
$MACRO REFERENCE_LOAD_BLOCK FROM EDIT;
    LOAD_BLOCK('reference');
END_MACRO;
\end{verbatim}
\end{mrmacextension}

\begin{mrmacextension}
\phantomsection
\label{prim:procedure-save-block}
\subsection*{SAVE\_BLOCK}
\index{SAVE\_BLOCK@\texttt{SAVE\_BLOCK}}
\begin{verbatim}
SAVE_BLOCK(path);
\end{verbatim}
Parameters:
\begin{description}
\item[path] A string-like file-system path.
\end{description}

Saves the current block to path.

\paragraph{Example macro.}
\begin{verbatim}
$MACRO REFERENCE_SAVE_BLOCK FROM EDIT;
    SAVE_BLOCK('reference');
END_MACRO;
\end{verbatim}
\end{mrmacextension}


\chapter{Files, Search and Windows}

File primitives operate on named files or the active editor file. Search primitives use editor search state. Window primitives select or arrange editor windows and workspaces. Their effects are on editor resources, rather than a local macro value.

\section{Primitive Map}
The map is a local index for this chapter. The final column is a generated page
reference and hyperlink to the full entry, where the source form, parameters,
semantics and a complete example macro appear.

\begin{longtable}{@{}L{0.38\textwidth}L{0.14\textwidth}L{0.27\textwidth}L{0.10\textwidth}@{}}
\toprule
Primitive & Role & Purpose & Reference \\
\midrule
\endhead
\toprule
Primitive & Role & Purpose & Reference \\
\midrule
\endhead
\primitiveoverview{COPY\_FILE}{procedure-copy-file}{procedure}{File, search or window.}
\primitiveoverview{FILE\_ATTR}{intrinsic-file-attr}{intrinsic}{File, search or window.}
\primitiveoverview{FILE\_EXISTS}{intrinsic-file-exists}{intrinsic}{File, search or window.}
\primitiveoverview{FIRST\_FILE}{intrinsic-first-file}{intrinsic}{File, search or window.}
\primitiveoverview{NEXT\_FILE}{intrinsic-next-file}{intrinsic}{File, search or window.}
\primitiveoverview{RENAME\_FILE}{intrinsic-rename-file}{intrinsic}{File, search or window.}
\primitiveoverview{SEARCH\_BWD}{intrinsic-search-bwd}{intrinsic}{File, search or window.}
\primitiveoverview{SEARCH\_FWD}{intrinsic-search-fwd}{intrinsic}{File, search or window.}
\primitiveoverview{SWITCH\_FILE}{intrinsic-switch-file}{intrinsic}{File, search or window.}
\primitiveoverview{CREATE\_WINDOW}{procedure-create-window}{procedure}{File, search or window.}
\primitiveoverview{DELETE\_WINDOW}{procedure-delete-window}{procedure}{File, search or window.}
\primitiveoverview{DESKTOP\_MOVE\_WINDOW\_LEFT}{procedure-desktop-move-window-left}{procedure}{File, search or window.}
\primitiveoverview{DESKTOP\_MOVE\_WINDOW\_RIGHT}{procedure-desktop-move-window-right}{procedure}{File, search or window.}
\primitiveoverview{DESKTOP\_VIEWPORT\_LEFT}{procedure-desktop-viewport-left}{procedure}{File, search or window.}
\primitiveoverview{DESKTOP\_VIEWPORT\_RIGHT}{procedure-desktop-viewport-right}{procedure}{File, search or window.}
\primitiveoverview{ERASE\_WINDOW}{procedure-erase-window}{procedure}{File, search or window.}
\primitiveoverview{FORCE\_SAVE}{procedure-force-save}{procedure}{File, search or window.}
\primitiveoverview{LINK\_WINDOW}{procedure-link-window}{procedure}{File, search or window.}
\primitiveoverview{LIST\_MATCHED\_FILES}{procedure-list-matched-files}{procedure}{File, search or window.}
\primitiveoverview{LOAD\_WORKSPACE}{procedure-load-workspace}{procedure}{File, search or window.}
\primitiveoverview{MODIFY\_WINDOW}{procedure-modify-window}{procedure}{File, search or window.}
\primitiveoverview{MOVE\_WIN\_TO\_NEXT\_DESKTOP}{procedure-move-win-to-next-desktop}{procedure}{File, search or window.}
\primitiveoverview{MOVE\_WIN\_TO\_PREV\_DESKTOP}{procedure-move-win-to-prev-desktop}{procedure}{File, search or window.}
\primitiveoverview{MULTI\_FILE\_SEARCH}{procedure-multi-file-search}{procedure}{File, search or window.}
\primitiveoverview{MULTI\_FILE\_SEARCH\_REPLACE}{procedure-multi-file-search-replace}{procedure}{File, search or window.}
\primitiveoverview{NEW\_SCREEN}{procedure-new-screen}{procedure}{File, search or window.}
\primitiveoverview{NEXT\_SEARCH\_RESULT}{procedure-next-search-result}{procedure}{File, search or window.}
\primitiveoverview{REDRAW}{procedure-redraw}{procedure}{File, search or window.}
\primitiveoverview{REPEAT\_SEARCH}{procedure-repeat-search}{procedure}{File, search or window.}
\primitiveoverview{REVERT\_FILE}{procedure-revert-file}{procedure}{File, search or window.}
\primitiveoverview{SAVE\_ALL}{procedure-save-all}{procedure}{File, search or window.}
\primitiveoverview{SAVE\_WORKSPACE}{procedure-save-workspace}{procedure}{File, search or window.}
\primitiveoverview{SEARCH}{procedure-search}{procedure}{File, search or window.}
\primitiveoverview{SEARCH\_REPLACE}{procedure-search-replace}{procedure}{File, search or window.}
\primitiveoverview{SWITCH\_WINDOW}{procedure-switch-window}{procedure}{File, search or window.}
\primitiveoverview{UNLINK\_WINDOW}{procedure-unlink-window}{procedure}{File, search or window.}
\primitiveoverview{WINDOW\_CASCADE}{procedure-window-cascade}{procedure}{File, search or window.}
\primitiveoverview{WINDOW\_CLOSE}{procedure-window-close}{procedure}{File, search or window.}
\primitiveoverview{WINDOW\_COPY\_BLOCK}{procedure-window-copy-block}{procedure}{File, search or window.}
\primitiveoverview{WINDOW\_LIST}{procedure-window-list}{procedure}{File, search or window.}
\primitiveoverview{WINDOW\_MOVE\_BLOCK}{procedure-window-move-block}{procedure}{File, search or window.}
\primitiveoverview{WINDOW\_NEXT}{procedure-window-next}{procedure}{File, search or window.}
\primitiveoverview{WINDOW\_OPEN}{procedure-window-open}{procedure}{File, search or window.}
\primitiveoverview{WINDOW\_PREVIOUS}{procedure-window-previous}{procedure}{File, search or window.}
\primitiveoverview{WINDOW\_SPLIT\_HORIZONTAL}{procedure-window-split-horizontal}{procedure}{File, search or window.}
\primitiveoverview{WINDOW\_SPLIT\_VERTICAL}{procedure-window-split-vertical}{procedure}{File, search or window.}
\primitiveoverview{WINDOW\_TILE}{procedure-window-tile}{procedure}{File, search or window.}
\primitiveoverview{WINDOW\_ZOOM}{procedure-window-zoom}{procedure}{File, search or window.}
\primitiveoverview{ZOOM}{procedure-zoom}{procedure}{File, search or window.}
\primitiveoverview{DEL\_FILE}{procedure-del-file}{procedure}{File, search or window.}
\primitiveoverview{LOAD\_FILE}{procedure-load-file}{procedure}{File, search or window.}
\primitiveoverview{REPLACE}{procedure-replace}{procedure}{File, search or window.}
\primitiveoverview{SAVE\_FILE}{procedure-save-file}{procedure}{File, search or window.}
\primitiveoverview{SET\_FILE\_ATTR}{procedure-set-file-attr}{procedure}{File, search or window.}
\primitiveoverview{SIZE\_WINDOW}{procedure-size-window}{procedure}{File, search or window.}
\primitiveoverview{WINDOW\_COPY}{procedure-window-copy}{procedure}{File, search or window.}
\primitiveoverview{WINDOW\_MOVE}{procedure-window-move}{procedure}{File, search or window.}
\bottomrule
\end{longtable}

\section{Reference Entries}

\phantomsection
\label{prim:procedure-copy-file}
\subsection*{COPY\_FILE}
\index{COPY\_FILE@\texttt{COPY\_FILE}}
\begin{verbatim}
COPY_FILE(sourcePath, targetPath);
COPY_FILE(sourcePath, targetPath, overwrite);
\end{verbatim}
Parameters:
\begin{description}
\item[sourcePath] Source file path.
\item[targetPath] Destination file path.
\item[overwrite] Non-zero permits overwrite.
\end{description}

Copies a file and returns a status value.

Returns an integer.

\paragraph{Example macro.}
\begin{verbatim}
$MACRO REFERENCE_COPY_FILE FROM EDIT;
    DEF_INT(Result);
    Result := COPY_FILE('reference', 'reference', 0);
END_MACRO;
\end{verbatim}

\phantomsection
\label{prim:intrinsic-file-attr}
\subsection*{FILE\_ATTR}
\index{FILE\_ATTR@\texttt{FILE\_ATTR}}
\begin{verbatim}
FILE_ATTR(path)
\end{verbatim}
Parameters:
\begin{description}
\item[path] A string-like file-system path.
\end{description}

Returns the file attribute mask.

Returns an integer.

\paragraph{Example macro.}
\begin{verbatim}
$MACRO REFERENCE_FILE_ATTR FROM EDIT;
    DEF_INT(Result);
    Result := FILE_ATTR('reference');
END_MACRO;
\end{verbatim}

\phantomsection
\label{prim:intrinsic-file-exists}
\subsection*{FILE\_EXISTS}
\index{FILE\_EXISTS@\texttt{FILE\_EXISTS}}
\begin{verbatim}
FILE_EXISTS(path)
\end{verbatim}
Parameters:
\begin{description}
\item[path] A string-like file-system path.
\end{description}

Returns non-zero when the path exists.

Returns an integer.

\paragraph{Example macro.}
\begin{verbatim}
$MACRO REFERENCE_FILE_EXISTS FROM EDIT;
    DEF_INT(Result);
    Result := FILE_EXISTS('reference');
END_MACRO;
\end{verbatim}

\phantomsection
\label{prim:intrinsic-first-file}
\subsection*{FIRST\_FILE}
\index{FIRST\_FILE@\texttt{FIRST\_FILE}}
\begin{verbatim}
FIRST_FILE(pattern)
\end{verbatim}
Parameters:
\begin{description}
\item[pattern] File-name pattern.
\end{description}

Starts pattern enumeration and returns the first status value.

Returns an integer.

\paragraph{Example macro.}
\begin{verbatim}
$MACRO REFERENCE_FIRST_FILE FROM EDIT;
    DEF_INT(Result);
    Result := FIRST_FILE('reference');
END_MACRO;
\end{verbatim}

\phantomsection
\label{prim:intrinsic-next-file}
\subsection*{NEXT\_FILE}
\index{NEXT\_FILE@\texttt{NEXT\_FILE}}
\begin{verbatim}
NEXT_FILE
\end{verbatim}
Parameters: None.

Advances the active file enumeration.

Returns an integer.

\paragraph{Example macro.}
\begin{verbatim}
$MACRO REFERENCE_NEXT_FILE FROM EDIT;
    DEF_INT(Result);
    Result := NEXT_FILE;
END_MACRO;
\end{verbatim}

\phantomsection
\label{prim:intrinsic-rename-file}
\subsection*{RENAME\_FILE}
\index{RENAME\_FILE@\texttt{RENAME\_FILE}}
\begin{verbatim}
RENAME_FILE(sourcePath, targetPath)
\end{verbatim}
Parameters:
\begin{description}
\item[sourcePath] Source file path.
\item[targetPath] Destination file path.
\end{description}

Renames a file and returns a status value.

Returns an integer.

\paragraph{Example macro.}
\begin{verbatim}
$MACRO REFERENCE_RENAME_FILE FROM EDIT;
    DEF_INT(Result);
    Result := RENAME_FILE('reference', 'reference');
END_MACRO;
\end{verbatim}

\phantomsection
\label{prim:intrinsic-search-bwd}
\subsection*{SEARCH\_BWD}
\index{SEARCH\_BWD@\texttt{SEARCH\_BWD}}
\begin{verbatim}
SEARCH_BWD(searchText, options)
\end{verbatim}
Parameters:
\begin{description}
\item[searchText] Text to search for.
\item[options] Integer search-option mask.
\end{description}

Searches backward and returns a status value.

Returns an integer.

\paragraph{Example macro.}
\begin{verbatim}
$MACRO REFERENCE_SEARCH_BWD FROM EDIT;
    DEF_INT(Result);
    Result := SEARCH_BWD('reference', 0);
END_MACRO;
\end{verbatim}

\phantomsection
\label{prim:intrinsic-search-fwd}
\subsection*{SEARCH\_FWD}
\index{SEARCH\_FWD@\texttt{SEARCH\_FWD}}
\begin{verbatim}
SEARCH_FWD(searchText, options)
\end{verbatim}
Parameters:
\begin{description}
\item[searchText] Text to search for.
\item[options] Integer search-option mask.
\end{description}

Searches forward and returns a status value.

Returns an integer.

\paragraph{Example macro.}
\begin{verbatim}
$MACRO REFERENCE_SEARCH_FWD FROM EDIT;
    DEF_INT(Result);
    Result := SEARCH_FWD('reference', 0);
END_MACRO;
\end{verbatim}

\phantomsection
\label{prim:intrinsic-switch-file}
\subsection*{SWITCH\_FILE}
\index{SWITCH\_FILE@\texttt{SWITCH\_FILE}}
\begin{verbatim}
SWITCH_FILE(fileName)
\end{verbatim}
Parameters:
\begin{description}
\item[fileName] Name of an open file.
\end{description}

Selects an already open file and returns a status value.

Returns an integer.

\paragraph{Example macro.}
\begin{verbatim}
$MACRO REFERENCE_SWITCH_FILE FROM EDIT;
    DEF_INT(Result);
    Result := SWITCH_FILE('reference');
END_MACRO;
\end{verbatim}

\phantomsection
\label{prim:procedure-create-window}
\subsection*{CREATE\_WINDOW}
\index{CREATE\_WINDOW@\texttt{CREATE\_WINDOW}}
\begin{verbatim}
CREATE_WINDOW;
\end{verbatim}
Parameters: None.

Performs the named MR window, screen or desktop operation in the active editor context.

\paragraph{Example macro.}
\begin{verbatim}
$MACRO REFERENCE_CREATE_WINDOW FROM EDIT;
    CREATE_WINDOW;
END_MACRO;
\end{verbatim}

\phantomsection
\label{prim:procedure-delete-window}
\subsection*{DELETE\_WINDOW}
\index{DELETE\_WINDOW@\texttt{DELETE\_WINDOW}}
\begin{verbatim}
DELETE_WINDOW;
\end{verbatim}
Parameters: None.

Performs the named MR window, screen or desktop operation in the active editor context.

\paragraph{Example macro.}
\begin{verbatim}
$MACRO REFERENCE_DELETE_WINDOW FROM EDIT;
    DELETE_WINDOW;
END_MACRO;
\end{verbatim}

\begin{mrmacextension}
\phantomsection
\label{prim:procedure-desktop-move-window-left}
\subsection*{DESKTOP\_MOVE\_WINDOW\_LEFT}
\index{DESKTOP\_MOVE\_WINDOW\_LEFT@\texttt{DESKTOP\_MOVE\_WINDOW\_LEFT}}
\begin{verbatim}
DESKTOP_MOVE_WINDOW_LEFT;
\end{verbatim}
Parameters: None.

Performs the named MR window, screen or desktop operation in the active editor context.

\paragraph{Example macro.}
\begin{verbatim}
$MACRO REFERENCE_DESKTOP_MOVE_WINDOW_LEFT FROM EDIT;
    DESKTOP_MOVE_WINDOW_LEFT;
END_MACRO;
\end{verbatim}
\end{mrmacextension}

\begin{mrmacextension}
\phantomsection
\label{prim:procedure-desktop-move-window-right}
\subsection*{DESKTOP\_MOVE\_WINDOW\_RIGHT}
\index{DESKTOP\_MOVE\_WINDOW\_RIGHT@\texttt{DESKTOP\_MOVE\_WINDOW\_RIGHT}}
\begin{verbatim}
DESKTOP_MOVE_WINDOW_RIGHT;
\end{verbatim}
Parameters: None.

Performs the named MR window, screen or desktop operation in the active editor context.

\paragraph{Example macro.}
\begin{verbatim}
$MACRO REFERENCE_DESKTOP_MOVE_WINDOW_RIGHT FROM EDIT;
    DESKTOP_MOVE_WINDOW_RIGHT;
END_MACRO;
\end{verbatim}
\end{mrmacextension}

\begin{mrmacextension}
\phantomsection
\label{prim:procedure-desktop-viewport-left}
\subsection*{DESKTOP\_VIEWPORT\_LEFT}
\index{DESKTOP\_VIEWPORT\_LEFT@\texttt{DESKTOP\_VIEWPORT\_LEFT}}
\begin{verbatim}
DESKTOP_VIEWPORT_LEFT;
\end{verbatim}
Parameters: None.

Performs the named MR window, screen or desktop operation in the active editor context.

\paragraph{Example macro.}
\begin{verbatim}
$MACRO REFERENCE_DESKTOP_VIEWPORT_LEFT FROM EDIT;
    DESKTOP_VIEWPORT_LEFT;
END_MACRO;
\end{verbatim}
\end{mrmacextension}

\begin{mrmacextension}
\phantomsection
\label{prim:procedure-desktop-viewport-right}
\subsection*{DESKTOP\_VIEWPORT\_RIGHT}
\index{DESKTOP\_VIEWPORT\_RIGHT@\texttt{DESKTOP\_VIEWPORT\_RIGHT}}
\begin{verbatim}
DESKTOP_VIEWPORT_RIGHT;
\end{verbatim}
Parameters: None.

Performs the named MR window, screen or desktop operation in the active editor context.

\paragraph{Example macro.}
\begin{verbatim}
$MACRO REFERENCE_DESKTOP_VIEWPORT_RIGHT FROM EDIT;
    DESKTOP_VIEWPORT_RIGHT;
END_MACRO;
\end{verbatim}
\end{mrmacextension}

\phantomsection
\label{prim:procedure-erase-window}
\subsection*{ERASE\_WINDOW}
\index{ERASE\_WINDOW@\texttt{ERASE\_WINDOW}}
\begin{verbatim}
ERASE_WINDOW;
\end{verbatim}
Parameters: None.

Performs the named MR window, screen or desktop operation in the active editor context.

\paragraph{Example macro.}
\begin{verbatim}
$MACRO REFERENCE_ERASE_WINDOW FROM EDIT;
    ERASE_WINDOW;
END_MACRO;
\end{verbatim}

\begin{mrmacextension}
\phantomsection
\label{prim:procedure-force-save}
\subsection*{FORCE\_SAVE}
\index{FORCE\_SAVE@\texttt{FORCE\_SAVE}}
\begin{verbatim}
FORCE_SAVE;
\end{verbatim}
Parameters: None.

Performs the named editor action through the active MR editor command route.

\paragraph{Example macro.}
\begin{verbatim}
$MACRO REFERENCE_FORCE_SAVE FROM EDIT;
    FORCE_SAVE;
END_MACRO;
\end{verbatim}
\end{mrmacextension}

\phantomsection
\label{prim:procedure-link-window}
\subsection*{LINK\_WINDOW}
\index{LINK\_WINDOW@\texttt{LINK\_WINDOW}}
\begin{verbatim}
LINK_WINDOW;
\end{verbatim}
Parameters: None.

Performs the named MR window, screen or desktop operation in the active editor context.

\paragraph{Example macro.}
\begin{verbatim}
$MACRO REFERENCE_LINK_WINDOW FROM EDIT;
    LINK_WINDOW;
END_MACRO;
\end{verbatim}

\begin{mrmacextension}
\phantomsection
\label{prim:procedure-list-matched-files}
\subsection*{LIST\_MATCHED\_FILES}
\index{LIST\_MATCHED\_FILES@\texttt{LIST\_MATCHED\_FILES}}
\begin{verbatim}
LIST_MATCHED_FILES;
\end{verbatim}
Parameters: None.

Invokes the named MR search operation in the active editor context.

\paragraph{Example macro.}
\begin{verbatim}
$MACRO REFERENCE_LIST_MATCHED_FILES FROM EDIT;
    LIST_MATCHED_FILES;
END_MACRO;
\end{verbatim}
\end{mrmacextension}

\begin{mrmacextension}
\phantomsection
\label{prim:procedure-load-workspace}
\subsection*{LOAD\_WORKSPACE}
\index{LOAD\_WORKSPACE@\texttt{LOAD\_WORKSPACE}}
\begin{verbatim}
LOAD_WORKSPACE;
\end{verbatim}
Parameters: None.

Performs the documented MR operation named by this primitive in the active execution context.

\paragraph{Example macro.}
\begin{verbatim}
$MACRO REFERENCE_LOAD_WORKSPACE FROM EDIT;
    LOAD_WORKSPACE;
END_MACRO;
\end{verbatim}
\end{mrmacextension}

\phantomsection
\label{prim:procedure-modify-window}
\subsection*{MODIFY\_WINDOW}
\index{MODIFY\_WINDOW@\texttt{MODIFY\_WINDOW}}
\begin{verbatim}
MODIFY_WINDOW;
\end{verbatim}
Parameters: None.

Performs the named MR window, screen or desktop operation in the active editor context.

\paragraph{Example macro.}
\begin{verbatim}
$MACRO REFERENCE_MODIFY_WINDOW FROM EDIT;
    MODIFY_WINDOW;
END_MACRO;
\end{verbatim}

\begin{mrmacextension}
\phantomsection
\label{prim:procedure-move-win-to-next-desktop}
\subsection*{MOVE\_WIN\_TO\_NEXT\_DESKTOP}
\index{MOVE\_WIN\_TO\_NEXT\_DESKTOP@\texttt{MOVE\_WIN\_TO\_NEXT\_DESKTOP}}
\begin{verbatim}
MOVE_WIN_TO_NEXT_DESKTOP;
\end{verbatim}
Parameters: None.

Performs the named MR window, screen or desktop operation in the active editor context.

\paragraph{Example macro.}
\begin{verbatim}
$MACRO REFERENCE_MOVE_WIN_TO_NEXT_DESKTOP FROM EDIT;
    MOVE_WIN_TO_NEXT_DESKTOP;
END_MACRO;
\end{verbatim}
\end{mrmacextension}

\begin{mrmacextension}
\phantomsection
\label{prim:procedure-move-win-to-prev-desktop}
\subsection*{MOVE\_WIN\_TO\_PREV\_DESKTOP}
\index{MOVE\_WIN\_TO\_PREV\_DESKTOP@\texttt{MOVE\_WIN\_TO\_PREV\_DESKTOP}}
\begin{verbatim}
MOVE_WIN_TO_PREV_DESKTOP;
\end{verbatim}
Parameters: None.

Performs the named MR window, screen or desktop operation in the active editor context.

\paragraph{Example macro.}
\begin{verbatim}
$MACRO REFERENCE_MOVE_WIN_TO_PREV_DESKTOP FROM EDIT;
    MOVE_WIN_TO_PREV_DESKTOP;
END_MACRO;
\end{verbatim}
\end{mrmacextension}

\begin{mrmacextension}
\phantomsection
\label{prim:procedure-multi-file-search}
\subsection*{MULTI\_FILE\_SEARCH}
\index{MULTI\_FILE\_SEARCH@\texttt{MULTI\_FILE\_SEARCH}}
\begin{verbatim}
MULTI_FILE_SEARCH;
\end{verbatim}
Parameters: None.

Invokes the named MR search operation in the active editor context.

\paragraph{Example macro.}
\begin{verbatim}
$MACRO REFERENCE_MULTI_FILE_SEARCH FROM EDIT;
    MULTI_FILE_SEARCH;
END_MACRO;
\end{verbatim}
\end{mrmacextension}

\begin{mrmacextension}
\phantomsection
\label{prim:procedure-multi-file-search-replace}
\subsection*{MULTI\_FILE\_SEARCH\_REPLACE}
\index{MULTI\_FILE\_SEARCH\_REPLACE@\texttt{MULTI\_FILE\_SEARCH\_REPLACE}}
\begin{verbatim}
MULTI_FILE_SEARCH_REPLACE;
\end{verbatim}
Parameters: None.

Invokes the named MR search operation in the active editor context.

\paragraph{Example macro.}
\begin{verbatim}
$MACRO REFERENCE_MULTI_FILE_SEARCH_REPLACE FROM EDIT;
    MULTI_FILE_SEARCH_REPLACE;
END_MACRO;
\end{verbatim}
\end{mrmacextension}

\phantomsection
\label{prim:procedure-new-screen}
\subsection*{NEW\_SCREEN}
\index{NEW\_SCREEN@\texttt{NEW\_SCREEN}}
\begin{verbatim}
NEW_SCREEN;
\end{verbatim}
Parameters: None.

Performs the named MR window, screen or desktop operation in the active editor context.

\paragraph{Example macro.}
\begin{verbatim}
$MACRO REFERENCE_NEW_SCREEN FROM EDIT;
    NEW_SCREEN;
END_MACRO;
\end{verbatim}

\begin{mrmacextension}
\phantomsection
\label{prim:procedure-next-search-result}
\subsection*{NEXT\_SEARCH\_RESULT}
\index{NEXT\_SEARCH\_RESULT@\texttt{NEXT\_SEARCH\_RESULT}}
\begin{verbatim}
NEXT_SEARCH_RESULT;
\end{verbatim}
Parameters: None.

Invokes the named MR search operation in the active editor context.

\paragraph{Example macro.}
\begin{verbatim}
$MACRO REFERENCE_NEXT_SEARCH_RESULT FROM EDIT;
    NEXT_SEARCH_RESULT;
END_MACRO;
\end{verbatim}
\end{mrmacextension}

\phantomsection
\label{prim:procedure-redraw}
\subsection*{REDRAW}
\index{REDRAW@\texttt{REDRAW}}
\begin{verbatim}
REDRAW;
\end{verbatim}
Parameters: None.

Performs the named MR window, screen or desktop operation in the active editor context.

\paragraph{Example macro.}
\begin{verbatim}
$MACRO REFERENCE_REDRAW FROM EDIT;
    REDRAW;
END_MACRO;
\end{verbatim}

\begin{mrmacextension}
\phantomsection
\label{prim:procedure-repeat-search}
\subsection*{REPEAT\_SEARCH}
\index{REPEAT\_SEARCH@\texttt{REPEAT\_SEARCH}}
\begin{verbatim}
REPEAT_SEARCH;
\end{verbatim}
Parameters: None.

Invokes the named MR search operation in the active editor context.

\paragraph{Example macro.}
\begin{verbatim}
$MACRO REFERENCE_REPEAT_SEARCH FROM EDIT;
    REPEAT_SEARCH;
END_MACRO;
\end{verbatim}
\end{mrmacextension}

\begin{mrmacextension}
\phantomsection
\label{prim:procedure-revert-file}
\subsection*{REVERT\_FILE}
\index{REVERT\_FILE@\texttt{REVERT\_FILE}}
\begin{verbatim}
REVERT_FILE;
\end{verbatim}
Parameters: None.

Performs the named editor action through the active MR editor command route.

\paragraph{Example macro.}
\begin{verbatim}
$MACRO REFERENCE_REVERT_FILE FROM EDIT;
    REVERT_FILE;
END_MACRO;
\end{verbatim}
\end{mrmacextension}

\begin{mrmacextension}
\phantomsection
\label{prim:procedure-save-all}
\subsection*{SAVE\_ALL}
\index{SAVE\_ALL@\texttt{SAVE\_ALL}}
\begin{verbatim}
SAVE_ALL;
\end{verbatim}
Parameters: None.

Performs the named editor action through the active MR editor command route.

\paragraph{Example macro.}
\begin{verbatim}
$MACRO REFERENCE_SAVE_ALL FROM EDIT;
    SAVE_ALL;
END_MACRO;
\end{verbatim}
\end{mrmacextension}

\begin{mrmacextension}
\phantomsection
\label{prim:procedure-save-workspace}
\subsection*{SAVE\_WORKSPACE}
\index{SAVE\_WORKSPACE@\texttt{SAVE\_WORKSPACE}}
\begin{verbatim}
SAVE_WORKSPACE;
\end{verbatim}
Parameters: None.

Performs the documented MR operation named by this primitive in the active execution context.

\paragraph{Example macro.}
\begin{verbatim}
$MACRO REFERENCE_SAVE_WORKSPACE FROM EDIT;
    SAVE_WORKSPACE;
END_MACRO;
\end{verbatim}
\end{mrmacextension}

\begin{mrmacextension}
\phantomsection
\label{prim:procedure-search}
\subsection*{SEARCH}
\index{SEARCH@\texttt{SEARCH}}
\begin{verbatim}
SEARCH;
\end{verbatim}
Parameters: None.

Invokes the named MR search operation in the active editor context.

\paragraph{Example macro.}
\begin{verbatim}
$MACRO REFERENCE_SEARCH FROM EDIT;
    SEARCH;
END_MACRO;
\end{verbatim}
\end{mrmacextension}

\begin{mrmacextension}
\phantomsection
\label{prim:procedure-search-replace}
\subsection*{SEARCH\_REPLACE}
\index{SEARCH\_REPLACE@\texttt{SEARCH\_REPLACE}}
\begin{verbatim}
SEARCH_REPLACE;
\end{verbatim}
Parameters: None.

Invokes the named MR search operation in the active editor context.

\paragraph{Example macro.}
\begin{verbatim}
$MACRO REFERENCE_SEARCH_REPLACE FROM EDIT;
    SEARCH_REPLACE;
END_MACRO;
\end{verbatim}
\end{mrmacextension}

\phantomsection
\label{prim:procedure-switch-window}
\subsection*{SWITCH\_WINDOW}
\index{SWITCH\_WINDOW@\texttt{SWITCH\_WINDOW}}
\begin{verbatim}
SWITCH_WINDOW(windowIndex);
\end{verbatim}
Parameters:
\begin{description}
\item[windowIndex] Integer positional argument.
\end{description}

Activates windowIndex.

\paragraph{Example macro.}
\begin{verbatim}
$MACRO REFERENCE_SWITCH_WINDOW FROM EDIT;
    SWITCH_WINDOW(0);
END_MACRO;
\end{verbatim}

\phantomsection
\label{prim:procedure-unlink-window}
\subsection*{UNLINK\_WINDOW}
\index{UNLINK\_WINDOW@\texttt{UNLINK\_WINDOW}}
\begin{verbatim}
UNLINK_WINDOW;
\end{verbatim}
Parameters: None.

Performs the named MR window, screen or desktop operation in the active editor context.

\paragraph{Example macro.}
\begin{verbatim}
$MACRO REFERENCE_UNLINK_WINDOW FROM EDIT;
    UNLINK_WINDOW;
END_MACRO;
\end{verbatim}

\begin{mrmacextension}
\phantomsection
\label{prim:procedure-window-cascade}
\subsection*{WINDOW\_CASCADE}
\index{WINDOW\_CASCADE@\texttt{WINDOW\_CASCADE}}
\begin{verbatim}
WINDOW_CASCADE;
\end{verbatim}
Parameters: None.

Performs the named MR window, screen or desktop operation in the active editor context.

\paragraph{Example macro.}
\begin{verbatim}
$MACRO REFERENCE_WINDOW_CASCADE FROM EDIT;
    WINDOW_CASCADE;
END_MACRO;
\end{verbatim}
\end{mrmacextension}

\begin{mrmacextension}
\phantomsection
\label{prim:procedure-window-close}
\subsection*{WINDOW\_CLOSE}
\index{WINDOW\_CLOSE@\texttt{WINDOW\_CLOSE}}
\begin{verbatim}
WINDOW_CLOSE;
\end{verbatim}
Parameters: None.

Performs the named MR window, screen or desktop operation in the active editor context.

\paragraph{Example macro.}
\begin{verbatim}
$MACRO REFERENCE_WINDOW_CLOSE FROM EDIT;
    WINDOW_CLOSE;
END_MACRO;
\end{verbatim}
\end{mrmacextension}

\begin{mrmacextension}
\phantomsection
\label{prim:procedure-window-copy-block}
\subsection*{WINDOW\_COPY\_BLOCK}
\index{WINDOW\_COPY\_BLOCK@\texttt{WINDOW\_COPY\_BLOCK}}
\begin{verbatim}
WINDOW_COPY_BLOCK;
\end{verbatim}
Parameters: None.

Performs the named MR window, screen or desktop operation in the active editor context.

\paragraph{Example macro.}
\begin{verbatim}
$MACRO REFERENCE_WINDOW_COPY_BLOCK FROM EDIT;
    WINDOW_COPY_BLOCK;
END_MACRO;
\end{verbatim}
\end{mrmacextension}

\begin{mrmacextension}
\phantomsection
\label{prim:procedure-window-list}
\subsection*{WINDOW\_LIST}
\index{WINDOW\_LIST@\texttt{WINDOW\_LIST}}
\begin{verbatim}
WINDOW_LIST;
\end{verbatim}
Parameters: None.

Performs the named MR window, screen or desktop operation in the active editor context.

\paragraph{Example macro.}
\begin{verbatim}
$MACRO REFERENCE_WINDOW_LIST FROM EDIT;
    WINDOW_LIST;
END_MACRO;
\end{verbatim}
\end{mrmacextension}

\begin{mrmacextension}
\phantomsection
\label{prim:procedure-window-move-block}
\subsection*{WINDOW\_MOVE\_BLOCK}
\index{WINDOW\_MOVE\_BLOCK@\texttt{WINDOW\_MOVE\_BLOCK}}
\begin{verbatim}
WINDOW_MOVE_BLOCK;
\end{verbatim}
Parameters: None.

Performs the named MR window, screen or desktop operation in the active editor context.

\paragraph{Example macro.}
\begin{verbatim}
$MACRO REFERENCE_WINDOW_MOVE_BLOCK FROM EDIT;
    WINDOW_MOVE_BLOCK;
END_MACRO;
\end{verbatim}
\end{mrmacextension}

\begin{mrmacextension}
\phantomsection
\label{prim:procedure-window-next}
\subsection*{WINDOW\_NEXT}
\index{WINDOW\_NEXT@\texttt{WINDOW\_NEXT}}
\begin{verbatim}
WINDOW_NEXT;
\end{verbatim}
Parameters: None.

Performs the named MR window, screen or desktop operation in the active editor context.

\paragraph{Example macro.}
\begin{verbatim}
$MACRO REFERENCE_WINDOW_NEXT FROM EDIT;
    WINDOW_NEXT;
END_MACRO;
\end{verbatim}
\end{mrmacextension}

\begin{mrmacextension}
\phantomsection
\label{prim:procedure-window-open}
\subsection*{WINDOW\_OPEN}
\index{WINDOW\_OPEN@\texttt{WINDOW\_OPEN}}
\begin{verbatim}
WINDOW_OPEN;
\end{verbatim}
Parameters: None.

Performs the named MR window, screen or desktop operation in the active editor context.

\paragraph{Example macro.}
\begin{verbatim}
$MACRO REFERENCE_WINDOW_OPEN FROM EDIT;
    WINDOW_OPEN;
END_MACRO;
\end{verbatim}
\end{mrmacextension}

\begin{mrmacextension}
\phantomsection
\label{prim:procedure-window-previous}
\subsection*{WINDOW\_PREVIOUS}
\index{WINDOW\_PREVIOUS@\texttt{WINDOW\_PREVIOUS}}
\begin{verbatim}
WINDOW_PREVIOUS;
\end{verbatim}
Parameters: None.

Performs the named MR window, screen or desktop operation in the active editor context.

\paragraph{Example macro.}
\begin{verbatim}
$MACRO REFERENCE_WINDOW_PREVIOUS FROM EDIT;
    WINDOW_PREVIOUS;
END_MACRO;
\end{verbatim}
\end{mrmacextension}

\begin{mrmacextension}
\phantomsection
\label{prim:procedure-window-split-horizontal}
\subsection*{WINDOW\_SPLIT\_HORIZONTAL}
\index{WINDOW\_SPLIT\_HORIZONTAL@\texttt{WINDOW\_SPLIT\_HORIZONTAL}}
\begin{verbatim}
WINDOW_SPLIT_HORIZONTAL;
\end{verbatim}
Parameters: None.

Performs the named MR window, screen or desktop operation in the active editor context.

\paragraph{Example macro.}
\begin{verbatim}
$MACRO REFERENCE_WINDOW_SPLIT_HORIZONTAL FROM EDIT;
    WINDOW_SPLIT_HORIZONTAL;
END_MACRO;
\end{verbatim}
\end{mrmacextension}

\begin{mrmacextension}
\phantomsection
\label{prim:procedure-window-split-vertical}
\subsection*{WINDOW\_SPLIT\_VERTICAL}
\index{WINDOW\_SPLIT\_VERTICAL@\texttt{WINDOW\_SPLIT\_VERTICAL}}
\begin{verbatim}
WINDOW_SPLIT_VERTICAL;
\end{verbatim}
Parameters: None.

Performs the named MR window, screen or desktop operation in the active editor context.

\paragraph{Example macro.}
\begin{verbatim}
$MACRO REFERENCE_WINDOW_SPLIT_VERTICAL FROM EDIT;
    WINDOW_SPLIT_VERTICAL;
END_MACRO;
\end{verbatim}
\end{mrmacextension}

\begin{mrmacextension}
\phantomsection
\label{prim:procedure-window-tile}
\subsection*{WINDOW\_TILE}
\index{WINDOW\_TILE@\texttt{WINDOW\_TILE}}
\begin{verbatim}
WINDOW_TILE;
\end{verbatim}
Parameters: None.

Performs the named MR window, screen or desktop operation in the active editor context.

\paragraph{Example macro.}
\begin{verbatim}
$MACRO REFERENCE_WINDOW_TILE FROM EDIT;
    WINDOW_TILE;
END_MACRO;
\end{verbatim}
\end{mrmacextension}

\begin{mrmacextension}
\phantomsection
\label{prim:procedure-window-zoom}
\subsection*{WINDOW\_ZOOM}
\index{WINDOW\_ZOOM@\texttt{WINDOW\_ZOOM}}
\begin{verbatim}
WINDOW_ZOOM;
\end{verbatim}
Parameters: None.

Performs the named MR window, screen or desktop operation in the active editor context.

\paragraph{Example macro.}
\begin{verbatim}
$MACRO REFERENCE_WINDOW_ZOOM FROM EDIT;
    WINDOW_ZOOM;
END_MACRO;
\end{verbatim}
\end{mrmacextension}

\phantomsection
\label{prim:procedure-zoom}
\subsection*{ZOOM}
\index{ZOOM@\texttt{ZOOM}}
\begin{verbatim}
ZOOM;
\end{verbatim}
Parameters: None.

Performs the named MR window, screen or desktop operation in the active editor context.

\paragraph{Example macro.}
\begin{verbatim}
$MACRO REFERENCE_ZOOM FROM EDIT;
    ZOOM;
END_MACRO;
\end{verbatim}

\phantomsection
\label{prim:procedure-del-file}
\subsection*{DEL\_FILE}
\index{DEL\_FILE@\texttt{DEL\_FILE}}
\begin{verbatim}
DEL_FILE(path);
\end{verbatim}
Parameters:
\begin{description}
\item[path] A string-like file-system path.
\end{description}

Deletes path through the editor file route.

\paragraph{Example macro.}
\begin{verbatim}
$MACRO REFERENCE_DEL_FILE FROM EDIT;
    DEL_FILE('reference');
END_MACRO;
\end{verbatim}

\phantomsection
\label{prim:procedure-load-file}
\subsection*{LOAD\_FILE}
\index{LOAD\_FILE@\texttt{LOAD\_FILE}}
\begin{verbatim}
LOAD_FILE(path);
\end{verbatim}
Parameters:
\begin{description}
\item[path] A string-like file-system path.
\end{description}

Loads path into the editor.

\paragraph{Example macro.}
\begin{verbatim}
$MACRO REFERENCE_LOAD_FILE FROM EDIT;
    LOAD_FILE('reference');
END_MACRO;
\end{verbatim}

\phantomsection
\label{prim:procedure-replace}
\subsection*{REPLACE}
\index{REPLACE@\texttt{REPLACE}}
\begin{verbatim}
REPLACE(replacement);
\end{verbatim}
Parameters:
\begin{description}
\item[replacement] Replacement text.
\end{description}

Performs replacement using the current search state.

\paragraph{Example macro.}
\begin{verbatim}
$MACRO REFERENCE_REPLACE FROM EDIT;
    REPLACE('reference');
END_MACRO;
\end{verbatim}

\phantomsection
\label{prim:procedure-save-file}
\subsection*{SAVE\_FILE}
\index{SAVE\_FILE@\texttt{SAVE\_FILE}}
\begin{verbatim}
SAVE_FILE;
\end{verbatim}
Parameters: None.

Saves the active editor file.

\paragraph{Example macro.}
\begin{verbatim}
$MACRO REFERENCE_SAVE_FILE FROM EDIT;
    SAVE_FILE;
END_MACRO;
\end{verbatim}

\phantomsection
\label{prim:procedure-set-file-attr}
\subsection*{SET\_FILE\_ATTR}
\index{SET\_FILE\_ATTR@\texttt{SET\_FILE\_ATTR}}
\begin{verbatim}
SET_FILE_ATTR(path, attributes);
\end{verbatim}
Parameters:
\begin{description}
\item[path] A string-like file-system path.
\item[attributes] Integer file attribute mask.
\end{description}

Applies the requested file attribute mask.

\paragraph{Example macro.}
\begin{verbatim}
$MACRO REFERENCE_SET_FILE_ATTR FROM EDIT;
    SET_FILE_ATTR('reference', 0);
END_MACRO;
\end{verbatim}

\phantomsection
\label{prim:procedure-size-window}
\subsection*{SIZE\_WINDOW}
\index{SIZE\_WINDOW@\texttt{SIZE\_WINDOW}}
\begin{verbatim}
SIZE_WINDOW(left, top, right, bottom);
\end{verbatim}
Parameters:
\begin{description}
\item[left] Left screen column.
\item[top] Top screen row.
\item[right] Right screen column.
\item[bottom] Bottom screen row.
\end{description}

Sets current window bounds.

\paragraph{Example macro.}
\begin{verbatim}
$MACRO REFERENCE_SIZE_WINDOW FROM EDIT;
    SIZE_WINDOW(0, 0, 0, 0);
END_MACRO;
\end{verbatim}

\phantomsection
\label{prim:procedure-window-copy}
\subsection*{WINDOW\_COPY}
\index{WINDOW\_COPY@\texttt{WINDOW\_COPY}}
\begin{verbatim}
WINDOW_COPY(windowIndex);
\end{verbatim}
Parameters:
\begin{description}
\item[windowIndex] Integer positional argument.
\end{description}

Copies the current window to windowIndex.

\paragraph{Example macro.}
\begin{verbatim}
$MACRO REFERENCE_WINDOW_COPY FROM EDIT;
    WINDOW_COPY(0);
END_MACRO;
\end{verbatim}

\phantomsection
\label{prim:procedure-window-move}
\subsection*{WINDOW\_MOVE}
\index{WINDOW\_MOVE@\texttt{WINDOW\_MOVE}}
\begin{verbatim}
WINDOW_MOVE(windowIndex);
\end{verbatim}
Parameters:
\begin{description}
\item[windowIndex] Integer positional argument.
\end{description}

Moves the current window to windowIndex.

\paragraph{Example macro.}
\begin{verbatim}
$MACRO REFERENCE_WINDOW_MOVE FROM EDIT;
    WINDOW_MOVE(0);
END_MACRO;
\end{verbatim}


\part{Interaction and UI}

\chapter{Screen, Input and Key Dispatch}

Interaction primitives request or stage user-facing work through the VM and UI boundary. They do not give macro source a drawing context, event object or direct keymap store. Use this chapter for screen output, input, feedback and key bindings.

\section{Primitive Map}
The map is a local index for this chapter. The final column is a generated page
reference and hyperlink to the full entry, where the source form, parameters,
semantics and a complete example macro appear.

\begin{longtable}{@{}L{0.38\textwidth}L{0.14\textwidth}L{0.27\textwidth}L{0.10\textwidth}@{}}
\toprule
Primitive & Role & Purpose & Reference \\
\midrule
\endhead
\toprule
Primitive & Role & Purpose & Reference \\
\midrule
\endhead
\primitiveoverview{BEEP}{procedure-beep}{procedure}{User interaction.}
\primitiveoverview{BRAIN}{procedure-brain}{procedure}{User interaction.}
\primitiveoverview{CLEAR\_SCREEN}{procedure-clear-screen}{procedure}{User interaction.}
\primitiveoverview{CLR\_LINE}{procedure-clr-line}{procedure}{User interaction.}
\primitiveoverview{CMD\_TO\_KEY}{procedure-cmd-to-key}{procedure}{User interaction.}
\primitiveoverview{FLABEL}{procedure-flabel}{procedure}{User interaction.}
\primitiveoverview{GOTOXY}{procedure-gotoxy}{procedure}{User interaction.}
\primitiveoverview{KEY\_IN}{procedure-key-in}{procedure}{User interaction.}
\primitiveoverview{KEY\_RECORD}{procedure-key-record}{procedure}{User interaction.}
\primitiveoverview{KILL\_BOX}{procedure-kill-box}{procedure}{User interaction.}
\primitiveoverview{MACRO\_TO\_KEY}{procedure-macro-to-key}{procedure}{User interaction.}
\primitiveoverview{MAKE\_MESSAGE}{procedure-make-message}{procedure}{User interaction.}
\primitiveoverview{MARQUEE}{procedure-marquee}{procedure}{User interaction.}
\primitiveoverview{MARQUEE\_ERROR}{procedure-marquee-error}{procedure}{User interaction.}
\primitiveoverview{MARQUEE\_WARNING}{procedure-marquee-warning}{procedure}{User interaction.}
\primitiveoverview{PASS\_KEY}{procedure-pass-key}{procedure}{User interaction.}
\primitiveoverview{PLAY\_KEY\_MACRO}{procedure-play-key-macro}{procedure}{User interaction.}
\primitiveoverview{POP\_LABELS}{procedure-pop-labels}{procedure}{User interaction.}
\primitiveoverview{PUSH\_KEY}{procedure-push-key}{procedure}{User interaction.}
\primitiveoverview{PUSH\_LABELS}{procedure-push-labels}{procedure}{User interaction.}
\primitiveoverview{PUT\_BOX}{procedure-put-box}{procedure}{User interaction.}
\primitiveoverview{PUT\_COL\_NUM}{procedure-put-col-num}{procedure}{User interaction.}
\primitiveoverview{PUT\_LINE\_NUM}{procedure-put-line-num}{procedure}{User interaction.}
\primitiveoverview{READ\_KEY}{procedure-read-key}{procedure}{User interaction.}
\primitiveoverview{REGISTER\_MENU\_ITEM}{procedure-register-menu-item}{procedure}{User interaction.}
\primitiveoverview{REMOVE\_MENU\_ITEM}{procedure-remove-menu-item}{procedure}{User interaction.}
\primitiveoverview{SCROLL\_BOX\_DN}{procedure-scroll-box-dn}{procedure}{User interaction.}
\primitiveoverview{SCROLL\_BOX\_UP}{procedure-scroll-box-up}{procedure}{User interaction.}
\primitiveoverview{SET\_CLIPBOARD\_TEXT}{procedure-set-clipboard-text}{procedure}{User interaction.}
\primitiveoverview{UNASSIGN\_ALL\_KEYS}{procedure-unassign-all-keys}{procedure}{User interaction.}
\primitiveoverview{UNASSIGN\_KEY}{procedure-unassign-key}{procedure}{User interaction.}
\primitiveoverview{WORKING}{procedure-working}{procedure}{User interaction.}
\primitiveoverview{WRITE}{procedure-write}{procedure}{User interaction.}
\bottomrule
\end{longtable}

\section{Reference Entries}

\begin{mrmacextension}
\phantomsection
\label{prim:procedure-beep}
\subsection*{BEEP}
\index{BEEP@\texttt{BEEP}}
\begin{verbatim}
BEEP;
\end{verbatim}
Parameters: None.

Requests the standard editor beep.

\paragraph{Example macro.}
\begin{verbatim}
$MACRO REFERENCE_BEEP FROM EDIT;
    BEEP;
END_MACRO;
\end{verbatim}
\end{mrmacextension}

\begin{mrmacextension}
\phantomsection
\label{prim:procedure-brain}
\subsection*{BRAIN}
\index{BRAIN@\texttt{BRAIN}}
\begin{verbatim}
BRAIN(enabled);
\end{verbatim}
Parameters:
\begin{description}
\item[enabled] Integer positional argument.
\end{description}

Enables or disables the typed busy indicator.

\paragraph{Example macro.}
\begin{verbatim}
$MACRO REFERENCE_BRAIN FROM EDIT;
    BRAIN(0);
END_MACRO;
\end{verbatim}
\end{mrmacextension}

\phantomsection
\label{prim:procedure-clear-screen}
\subsection*{CLEAR\_SCREEN}
\index{CLEAR\_SCREEN@\texttt{CLEAR\_SCREEN}}
\begin{verbatim}
CLEAR_SCREEN();
CLEAR_SCREEN(attribute);
\end{verbatim}
Parameters:
\begin{description}
\item[attribute] Screen attribute code.
\end{description}

Clears the staged screen using the optional attribute.

\paragraph{Example macro.}
\begin{verbatim}
$MACRO REFERENCE_CLEAR_SCREEN FROM EDIT;
    CLEAR_SCREEN(0);
END_MACRO;
\end{verbatim}

\phantomsection
\label{prim:procedure-clr-line}
\subsection*{CLR\_LINE}
\index{CLR\_LINE@\texttt{CLR\_LINE}}
\begin{verbatim}
CLR_LINE();
CLR_LINE(column, row, attribute);
\end{verbatim}
Parameters:
\begin{description}
\item[column] Screen or editor column, as required by the primitive.
\item[row] Zero-based screen row.
\item[attribute] Screen attribute code.
\end{description}

Clears the current line or the specified screen range.

\paragraph{Example macro.}
\begin{verbatim}
$MACRO REFERENCE_CLR_LINE FROM EDIT;
    CLR_LINE(0, 0, 0);
END_MACRO;
\end{verbatim}

\phantomsection
\label{prim:procedure-cmd-to-key}
\subsection*{CMD\_TO\_KEY}
\index{CMD\_TO\_KEY@\texttt{CMD\_TO\_KEY}}
\begin{verbatim}
CMD_TO_KEY(keySpec, commandId, mode);
\end{verbatim}
Parameters:
\begin{description}
\item[keySpec] MRMAC key specification.
\item[commandId] MR application command identifier.
\item[mode] Editor mode constant.
\end{description}

Binds an MR command id to keySpec in mode.

\paragraph{Example macro.}
\begin{verbatim}
$MACRO REFERENCE_CMD_TO_KEY FROM EDIT;
    CMD_TO_KEY('reference', 0, 0);
END_MACRO;
\end{verbatim}

\phantomsection
\label{prim:procedure-flabel}
\subsection*{FLABEL}
\index{FLABEL@\texttt{FLABEL}}
\begin{verbatim}
FLABEL(label, keyCode, mode);
\end{verbatim}
Parameters:
\begin{description}
\item[label] User-visible field label.
\item[keyCode] Integer key code.
\item[mode] Editor mode constant.
\end{description}

Assigns a function-label text to a key and mode.

\paragraph{Example macro.}
\begin{verbatim}
$MACRO REFERENCE_FLABEL FROM EDIT;
    FLABEL('reference', 0, 0);
END_MACRO;
\end{verbatim}

\phantomsection
\label{prim:procedure-gotoxy}
\subsection*{GOTOXY}
\index{GOTOXY@\texttt{GOTOXY}}
\begin{verbatim}
GOTOXY(column, row);
\end{verbatim}
Parameters:
\begin{description}
\item[column] Screen or editor column, as required by the primitive.
\item[row] Zero-based screen row.
\end{description}

Moves the staged screen cursor.

\paragraph{Example macro.}
\begin{verbatim}
$MACRO REFERENCE_GOTOXY FROM EDIT;
    GOTOXY(0, 0);
END_MACRO;
\end{verbatim}

\phantomsection
\label{prim:procedure-key-in}
\subsection*{KEY\_IN}
\index{KEY\_IN@\texttt{KEY\_IN}}
\begin{verbatim}
KEY_IN(keySpec);
\end{verbatim}
Parameters:
\begin{description}
\item[keySpec] MRMAC key specification.
\end{description}

Tests the supplied key specification against current input.

\paragraph{Example macro.}
\begin{verbatim}
$MACRO REFERENCE_KEY_IN FROM EDIT;
    KEY_IN('reference');
END_MACRO;
\end{verbatim}

\phantomsection
\label{prim:procedure-key-record}
\subsection*{KEY\_RECORD}
\index{KEY\_RECORD@\texttt{KEY\_RECORD}}
\begin{verbatim}
KEY_RECORD;
\end{verbatim}
Parameters: None.

Toggles macro-key recording.

\paragraph{Example macro.}
\begin{verbatim}
$MACRO REFERENCE_KEY_RECORD FROM EDIT;
    KEY_RECORD;
END_MACRO;
\end{verbatim}

\phantomsection
\label{prim:procedure-kill-box}
\subsection*{KILL\_BOX}
\index{KILL\_BOX@\texttt{KILL\_BOX}}
\begin{verbatim}
KILL_BOX;
\end{verbatim}
Parameters: None.

Removes the current staged screen box.

\paragraph{Example macro.}
\begin{verbatim}
$MACRO REFERENCE_KILL_BOX FROM EDIT;
    KILL_BOX;
END_MACRO;
\end{verbatim}

\phantomsection
\label{prim:procedure-macro-to-key}
\subsection*{MACRO\_TO\_KEY}
\index{MACRO\_TO\_KEY@\texttt{MACRO\_TO\_KEY}}
\begin{verbatim}
MACRO_TO_KEY(keySpec, macroSpec, mode);
\end{verbatim}
Parameters:
\begin{description}
\item[keySpec] MRMAC key specification.
\item[macroSpec] Macro specification or callback target.
\item[mode] Editor mode constant.
\end{description}

Binds macroSpec to keySpec in mode.

\paragraph{Example macro.}
\begin{verbatim}
$MACRO REFERENCE_MACRO_TO_KEY FROM EDIT;
    MACRO_TO_KEY('reference', 'reference', 0);
END_MACRO;
\end{verbatim}

\begin{mrmacextension}
\phantomsection
\label{prim:procedure-make-message}
\subsection*{MAKE\_MESSAGE}
\index{MAKE\_MESSAGE@\texttt{MAKE\_MESSAGE}}
\begin{verbatim}
MAKE_MESSAGE(message);
\end{verbatim}
Parameters:
\begin{description}
\item[message] String-like user-visible message.
\end{description}

Writes a message through the MR message-line path.

\paragraph{Example macro.}
\begin{verbatim}
$MACRO REFERENCE_MAKE_MESSAGE FROM EDIT;
    MAKE_MESSAGE('reference');
END_MACRO;
\end{verbatim}
\end{mrmacextension}

\begin{mrmacextension}
\phantomsection
\label{prim:procedure-marquee}
\subsection*{MARQUEE}
\index{MARQUEE@\texttt{MARQUEE}}
\begin{verbatim}
MARQUEE(message);
\end{verbatim}
Parameters:
\begin{description}
\item[message] String-like user-visible message.
\end{description}

Displays an informational marquee message.

\paragraph{Example macro.}
\begin{verbatim}
$MACRO REFERENCE_MARQUEE FROM EDIT;
    MARQUEE('reference');
END_MACRO;
\end{verbatim}
\end{mrmacextension}

\begin{mrmacextension}
\phantomsection
\label{prim:procedure-marquee-error}
\subsection*{MARQUEE\_ERROR}
\index{MARQUEE\_ERROR@\texttt{MARQUEE\_ERROR}}
\begin{verbatim}
MARQUEE_ERROR(message);
\end{verbatim}
Parameters:
\begin{description}
\item[message] String-like user-visible message.
\end{description}

Displays an error marquee message.

\paragraph{Example macro.}
\begin{verbatim}
$MACRO REFERENCE_MARQUEE_ERROR FROM EDIT;
    MARQUEE_ERROR('reference');
END_MACRO;
\end{verbatim}
\end{mrmacextension}

\begin{mrmacextension}
\phantomsection
\label{prim:procedure-marquee-warning}
\subsection*{MARQUEE\_WARNING}
\index{MARQUEE\_WARNING@\texttt{MARQUEE\_WARNING}}
\begin{verbatim}
MARQUEE_WARNING(message);
\end{verbatim}
Parameters:
\begin{description}
\item[message] String-like user-visible message.
\end{description}

Displays a warning marquee message.

\paragraph{Example macro.}
\begin{verbatim}
$MACRO REFERENCE_MARQUEE_WARNING FROM EDIT;
    MARQUEE_WARNING('reference');
END_MACRO;
\end{verbatim}
\end{mrmacextension}

\phantomsection
\label{prim:procedure-pass-key}
\subsection*{PASS\_KEY}
\index{PASS\_KEY@\texttt{PASS\_KEY}}
\begin{verbatim}
PASS_KEY(keyCode, keyFlags);
\end{verbatim}
Parameters:
\begin{description}
\item[keyCode] Integer key code.
\item[keyFlags] Integer key modifier mask.
\end{description}

Passes a key with its modifier mask to the editor.

\paragraph{Example macro.}
\begin{verbatim}
$MACRO REFERENCE_PASS_KEY FROM EDIT;
    PASS_KEY(0, 0);
END_MACRO;
\end{verbatim}

\phantomsection
\label{prim:procedure-play-key-macro}
\subsection*{PLAY\_KEY\_MACRO}
\index{PLAY\_KEY\_MACRO@\texttt{PLAY\_KEY\_MACRO}}
\begin{verbatim}
PLAY_KEY_MACRO(keyCode, repeatCount);
PLAY_KEY_MACRO(keyCode, repeatCount, mode);
\end{verbatim}
Parameters:
\begin{description}
\item[keyCode] Integer key code.
\item[repeatCount] Number of key macro repetitions.
\end{description}

Replays a recorded key macro.

\paragraph{Example macro.}
\begin{verbatim}
$MACRO REFERENCE_PLAY_KEY_MACRO FROM EDIT;
    PLAY_KEY_MACRO(0, 0);
END_MACRO;
\end{verbatim}

\phantomsection
\label{prim:procedure-pop-labels}
\subsection*{POP\_LABELS}
\index{POP\_LABELS@\texttt{POP\_LABELS}}
\begin{verbatim}
POP_LABELS;
\end{verbatim}
Parameters: None.

Restores the saved function-label state.

\paragraph{Example macro.}
\begin{verbatim}
$MACRO REFERENCE_POP_LABELS FROM EDIT;
    POP_LABELS;
END_MACRO;
\end{verbatim}

\phantomsection
\label{prim:procedure-push-key}
\subsection*{PUSH\_KEY}
\index{PUSH\_KEY@\texttt{PUSH\_KEY}}
\begin{verbatim}
PUSH_KEY(keyCode, keyFlags);
\end{verbatim}
Parameters:
\begin{description}
\item[keyCode] Integer key code.
\item[keyFlags] Integer key modifier mask.
\end{description}

Queues a key with its modifier mask.

\paragraph{Example macro.}
\begin{verbatim}
$MACRO REFERENCE_PUSH_KEY FROM EDIT;
    PUSH_KEY(0, 0);
END_MACRO;
\end{verbatim}

\phantomsection
\label{prim:procedure-push-labels}
\subsection*{PUSH\_LABELS}
\index{PUSH\_LABELS@\texttt{PUSH\_LABELS}}
\begin{verbatim}
PUSH_LABELS;
\end{verbatim}
Parameters: None.

Saves the current function-label state.

\paragraph{Example macro.}
\begin{verbatim}
$MACRO REFERENCE_PUSH_LABELS FROM EDIT;
    PUSH_LABELS;
END_MACRO;
\end{verbatim}

\phantomsection
\label{prim:procedure-put-box}
\subsection*{PUT\_BOX}
\index{PUT\_BOX@\texttt{PUT\_BOX}}
\begin{verbatim}
PUT_BOX(left, top, right, bottom, foreground, background, text, border);
\end{verbatim}
Parameters:
\begin{description}
\item[left] Left screen column.
\item[top] Top screen row.
\item[right] Right screen column.
\item[bottom] Bottom screen row.
\item[foreground] Foreground colour code.
\item[background] Background colour code.
\item[text] String-like text.
\item[border] Integer border style or attribute.
\end{description}

Draws a staged text box at the supplied screen bounds.

\paragraph{Example macro.}
\begin{verbatim}
$MACRO REFERENCE_PUT_BOX FROM EDIT;
    PUT_BOX(0, 0, 0, 0, 0, 0, 'reference', 0);
END_MACRO;
\end{verbatim}

\phantomsection
\label{prim:procedure-put-col-num}
\subsection*{PUT\_COL\_NUM}
\index{PUT\_COL\_NUM@\texttt{PUT\_COL\_NUM}}
\begin{verbatim}
PUT_COL_NUM(column);
\end{verbatim}
Parameters:
\begin{description}
\item[column] Screen or editor column, as required by the primitive.
\end{description}

Displays the supplied column number.

\paragraph{Example macro.}
\begin{verbatim}
$MACRO REFERENCE_PUT_COL_NUM FROM EDIT;
    PUT_COL_NUM(0);
END_MACRO;
\end{verbatim}

\phantomsection
\label{prim:procedure-put-line-num}
\subsection*{PUT\_LINE\_NUM}
\index{PUT\_LINE\_NUM@\texttt{PUT\_LINE\_NUM}}
\begin{verbatim}
PUT_LINE_NUM(line);
\end{verbatim}
Parameters:
\begin{description}
\item[line] One-based editor line number.
\end{description}

Displays the supplied line number.

\paragraph{Example macro.}
\begin{verbatim}
$MACRO REFERENCE_PUT_LINE_NUM FROM EDIT;
    PUT_LINE_NUM(0);
END_MACRO;
\end{verbatim}

\phantomsection
\label{prim:procedure-read-key}
\subsection*{READ\_KEY}
\index{READ\_KEY@\texttt{READ\_KEY}}
\begin{verbatim}
READ_KEY;
\end{verbatim}
Parameters: None.

Waits for and dispatches the next input key in the active interaction context.

\paragraph{Example macro.}
\begin{verbatim}
$MACRO REFERENCE_READ_KEY FROM EDIT;
    READ_KEY;
END_MACRO;
\end{verbatim}

\begin{mrmacextension}
\phantomsection
\label{prim:procedure-register-menu-item}
\subsection*{REGISTER\_MENU\_ITEM}
\index{REGISTER\_MENU\_ITEM@\texttt{REGISTER\_MENU\_ITEM}}
\begin{verbatim}
REGISTER_MENU_ITEM(menuName, itemLabel);
REGISTER_MENU_ITEM(menuName, itemLabel, macroSpec);
\end{verbatim}
Parameters:
\begin{description}
\item[menuName] Named menu definition.
\item[itemLabel] String or character positional argument.
\item[macroSpec] Macro specification or callback target.
\end{description}

Registers an owned runtime menu item and optional macro callback.

\paragraph{Example macro.}
\begin{verbatim}
$MACRO REFERENCE_REGISTER_MENU_ITEM FROM EDIT;
    REGISTER_MENU_ITEM('reference', 'reference', 'reference');
END_MACRO;
\end{verbatim}
\end{mrmacextension}

\begin{mrmacextension}
\phantomsection
\label{prim:procedure-remove-menu-item}
\subsection*{REMOVE\_MENU\_ITEM}
\index{REMOVE\_MENU\_ITEM@\texttt{REMOVE\_MENU\_ITEM}}
\begin{verbatim}
REMOVE_MENU_ITEM(menuName, itemLabel);
\end{verbatim}
Parameters:
\begin{description}
\item[menuName] Named menu definition.
\item[itemLabel] String or character positional argument.
\end{description}

Removes an owned runtime menu item.

\paragraph{Example macro.}
\begin{verbatim}
$MACRO REFERENCE_REMOVE_MENU_ITEM FROM EDIT;
    REMOVE_MENU_ITEM('reference', 'reference');
END_MACRO;
\end{verbatim}
\end{mrmacextension}

\phantomsection
\label{prim:procedure-scroll-box-dn}
\subsection*{SCROLL\_BOX\_DN}
\index{SCROLL\_BOX\_DN@\texttt{SCROLL\_BOX\_DN}}
\begin{verbatim}
SCROLL_BOX_DN(left, top, right, bottom, attribute);
\end{verbatim}
Parameters:
\begin{description}
\item[left] Left screen column.
\item[top] Top screen row.
\item[right] Right screen column.
\item[bottom] Bottom screen row.
\item[attribute] Screen attribute code.
\end{description}

Scrolls the supplied staged screen rectangle down.

\paragraph{Example macro.}
\begin{verbatim}
$MACRO REFERENCE_SCROLL_BOX_DN FROM EDIT;
    SCROLL_BOX_DN(0, 0, 0, 0, 0);
END_MACRO;
\end{verbatim}

\phantomsection
\label{prim:procedure-scroll-box-up}
\subsection*{SCROLL\_BOX\_UP}
\index{SCROLL\_BOX\_UP@\texttt{SCROLL\_BOX\_UP}}
\begin{verbatim}
SCROLL_BOX_UP(left, top, right, bottom, attribute);
\end{verbatim}
Parameters:
\begin{description}
\item[left] Left screen column.
\item[top] Top screen row.
\item[right] Right screen column.
\item[bottom] Bottom screen row.
\item[attribute] Screen attribute code.
\end{description}

Scrolls the supplied staged screen rectangle up.

\paragraph{Example macro.}
\begin{verbatim}
$MACRO REFERENCE_SCROLL_BOX_UP FROM EDIT;
    SCROLL_BOX_UP(0, 0, 0, 0, 0);
END_MACRO;
\end{verbatim}

\begin{mrmacextension}
\phantomsection
\label{prim:procedure-set-clipboard-text}
\subsection*{SET\_CLIPBOARD\_TEXT}
\index{SET\_CLIPBOARD\_TEXT@\texttt{SET\_CLIPBOARD\_TEXT}}
\begin{verbatim}
SET_CLIPBOARD_TEXT(text);
\end{verbatim}
Parameters:
\begin{description}
\item[text] String-like text.
\end{description}

Replaces system clipboard text.

\paragraph{Example macro.}
\begin{verbatim}
$MACRO REFERENCE_SET_CLIPBOARD_TEXT FROM EDIT;
    SET_CLIPBOARD_TEXT('reference');
END_MACRO;
\end{verbatim}
\end{mrmacextension}

\phantomsection
\label{prim:procedure-unassign-all-keys}
\subsection*{UNASSIGN\_ALL\_KEYS}
\index{UNASSIGN\_ALL\_KEYS@\texttt{UNASSIGN\_ALL\_KEYS}}
\begin{verbatim}
UNASSIGN_ALL_KEYS;
\end{verbatim}
Parameters: None.

Removes all active macro key bindings.

\paragraph{Example macro.}
\begin{verbatim}
$MACRO REFERENCE_UNASSIGN_ALL_KEYS FROM EDIT;
    UNASSIGN_ALL_KEYS;
END_MACRO;
\end{verbatim}

\phantomsection
\label{prim:procedure-unassign-key}
\subsection*{UNASSIGN\_KEY}
\index{UNASSIGN\_KEY@\texttt{UNASSIGN\_KEY}}
\begin{verbatim}
UNASSIGN_KEY(keySpec, mode);
\end{verbatim}
Parameters:
\begin{description}
\item[keySpec] MRMAC key specification.
\item[mode] Editor mode constant.
\end{description}

Removes the binding for keySpec in mode.

\paragraph{Example macro.}
\begin{verbatim}
$MACRO REFERENCE_UNASSIGN_KEY FROM EDIT;
    UNASSIGN_KEY('reference', 0);
END_MACRO;
\end{verbatim}

\begin{mrmacextension}
\phantomsection
\label{prim:procedure-working}
\subsection*{WORKING}
\index{WORKING@\texttt{WORKING}}
\begin{verbatim}
WORKING;
\end{verbatim}
Parameters: None.

Shows the standard working indication.

\paragraph{Example macro.}
\begin{verbatim}
$MACRO REFERENCE_WORKING FROM EDIT;
    WORKING;
END_MACRO;
\end{verbatim}
\end{mrmacextension}

\phantomsection
\label{prim:procedure-write}
\subsection*{WRITE}
\index{WRITE@\texttt{WRITE}}
\begin{verbatim}
WRITE(text, column, row, foreground, background);
\end{verbatim}
Parameters:
\begin{description}
\item[text] String-like text.
\item[column] Screen or editor column, as required by the primitive.
\item[row] Zero-based screen row.
\item[foreground] Foreground colour code.
\item[background] Background colour code.
\end{description}

Writes staged screen text with foreground and background attributes.

\paragraph{Example macro.}
\begin{verbatim}
$MACRO REFERENCE_WRITE FROM EDIT;
    WRITE('reference', 0, 0, 0, 0);
END_MACRO;
\end{verbatim}


\chapter{Typed Dialogs and Collections}

Typed dialogs are built from explicit controls and named collections. A macro declares a dialog, adds controls, fills list, tree or table data, then executes it. The resulting control id and values remain ordinary MRMAC values.

\section{Primitive Map}
The map is a local index for this chapter. The final column is a generated page
reference and hyperlink to the full entry, where the source form, parameters,
semantics and a complete example macro appear.

\begin{longtable}{@{}L{0.38\textwidth}L{0.14\textwidth}L{0.27\textwidth}L{0.10\textwidth}@{}}
\toprule
Primitive & Role & Purpose & Reference \\
\midrule
\endhead
\toprule
Primitive & Role & Purpose & Reference \\
\midrule
\endhead
\primitiveoverview{UI\_BUTTON}{procedure-ui-button}{procedure}{Dialog construction.}
\primitiveoverview{UI\_DIALOG}{procedure-ui-dialog}{procedure}{Dialog construction.}
\primitiveoverview{UI\_DISPLAY}{procedure-ui-display}{procedure}{Dialog construction.}
\primitiveoverview{UI\_EXEC}{intrinsic-ui-exec}{intrinsic}{Dialog construction.}
\primitiveoverview{UI\_GRID}{procedure-ui-grid}{procedure}{Dialog construction.}
\primitiveoverview{UI\_INDEX}{intrinsic-ui-index}{intrinsic}{Dialog construction.}
\primitiveoverview{UI\_INPUT}{procedure-ui-input}{procedure}{Dialog construction.}
\primitiveoverview{UI\_LABEL}{procedure-ui-label}{procedure}{Dialog construction.}
\primitiveoverview{UI\_LISTBOX}{procedure-ui-listbox}{procedure}{Dialog construction.}
\primitiveoverview{UI\_LIST\_ADD}{procedure-ui-list-add}{procedure}{Dialog construction.}
\primitiveoverview{UI\_LIST\_CLEAR}{procedure-ui-list-clear}{procedure}{Dialog construction.}
\primitiveoverview{UI\_MESSAGEBOX}{procedure-ui-messagebox}{procedure}{Dialog construction.}
\primitiveoverview{UI\_TABLE}{procedure-ui-table}{procedure}{Dialog construction.}
\primitiveoverview{UI\_TABLE\_CLEAR}{procedure-ui-table-clear}{procedure}{Dialog construction.}
\primitiveoverview{UI\_TABLE\_COLUMN}{procedure-ui-table-column}{procedure}{Dialog construction.}
\primitiveoverview{UI\_TABLE\_ROW}{procedure-ui-table-row}{procedure}{Dialog construction.}
\primitiveoverview{UI\_TEXT}{intrinsic-ui-text}{intrinsic}{Dialog construction.}
\primitiveoverview{UI\_TREE}{procedure-ui-tree}{procedure}{Dialog construction.}
\primitiveoverview{UI\_TREE\_CLEAR}{procedure-ui-tree-clear}{procedure}{Dialog construction.}
\primitiveoverview{UI\_TREE\_NODE}{procedure-ui-tree-node}{procedure}{Dialog construction.}
\bottomrule
\end{longtable}

\section{Reference Entries}

\begin{mrmacextension}
\phantomsection
\label{prim:procedure-ui-button}
\subsection*{UI\_BUTTON}
\index{UI\_BUTTON@\texttt{UI\_BUTTON}}
\begin{verbatim}
UI_BUTTON(column, row, width, id, text);
\end{verbatim}
Parameters:
\begin{description}
\item[column] Screen or editor column, as required by the primitive.
\item[row] Zero-based screen row.
\item[width] Control or region width.
\item[id] Typed-dialog control identifier.
\item[text] String-like text.
\end{description}

Adds a typed dialog button with the supplied control id.

\paragraph{Example macro.}
\begin{verbatim}
$MACRO REFERENCE_UI_BUTTON FROM EDIT;
    UI_BUTTON(0, 0, 0, 0, 'reference');
END_MACRO;
\end{verbatim}
\end{mrmacextension}

\begin{mrmacextension}
\phantomsection
\label{prim:procedure-ui-dialog}
\subsection*{UI\_DIALOG}
\index{UI\_DIALOG@\texttt{UI\_DIALOG}}
\begin{verbatim}
UI_DIALOG(left, top, width, height, title);
\end{verbatim}
Parameters:
\begin{description}
\item[left] Left screen column.
\item[top] Top screen row.
\item[width] Control or region width.
\item[height] Control or region height.
\item[title] User-visible title.
\end{description}

Starts a typed dialog definition at the supplied rectangle and gives it title.

\paragraph{Example macro.}
\begin{verbatim}
$MACRO REFERENCE_UI_DIALOG FROM EDIT;
    UI_DIALOG(0, 0, 0, 0, 'reference');
END_MACRO;
\end{verbatim}
\end{mrmacextension}

\begin{mrmacextension}
\phantomsection
\label{prim:procedure-ui-display}
\subsection*{UI\_DISPLAY}
\index{UI\_DISPLAY@\texttt{UI\_DISPLAY}}
\begin{verbatim}
UI_DISPLAY(column, row, width, text);
\end{verbatim}
Parameters:
\begin{description}
\item[column] Screen or editor column, as required by the primitive.
\item[row] Zero-based screen row.
\item[width] Control or region width.
\item[text] String-like text.
\end{description}

Adds a read-only display field to the current typed dialog.

\paragraph{Example macro.}
\begin{verbatim}
$MACRO REFERENCE_UI_DISPLAY FROM EDIT;
    UI_DISPLAY(0, 0, 0, 'reference');
END_MACRO;
\end{verbatim}
\end{mrmacextension}

\begin{mrmacextension}
\phantomsection
\label{prim:intrinsic-ui-exec}
\subsection*{UI\_EXEC}
\index{UI\_EXEC@\texttt{UI\_EXEC}}
\begin{verbatim}
UI_EXEC
\end{verbatim}
Parameters: None.

Executes the current typed dialog and returns its selected control id.

Returns an integer.

\paragraph{Example macro.}
\begin{verbatim}
$MACRO REFERENCE_UI_EXEC FROM EDIT;
    DEF_INT(Result);
    Result := UI_EXEC;
END_MACRO;
\end{verbatim}
\end{mrmacextension}

\begin{mrmacextension}
\phantomsection
\label{prim:procedure-ui-grid}
\subsection*{UI\_GRID}
\index{UI\_GRID@\texttt{UI\_GRID}}
\begin{verbatim}
UI_GRID(column, row, width, height, id, tableName, initialValue, flags);
\end{verbatim}
Parameters:
\begin{description}
\item[column] Screen or editor column, as required by the primitive.
\item[row] Zero-based screen row.
\item[width] Control or region width.
\item[height] Control or region height.
\item[id] Typed-dialog control identifier.
\item[tableName] Named table definition.
\item[initialValue] Initial selected value.
\item[flags] Integer behaviour flags defined for the control.
\end{description}

Adds a typed grid control backed by tableName.

\paragraph{Example macro.}
\begin{verbatim}
$MACRO REFERENCE_UI_GRID FROM EDIT;
    UI_GRID(0, 0, 0, 0, 0, 'reference', 'reference', 0);
END_MACRO;
\end{verbatim}
\end{mrmacextension}

\begin{mrmacextension}
\phantomsection
\label{prim:intrinsic-ui-index}
\subsection*{UI\_INDEX}
\index{UI\_INDEX@\texttt{UI\_INDEX}}
\begin{verbatim}
UI_INDEX(controlId)
\end{verbatim}
Parameters:
\begin{description}
\item[controlId] Typed-dialog control identifier.
\end{description}

Returns the selection index of a typed-dialog control.

Returns an integer.

\paragraph{Example macro.}
\begin{verbatim}
$MACRO REFERENCE_UI_INDEX FROM EDIT;
    DEF_INT(Result);
    Result := UI_INDEX(0);
END_MACRO;
\end{verbatim}
\end{mrmacextension}

\begin{mrmacextension}
\phantomsection
\label{prim:procedure-ui-input}
\subsection*{UI\_INPUT}
\index{UI\_INPUT@\texttt{UI\_INPUT}}
\begin{verbatim}
UI_INPUT(column, row, width, id, fieldName, initialText);
\end{verbatim}
Parameters:
\begin{description}
\item[column] Screen or editor column, as required by the primitive.
\item[row] Zero-based screen row.
\item[width] Control or region width.
\item[id] Typed-dialog control identifier.
\item[fieldName] Named retained field.
\item[initialText] Initial text value.
\end{description}

Adds a named editable text field to the current typed dialog.

\paragraph{Example macro.}
\begin{verbatim}
$MACRO REFERENCE_UI_INPUT FROM EDIT;
    UI_INPUT(0, 0, 0, 0, 'reference', 'reference');
END_MACRO;
\end{verbatim}
\end{mrmacextension}

\begin{mrmacextension}
\phantomsection
\label{prim:procedure-ui-label}
\subsection*{UI\_LABEL}
\index{UI\_LABEL@\texttt{UI\_LABEL}}
\begin{verbatim}
UI_LABEL(column, row, text);
\end{verbatim}
Parameters:
\begin{description}
\item[column] Screen or editor column, as required by the primitive.
\item[row] Zero-based screen row.
\item[text] String-like text.
\end{description}

Adds static label text to the current typed dialog.

\paragraph{Example macro.}
\begin{verbatim}
$MACRO REFERENCE_UI_LABEL FROM EDIT;
    UI_LABEL(0, 0, 'reference');
END_MACRO;
\end{verbatim}
\end{mrmacextension}

\begin{mrmacextension}
\phantomsection
\label{prim:procedure-ui-listbox}
\subsection*{UI\_LISTBOX}
\index{UI\_LISTBOX@\texttt{UI\_LISTBOX}}
\begin{verbatim}
UI_LISTBOX(column, row, width, height, id, listName, initialValue, flags);
\end{verbatim}
Parameters:
\begin{description}
\item[column] Screen or editor column, as required by the primitive.
\item[row] Zero-based screen row.
\item[width] Control or region width.
\item[height] Control or region height.
\item[id] Typed-dialog control identifier.
\item[listName] Named list definition.
\item[initialValue] Initial selected value.
\item[flags] Integer behaviour flags defined for the control.
\end{description}

Adds a typed list-box control backed by listName.

\paragraph{Example macro.}
\begin{verbatim}
$MACRO REFERENCE_UI_LISTBOX FROM EDIT;
    UI_LISTBOX(0, 0, 0, 0, 0, 'reference', 'reference', 0);
END_MACRO;
\end{verbatim}
\end{mrmacextension}

\begin{mrmacextension}
\phantomsection
\label{prim:procedure-ui-list-add}
\subsection*{UI\_LIST\_ADD}
\index{UI\_LIST\_ADD@\texttt{UI\_LIST\_ADD}}
\begin{verbatim}
UI_LIST_ADD(listName, value);
\end{verbatim}
Parameters:
\begin{description}
\item[listName] Named list definition.
\item[value] Value written or selected by the primitive.
\end{description}

Appends one selectable value to the named typed-dialog list.

\paragraph{Example macro.}
\begin{verbatim}
$MACRO REFERENCE_UI_LIST_ADD FROM EDIT;
    UI_LIST_ADD('reference', 'reference');
END_MACRO;
\end{verbatim}
\end{mrmacextension}

\begin{mrmacextension}
\phantomsection
\label{prim:procedure-ui-list-clear}
\subsection*{UI\_LIST\_CLEAR}
\index{UI\_LIST\_CLEAR@\texttt{UI\_LIST\_CLEAR}}
\begin{verbatim}
UI_LIST_CLEAR(listName);
\end{verbatim}
Parameters:
\begin{description}
\item[listName] Named list definition.
\end{description}

Clears all values in the named typed-dialog list.

\paragraph{Example macro.}
\begin{verbatim}
$MACRO REFERENCE_UI_LIST_CLEAR FROM EDIT;
    UI_LIST_CLEAR('reference');
END_MACRO;
\end{verbatim}
\end{mrmacextension}

\begin{mrmacextension}
\phantomsection
\label{prim:procedure-ui-messagebox}
\subsection*{UI\_MESSAGEBOX}
\index{UI\_MESSAGEBOX@\texttt{UI\_MESSAGEBOX}}
\begin{verbatim}
UI_MESSAGEBOX(message);
\end{verbatim}
Parameters:
\begin{description}
\item[message] String-like user-visible message.
\end{description}

Shows a typed modal message box.

\paragraph{Example macro.}
\begin{verbatim}
$MACRO REFERENCE_UI_MESSAGEBOX FROM EDIT;
    UI_MESSAGEBOX('reference');
END_MACRO;
\end{verbatim}
\end{mrmacextension}

\begin{mrmacextension}
\phantomsection
\label{prim:procedure-ui-table}
\subsection*{UI\_TABLE}
\index{UI\_TABLE@\texttt{UI\_TABLE}}
\begin{verbatim}
UI_TABLE(column, row, width, height, id, tableName, initialValue, flags);
\end{verbatim}
Parameters:
\begin{description}
\item[column] Screen or editor column, as required by the primitive.
\item[row] Zero-based screen row.
\item[width] Control or region width.
\item[height] Control or region height.
\item[id] Typed-dialog control identifier.
\item[tableName] Named table definition.
\item[initialValue] Initial selected value.
\item[flags] Integer behaviour flags defined for the control.
\end{description}

Adds a typed table control backed by tableName.

\paragraph{Example macro.}
\begin{verbatim}
$MACRO REFERENCE_UI_TABLE FROM EDIT;
    UI_TABLE(0, 0, 0, 0, 0, 'reference', 'reference', 0);
END_MACRO;
\end{verbatim}
\end{mrmacextension}

\begin{mrmacextension}
\phantomsection
\label{prim:procedure-ui-table-clear}
\subsection*{UI\_TABLE\_CLEAR}
\index{UI\_TABLE\_CLEAR@\texttt{UI\_TABLE\_CLEAR}}
\begin{verbatim}
UI_TABLE_CLEAR(tableName);
\end{verbatim}
Parameters:
\begin{description}
\item[tableName] Named table definition.
\end{description}

Clears all rows and columns of the named typed-dialog table.

\paragraph{Example macro.}
\begin{verbatim}
$MACRO REFERENCE_UI_TABLE_CLEAR FROM EDIT;
    UI_TABLE_CLEAR('reference');
END_MACRO;
\end{verbatim}
\end{mrmacextension}

\begin{mrmacextension}
\phantomsection
\label{prim:procedure-ui-table-column}
\subsection*{UI\_TABLE\_COLUMN}
\index{UI\_TABLE\_COLUMN@\texttt{UI\_TABLE\_COLUMN}}
\begin{verbatim}
UI_TABLE_COLUMN(tableName, columnName, width);
\end{verbatim}
Parameters:
\begin{description}
\item[tableName] Named table definition.
\item[columnName] String or character positional argument.
\item[width] Control or region width.
\end{description}

Adds one named column to the typed-dialog table.

\paragraph{Example macro.}
\begin{verbatim}
$MACRO REFERENCE_UI_TABLE_COLUMN FROM EDIT;
    UI_TABLE_COLUMN('reference', 'reference', 0);
END_MACRO;
\end{verbatim}
\end{mrmacextension}

\begin{mrmacextension}
\phantomsection
\label{prim:procedure-ui-table-row}
\subsection*{UI\_TABLE\_ROW}
\index{UI\_TABLE\_ROW@\texttt{UI\_TABLE\_ROW}}
\begin{verbatim}
UI_TABLE_ROW(tableName, rowId, values);
\end{verbatim}
Parameters:
\begin{description}
\item[tableName] Named table definition.
\item[rowId] String-like stable row identifier.
\item[values] String-like encoded cell values for the row.
\end{description}

Adds or replaces one identified typed-dialog table row.

\paragraph{Example macro.}
\begin{verbatim}
$MACRO REFERENCE_UI_TABLE_ROW FROM EDIT;
    UI_TABLE_ROW('reference', 'reference', 'reference');
END_MACRO;
\end{verbatim}
\end{mrmacextension}

\begin{mrmacextension}
\phantomsection
\label{prim:intrinsic-ui-text}
\subsection*{UI\_TEXT}
\index{UI\_TEXT@\texttt{UI\_TEXT}}
\begin{verbatim}
UI_TEXT(controlId)
\end{verbatim}
Parameters:
\begin{description}
\item[controlId] Typed-dialog control identifier.
\end{description}

Returns the text value of a typed-dialog control.

Returns a string.

\paragraph{Example macro.}
\begin{verbatim}
$MACRO REFERENCE_UI_TEXT FROM EDIT;
    DEF_STR(Result);
    Result := UI_TEXT(0);
END_MACRO;
\end{verbatim}
\end{mrmacextension}

\begin{mrmacextension}
\phantomsection
\label{prim:procedure-ui-tree}
\subsection*{UI\_TREE}
\index{UI\_TREE@\texttt{UI\_TREE}}
\begin{verbatim}
UI_TREE(column, row, width, height, id, treeName, initialValue, flags);
\end{verbatim}
Parameters:
\begin{description}
\item[column] Screen or editor column, as required by the primitive.
\item[row] Zero-based screen row.
\item[width] Control or region width.
\item[height] Control or region height.
\item[id] Typed-dialog control identifier.
\item[treeName] Named tree definition.
\item[initialValue] Initial selected value.
\item[flags] Integer behaviour flags defined for the control.
\end{description}

Adds a typed tree control backed by treeName.

\paragraph{Example macro.}
\begin{verbatim}
$MACRO REFERENCE_UI_TREE FROM EDIT;
    UI_TREE(0, 0, 0, 0, 0, 'reference', 'reference', 0);
END_MACRO;
\end{verbatim}
\end{mrmacextension}

\begin{mrmacextension}
\phantomsection
\label{prim:procedure-ui-tree-clear}
\subsection*{UI\_TREE\_CLEAR}
\index{UI\_TREE\_CLEAR@\texttt{UI\_TREE\_CLEAR}}
\begin{verbatim}
UI_TREE_CLEAR(treeName);
\end{verbatim}
Parameters:
\begin{description}
\item[treeName] Named tree definition.
\end{description}

Clears the named typed-dialog tree data source.

\paragraph{Example macro.}
\begin{verbatim}
$MACRO REFERENCE_UI_TREE_CLEAR FROM EDIT;
    UI_TREE_CLEAR('reference');
END_MACRO;
\end{verbatim}
\end{mrmacextension}

\begin{mrmacextension}
\phantomsection
\label{prim:procedure-ui-tree-node}
\subsection*{UI\_TREE\_NODE}
\index{UI\_TREE\_NODE@\texttt{UI\_TREE\_NODE}}
\begin{verbatim}
UI_TREE_NODE(treeName, nodeId, parentId, text, flags);
\end{verbatim}
Parameters:
\begin{description}
\item[treeName] Named tree definition.
\item[nodeId] String-like tree-node identifier.
\item[parentId] String-like parent tree-node identifier.
\item[text] String-like text.
\item[flags] Integer behaviour flags defined for the control.
\end{description}

Adds or replaces one typed-dialog tree node.

\paragraph{Example macro.}
\begin{verbatim}
$MACRO REFERENCE_UI_TREE_NODE FROM EDIT;
    UI_TREE_NODE('reference', 'reference', 'reference', 'reference', 0);
END_MACRO;
\end{verbatim}
\end{mrmacextension}


\chapter{Modeless MMP Interfaces}

Modeless MRMAC Primitives, abbreviated MMP, build retained modeless interfaces from typed MRMAC values. A macro declares controls, shows a named window, and later updates its retained fields. MMP never exposes a native view, pointer or event object.

\section{Primitive Map}
The map is a local index for this chapter. The final column is a generated page
reference and hyperlink to the full entry, where the source form, parameters,
semantics and a complete example macro appear.

\begin{longtable}{@{}L{0.38\textwidth}L{0.14\textwidth}L{0.27\textwidth}L{0.10\textwidth}@{}}
\toprule
Primitive & Role & Purpose & Reference \\
\midrule
\endhead
\toprule
Primitive & Role & Purpose & Reference \\
\midrule
\endhead
\primitiveoverview{UI\_MODELESS\_CLOSE}{procedure-ui-modeless-close}{procedure}{Modeless UI.}
\primitiveoverview{UI\_MODELESS\_DISPLAY}{procedure-ui-modeless-display}{procedure}{Modeless UI.}
\primitiveoverview{UI\_MODELESS\_ON}{procedure-ui-modeless-on}{procedure}{Modeless UI.}
\primitiveoverview{UI\_MODELESS\_SHOW}{procedure-ui-modeless-show}{procedure}{Modeless UI.}
\primitiveoverview{UI\_MODELESS\_UPDATE}{procedure-ui-modeless-update}{procedure}{Modeless UI.}
\primitiveoverview{MMP\_ACTION\_BUTTON}{procedure-mmp-action-button}{procedure}{Modeless UI.}
\primitiveoverview{MMP\_ACTION\_MENU}{procedure-mmp-action-menu}{procedure}{Modeless UI.}
\primitiveoverview{MMP\_BOOL\_FIELD}{procedure-mmp-bool-field}{procedure}{Modeless UI.}
\primitiveoverview{MMP\_BOOL\_SET}{procedure-mmp-bool-set}{procedure}{Modeless UI.}
\primitiveoverview{MMP\_BOOL\_VALUE}{intrinsic-mmp-bool-value}{intrinsic}{Modeless UI.}
\primitiveoverview{MMP\_CANVAS}{procedure-mmp-canvas}{procedure}{Modeless UI.}
\primitiveoverview{MMP\_CANVAS\_BOX}{procedure-mmp-canvas-box}{procedure}{Modeless UI.}
\primitiveoverview{MMP\_CANVAS\_CLEAR}{procedure-mmp-canvas-clear}{procedure}{Modeless UI.}
\primitiveoverview{MMP\_CANVAS\_COMMIT}{procedure-mmp-canvas-commit}{procedure}{Modeless UI.}
\primitiveoverview{MMP\_CANVAS\_FILL}{procedure-mmp-canvas-fill}{procedure}{Modeless UI.}
\primitiveoverview{MMP\_CANVAS\_GLYPH}{procedure-mmp-canvas-glyph}{procedure}{Modeless UI.}
\primitiveoverview{MMP\_CANVAS\_HOTSPOT}{procedure-mmp-canvas-hotspot}{procedure}{Modeless UI.}
\primitiveoverview{MMP\_CANVAS\_LINE}{procedure-mmp-canvas-line}{procedure}{Modeless UI.}
\primitiveoverview{MMP\_CANVAS\_TEXT}{procedure-mmp-canvas-text}{procedure}{Modeless UI.}
\primitiveoverview{MMP\_INT\_FIELD}{procedure-mmp-int-field}{procedure}{Modeless UI.}
\primitiveoverview{MMP\_INT\_SET}{procedure-mmp-int-set}{procedure}{Modeless UI.}
\primitiveoverview{MMP\_INT\_VALUE}{intrinsic-mmp-int-value}{intrinsic}{Modeless UI.}
\primitiveoverview{MMP\_LOG\_APPEND}{procedure-mmp-log-append}{procedure}{Modeless UI.}
\primitiveoverview{MMP\_LOG\_CLEAR}{procedure-mmp-log-clear}{procedure}{Modeless UI.}
\primitiveoverview{MMP\_LOG\_COUNT}{intrinsic-mmp-log-count}{intrinsic}{Modeless UI.}
\primitiveoverview{MMP\_LOG\_FIELD}{procedure-mmp-log-field}{procedure}{Modeless UI.}
\primitiveoverview{MMP\_MENU\_CLEAR}{procedure-mmp-menu-clear}{procedure}{Modeless UI.}
\primitiveoverview{MMP\_MENU\_ITEM}{procedure-mmp-menu-item}{procedure}{Modeless UI.}
\primitiveoverview{MMP\_PROGRESS\_FIELD}{procedure-mmp-progress-field}{procedure}{Modeless UI.}
\primitiveoverview{MMP\_PROGRESS\_SET}{procedure-mmp-progress-set}{procedure}{Modeless UI.}
\primitiveoverview{MMP\_PROGRESS\_VALUE}{intrinsic-mmp-progress-value}{intrinsic}{Modeless UI.}
\primitiveoverview{MMP\_SELECT\_FIELD}{procedure-mmp-select-field}{procedure}{Modeless UI.}
\primitiveoverview{MMP\_SELECT\_OPTION}{procedure-mmp-select-option}{procedure}{Modeless UI.}
\primitiveoverview{MMP\_SELECT\_SET}{procedure-mmp-select-set}{procedure}{Modeless UI.}
\primitiveoverview{MMP\_SELECT\_VALUE}{intrinsic-mmp-select-value}{intrinsic}{Modeless UI.}
\primitiveoverview{MMP\_STATUS\_FIELD}{procedure-mmp-status-field}{procedure}{Modeless UI.}
\primitiveoverview{MMP\_STATUS\_SET}{procedure-mmp-status-set}{procedure}{Modeless UI.}
\primitiveoverview{MMP\_STYLE\_ACCENT}{intrinsic-mmp-style-accent}{intrinsic}{Modeless UI.}
\primitiveoverview{MMP\_STYLE\_MUTED}{intrinsic-mmp-style-muted}{intrinsic}{Modeless UI.}
\primitiveoverview{MMP\_STYLE\_SURFACE}{intrinsic-mmp-style-surface}{intrinsic}{Modeless UI.}
\primitiveoverview{MMP\_STYLE\_TEXT}{intrinsic-mmp-style-text}{intrinsic}{Modeless UI.}
\primitiveoverview{MMP\_TEXT\_FIELD}{procedure-mmp-text-field}{procedure}{Modeless UI.}
\primitiveoverview{MMP\_TEXT\_SET}{procedure-mmp-text-set}{procedure}{Modeless UI.}
\primitiveoverview{MMP\_TEXT\_VALUE}{intrinsic-mmp-text-value}{intrinsic}{Modeless UI.}
\primitiveoverview{MMP\_TIMER\_START}{procedure-mmp-timer-start}{procedure}{Modeless UI.}
\primitiveoverview{MMP\_TIMER\_STOP}{procedure-mmp-timer-stop}{procedure}{Modeless UI.}
\primitiveoverview{MMP\_WINDOW\_EXISTS}{intrinsic-mmp-window-exists}{intrinsic}{Modeless UI.}
\primitiveoverview{MMP\_WINDOW\_HEIGHT}{intrinsic-mmp-window-height}{intrinsic}{Modeless UI.}
\primitiveoverview{MMP\_WINDOW\_INSTANCE}{intrinsic-mmp-window-instance}{intrinsic}{Modeless UI.}
\primitiveoverview{MMP\_WINDOW\_WIDTH}{intrinsic-mmp-window-width}{intrinsic}{Modeless UI.}
\primitiveoverview{MMP\_WINDOW\_X}{intrinsic-mmp-window-x}{intrinsic}{Modeless UI.}
\primitiveoverview{MMP\_WINDOW\_Y}{intrinsic-mmp-window-y}{intrinsic}{Modeless UI.}
\bottomrule
\end{longtable}

\section{Reference Entries}

\begin{mrmacextension}
\phantomsection
\label{prim:procedure-ui-modeless-close}
\subsection*{UI\_MODELESS\_CLOSE}
\index{UI\_MODELESS\_CLOSE@\texttt{UI\_MODELESS\_CLOSE}}
\begin{verbatim}
UI_MODELESS_CLOSE(windowName);
\end{verbatim}
Parameters:
\begin{description}
\item[windowName] Logical modeless window name.
\end{description}

Requests that the named modeless window close through the UI event path.

\paragraph{Example macro.}
\begin{verbatim}
$MACRO REFERENCE_UI_MODELESS_CLOSE FROM EDIT;
    UI_MODELESS_CLOSE('reference');
END_MACRO;
\end{verbatim}
\end{mrmacextension}

\begin{mrmacextension}
\phantomsection
\label{prim:procedure-ui-modeless-display}
\subsection*{UI\_MODELESS\_DISPLAY}
\index{UI\_MODELESS\_DISPLAY@\texttt{UI\_MODELESS\_DISPLAY}}
\begin{verbatim}
UI_MODELESS_DISPLAY(windowName, displayId, text);
\end{verbatim}
Parameters:
\begin{description}
\item[windowName] Logical modeless window name.
\item[displayId] Integer retained display-line identifier.
\item[text] String-like text.
\end{description}

Replaces one retained modeless display line.

\paragraph{Example macro.}
\begin{verbatim}
$MACRO REFERENCE_UI_MODELESS_DISPLAY FROM EDIT;
    UI_MODELESS_DISPLAY('reference', 0, 'reference');
END_MACRO;
\end{verbatim}
\end{mrmacextension}

\begin{mrmacextension}
\phantomsection
\label{prim:procedure-ui-modeless-on}
\subsection*{UI\_MODELESS\_ON}
\index{UI\_MODELESS\_ON@\texttt{UI\_MODELESS\_ON}}
\begin{verbatim}
UI_MODELESS_ON(controlId, macroSpec);
\end{verbatim}
Parameters:
\begin{description}
\item[controlId] Typed-dialog control identifier.
\item[macroSpec] Macro specification or callback target.
\end{description}

Binds a typed-dialog control id to a macro callback.

\paragraph{Example macro.}
\begin{verbatim}
$MACRO REFERENCE_UI_MODELESS_ON FROM EDIT;
    UI_MODELESS_ON(0, 'reference');
END_MACRO;
\end{verbatim}
\end{mrmacextension}

\begin{mrmacextension}
\phantomsection
\label{prim:procedure-ui-modeless-show}
\subsection*{UI\_MODELESS\_SHOW}
\index{UI\_MODELESS\_SHOW@\texttt{UI\_MODELESS\_SHOW}}
\begin{verbatim}
UI_MODELESS_SHOW(windowName);
\end{verbatim}
Parameters:
\begin{description}
\item[windowName] Logical modeless window name.
\end{description}

Shows the named retained modeless window.

\paragraph{Example macro.}
\begin{verbatim}
$MACRO REFERENCE_UI_MODELESS_SHOW FROM EDIT;
    UI_MODELESS_SHOW('reference');
END_MACRO;
\end{verbatim}
\end{mrmacextension}

\begin{mrmacextension}
\phantomsection
\label{prim:procedure-ui-modeless-update}
\subsection*{UI\_MODELESS\_UPDATE}
\index{UI\_MODELESS\_UPDATE@\texttt{UI\_MODELESS\_UPDATE}}
\begin{verbatim}
UI_MODELESS_UPDATE(windowName);
\end{verbatim}
Parameters:
\begin{description}
\item[windowName] Logical modeless window name.
\end{description}

Requests a retained modeless window refresh without taking focus.

\paragraph{Example macro.}
\begin{verbatim}
$MACRO REFERENCE_UI_MODELESS_UPDATE FROM EDIT;
    UI_MODELESS_UPDATE('reference');
END_MACRO;
\end{verbatim}
\end{mrmacextension}

\begin{mrmacextension}
\phantomsection
\label{prim:procedure-mmp-action-button}
\subsection*{MMP\_ACTION\_BUTTON}
\index{MMP\_ACTION\_BUTTON@\texttt{MMP\_ACTION\_BUTTON}}
\begin{verbatim}
MMP_ACTION_BUTTON(column, row, width, height, label, macroSpec);
\end{verbatim}
Parameters:
\begin{description}
\item[column] Screen or editor column, as required by the primitive.
\item[row] Zero-based screen row.
\item[width] Control or region width.
\item[height] Control or region height.
\item[label] User-visible field label.
\item[macroSpec] Macro specification or callback target.
\end{description}

Adds a modeless action button and its macro callback.

\paragraph{Example macro.}
\begin{verbatim}
$MACRO REFERENCE_MMP_ACTION_BUTTON FROM EDIT;
    MMP_ACTION_BUTTON(0, 0, 0, 0, 'reference', 'reference');
END_MACRO;
\end{verbatim}
\end{mrmacextension}

\begin{mrmacextension}
\phantomsection
\label{prim:procedure-mmp-action-menu}
\subsection*{MMP\_ACTION\_MENU}
\index{MMP\_ACTION\_MENU@\texttt{MMP\_ACTION\_MENU}}
\begin{verbatim}
MMP_ACTION_MENU(column, row, width, height, id, menuName, initialValue, macroSpec);
\end{verbatim}
Parameters:
\begin{description}
\item[column] Screen or editor column, as required by the primitive.
\item[row] Zero-based screen row.
\item[width] Control or region width.
\item[height] Control or region height.
\item[id] Typed-dialog control identifier.
\item[menuName] Named menu definition.
\item[initialValue] Initial selected value.
\item[macroSpec] Macro specification or callback target.
\end{description}

Adds a modeless action menu backed by menuName.

\paragraph{Example macro.}
\begin{verbatim}
$MACRO REFERENCE_MMP_ACTION_MENU FROM EDIT;
    MMP_ACTION_MENU(0, 0, 0, 0, 0, 'reference', 'reference', 'reference');
END_MACRO;
\end{verbatim}
\end{mrmacextension}

\begin{mrmacextension}
\phantomsection
\label{prim:procedure-mmp-bool-field}
\subsection*{MMP\_BOOL\_FIELD}
\index{MMP\_BOOL\_FIELD@\texttt{MMP\_BOOL\_FIELD}}
\begin{verbatim}
MMP_BOOL_FIELD(column, row, fieldName, label, initialValue);
\end{verbatim}
Parameters:
\begin{description}
\item[column] Screen or editor column, as required by the primitive.
\item[row] Zero-based screen row.
\item[fieldName] Named retained field.
\item[label] User-visible field label.
\item[initialValue] Initial selected value.
\end{description}

Adds a retained named boolean field.

\paragraph{Example macro.}
\begin{verbatim}
$MACRO REFERENCE_MMP_BOOL_FIELD FROM EDIT;
    MMP_BOOL_FIELD(0, 0, 'reference', 'reference', 0);
END_MACRO;
\end{verbatim}
\end{mrmacextension}

\begin{mrmacextension}
\phantomsection
\label{prim:procedure-mmp-bool-set}
\subsection*{MMP\_BOOL\_SET}
\index{MMP\_BOOL\_SET@\texttt{MMP\_BOOL\_SET}}
\begin{verbatim}
MMP_BOOL_SET(windowName, fieldName, value);
\end{verbatim}
Parameters:
\begin{description}
\item[windowName] Logical modeless window name.
\item[fieldName] Named retained field.
\item[value] Value written or selected by the primitive.
\end{description}

Writes a retained modeless boolean field.

\paragraph{Example macro.}
\begin{verbatim}
$MACRO REFERENCE_MMP_BOOL_SET FROM EDIT;
    MMP_BOOL_SET('reference', 'reference', 0);
END_MACRO;
\end{verbatim}
\end{mrmacextension}

\begin{mrmacextension}
\phantomsection
\label{prim:intrinsic-mmp-bool-value}
\subsection*{MMP\_BOOL\_VALUE}
\index{MMP\_BOOL\_VALUE@\texttt{MMP\_BOOL\_VALUE}}
\begin{verbatim}
MMP_BOOL_VALUE(windowName, fieldName)
\end{verbatim}
Parameters:
\begin{description}
\item[windowName] Logical modeless window name.
\item[fieldName] Named retained field.
\end{description}

Reads the retained boolean value of a modeless field.

Returns an integer.

\paragraph{Example macro.}
\begin{verbatim}
$MACRO REFERENCE_MMP_BOOL_VALUE FROM EDIT;
    DEF_INT(Result);
    Result := MMP_BOOL_VALUE('reference', 'reference');
END_MACRO;
\end{verbatim}
\end{mrmacextension}

\begin{mrmacextension}
\phantomsection
\label{prim:procedure-mmp-canvas}
\subsection*{MMP\_CANVAS}
\index{MMP\_CANVAS@\texttt{MMP\_CANVAS}}
\begin{verbatim}
MMP_CANVAS(left, top, width, height, canvasName);
\end{verbatim}
Parameters:
\begin{description}
\item[left] Left screen column.
\item[top] Top screen row.
\item[width] Control or region width.
\item[height] Control or region height.
\item[canvasName] Named retained canvas.
\end{description}

Adds a named retained canvas to the current modeless definition.

\paragraph{Example macro.}
\begin{verbatim}
$MACRO REFERENCE_MMP_CANVAS FROM EDIT;
    MMP_CANVAS(0, 0, 0, 0, 'reference');
END_MACRO;
\end{verbatim}
\end{mrmacextension}

\begin{mrmacextension}
\phantomsection
\label{prim:procedure-mmp-canvas-box}
\subsection*{MMP\_CANVAS\_BOX}
\index{MMP\_CANVAS\_BOX@\texttt{MMP\_CANVAS\_BOX}}
\begin{verbatim}
MMP_CANVAS_BOX(windowName, canvasName, left, top, width, height, style);
\end{verbatim}
Parameters:
\begin{description}
\item[windowName] Logical modeless window name.
\item[canvasName] Named retained canvas.
\item[left] Left screen column.
\item[top] Top screen row.
\item[width] Control or region width.
\item[height] Control or region height.
\item[style] MRMAC style or attribute code.
\end{description}

Draws a retained outlined canvas rectangle.

\paragraph{Example macro.}
\begin{verbatim}
$MACRO REFERENCE_MMP_CANVAS_BOX FROM EDIT;
    MMP_CANVAS_BOX('reference', 'reference', 0, 0, 0, 0, 0);
END_MACRO;
\end{verbatim}
\end{mrmacextension}

\begin{mrmacextension}
\phantomsection
\label{prim:procedure-mmp-canvas-clear}
\subsection*{MMP\_CANVAS\_CLEAR}
\index{MMP\_CANVAS\_CLEAR@\texttt{MMP\_CANVAS\_CLEAR}}
\begin{verbatim}
MMP_CANVAS_CLEAR(windowName, canvasName, style);
\end{verbatim}
Parameters:
\begin{description}
\item[windowName] Logical modeless window name.
\item[canvasName] Named retained canvas.
\item[style] MRMAC style or attribute code.
\end{description}

Clears the retained scene of the named canvas using style.

\paragraph{Example macro.}
\begin{verbatim}
$MACRO REFERENCE_MMP_CANVAS_CLEAR FROM EDIT;
    MMP_CANVAS_CLEAR('reference', 'reference', 0);
END_MACRO;
\end{verbatim}
\end{mrmacextension}

\begin{mrmacextension}
\phantomsection
\label{prim:procedure-mmp-canvas-commit}
\subsection*{MMP\_CANVAS\_COMMIT}
\index{MMP\_CANVAS\_COMMIT@\texttt{MMP\_CANVAS\_COMMIT}}
\begin{verbatim}
MMP_CANVAS_COMMIT(windowName, canvasName);
\end{verbatim}
Parameters:
\begin{description}
\item[windowName] Logical modeless window name.
\item[canvasName] Named retained canvas.
\end{description}

Redraws only the named retained canvas.

\paragraph{Example macro.}
\begin{verbatim}
$MACRO REFERENCE_MMP_CANVAS_COMMIT FROM EDIT;
    MMP_CANVAS_COMMIT('reference', 'reference');
END_MACRO;
\end{verbatim}
\end{mrmacextension}

\begin{mrmacextension}
\phantomsection
\label{prim:procedure-mmp-canvas-fill}
\subsection*{MMP\_CANVAS\_FILL}
\index{MMP\_CANVAS\_FILL@\texttt{MMP\_CANVAS\_FILL}}
\begin{verbatim}
MMP_CANVAS_FILL(windowName, canvasName, left, top, width, height, style);
\end{verbatim}
Parameters:
\begin{description}
\item[windowName] Logical modeless window name.
\item[canvasName] Named retained canvas.
\item[left] Left screen column.
\item[top] Top screen row.
\item[width] Control or region width.
\item[height] Control or region height.
\item[style] MRMAC style or attribute code.
\end{description}

Fills a retained canvas rectangle.

\paragraph{Example macro.}
\begin{verbatim}
$MACRO REFERENCE_MMP_CANVAS_FILL FROM EDIT;
    MMP_CANVAS_FILL('reference', 'reference', 0, 0, 0, 0, 0);
END_MACRO;
\end{verbatim}
\end{mrmacextension}

\begin{mrmacextension}
\phantomsection
\label{prim:procedure-mmp-canvas-glyph}
\subsection*{MMP\_CANVAS\_GLYPH}
\index{MMP\_CANVAS\_GLYPH@\texttt{MMP\_CANVAS\_GLYPH}}
\begin{verbatim}
MMP_CANVAS_GLYPH(windowName, canvasName, column, row, style, glyph);
\end{verbatim}
Parameters:
\begin{description}
\item[windowName] Logical modeless window name.
\item[canvasName] Named retained canvas.
\item[column] Screen or editor column, as required by the primitive.
\item[row] Zero-based screen row.
\item[style] MRMAC style or attribute code.
\item[glyph] String-like glyph text used by the retained canvas.
\end{description}

Places one retained glyph at a canvas coordinate.

\paragraph{Example macro.}
\begin{verbatim}
$MACRO REFERENCE_MMP_CANVAS_GLYPH FROM EDIT;
    MMP_CANVAS_GLYPH('reference', 'reference', 0, 0, 0, 'reference');
END_MACRO;
\end{verbatim}
\end{mrmacextension}

\begin{mrmacextension}
\phantomsection
\label{prim:procedure-mmp-canvas-hotspot}
\subsection*{MMP\_CANVAS\_HOTSPOT}
\index{MMP\_CANVAS\_HOTSPOT@\texttt{MMP\_CANVAS\_HOTSPOT}}
\begin{verbatim}
MMP_CANVAS_HOTSPOT(windowName, left, top, width, height, button, macroSpec);
\end{verbatim}
Parameters:
\begin{description}
\item[windowName] Logical modeless window name.
\item[left] Left screen column.
\item[top] Top screen row.
\item[width] Control or region width.
\item[height] Control or region height.
\item[button] Integer mouse-button identifier.
\item[macroSpec] Macro specification or callback target.
\end{description}

Registers a canvas-local left-click region and macro callback.

\paragraph{Example macro.}
\begin{verbatim}
$MACRO REFERENCE_MMP_CANVAS_HOTSPOT FROM EDIT;
    MMP_CANVAS_HOTSPOT('reference', 0, 0, 0, 0, 0, 'reference');
END_MACRO;
\end{verbatim}
\end{mrmacextension}

\begin{mrmacextension}
\phantomsection
\label{prim:procedure-mmp-canvas-line}
\subsection*{MMP\_CANVAS\_LINE}
\index{MMP\_CANVAS\_LINE@\texttt{MMP\_CANVAS\_LINE}}
\begin{verbatim}
MMP_CANVAS_LINE(windowName, canvasName, x1, y1, x2, y2, style, glyph);
\end{verbatim}
Parameters:
\begin{description}
\item[windowName] Logical modeless window name.
\item[canvasName] Named retained canvas.
\item[x1] Starting canvas column.
\item[y1] Starting canvas row.
\item[x2] Ending canvas column.
\item[y2] Ending canvas row.
\item[style] MRMAC style or attribute code.
\item[glyph] String-like glyph text used by the retained canvas.
\end{description}

Draws a retained line between two canvas coordinates.

\paragraph{Example macro.}
\begin{verbatim}
$MACRO REFERENCE_MMP_CANVAS_LINE FROM EDIT;
    MMP_CANVAS_LINE('reference', 'reference', 0, 0, 0, 0, 0, 'reference');
END_MACRO;
\end{verbatim}
\end{mrmacextension}

\begin{mrmacextension}
\phantomsection
\label{prim:procedure-mmp-canvas-text}
\subsection*{MMP\_CANVAS\_TEXT}
\index{MMP\_CANVAS\_TEXT@\texttt{MMP\_CANVAS\_TEXT}}
\begin{verbatim}
MMP_CANVAS_TEXT(windowName, canvasName, column, row, style, text);
\end{verbatim}
Parameters:
\begin{description}
\item[windowName] Logical modeless window name.
\item[canvasName] Named retained canvas.
\item[column] Screen or editor column, as required by the primitive.
\item[row] Zero-based screen row.
\item[style] MRMAC style or attribute code.
\item[text] String-like text.
\end{description}

Places retained text at a canvas coordinate.

\paragraph{Example macro.}
\begin{verbatim}
$MACRO REFERENCE_MMP_CANVAS_TEXT FROM EDIT;
    MMP_CANVAS_TEXT('reference', 'reference', 0, 0, 0, 'reference');
END_MACRO;
\end{verbatim}
\end{mrmacextension}

\begin{mrmacextension}
\phantomsection
\label{prim:procedure-mmp-int-field}
\subsection*{MMP\_INT\_FIELD}
\index{MMP\_INT\_FIELD@\texttt{MMP\_INT\_FIELD}}
\begin{verbatim}
MMP_INT_FIELD(column, row, width, fieldName, label, minimum, maximum, initialValue);
\end{verbatim}
Parameters:
\begin{description}
\item[column] Screen or editor column, as required by the primitive.
\item[row] Zero-based screen row.
\item[width] Control or region width.
\item[fieldName] Named retained field.
\item[label] User-visible field label.
\item[minimum] Inclusive minimum integer value.
\item[maximum] Inclusive maximum integer value.
\item[initialValue] Initial selected value.
\end{description}

Adds a range-bounded retained integer field.

\paragraph{Example macro.}
\begin{verbatim}
$MACRO REFERENCE_MMP_INT_FIELD FROM EDIT;
    MMP_INT_FIELD(0, 0, 0, 'reference', 'reference', 0, 0, 0);
END_MACRO;
\end{verbatim}
\end{mrmacextension}

\begin{mrmacextension}
\phantomsection
\label{prim:procedure-mmp-int-set}
\subsection*{MMP\_INT\_SET}
\index{MMP\_INT\_SET@\texttt{MMP\_INT\_SET}}
\begin{verbatim}
MMP_INT_SET(windowName, fieldName, value);
\end{verbatim}
Parameters:
\begin{description}
\item[windowName] Logical modeless window name.
\item[fieldName] Named retained field.
\item[value] Value written or selected by the primitive.
\end{description}

Writes a retained modeless integer field.

\paragraph{Example macro.}
\begin{verbatim}
$MACRO REFERENCE_MMP_INT_SET FROM EDIT;
    MMP_INT_SET('reference', 'reference', 0);
END_MACRO;
\end{verbatim}
\end{mrmacextension}

\begin{mrmacextension}
\phantomsection
\label{prim:intrinsic-mmp-int-value}
\subsection*{MMP\_INT\_VALUE}
\index{MMP\_INT\_VALUE@\texttt{MMP\_INT\_VALUE}}
\begin{verbatim}
MMP_INT_VALUE(windowName, fieldName)
\end{verbatim}
Parameters:
\begin{description}
\item[windowName] Logical modeless window name.
\item[fieldName] Named retained field.
\end{description}

Reads the retained integer value of a modeless field.

Returns an integer.

\paragraph{Example macro.}
\begin{verbatim}
$MACRO REFERENCE_MMP_INT_VALUE FROM EDIT;
    DEF_INT(Result);
    Result := MMP_INT_VALUE('reference', 'reference');
END_MACRO;
\end{verbatim}
\end{mrmacextension}

\begin{mrmacextension}
\phantomsection
\label{prim:procedure-mmp-log-append}
\subsection*{MMP\_LOG\_APPEND}
\index{MMP\_LOG\_APPEND@\texttt{MMP\_LOG\_APPEND}}
\begin{verbatim}
MMP_LOG_APPEND(windowName, fieldName, text);
\end{verbatim}
Parameters:
\begin{description}
\item[windowName] Logical modeless window name.
\item[fieldName] Named retained field.
\item[text] String-like text.
\end{description}

Appends text to a retained modeless log.

\paragraph{Example macro.}
\begin{verbatim}
$MACRO REFERENCE_MMP_LOG_APPEND FROM EDIT;
    MMP_LOG_APPEND('reference', 'reference', 'reference');
END_MACRO;
\end{verbatim}
\end{mrmacextension}

\begin{mrmacextension}
\phantomsection
\label{prim:procedure-mmp-log-clear}
\subsection*{MMP\_LOG\_CLEAR}
\index{MMP\_LOG\_CLEAR@\texttt{MMP\_LOG\_CLEAR}}
\begin{verbatim}
MMP_LOG_CLEAR(windowName, fieldName);
\end{verbatim}
Parameters:
\begin{description}
\item[windowName] Logical modeless window name.
\item[fieldName] Named retained field.
\end{description}

Clears a retained modeless log.

\paragraph{Example macro.}
\begin{verbatim}
$MACRO REFERENCE_MMP_LOG_CLEAR FROM EDIT;
    MMP_LOG_CLEAR('reference', 'reference');
END_MACRO;
\end{verbatim}
\end{mrmacextension}

\begin{mrmacextension}
\phantomsection
\label{prim:intrinsic-mmp-log-count}
\subsection*{MMP\_LOG\_COUNT}
\index{MMP\_LOG\_COUNT@\texttt{MMP\_LOG\_COUNT}}
\begin{verbatim}
MMP_LOG_COUNT(windowName, fieldName)
\end{verbatim}
Parameters:
\begin{description}
\item[windowName] Logical modeless window name.
\item[fieldName] Named retained field.
\end{description}

Returns the retained entry count of a modeless log field.

Returns an integer.

\paragraph{Example macro.}
\begin{verbatim}
$MACRO REFERENCE_MMP_LOG_COUNT FROM EDIT;
    DEF_INT(Result);
    Result := MMP_LOG_COUNT('reference', 'reference');
END_MACRO;
\end{verbatim}
\end{mrmacextension}

\begin{mrmacextension}
\phantomsection
\label{prim:procedure-mmp-log-field}
\subsection*{MMP\_LOG\_FIELD}
\index{MMP\_LOG\_FIELD@\texttt{MMP\_LOG\_FIELD}}
\begin{verbatim}
MMP_LOG_FIELD(column, row, width, height, fieldName, label, capacity);
\end{verbatim}
Parameters:
\begin{description}
\item[column] Screen or editor column, as required by the primitive.
\item[row] Zero-based screen row.
\item[width] Control or region width.
\item[height] Control or region height.
\item[fieldName] Named retained field.
\item[label] User-visible field label.
\item[capacity] Maximum retained log entry count.
\end{description}

Adds a bounded retained log field.

\paragraph{Example macro.}
\begin{verbatim}
$MACRO REFERENCE_MMP_LOG_FIELD FROM EDIT;
    MMP_LOG_FIELD(0, 0, 0, 0, 'reference', 'reference', 0);
END_MACRO;
\end{verbatim}
\end{mrmacextension}

\begin{mrmacextension}
\phantomsection
\label{prim:procedure-mmp-menu-clear}
\subsection*{MMP\_MENU\_CLEAR}
\index{MMP\_MENU\_CLEAR@\texttt{MMP\_MENU\_CLEAR}}
\begin{verbatim}
MMP_MENU_CLEAR(menuName);
\end{verbatim}
Parameters:
\begin{description}
\item[menuName] Named menu definition.
\end{description}

Clears a named retained action menu.

\paragraph{Example macro.}
\begin{verbatim}
$MACRO REFERENCE_MMP_MENU_CLEAR FROM EDIT;
    MMP_MENU_CLEAR('reference');
END_MACRO;
\end{verbatim}
\end{mrmacextension}

\begin{mrmacextension}
\phantomsection
\label{prim:procedure-mmp-menu-item}
\subsection*{MMP\_MENU\_ITEM}
\index{MMP\_MENU\_ITEM@\texttt{MMP\_MENU\_ITEM}}
\begin{verbatim}
MMP_MENU_ITEM(menuName, label, value, detail);
\end{verbatim}
Parameters:
\begin{description}
\item[menuName] Named menu definition.
\item[label] User-visible field label.
\item[value] Value written or selected by the primitive.
\item[detail] String or character positional argument.
\end{description}

Adds one label, value and detail row to a retained action menu.

\paragraph{Example macro.}
\begin{verbatim}
$MACRO REFERENCE_MMP_MENU_ITEM FROM EDIT;
    MMP_MENU_ITEM('reference', 'reference', 'reference', 'reference');
END_MACRO;
\end{verbatim}
\end{mrmacextension}

\begin{mrmacextension}
\phantomsection
\label{prim:procedure-mmp-progress-field}
\subsection*{MMP\_PROGRESS\_FIELD}
\index{MMP\_PROGRESS\_FIELD@\texttt{MMP\_PROGRESS\_FIELD}}
\begin{verbatim}
MMP_PROGRESS_FIELD(column, row, width, fieldName, label, maximum, initialValue);
\end{verbatim}
Parameters:
\begin{description}
\item[column] Screen or editor column, as required by the primitive.
\item[row] Zero-based screen row.
\item[width] Control or region width.
\item[fieldName] Named retained field.
\item[label] User-visible field label.
\item[maximum] Inclusive maximum integer value.
\item[initialValue] Initial selected value.
\end{description}

Adds a retained progress display with maximum and initial value.

\paragraph{Example macro.}
\begin{verbatim}
$MACRO REFERENCE_MMP_PROGRESS_FIELD FROM EDIT;
    MMP_PROGRESS_FIELD(0, 0, 0, 'reference', 'reference', 0, 0);
END_MACRO;
\end{verbatim}
\end{mrmacextension}

\begin{mrmacextension}
\phantomsection
\label{prim:procedure-mmp-progress-set}
\subsection*{MMP\_PROGRESS\_SET}
\index{MMP\_PROGRESS\_SET@\texttt{MMP\_PROGRESS\_SET}}
\begin{verbatim}
MMP_PROGRESS_SET(windowName, fieldName, value);
\end{verbatim}
Parameters:
\begin{description}
\item[windowName] Logical modeless window name.
\item[fieldName] Named retained field.
\item[value] Value written or selected by the primitive.
\end{description}

Writes a retained modeless progress value.

\paragraph{Example macro.}
\begin{verbatim}
$MACRO REFERENCE_MMP_PROGRESS_SET FROM EDIT;
    MMP_PROGRESS_SET('reference', 'reference', 0);
END_MACRO;
\end{verbatim}
\end{mrmacextension}

\begin{mrmacextension}
\phantomsection
\label{prim:intrinsic-mmp-progress-value}
\subsection*{MMP\_PROGRESS\_VALUE}
\index{MMP\_PROGRESS\_VALUE@\texttt{MMP\_PROGRESS\_VALUE}}
\begin{verbatim}
MMP_PROGRESS_VALUE(windowName, fieldName)
\end{verbatim}
Parameters:
\begin{description}
\item[windowName] Logical modeless window name.
\item[fieldName] Named retained field.
\end{description}

Reads the retained progress value of a modeless field.

Returns an integer.

\paragraph{Example macro.}
\begin{verbatim}
$MACRO REFERENCE_MMP_PROGRESS_VALUE FROM EDIT;
    DEF_INT(Result);
    Result := MMP_PROGRESS_VALUE('reference', 'reference');
END_MACRO;
\end{verbatim}
\end{mrmacextension}

\begin{mrmacextension}
\phantomsection
\label{prim:procedure-mmp-select-field}
\subsection*{MMP\_SELECT\_FIELD}
\index{MMP\_SELECT\_FIELD@\texttt{MMP\_SELECT\_FIELD}}
\begin{verbatim}
MMP_SELECT_FIELD(column, row, width, height, fieldName, label, initialValue);
\end{verbatim}
Parameters:
\begin{description}
\item[column] Screen or editor column, as required by the primitive.
\item[row] Zero-based screen row.
\item[width] Control or region width.
\item[height] Control or region height.
\item[fieldName] Named retained field.
\item[label] User-visible field label.
\item[initialValue] Initial selected value.
\end{description}

Adds a retained selection field.

\paragraph{Example macro.}
\begin{verbatim}
$MACRO REFERENCE_MMP_SELECT_FIELD FROM EDIT;
    MMP_SELECT_FIELD(0, 0, 0, 0, 'reference', 'reference', 'reference');
END_MACRO;
\end{verbatim}
\end{mrmacextension}

\begin{mrmacextension}
\phantomsection
\label{prim:procedure-mmp-select-option}
\subsection*{MMP\_SELECT\_OPTION}
\index{MMP\_SELECT\_OPTION@\texttt{MMP\_SELECT\_OPTION}}
\begin{verbatim}
MMP_SELECT_OPTION(fieldName, value);
\end{verbatim}
Parameters:
\begin{description}
\item[fieldName] Named retained field.
\item[value] Value written or selected by the primitive.
\end{description}

Adds one selectable option to a retained selection field.

\paragraph{Example macro.}
\begin{verbatim}
$MACRO REFERENCE_MMP_SELECT_OPTION FROM EDIT;
    MMP_SELECT_OPTION('reference', 'reference');
END_MACRO;
\end{verbatim}
\end{mrmacextension}

\begin{mrmacextension}
\phantomsection
\label{prim:procedure-mmp-select-set}
\subsection*{MMP\_SELECT\_SET}
\index{MMP\_SELECT\_SET@\texttt{MMP\_SELECT\_SET}}
\begin{verbatim}
MMP_SELECT_SET(windowName, fieldName, value);
\end{verbatim}
Parameters:
\begin{description}
\item[windowName] Logical modeless window name.
\item[fieldName] Named retained field.
\item[value] Value written or selected by the primitive.
\end{description}

Writes the selected value of a retained selection field.

\paragraph{Example macro.}
\begin{verbatim}
$MACRO REFERENCE_MMP_SELECT_SET FROM EDIT;
    MMP_SELECT_SET('reference', 'reference', 'reference');
END_MACRO;
\end{verbatim}
\end{mrmacextension}

\begin{mrmacextension}
\phantomsection
\label{prim:intrinsic-mmp-select-value}
\subsection*{MMP\_SELECT\_VALUE}
\index{MMP\_SELECT\_VALUE@\texttt{MMP\_SELECT\_VALUE}}
\begin{verbatim}
MMP_SELECT_VALUE(windowName, fieldName)
\end{verbatim}
Parameters:
\begin{description}
\item[windowName] Logical modeless window name.
\item[fieldName] Named retained field.
\end{description}

Reads the selected value of a modeless selection field.

Returns a string.

\paragraph{Example macro.}
\begin{verbatim}
$MACRO REFERENCE_MMP_SELECT_VALUE FROM EDIT;
    DEF_STR(Result);
    Result := MMP_SELECT_VALUE('reference', 'reference');
END_MACRO;
\end{verbatim}
\end{mrmacextension}

\begin{mrmacextension}
\phantomsection
\label{prim:procedure-mmp-status-field}
\subsection*{MMP\_STATUS\_FIELD}
\index{MMP\_STATUS\_FIELD@\texttt{MMP\_STATUS\_FIELD}}
\begin{verbatim}
MMP_STATUS_FIELD(column, row, width, fieldName, initialText);
\end{verbatim}
Parameters:
\begin{description}
\item[column] Screen or editor column, as required by the primitive.
\item[row] Zero-based screen row.
\item[width] Control or region width.
\item[fieldName] Named retained field.
\item[initialText] Initial text value.
\end{description}

Adds a retained status line.

\paragraph{Example macro.}
\begin{verbatim}
$MACRO REFERENCE_MMP_STATUS_FIELD FROM EDIT;
    MMP_STATUS_FIELD(0, 0, 0, 'reference', 'reference');
END_MACRO;
\end{verbatim}
\end{mrmacextension}

\begin{mrmacextension}
\phantomsection
\label{prim:procedure-mmp-status-set}
\subsection*{MMP\_STATUS\_SET}
\index{MMP\_STATUS\_SET@\texttt{MMP\_STATUS\_SET}}
\begin{verbatim}
MMP_STATUS_SET(windowName, fieldName, text);
\end{verbatim}
Parameters:
\begin{description}
\item[windowName] Logical modeless window name.
\item[fieldName] Named retained field.
\item[text] String-like text.
\end{description}

Writes a retained status line.

\paragraph{Example macro.}
\begin{verbatim}
$MACRO REFERENCE_MMP_STATUS_SET FROM EDIT;
    MMP_STATUS_SET('reference', 'reference', 'reference');
END_MACRO;
\end{verbatim}
\end{mrmacextension}

\begin{mrmacextension}
\phantomsection
\label{prim:intrinsic-mmp-style-accent}
\subsection*{MMP\_STYLE\_ACCENT}
\index{MMP\_STYLE\_ACCENT@\texttt{MMP\_STYLE\_ACCENT}}
\begin{verbatim}
MMP_STYLE_ACCENT()
\end{verbatim}
Parameters: None.

Returns the MMP accent style code.

Returns an integer.

\paragraph{Example macro.}
\begin{verbatim}
$MACRO REFERENCE_MMP_STYLE_ACCENT FROM EDIT;
    DEF_INT(Result);
    Result := MMP_STYLE_ACCENT();
END_MACRO;
\end{verbatim}
\end{mrmacextension}

\begin{mrmacextension}
\phantomsection
\label{prim:intrinsic-mmp-style-muted}
\subsection*{MMP\_STYLE\_MUTED}
\index{MMP\_STYLE\_MUTED@\texttt{MMP\_STYLE\_MUTED}}
\begin{verbatim}
MMP_STYLE_MUTED()
\end{verbatim}
Parameters: None.

Returns the MMP muted style code.

Returns an integer.

\paragraph{Example macro.}
\begin{verbatim}
$MACRO REFERENCE_MMP_STYLE_MUTED FROM EDIT;
    DEF_INT(Result);
    Result := MMP_STYLE_MUTED();
END_MACRO;
\end{verbatim}
\end{mrmacextension}

\begin{mrmacextension}
\phantomsection
\label{prim:intrinsic-mmp-style-surface}
\subsection*{MMP\_STYLE\_SURFACE}
\index{MMP\_STYLE\_SURFACE@\texttt{MMP\_STYLE\_SURFACE}}
\begin{verbatim}
MMP_STYLE_SURFACE()
\end{verbatim}
Parameters: None.

Returns the MMP surface style code.

Returns an integer.

\paragraph{Example macro.}
\begin{verbatim}
$MACRO REFERENCE_MMP_STYLE_SURFACE FROM EDIT;
    DEF_INT(Result);
    Result := MMP_STYLE_SURFACE();
END_MACRO;
\end{verbatim}
\end{mrmacextension}

\begin{mrmacextension}
\phantomsection
\label{prim:intrinsic-mmp-style-text}
\subsection*{MMP\_STYLE\_TEXT}
\index{MMP\_STYLE\_TEXT@\texttt{MMP\_STYLE\_TEXT}}
\begin{verbatim}
MMP_STYLE_TEXT()
\end{verbatim}
Parameters: None.

Returns the MMP text style code.

Returns an integer.

\paragraph{Example macro.}
\begin{verbatim}
$MACRO REFERENCE_MMP_STYLE_TEXT FROM EDIT;
    DEF_INT(Result);
    Result := MMP_STYLE_TEXT();
END_MACRO;
\end{verbatim}
\end{mrmacextension}

\begin{mrmacextension}
\phantomsection
\label{prim:procedure-mmp-text-field}
\subsection*{MMP\_TEXT\_FIELD}
\index{MMP\_TEXT\_FIELD@\texttt{MMP\_TEXT\_FIELD}}
\begin{verbatim}
MMP_TEXT_FIELD(column, row, width, fieldName, label, initialText);
\end{verbatim}
Parameters:
\begin{description}
\item[column] Screen or editor column, as required by the primitive.
\item[row] Zero-based screen row.
\item[width] Control or region width.
\item[fieldName] Named retained field.
\item[label] User-visible field label.
\item[initialText] Initial text value.
\end{description}

Adds a retained named text field.

\paragraph{Example macro.}
\begin{verbatim}
$MACRO REFERENCE_MMP_TEXT_FIELD FROM EDIT;
    MMP_TEXT_FIELD(0, 0, 0, 'reference', 'reference', 'reference');
END_MACRO;
\end{verbatim}
\end{mrmacextension}

\begin{mrmacextension}
\phantomsection
\label{prim:procedure-mmp-text-set}
\subsection*{MMP\_TEXT\_SET}
\index{MMP\_TEXT\_SET@\texttt{MMP\_TEXT\_SET}}
\begin{verbatim}
MMP_TEXT_SET(windowName, fieldName, text);
\end{verbatim}
Parameters:
\begin{description}
\item[windowName] Logical modeless window name.
\item[fieldName] Named retained field.
\item[text] String-like text.
\end{description}

Writes a retained modeless text field.

\paragraph{Example macro.}
\begin{verbatim}
$MACRO REFERENCE_MMP_TEXT_SET FROM EDIT;
    MMP_TEXT_SET('reference', 'reference', 'reference');
END_MACRO;
\end{verbatim}
\end{mrmacextension}

\begin{mrmacextension}
\phantomsection
\label{prim:intrinsic-mmp-text-value}
\subsection*{MMP\_TEXT\_VALUE}
\index{MMP\_TEXT\_VALUE@\texttt{MMP\_TEXT\_VALUE}}
\begin{verbatim}
MMP_TEXT_VALUE(windowName, fieldName)
\end{verbatim}
Parameters:
\begin{description}
\item[windowName] Logical modeless window name.
\item[fieldName] Named retained field.
\end{description}

Reads the retained text value of a modeless text field.

Returns a string.

\paragraph{Example macro.}
\begin{verbatim}
$MACRO REFERENCE_MMP_TEXT_VALUE FROM EDIT;
    DEF_STR(Result);
    Result := MMP_TEXT_VALUE('reference', 'reference');
END_MACRO;
\end{verbatim}
\end{mrmacextension}

\begin{mrmacextension}
\phantomsection
\label{prim:procedure-mmp-timer-start}
\subsection*{MMP\_TIMER\_START}
\index{MMP\_TIMER\_START@\texttt{MMP\_TIMER\_START}}
\begin{verbatim}
MMP_TIMER_START(windowName, timerName, milliseconds, macroSpec);
\end{verbatim}
Parameters:
\begin{description}
\item[windowName] Logical modeless window name.
\item[timerName] String or character positional argument.
\item[milliseconds] Positive delay or timer interval in milliseconds.
\item[macroSpec] Macro specification or callback target.
\end{description}

Starts a named logical modeless window timer and callback.

\paragraph{Example macro.}
\begin{verbatim}
$MACRO REFERENCE_MMP_TIMER_START FROM EDIT;
    MMP_TIMER_START('reference', 'reference', 0, 'reference');
END_MACRO;
\end{verbatim}
\end{mrmacextension}

\begin{mrmacextension}
\phantomsection
\label{prim:procedure-mmp-timer-stop}
\subsection*{MMP\_TIMER\_STOP}
\index{MMP\_TIMER\_STOP@\texttt{MMP\_TIMER\_STOP}}
\begin{verbatim}
MMP_TIMER_STOP(windowName, timerName);
\end{verbatim}
Parameters:
\begin{description}
\item[windowName] Logical modeless window name.
\item[timerName] String or character positional argument.
\end{description}

Stops a named logical modeless window timer.

\paragraph{Example macro.}
\begin{verbatim}
$MACRO REFERENCE_MMP_TIMER_STOP FROM EDIT;
    MMP_TIMER_STOP('reference', 'reference');
END_MACRO;
\end{verbatim}
\end{mrmacextension}

\begin{mrmacextension}
\phantomsection
\label{prim:intrinsic-mmp-window-exists}
\subsection*{MMP\_WINDOW\_EXISTS}
\index{MMP\_WINDOW\_EXISTS@\texttt{MMP\_WINDOW\_EXISTS}}
\begin{verbatim}
MMP_WINDOW_EXISTS(windowName)
\end{verbatim}
Parameters:
\begin{description}
\item[windowName] Logical modeless window name.
\end{description}

Reports whether a named logical modeless window exists.

Returns an integer.

\paragraph{Example macro.}
\begin{verbatim}
$MACRO REFERENCE_MMP_WINDOW_EXISTS FROM EDIT;
    DEF_INT(Result);
    Result := MMP_WINDOW_EXISTS('reference');
END_MACRO;
\end{verbatim}
\end{mrmacextension}

\begin{mrmacextension}
\phantomsection
\label{prim:intrinsic-mmp-window-height}
\subsection*{MMP\_WINDOW\_HEIGHT}
\index{MMP\_WINDOW\_HEIGHT@\texttt{MMP\_WINDOW\_HEIGHT}}
\begin{verbatim}
MMP_WINDOW_HEIGHT(windowName)
\end{verbatim}
Parameters:
\begin{description}
\item[windowName] Logical modeless window name.
\end{description}

Returns the logical height of the named modeless window.

Returns an integer.

\paragraph{Example macro.}
\begin{verbatim}
$MACRO REFERENCE_MMP_WINDOW_HEIGHT FROM EDIT;
    DEF_INT(Result);
    Result := MMP_WINDOW_HEIGHT('reference');
END_MACRO;
\end{verbatim}
\end{mrmacextension}

\begin{mrmacextension}
\phantomsection
\label{prim:intrinsic-mmp-window-instance}
\subsection*{MMP\_WINDOW\_INSTANCE}
\index{MMP\_WINDOW\_INSTANCE@\texttt{MMP\_WINDOW\_INSTANCE}}
\begin{verbatim}
MMP_WINDOW_INSTANCE(windowName)
\end{verbatim}
Parameters:
\begin{description}
\item[windowName] Logical modeless window name.
\end{description}

Returns the current instance identity of the named modeless window.

Returns a string.

\paragraph{Example macro.}
\begin{verbatim}
$MACRO REFERENCE_MMP_WINDOW_INSTANCE FROM EDIT;
    DEF_STR(Result);
    Result := MMP_WINDOW_INSTANCE('reference');
END_MACRO;
\end{verbatim}
\end{mrmacextension}

\begin{mrmacextension}
\phantomsection
\label{prim:intrinsic-mmp-window-width}
\subsection*{MMP\_WINDOW\_WIDTH}
\index{MMP\_WINDOW\_WIDTH@\texttt{MMP\_WINDOW\_WIDTH}}
\begin{verbatim}
MMP_WINDOW_WIDTH(windowName)
\end{verbatim}
Parameters:
\begin{description}
\item[windowName] Logical modeless window name.
\end{description}

Returns the logical width of the named modeless window.

Returns an integer.

\paragraph{Example macro.}
\begin{verbatim}
$MACRO REFERENCE_MMP_WINDOW_WIDTH FROM EDIT;
    DEF_INT(Result);
    Result := MMP_WINDOW_WIDTH('reference');
END_MACRO;
\end{verbatim}
\end{mrmacextension}

\begin{mrmacextension}
\phantomsection
\label{prim:intrinsic-mmp-window-x}
\subsection*{MMP\_WINDOW\_X}
\index{MMP\_WINDOW\_X@\texttt{MMP\_WINDOW\_X}}
\begin{verbatim}
MMP_WINDOW_X(windowName)
\end{verbatim}
Parameters:
\begin{description}
\item[windowName] Logical modeless window name.
\end{description}

Returns the logical x coordinate of the named modeless window.

Returns an integer.

\paragraph{Example macro.}
\begin{verbatim}
$MACRO REFERENCE_MMP_WINDOW_X FROM EDIT;
    DEF_INT(Result);
    Result := MMP_WINDOW_X('reference');
END_MACRO;
\end{verbatim}
\end{mrmacextension}

\begin{mrmacextension}
\phantomsection
\label{prim:intrinsic-mmp-window-y}
\subsection*{MMP\_WINDOW\_Y}
\index{MMP\_WINDOW\_Y@\texttt{MMP\_WINDOW\_Y}}
\begin{verbatim}
MMP_WINDOW_Y(windowName)
\end{verbatim}
Parameters:
\begin{description}
\item[windowName] Logical modeless window name.
\end{description}

Returns the logical y coordinate of the named modeless window.

Returns an integer.

\paragraph{Example macro.}
\begin{verbatim}
$MACRO REFERENCE_MMP_WINDOW_Y FROM EDIT;
    DEF_INT(Result);
    Result := MMP_WINDOW_Y('reference');
END_MACRO;
\end{verbatim}
\end{mrmacextension}


\part{Runtime and Integration}

\chapter{Globals, Catalogs and Execution Sessions}

Local declarations belong to one running macro. Explicit macro globals, the loaded macro catalogue and execution sessions are different runtime concepts. This chapter covers their controlled primitive interfaces and enumerators.

\section{Primitive Map}
The map is a local index for this chapter. The final column is a generated page
reference and hyperlink to the full entry, where the source form, parameters,
semantics and a complete example macro appear.

\begin{longtable}{@{}L{0.38\textwidth}L{0.14\textwidth}L{0.27\textwidth}L{0.10\textwidth}@{}}
\toprule
Primitive & Role & Purpose & Reference \\
\midrule
\endhead
\toprule
Primitive & Role & Purpose & Reference \\
\midrule
\endhead
\primitiveoverview{FIRST\_GLOBAL}{intrinsic-first-global}{intrinsic}{Macro runtime state.}
\primitiveoverview{NEXT\_GLOBAL}{intrinsic-next-global}{intrinsic}{Macro runtime state.}
\primitiveoverview{GLOBAL\_HASH}{intrinsic-global-hash}{intrinsic}{Macro runtime state.}
\primitiveoverview{GLOBAL\_INT}{intrinsic-global-int}{intrinsic}{Macro runtime state.}
\primitiveoverview{GLOBAL\_STR}{intrinsic-global-str}{intrinsic}{Macro runtime state.}
\primitiveoverview{INQ\_MACRO}{intrinsic-inq-macro}{intrinsic}{Macro runtime state.}
\primitiveoverview{CREATE\_GLOBAL\_STR}{procedure-create-global-str}{procedure}{Macro runtime state.}
\primitiveoverview{DELAY}{procedure-delay}{procedure}{Macro runtime state.}
\primitiveoverview{EXEC\_ASSIGN}{procedure-exec-assign}{procedure}{Macro runtime state.}
\primitiveoverview{EXEC\_SESSION\_LIST}{procedure-exec-session-list}{procedure}{Macro runtime state.}
\primitiveoverview{EXEC\_SESSION\_STOP}{procedure-exec-session-stop}{procedure}{Macro runtime state.}
\primitiveoverview{LOAD\_MACRO\_FILE}{procedure-load-macro-file}{procedure}{Macro runtime state.}
\primitiveoverview{RUN\_MACRO}{procedure-run-macro}{procedure}{Macro runtime state.}
\primitiveoverview{SET\_GLOBAL\_HASH}{procedure-set-global-hash}{procedure}{Macro runtime state.}
\primitiveoverview{SET\_GLOBAL\_INT}{procedure-set-global-int}{procedure}{Macro runtime state.}
\primitiveoverview{SET\_GLOBAL\_STR}{procedure-set-global-str}{procedure}{Macro runtime state.}
\primitiveoverview{UNLOAD\_MACRO}{procedure-unload-macro}{procedure}{Macro runtime state.}
\bottomrule
\end{longtable}

\section{Reference Entries}

\begin{mrmacextension}
\phantomsection
\label{prim:intrinsic-first-global}
\subsection*{FIRST\_GLOBAL}
\index{FIRST\_GLOBAL@\texttt{FIRST\_GLOBAL}}
\begin{verbatim}
FIRST_GLOBAL(cursor)
\end{verbatim}
Parameters:
\begin{description}
\item[cursor] Integer enumeration cursor variable.
\end{description}

Starts global-name enumeration with cursor and returns the first name.

Returns a string.

\end{mrmacextension}

\paragraph{Example macro.}
\begin{verbatim}
$MACRO REFERENCE_FIRST_GLOBAL FROM EDIT;
    DEF_INT(Cursor);
    DEF_STR(Result);
    Result := FIRST_GLOBAL(Cursor);
END_MACRO;
\end{verbatim}

\begin{mrmacextension}
\phantomsection
\label{prim:intrinsic-next-global}
\subsection*{NEXT\_GLOBAL}
\index{NEXT\_GLOBAL@\texttt{NEXT\_GLOBAL}}
\begin{verbatim}
NEXT_GLOBAL(cursor)
\end{verbatim}
Parameters:
\begin{description}
\item[cursor] Integer enumeration cursor variable.
\end{description}

Continues global-name enumeration with cursor and returns the next name.

Returns a string.

\paragraph{Example macro.}
\begin{verbatim}
$MACRO REFERENCE_NEXT_GLOBAL FROM EDIT;
    DEF_INT(Cursor);
    DEF_STR(Result);
    Result := NEXT_GLOBAL(Cursor);
END_MACRO;
\end{verbatim}
\end{mrmacextension}

\begin{mrmacextension}
\phantomsection
\label{prim:intrinsic-global-hash}
\subsection*{GLOBAL\_HASH}
\index{GLOBAL\_HASH@\texttt{GLOBAL\_HASH}}
\begin{verbatim}
GLOBAL_HASH(name)
\end{verbatim}
Parameters:
\begin{description}
\item[name] String-like global name.
\end{description}

Reads a named hash global.

Returns a hash.

\paragraph{Example macro.}
\begin{verbatim}
$MACRO REFERENCE_GLOBAL_HASH FROM EDIT;
    DEF_HASH(State);
    DEF_HASH(Result);
    Result := GLOBAL_HASH('reference');
END_MACRO;
\end{verbatim}
\end{mrmacextension}

\begin{mrmacextension}
\phantomsection
\label{prim:intrinsic-global-int}
\subsection*{GLOBAL\_INT}
\index{GLOBAL\_INT@\texttt{GLOBAL\_INT}}
\begin{verbatim}
GLOBAL_INT(name)
\end{verbatim}
Parameters:
\begin{description}
\item[name] String-like global name.
\end{description}

Reads a named integer global.

Returns an integer.

\paragraph{Example macro.}
\begin{verbatim}
$MACRO REFERENCE_GLOBAL_INT FROM EDIT;
    DEF_INT(Result);
    Result := GLOBAL_INT('reference');
END_MACRO;
\end{verbatim}
\end{mrmacextension}

\begin{mrmacextension}
\phantomsection
\label{prim:intrinsic-global-str}
\subsection*{GLOBAL\_STR}
\index{GLOBAL\_STR@\texttt{GLOBAL\_STR}}
\begin{verbatim}
GLOBAL_STR(name)
\end{verbatim}
Parameters:
\begin{description}
\item[name] String-like global name.
\end{description}

Reads a named string global.

Returns a string.

\paragraph{Example macro.}
\begin{verbatim}
$MACRO REFERENCE_GLOBAL_STR FROM EDIT;
    DEF_STR(Result);
    Result := GLOBAL_STR('reference');
END_MACRO;
\end{verbatim}
\end{mrmacextension}

\phantomsection
\label{prim:intrinsic-inq-macro}
\subsection*{INQ\_MACRO}
\index{INQ\_MACRO@\texttt{INQ\_MACRO}}
\begin{verbatim}
INQ_MACRO(macroSpec)
\end{verbatim}
Parameters:
\begin{description}
\item[macroSpec] Macro specification or callback target.
\end{description}

Reports whether the macro catalog contains macroSpec.

Returns an integer.

\paragraph{Example macro.}
\begin{verbatim}
$MACRO REFERENCE_INQ_MACRO FROM EDIT;
    DEF_INT(Result);
    Result := INQ_MACRO('reference');
END_MACRO;
\end{verbatim}


\begin{mrmacextension}
\phantomsection
\label{prim:procedure-create-global-str}
\subsection*{CREATE\_GLOBAL\_STR}
\index{CREATE\_GLOBAL\_STR@\texttt{CREATE\_GLOBAL\_STR}}
\begin{verbatim}
CREATE_GLOBAL_STR(name, value);
\end{verbatim}
Parameters:
\begin{description}
\item[name] String-like global name.
\item[value] Value written or selected by the primitive.
\end{description}

Creates or replaces a named process-wide string global.

\paragraph{Example macro.}
\begin{verbatim}
$MACRO REFERENCE_CREATE_GLOBAL_STR FROM EDIT;
    CREATE_GLOBAL_STR('reference', 'reference');
END_MACRO;
\end{verbatim}
\end{mrmacextension}

\begin{mrmacextension}
\phantomsection
\label{prim:procedure-delay}
\subsection*{DELAY}
\index{DELAY@\texttt{DELAY}}
\begin{verbatim}
DELAY(milliseconds);
\end{verbatim}
Parameters:
\begin{description}
\item[milliseconds] Positive delay or timer interval in milliseconds.
\end{description}

Yields macro execution for the requested number of milliseconds.

\paragraph{Example macro.}
\begin{verbatim}
$MACRO REFERENCE_DELAY FROM EDIT;
    DELAY(0);
END_MACRO;
\end{verbatim}
\end{mrmacextension}

\begin{mrmacextension}
\phantomsection
\label{prim:procedure-exec-assign}
\subsection*{EXEC\_ASSIGN}
\index{EXEC\_ASSIGN@\texttt{EXEC\_ASSIGN}}
\begin{verbatim}
EXEC_ASSIGN(actionName, macroSpec, label);
\end{verbatim}
Parameters:
\begin{description}
\item[actionName] Named execution action.
\item[macroSpec] Macro specification or callback target.
\item[label] User-visible field label.
\end{description}

Assigns macroSpec to a named execution action.

\paragraph{Example macro.}
\begin{verbatim}
$MACRO REFERENCE_EXEC_ASSIGN FROM EDIT;
    EXEC_ASSIGN('reference', 'reference', 'reference');
END_MACRO;
\end{verbatim}
\end{mrmacextension}

\begin{mrmacextension}
\phantomsection
\label{prim:procedure-exec-session-list}
\subsection*{EXEC\_SESSION\_LIST}
\index{EXEC\_SESSION\_LIST@\texttt{EXEC\_SESSION\_LIST}}
\begin{verbatim}
EXEC_SESSION_LIST(targetName);
\end{verbatim}
Parameters:
\begin{description}
\item[targetName] String or character positional argument.
\end{description}

Publishes execution-session information to targetName.

\paragraph{Example macro.}
\begin{verbatim}
$MACRO REFERENCE_EXEC_SESSION_LIST FROM EDIT;
    EXEC_SESSION_LIST('reference');
END_MACRO;
\end{verbatim}
\end{mrmacextension}

\begin{mrmacextension}
\phantomsection
\label{prim:procedure-exec-session-stop}
\subsection*{EXEC\_SESSION\_STOP}
\index{EXEC\_SESSION\_STOP@\texttt{EXEC\_SESSION\_STOP}}
\begin{verbatim}
EXEC_SESSION_STOP(sessionId);
\end{verbatim}
Parameters:
\begin{description}
\item[sessionId] Integer execution session identifier.
\end{description}

Requests cancellation of the named execution session.

\paragraph{Example macro.}
\begin{verbatim}
$MACRO REFERENCE_EXEC_SESSION_STOP FROM EDIT;
    EXEC_SESSION_STOP(0);
END_MACRO;
\end{verbatim}
\end{mrmacextension}

\phantomsection
\label{prim:procedure-load-macro-file}
\subsection*{LOAD\_MACRO\_FILE}
\index{LOAD\_MACRO\_FILE@\texttt{LOAD\_MACRO\_FILE}}
\begin{verbatim}
LOAD_MACRO_FILE(path);
\end{verbatim}
Parameters:
\begin{description}
\item[path] A string-like file-system path.
\end{description}

Loads and compiles a macro source file into the catalog.

\paragraph{Example macro.}
\begin{verbatim}
$MACRO REFERENCE_LOAD_MACRO_FILE FROM EDIT;
    LOAD_MACRO_FILE('reference');
END_MACRO;
\end{verbatim}


\phantomsection
\label{prim:procedure-run-macro}
\subsection*{RUN\_MACRO}
\index{RUN\_MACRO@\texttt{RUN\_MACRO}}
\begin{verbatim}
RUN_MACRO(macroSpec);
\end{verbatim}
Parameters:
\begin{description}
\item[macroSpec] Macro specification or callback target.
\end{description}

Runs a macro selected by macroSpec from the loaded catalog.

\paragraph{Example macro.}
\begin{verbatim}
$MACRO REFERENCE_RUN_MACRO FROM EDIT;
    RUN_MACRO('reference');
END_MACRO;
\end{verbatim}


\begin{mrmacextension}
\phantomsection
\label{prim:procedure-set-global-hash}
\subsection*{SET\_GLOBAL\_HASH}
\index{SET\_GLOBAL\_HASH@\texttt{SET\_GLOBAL\_HASH}}
\begin{verbatim}
SET_GLOBAL_HASH(name, value);
\end{verbatim}
Parameters:
\begin{description}
\item[name] String-like global name.
\item[value] Value written or selected by the primitive.
\end{description}

Writes a named process-wide hash global.

\paragraph{Example macro.}
\begin{verbatim}
$MACRO REFERENCE_SET_GLOBAL_HASH FROM EDIT;
    DEF_HASH(State);
    SET_GLOBAL_HASH('reference', State);
END_MACRO;
\end{verbatim}
\end{mrmacextension}

\begin{mrmacextension}
\phantomsection
\label{prim:procedure-set-global-int}
\subsection*{SET\_GLOBAL\_INT}
\index{SET\_GLOBAL\_INT@\texttt{SET\_GLOBAL\_INT}}
\begin{verbatim}
SET_GLOBAL_INT(name, value);
\end{verbatim}
Parameters:
\begin{description}
\item[name] String-like global name.
\item[value] Value written or selected by the primitive.
\end{description}

Writes a named process-wide integer global.

\paragraph{Example macro.}
\begin{verbatim}
$MACRO REFERENCE_SET_GLOBAL_INT FROM EDIT;
    SET_GLOBAL_INT('reference', 0);
END_MACRO;
\end{verbatim}
\end{mrmacextension}

\begin{mrmacextension}
\phantomsection
\label{prim:procedure-set-global-str}
\subsection*{SET\_GLOBAL\_STR}
\index{SET\_GLOBAL\_STR@\texttt{SET\_GLOBAL\_STR}}
\begin{verbatim}
SET_GLOBAL_STR(name, value);
\end{verbatim}
Parameters:
\begin{description}
\item[name] String-like global name.
\item[value] Value written or selected by the primitive.
\end{description}

Writes a named process-wide string global.

\paragraph{Example macro.}
\begin{verbatim}
$MACRO REFERENCE_SET_GLOBAL_STR FROM EDIT;
    SET_GLOBAL_STR('reference', 'reference');
END_MACRO;
\end{verbatim}
\end{mrmacextension}

\phantomsection
\label{prim:procedure-unload-macro}
\subsection*{UNLOAD\_MACRO}
\index{UNLOAD\_MACRO@\texttt{UNLOAD\_MACRO}}
\begin{verbatim}
UNLOAD_MACRO(macroSpec);
\end{verbatim}
Parameters:
\begin{description}
\item[macroSpec] Macro specification or callback target.
\end{description}

Removes macroSpec from the loaded macro catalog.

\paragraph{Example macro.}
\begin{verbatim}
$MACRO REFERENCE_UNLOAD_MACRO FROM EDIT;
    UNLOAD_MACRO('reference');
END_MACRO;
\end{verbatim}


\chapter{Settings, Keymaps and Themes}

Settings and serialized UI descriptions cross a controlled VM boundary. They are not a macro-owned configuration store. This chapter documents the accepted settings, keymap and theme source forms.

\section{Primitive Map}
The map is a local index for this chapter. The final column is a generated page
reference and hyperlink to the full entry, where the source form, parameters,
semantics and a complete example macro appear.

\begin{longtable}{@{}L{0.38\textwidth}L{0.14\textwidth}L{0.27\textwidth}L{0.10\textwidth}@{}}
\toprule
Primitive & Role & Purpose & Reference \\
\midrule
\endhead
\toprule
Primitive & Role & Purpose & Reference \\
\midrule
\endhead
\primitiveoverview{CODECOLORS}{procedure-codecolors}{procedure}{Settings source.}
\primitiveoverview{DEBUGGERCOLORS}{procedure-debuggercolors}{procedure}{Settings source.}
\primitiveoverview{FILECOMPARECOLORS}{procedure-filecomparecolors}{procedure}{Settings source.}
\primitiveoverview{FILECOMPAREMINIMAPCOLORS}{procedure-filecompareminimapcolors}{procedure}{Settings source.}
\primitiveoverview{HELPCOLORS}{procedure-helpcolors}{procedure}{Settings source.}
\primitiveoverview{KEYMAP\_BIND}{procedure-keymap-bind}{procedure}{Settings source.}
\primitiveoverview{KEYMAP\_PROFILE}{procedure-keymap-profile}{procedure}{Settings source.}
\primitiveoverview{KEYMAP\_RESET}{procedure-keymap-reset}{procedure}{Settings source.}
\primitiveoverview{KEYMAP\_VERSION}{procedure-keymap-version}{procedure}{Settings source.}
\primitiveoverview{MENUDIALOGCOLORS}{procedure-menudialogcolors}{procedure}{Settings source.}
\primitiveoverview{MINIMAPCOLORS}{procedure-minimapcolors}{procedure}{Settings source.}
\primitiveoverview{MRCOMPILERPROFILE}{procedure-mrcompilerprofile}{procedure}{Settings source.}
\primitiveoverview{MRFEPROFILE}{procedure-mrfeprofile}{procedure}{Settings source.}
\primitiveoverview{MRSETUP}{procedure-mrsetup}{procedure}{Settings source.}
\primitiveoverview{OTHERCOLORS}{procedure-othercolors}{procedure}{Settings source.}
\primitiveoverview{SAVE\_SETTINGS}{procedure-save-settings}{procedure}{Settings source.}
\primitiveoverview{THEME\_NAME}{procedure-theme-name}{procedure}{Settings source.}
\primitiveoverview{THEME\_RESET}{procedure-theme-reset}{procedure}{Settings source.}
\primitiveoverview{THEME\_VERSION}{procedure-theme-version}{procedure}{Settings source.}
\primitiveoverview{WINDOWCOLORS}{procedure-windowcolors}{procedure}{Settings source.}
\bottomrule
\end{longtable}

\section{Reference Entries}

\begin{mrmacextension}
\phantomsection
\label{prim:procedure-codecolors}
\subsection*{CODECOLORS}
\index{CODECOLORS@\texttt{CODECOLORS}}
\begin{verbatim}
CODECOLORS(definition);
\end{verbatim}
Parameters:
\begin{description}
\item[definition] Serialized keymap or theme definition.
\end{description}

Adds the named serialized theme colour definition.

\end{mrmacextension}

\paragraph{Example macro.}
\begin{verbatim}
$MACRO REFERENCE_CODECOLORS FROM EDIT;
    CODECOLORS('reference');
END_MACRO;
\end{verbatim}

\begin{mrmacextension}
\phantomsection
\label{prim:procedure-debuggercolors}
\subsection*{DEBUGGERCOLORS}
\index{DEBUGGERCOLORS@\texttt{DEBUGGERCOLORS}}
\begin{verbatim}
DEBUGGERCOLORS(definition);
\end{verbatim}
Parameters:
\begin{description}
\item[definition] Serialized keymap or theme definition.
\end{description}

Adds the named serialized theme colour definition.

\paragraph{Example macro.}
\begin{verbatim}
$MACRO REFERENCE_DEBUGGERCOLORS FROM EDIT;
    DEBUGGERCOLORS('reference');
END_MACRO;
\end{verbatim}
\end{mrmacextension}

\begin{mrmacextension}
\phantomsection
\label{prim:procedure-filecomparecolors}
\subsection*{FILECOMPARECOLORS}
\index{FILECOMPARECOLORS@\texttt{FILECOMPARECOLORS}}
\begin{verbatim}
FILECOMPARECOLORS(definition);
\end{verbatim}
Parameters:
\begin{description}
\item[definition] Serialized keymap or theme definition.
\end{description}

Adds the named serialized theme colour definition.

\paragraph{Example macro.}
\begin{verbatim}
$MACRO REFERENCE_FILECOMPARECOLORS FROM EDIT;
    FILECOMPARECOLORS('reference');
END_MACRO;
\end{verbatim}
\end{mrmacextension}

\begin{mrmacextension}
\phantomsection
\label{prim:procedure-filecompareminimapcolors}
\subsection*{FILECOMPAREMINIMAPCOLORS}
\index{FILECOMPAREMINIMAPCOLORS@\texttt{FILECOMPAREMINIMAPCOLORS}}
\begin{verbatim}
FILECOMPAREMINIMAPCOLORS(definition);
\end{verbatim}
Parameters:
\begin{description}
\item[definition] Serialized keymap or theme definition.
\end{description}

Adds the named serialized theme colour definition.

\paragraph{Example macro.}
\begin{verbatim}
$MACRO REFERENCE_FILECOMPAREMINIMAPCOLORS FROM EDIT;
    FILECOMPAREMINIMAPCOLORS('reference');
END_MACRO;
\end{verbatim}
\end{mrmacextension}

\begin{mrmacextension}
\phantomsection
\label{prim:procedure-helpcolors}
\subsection*{HELPCOLORS}
\index{HELPCOLORS@\texttt{HELPCOLORS}}
\begin{verbatim}
HELPCOLORS(definition);
\end{verbatim}
Parameters:
\begin{description}
\item[definition] Serialized keymap or theme definition.
\end{description}

Adds the named serialized theme colour definition.

\paragraph{Example macro.}
\begin{verbatim}
$MACRO REFERENCE_HELPCOLORS FROM EDIT;
    HELPCOLORS('reference');
END_MACRO;
\end{verbatim}
\end{mrmacextension}

\begin{mrmacextension}
\phantomsection
\label{prim:procedure-keymap-bind}
\subsection*{KEYMAP\_BIND}
\index{KEYMAP\_BIND@\texttt{KEYMAP\_BIND}}
\begin{verbatim}
KEYMAP_BIND(binding);
\end{verbatim}
Parameters:
\begin{description}
\item[binding] String or character positional argument.
\end{description}

Adds one serialized keymap binding definition.

\paragraph{Example macro.}
\begin{verbatim}
$MACRO REFERENCE_KEYMAP_BIND FROM EDIT;
    KEYMAP_BIND('reference');
END_MACRO;
\end{verbatim}
\end{mrmacextension}

\begin{mrmacextension}
\phantomsection
\label{prim:procedure-keymap-profile}
\subsection*{KEYMAP\_PROFILE}
\index{KEYMAP\_PROFILE@\texttt{KEYMAP\_PROFILE}}
\begin{verbatim}
KEYMAP_PROFILE(profileName);
\end{verbatim}
Parameters:
\begin{description}
\item[profileName] Serialized profile name.
\end{description}

Selects or defines the named serialized keymap profile.

\paragraph{Example macro.}
\begin{verbatim}
$MACRO REFERENCE_KEYMAP_PROFILE FROM EDIT;
    KEYMAP_PROFILE('reference');
END_MACRO;
\end{verbatim}
\end{mrmacextension}

\begin{mrmacextension}
\phantomsection
\label{prim:procedure-keymap-reset}
\subsection*{KEYMAP\_RESET}
\index{KEYMAP\_RESET@\texttt{KEYMAP\_RESET}}
\begin{verbatim}
KEYMAP_RESET();
\end{verbatim}
Parameters: None.

Resets the serialized keymap source currently being assembled.

\paragraph{Example macro.}
\begin{verbatim}
$MACRO REFERENCE_KEYMAP_RESET FROM EDIT;
    KEYMAP_RESET();
END_MACRO;
\end{verbatim}
\end{mrmacextension}

\begin{mrmacextension}
\phantomsection
\label{prim:procedure-keymap-version}
\subsection*{KEYMAP\_VERSION}
\index{KEYMAP\_VERSION@\texttt{KEYMAP\_VERSION}}
\begin{verbatim}
KEYMAP_VERSION(version);
\end{verbatim}
Parameters:
\begin{description}
\item[version] Serialized version text.
\end{description}

Sets the keymap source version.

\paragraph{Example macro.}
\begin{verbatim}
$MACRO REFERENCE_KEYMAP_VERSION FROM EDIT;
    KEYMAP_VERSION('reference');
END_MACRO;
\end{verbatim}
\end{mrmacextension}

\begin{mrmacextension}
\phantomsection
\label{prim:procedure-menudialogcolors}
\subsection*{MENUDIALOGCOLORS}
\index{MENUDIALOGCOLORS@\texttt{MENUDIALOGCOLORS}}
\begin{verbatim}
MENUDIALOGCOLORS(definition);
\end{verbatim}
Parameters:
\begin{description}
\item[definition] Serialized keymap or theme definition.
\end{description}

Adds the named serialized theme colour definition.

\paragraph{Example macro.}
\begin{verbatim}
$MACRO REFERENCE_MENUDIALOGCOLORS FROM EDIT;
    MENUDIALOGCOLORS('reference');
END_MACRO;
\end{verbatim}
\end{mrmacextension}

\begin{mrmacextension}
\phantomsection
\label{prim:procedure-minimapcolors}
\subsection*{MINIMAPCOLORS}
\index{MINIMAPCOLORS@\texttt{MINIMAPCOLORS}}
\begin{verbatim}
MINIMAPCOLORS(definition);
\end{verbatim}
Parameters:
\begin{description}
\item[definition] Serialized keymap or theme definition.
\end{description}

Adds the named serialized theme colour definition.

\paragraph{Example macro.}
\begin{verbatim}
$MACRO REFERENCE_MINIMAPCOLORS FROM EDIT;
    MINIMAPCOLORS('reference');
END_MACRO;
\end{verbatim}
\end{mrmacextension}

\begin{mrmacextension}
\phantomsection
\label{prim:procedure-mrcompilerprofile}
\subsection*{MRCOMPILERPROFILE}
\index{MRCOMPILERPROFILE@\texttt{MRCOMPILERPROFILE}}
\begin{verbatim}
MRCOMPILERPROFILE(profileName, compilerPath, arguments, workingDirectory);
\end{verbatim}
Parameters:
\begin{description}
\item[profileName] Serialized profile name.
\item[compilerPath] String or character positional argument.
\item[arguments] String or character positional argument.
\item[workingDirectory] String or character positional argument.
\end{description}

Defines one serialized compiler profile.

\paragraph{Example macro.}
\begin{verbatim}
$MACRO REFERENCE_MRCOMPILERPROFILE FROM EDIT;
    MRCOMPILERPROFILE('reference', 'reference', 'reference', 'reference');
END_MACRO;
\end{verbatim}
\end{mrmacextension}

\begin{mrmacextension}
\phantomsection
\label{prim:procedure-mrfeprofile}
\subsection*{MRFEPROFILE}
\index{MRFEPROFILE@\texttt{MRFEPROFILE}}
\begin{verbatim}
MRFEPROFILE(profileName, sourcePath, outputPath, arguments);
\end{verbatim}
Parameters:
\begin{description}
\item[profileName] Serialized profile name.
\item[sourcePath] Source file path.
\item[outputPath] String or character positional argument.
\item[arguments] String or character positional argument.
\end{description}

Defines one serialized editor profile.

\paragraph{Example macro.}
\begin{verbatim}
$MACRO REFERENCE_MRFEPROFILE FROM EDIT;
    MRFEPROFILE('reference', 'reference', 'reference', 'reference');
END_MACRO;
\end{verbatim}
\end{mrmacextension}

\begin{mrmacextension}
\phantomsection
\label{prim:procedure-mrsetup}
\subsection*{MRSETUP}
\index{MRSETUP@\texttt{MRSETUP}}
\begin{verbatim}
MRSETUP(key, value);
\end{verbatim}
Parameters:
\begin{description}
\item[key] String-like hash key.
\item[value] Value written or selected by the primitive.
\end{description}

Accepts one serialized settings key and value only at the controlled settings boundary.

\paragraph{Example macro.}
\begin{verbatim}
$MACRO REFERENCE_MRSETUP FROM EDIT;
    MRSETUP('reference', 'reference');
END_MACRO;
\end{verbatim}
\end{mrmacextension}

\begin{mrmacextension}
\phantomsection
\label{prim:procedure-othercolors}
\subsection*{OTHERCOLORS}
\index{OTHERCOLORS@\texttt{OTHERCOLORS}}
\begin{verbatim}
OTHERCOLORS(definition);
\end{verbatim}
Parameters:
\begin{description}
\item[definition] Serialized keymap or theme definition.
\end{description}

Adds the named serialized theme colour definition.

\paragraph{Example macro.}
\begin{verbatim}
$MACRO REFERENCE_OTHERCOLORS FROM EDIT;
    OTHERCOLORS('reference');
END_MACRO;
\end{verbatim}
\end{mrmacextension}

\begin{mrmacextension}
\phantomsection
\label{prim:procedure-save-settings}
\subsection*{SAVE\_SETTINGS}
\index{SAVE\_SETTINGS@\texttt{SAVE\_SETTINGS}}
\begin{verbatim}
SAVE_SETTINGS;
\end{verbatim}
Parameters: None.

Persists the current runtime settings model through the VM boundary.

\paragraph{Example macro.}
\begin{verbatim}
$MACRO REFERENCE_SAVE_SETTINGS FROM EDIT;
    SAVE_SETTINGS;
END_MACRO;
\end{verbatim}
\end{mrmacextension}

\begin{mrmacextension}
\phantomsection
\label{prim:procedure-theme-name}
\subsection*{THEME\_NAME}
\index{THEME\_NAME@\texttt{THEME\_NAME}}
\begin{verbatim}
THEME_NAME(themeName);
\end{verbatim}
Parameters:
\begin{description}
\item[themeName] String or character positional argument.
\end{description}

Sets the serialized theme name.

\paragraph{Example macro.}
\begin{verbatim}
$MACRO REFERENCE_THEME_NAME FROM EDIT;
    THEME_NAME('reference');
END_MACRO;
\end{verbatim}
\end{mrmacextension}

\begin{mrmacextension}
\phantomsection
\label{prim:procedure-theme-reset}
\subsection*{THEME\_RESET}
\index{THEME\_RESET@\texttt{THEME\_RESET}}
\begin{verbatim}
THEME_RESET();
\end{verbatim}
Parameters: None.

Resets the serialized theme source currently being assembled.

\paragraph{Example macro.}
\begin{verbatim}
$MACRO REFERENCE_THEME_RESET FROM EDIT;
    THEME_RESET();
END_MACRO;
\end{verbatim}
\end{mrmacextension}

\begin{mrmacextension}
\phantomsection
\label{prim:procedure-theme-version}
\subsection*{THEME\_VERSION}
\index{THEME\_VERSION@\texttt{THEME\_VERSION}}
\begin{verbatim}
THEME_VERSION(version);
\end{verbatim}
Parameters:
\begin{description}
\item[version] Serialized version text.
\end{description}

Sets the theme source version.

\paragraph{Example macro.}
\begin{verbatim}
$MACRO REFERENCE_THEME_VERSION FROM EDIT;
    THEME_VERSION('reference');
END_MACRO;
\end{verbatim}
\end{mrmacextension}

\begin{mrmacextension}
\phantomsection
\label{prim:procedure-windowcolors}
\subsection*{WINDOWCOLORS}
\index{WINDOWCOLORS@\texttt{WINDOWCOLORS}}
\begin{verbatim}
WINDOWCOLORS(definition);
\end{verbatim}
Parameters:
\begin{description}
\item[definition] Serialized keymap or theme definition.
\end{description}

Adds the named serialized theme colour definition.

\paragraph{Example macro.}
\begin{verbatim}
$MACRO REFERENCE_WINDOWCOLORS FROM EDIT;
    WINDOWCOLORS('reference');
END_MACRO;
\end{verbatim}
\end{mrmacextension}


\chapter{External Environment and Editor Commands}

This chapter contains the bounded process-environment bridge and procedures that invoke established MR application commands. They request existing behaviour; they do not provide unrestricted native handles or a separate DOS shell.

\section{Primitive Map}
The map is a local index for this chapter. The final column is a generated page
reference and hyperlink to the full entry, where the source form, parameters,
semantics and a complete example macro appear.

\begin{longtable}{@{}L{0.38\textwidth}L{0.14\textwidth}L{0.27\textwidth}L{0.10\textwidth}@{}}
\toprule
Primitive & Role & Purpose & Reference \\
\midrule
\endhead
\toprule
Primitive & Role & Purpose & Reference \\
\midrule
\endhead
\primitiveoverview{GET\_ENVIRONMENT}{intrinsic-get-environment}{intrinsic}{MR command bridge.}
\primitiveoverview{OS\_BACK}{intrinsic-os-back}{intrinsic}{MR command bridge.}
\primitiveoverview{OS\_COLOR}{intrinsic-os-color}{intrinsic}{MR command bridge.}
\primitiveoverview{SUBSHELL}{intrinsic-subshell}{intrinsic}{MR command bridge.}
\primitiveoverview{VERSION}{intrinsic-version}{intrinsic}{MR command bridge.}
\primitiveoverview{ACTIVE\_KEYMAP\_PROFILE}{procedure-active-keymap-profile}{procedure}{MR command bridge.}
\primitiveoverview{EXIT\_SAVE\_ALL}{procedure-exit-save-all}{procedure}{MR command bridge.}
\primitiveoverview{MACRO\_TOGGLE\_RECORDING}{procedure-macro-toggle-recording}{procedure}{MR command bridge.}
\primitiveoverview{MATCH\_PARENTHESIS}{procedure-match-parenthesis}{procedure}{MR command bridge.}
\primitiveoverview{SETUP\_BACKUPS\_AUTOSAVE}{procedure-setup-backups-autosave}{procedure}{MR command bridge.}
\primitiveoverview{SETUP\_COLOR}{procedure-setup-color}{procedure}{MR command bridge.}
\primitiveoverview{SETUP\_COMPILER\_PROFILES}{procedure-setup-compiler-profiles}{procedure}{MR command bridge.}
\primitiveoverview{SETUP\_EDIT\_SETTINGS}{procedure-setup-edit-settings}{procedure}{MR command bridge.}
\primitiveoverview{SETUP\_FILENAME\_EXTENSIONS}{procedure-setup-filename-extensions}{procedure}{MR command bridge.}
\primitiveoverview{SETUP\_KEYMAP}{procedure-setup-keymap}{procedure}{MR command bridge.}
\primitiveoverview{SETUP\_LIVE\_LOGS}{procedure-setup-live-logs}{procedure}{MR command bridge.}
\primitiveoverview{SETUP\_PATHS}{procedure-setup-paths}{procedure}{MR command bridge.}
\primitiveoverview{SETUP\_SEARCH\_REPLACE\_DEFAULTS}{procedure-setup-search-replace-defaults}{procedure}{MR command bridge.}
\primitiveoverview{SETUP\_USER\_INTERFACE}{procedure-setup-user-interface}{procedure}{MR command bridge.}
\primitiveoverview{QUIT}{procedure-quit}{procedure}{MR command bridge.}
\primitiveoverview{REST\_OS\_SCREEN}{procedure-rest-os-screen}{procedure}{MR command bridge.}
\primitiveoverview{SAVE\_OS\_SCREEN}{procedure-save-os-screen}{procedure}{MR command bridge.}
\primitiveoverview{CHANGE\_DIR}{procedure-change-dir}{procedure}{MR command bridge.}
\primitiveoverview{FORK}{procedure-fork}{procedure}{MR command bridge.}
\primitiveoverview{SHELL\_TO\_OS}{procedure-shell-to-os}{procedure}{MR command bridge.}
\primitiveoverview{WRITE\_SOD}{procedure-write-sod}{procedure}{MR command bridge.}
\bottomrule
\end{longtable}

\section{Reference Entries}

\phantomsection
\label{prim:intrinsic-get-environment}
\subsection*{GET\_ENVIRONMENT}
\index{GET\_ENVIRONMENT@\texttt{GET\_ENVIRONMENT}}
\begin{verbatim}
GET_ENVIRONMENT(variableName)
\end{verbatim}
Parameters:
\begin{description}
\item[variableName] String-like environment variable name.
\end{description}

Returns the named environment value.

Returns a string.

\paragraph{Example macro.}
\begin{verbatim}
$MACRO REFERENCE_GET_ENVIRONMENT FROM EDIT;
    DEF_STR(Result);
    Result := GET_ENVIRONMENT('reference');
END_MACRO;
\end{verbatim}

\phantomsection
\label{prim:intrinsic-os-back}
\subsection*{OS\_BACK}
\index{OS\_BACK@\texttt{OS\_BACK}}
\begin{verbatim}
OS_BACK
\end{verbatim}
Parameters: None.

Returns the supported system background attribute.

Returns an integer.

\paragraph{Example macro.}
\begin{verbatim}
$MACRO REFERENCE_OS_BACK FROM EDIT;
    DEF_INT(Result);
    Result := OS_BACK;
END_MACRO;
\end{verbatim}

\phantomsection
\label{prim:intrinsic-os-color}
\subsection*{OS\_COLOR}
\index{OS\_COLOR@\texttt{OS\_COLOR}}
\begin{verbatim}
OS_COLOR
\end{verbatim}
Parameters: None.

Returns the supported system colour attribute.

Returns an integer.

\paragraph{Example macro.}
\begin{verbatim}
$MACRO REFERENCE_OS_COLOR FROM EDIT;
    DEF_INT(Result);
    Result := OS_COLOR;
END_MACRO;
\end{verbatim}

\phantomsection
\label{prim:intrinsic-subshell}
\subsection*{SUBSHELL}
\index{SUBSHELL@\texttt{SUBSHELL}}
\begin{verbatim}
SUBSHELL(command, mode)
\end{verbatim}
Parameters:
\begin{description}
\item[command] Controlled command text.
\item[mode] Editor mode constant.
\end{description}

Runs the controlled command path and returns captured output.

Returns a string.

\paragraph{Example macro.}
\begin{verbatim}
$MACRO REFERENCE_SUBSHELL FROM EDIT;
    DEF_STR(Result);
    Result := SUBSHELL('reference', 0);
END_MACRO;
\end{verbatim}

\phantomsection
\label{prim:intrinsic-version}
\subsection*{VERSION}
\index{VERSION@\texttt{VERSION}}
\begin{verbatim}
VERSION
\end{verbatim}
Parameters: None.

Returns the MRMAC runtime version string.

Returns a string.

\paragraph{Example macro.}
\begin{verbatim}
$MACRO REFERENCE_VERSION FROM EDIT;
    DEF_STR(Result);
    Result := VERSION;
END_MACRO;
\end{verbatim}

\phantomsection
\label{prim:procedure-active-keymap-profile}
\subsection*{ACTIVE\_KEYMAP\_PROFILE}
\index{ACTIVE\_KEYMAP\_PROFILE@\texttt{ACTIVE\_KEYMAP\_PROFILE}}
\begin{verbatim}
ACTIVE_KEYMAP_PROFILE(profileName);
\end{verbatim}
Parameters:
\begin{description}
\item[profileName] Serialized profile name.
\end{description}

Sets the active serialized keymap profile.

\paragraph{Example macro.}
\begin{verbatim}
$MACRO REFERENCE_ACTIVE_KEYMAP_PROFILE FROM EDIT;
    ACTIVE_KEYMAP_PROFILE('reference');
END_MACRO;
\end{verbatim}

\begin{mrmacextension}
\phantomsection
\label{prim:procedure-exit-save-all}
\subsection*{EXIT\_SAVE\_ALL}
\index{EXIT\_SAVE\_ALL@\texttt{EXIT\_SAVE\_ALL}}
\begin{verbatim}
EXIT_SAVE_ALL;
\end{verbatim}
Parameters: None.

Performs the documented MR operation named by this primitive in the active execution context.

\paragraph{Example macro.}
\begin{verbatim}
$MACRO REFERENCE_EXIT_SAVE_ALL FROM EDIT;
    EXIT_SAVE_ALL;
END_MACRO;
\end{verbatim}
\end{mrmacextension}

\begin{mrmacextension}
\phantomsection
\label{prim:procedure-macro-toggle-recording}
\subsection*{MACRO\_TOGGLE\_RECORDING}
\index{MACRO\_TOGGLE\_RECORDING@\texttt{MACRO\_TOGGLE\_RECORDING}}
\begin{verbatim}
MACRO_TOGGLE_RECORDING;
\end{verbatim}
Parameters: None.

Performs the documented MR operation named by this primitive in the active execution context.

\paragraph{Example macro.}
\begin{verbatim}
$MACRO REFERENCE_MACRO_TOGGLE_RECORDING FROM EDIT;
    MACRO_TOGGLE_RECORDING;
END_MACRO;
\end{verbatim}
\end{mrmacextension}

\begin{mrmacextension}
\phantomsection
\label{prim:procedure-match-parenthesis}
\subsection*{MATCH\_PARENTHESIS}
\index{MATCH\_PARENTHESIS@\texttt{MATCH\_PARENTHESIS}}
\begin{verbatim}
MATCH_PARENTHESIS;
\end{verbatim}
Parameters: None.

Performs the documented MR operation named by this primitive in the active execution context.

\paragraph{Example macro.}
\begin{verbatim}
$MACRO REFERENCE_MATCH_PARENTHESIS FROM EDIT;
    MATCH_PARENTHESIS;
END_MACRO;
\end{verbatim}
\end{mrmacextension}

\begin{mrmacextension}
\phantomsection
\label{prim:procedure-setup-backups-autosave}
\subsection*{SETUP\_BACKUPS\_AUTOSAVE}
\index{SETUP\_BACKUPS\_AUTOSAVE@\texttt{SETUP\_BACKUPS\_AUTOSAVE}}
\begin{verbatim}
SETUP_BACKUPS_AUTOSAVE;
\end{verbatim}
Parameters: None.

Opens the corresponding MR setup dialog through the application command route.

\paragraph{Example macro.}
\begin{verbatim}
$MACRO REFERENCE_SETUP_BACKUPS_AUTOSAVE FROM EDIT;
    SETUP_BACKUPS_AUTOSAVE;
END_MACRO;
\end{verbatim}
\end{mrmacextension}

\begin{mrmacextension}
\phantomsection
\label{prim:procedure-setup-color}
\subsection*{SETUP\_COLOR}
\index{SETUP\_COLOR@\texttt{SETUP\_COLOR}}
\begin{verbatim}
SETUP_COLOR;
\end{verbatim}
Parameters: None.

Opens the corresponding MR setup dialog through the application command route.

\paragraph{Example macro.}
\begin{verbatim}
$MACRO REFERENCE_SETUP_COLOR FROM EDIT;
    SETUP_COLOR;
END_MACRO;
\end{verbatim}
\end{mrmacextension}

\begin{mrmacextension}
\phantomsection
\label{prim:procedure-setup-compiler-profiles}
\subsection*{SETUP\_COMPILER\_PROFILES}
\index{SETUP\_COMPILER\_PROFILES@\texttt{SETUP\_COMPILER\_PROFILES}}
\begin{verbatim}
SETUP_COMPILER_PROFILES;
\end{verbatim}
Parameters: None.

Opens the corresponding MR setup dialog through the application command route.

\paragraph{Example macro.}
\begin{verbatim}
$MACRO REFERENCE_SETUP_COMPILER_PROFILES FROM EDIT;
    SETUP_COMPILER_PROFILES;
END_MACRO;
\end{verbatim}
\end{mrmacextension}

\begin{mrmacextension}
\phantomsection
\label{prim:procedure-setup-edit-settings}
\subsection*{SETUP\_EDIT\_SETTINGS}
\index{SETUP\_EDIT\_SETTINGS@\texttt{SETUP\_EDIT\_SETTINGS}}
\begin{verbatim}
SETUP_EDIT_SETTINGS;
\end{verbatim}
Parameters: None.

Opens the corresponding MR setup dialog through the application command route.

\paragraph{Example macro.}
\begin{verbatim}
$MACRO REFERENCE_SETUP_EDIT_SETTINGS FROM EDIT;
    SETUP_EDIT_SETTINGS;
END_MACRO;
\end{verbatim}
\end{mrmacextension}

\begin{mrmacextension}
\phantomsection
\label{prim:procedure-setup-filename-extensions}
\subsection*{SETUP\_FILENAME\_EXTENSIONS}
\index{SETUP\_FILENAME\_EXTENSIONS@\texttt{SETUP\_FILENAME\_EXTENSIONS}}
\begin{verbatim}
SETUP_FILENAME_EXTENSIONS;
\end{verbatim}
Parameters: None.

Opens the corresponding MR setup dialog through the application command route.

\paragraph{Example macro.}
\begin{verbatim}
$MACRO REFERENCE_SETUP_FILENAME_EXTENSIONS FROM EDIT;
    SETUP_FILENAME_EXTENSIONS;
END_MACRO;
\end{verbatim}
\end{mrmacextension}

\begin{mrmacextension}
\phantomsection
\label{prim:procedure-setup-keymap}
\subsection*{SETUP\_KEYMAP}
\index{SETUP\_KEYMAP@\texttt{SETUP\_KEYMAP}}
\begin{verbatim}
SETUP_KEYMAP;
\end{verbatim}
Parameters: None.

Opens the corresponding MR setup dialog through the application command route.

\paragraph{Example macro.}
\begin{verbatim}
$MACRO REFERENCE_SETUP_KEYMAP FROM EDIT;
    SETUP_KEYMAP;
END_MACRO;
\end{verbatim}
\end{mrmacextension}

\begin{mrmacextension}
\phantomsection
\label{prim:procedure-setup-live-logs}
\subsection*{SETUP\_LIVE\_LOGS}
\index{SETUP\_LIVE\_LOGS@\texttt{SETUP\_LIVE\_LOGS}}
\begin{verbatim}
SETUP_LIVE_LOGS;
\end{verbatim}
Parameters: None.

Opens the corresponding MR setup dialog through the application command route.

\paragraph{Example macro.}
\begin{verbatim}
$MACRO REFERENCE_SETUP_LIVE_LOGS FROM EDIT;
    SETUP_LIVE_LOGS;
END_MACRO;
\end{verbatim}
\end{mrmacextension}

\begin{mrmacextension}
\phantomsection
\label{prim:procedure-setup-paths}
\subsection*{SETUP\_PATHS}
\index{SETUP\_PATHS@\texttt{SETUP\_PATHS}}
\begin{verbatim}
SETUP_PATHS;
\end{verbatim}
Parameters: None.

Opens the corresponding MR setup dialog through the application command route.

\paragraph{Example macro.}
\begin{verbatim}
$MACRO REFERENCE_SETUP_PATHS FROM EDIT;
    SETUP_PATHS;
END_MACRO;
\end{verbatim}
\end{mrmacextension}

\begin{mrmacextension}
\phantomsection
\label{prim:procedure-setup-search-replace-defaults}
\subsection*{SETUP\_SEARCH\_REPLACE\_DEFAULTS}
\index{SETUP\_SEARCH\_REPLACE\_DEFAULTS@\texttt{SETUP\_SEARCH\_REPLACE\_DEFAULTS}}
\begin{verbatim}
SETUP_SEARCH_REPLACE_DEFAULTS;
\end{verbatim}
Parameters: None.

Opens the corresponding MR setup dialog through the application command route.

\paragraph{Example macro.}
\begin{verbatim}
$MACRO REFERENCE_SETUP_SEARCH_REPLACE_DEFAULTS FROM EDIT;
    SETUP_SEARCH_REPLACE_DEFAULTS;
END_MACRO;
\end{verbatim}
\end{mrmacextension}

\begin{mrmacextension}
\phantomsection
\label{prim:procedure-setup-user-interface}
\subsection*{SETUP\_USER\_INTERFACE}
\index{SETUP\_USER\_INTERFACE@\texttt{SETUP\_USER\_INTERFACE}}
\begin{verbatim}
SETUP_USER_INTERFACE;
\end{verbatim}
Parameters: None.

Opens the corresponding MR setup dialog through the application command route.

\paragraph{Example macro.}
\begin{verbatim}
$MACRO REFERENCE_SETUP_USER_INTERFACE FROM EDIT;
    SETUP_USER_INTERFACE;
END_MACRO;
\end{verbatim}
\end{mrmacextension}

\phantomsection
\label{prim:procedure-quit}
\subsection*{QUIT}
\index{QUIT@\texttt{QUIT}}
\begin{verbatim}
QUIT();
QUIT(saveMode);
\end{verbatim}
Parameters:
\begin{description}
\item[saveMode] Optional integer quit/save policy.
\end{description}

Requests the editor exit route with optional save policy.

\paragraph{Example macro.}
\begin{verbatim}
$MACRO REFERENCE_QUIT FROM EDIT;
    QUIT(0);
END_MACRO;
\end{verbatim}

\phantomsection
\label{prim:procedure-rest-os-screen}
\subsection*{REST\_OS\_SCREEN}
\index{REST\_OS\_SCREEN@\texttt{REST\_OS\_SCREEN}}
\begin{verbatim}
REST_OS_SCREEN;
\end{verbatim}
Parameters: None.

Restores saved supported system-screen state.

\paragraph{Example macro.}
\begin{verbatim}
$MACRO REFERENCE_REST_OS_SCREEN FROM EDIT;
    REST_OS_SCREEN;
END_MACRO;
\end{verbatim}

\phantomsection
\label{prim:procedure-save-os-screen}
\subsection*{SAVE\_OS\_SCREEN}
\index{SAVE\_OS\_SCREEN@\texttt{SAVE\_OS\_SCREEN}}
\begin{verbatim}
SAVE_OS_SCREEN;
\end{verbatim}
Parameters: None.

Saves supported system-screen state.

\paragraph{Example macro.}
\begin{verbatim}
$MACRO REFERENCE_SAVE_OS_SCREEN FROM EDIT;
    SAVE_OS_SCREEN;
END_MACRO;
\end{verbatim}

\phantomsection
\label{prim:procedure-change-dir}
\subsection*{CHANGE\_DIR}
\index{CHANGE\_DIR@\texttt{CHANGE\_DIR}}
\begin{verbatim}
CHANGE_DIR(path);
\end{verbatim}
Parameters:
\begin{description}
\item[path] A string-like file-system path.
\end{description}

Changes the macro runtime working directory.

\paragraph{Example macro.}
\begin{verbatim}
$MACRO REFERENCE_CHANGE_DIR FROM EDIT;
    CHANGE_DIR('reference');
END_MACRO;
\end{verbatim}

\phantomsection
\label{prim:procedure-fork}
\subsection*{FORK}
\index{FORK@\texttt{FORK}}
\begin{verbatim}
FORK(command, argument1, argument2, argument3,
     argument4, argument5, argument6, argument7);
\end{verbatim}
Parameters:
\begin{description}
\item[command] Controlled command text.
\item[argument1] First controlled-process string argument.
\item[argument2] Second controlled-process string argument.
\item[argument3] Third controlled-process string argument.
\item[argument4] Fourth controlled-process string argument.
\item[argument5] Fifth controlled-process string argument.
\item[argument6] Sixth controlled-process string argument.
\item[argument7] Seventh controlled-process string argument.
\end{description}

Starts a controlled external process with its command and positional string arguments.

\paragraph{Example macro.}
\begin{verbatim}
$MACRO REFERENCE_FORK FROM EDIT;
    FORK('reference', 'reference', 'reference', 'reference',
         'reference', 'reference', 'reference', 'reference');
END_MACRO;
\end{verbatim}

\phantomsection
\label{prim:procedure-shell-to-os}
\subsection*{SHELL\_TO\_OS}
\index{SHELL\_TO\_OS@\texttt{SHELL\_TO\_OS}}
\begin{verbatim}
SHELL_TO_OS(command, mode);
\end{verbatim}
Parameters:
\begin{description}
\item[command] Controlled command text.
\item[mode] Editor mode constant.
\end{description}

Executes a controlled system command using mode.

\paragraph{Example macro.}
\begin{verbatim}
$MACRO REFERENCE_SHELL_TO_OS FROM EDIT;
    SHELL_TO_OS('reference', 0);
END_MACRO;
\end{verbatim}

\phantomsection
\label{prim:procedure-write-sod}
\subsection*{WRITE\_SOD}
\index{WRITE\_SOD@\texttt{WRITE\_SOD}}
\begin{verbatim}
WRITE_SOD(text);
\end{verbatim}
Parameters:
\begin{description}
\item[text] String-like text.
\end{description}

Writes text through the supported standard-output device path.

\paragraph{Example macro.}
\begin{verbatim}
$MACRO REFERENCE_WRITE_SOD FROM EDIT;
    WRITE_SOD('reference');
END_MACRO;
\end{verbatim}


\part{Debugging}

\chapter{Debugging MRMAC Macros}

\begin{mrmacextension}
\section{Starting a debug session}

The MRMAC debugger starts a loaded macro from the macro catalog and compiles it
with source-map information only for that debug run. Normal macro loading does
not create source maps. A debugger session runs real MRMAC semantics: file,
editor and UI side effects are not a preview.
\end{mrmacextension}

\debuggerscreenshot{DBG_overview.png}{fig:debugger-library}{The Macro Library holds loaded macros. Select the macro to inspect, then choose Debug to create a debugger session for that macro.}

\begin{mrmacextension}
The new session pauses at its first debuggable source span. Debugger Output
identifies the macro, session, state and stop reason; the Variables and Watches
panes show the state available at that pause.
\end{mrmacextension}

\debuggerscreenshot{DBG_session_starting.png}{fig:debugger-session}{A newly started MRMAC debug session paused at a source span. The debugger panes identify the session and expose its current state.}

\begin{mrmacextension}
\section{Breakpoints, running and stopping}

\texttt{F9} sets a breakpoint at the cursor source line or clears one already
there. The debugger maps the line to its debuggable source span, so a comment,
declaration or other non-executable line cannot receive a breakpoint.
\texttt{Shift-F9} enables or disables the breakpoint at the cursor. Its
\texttt{Alt-Shift-F9} and \texttt{Ctrl-Shift-F9} forms respectively enable or
disable all breakpoints, or clear all breakpoints, in the current macro file.

\texttt{F6} Run Here closes any live session and starts the selected macro
again, with a temporary breakpoint at the cursor line. The cursor line must be
debuggable; another enabled breakpoint reached first still stops execution
first. \texttt{F5} continues a paused session. While that session is running,
\texttt{F5} requests a pause. \texttt{F8} closes a live session, clears its
current-location marker and removes its variable state. With no live session,
\texttt{F8} resets the debugger by starting a new session which pauses at the
first debuggable source span.

Breakpoints and watches are debugger state, separate from the general settings.
Their definitions are stored with the debugger Bento's workspace configuration
and rebound when that workspace is restored.
\end{mrmacextension}

\debuggerscreenshot{DBG_BP_with_running_MMP.png}{fig:debugger-breakpoint}{A breakpoint pauses the macro at the marked source span while the displayed modeless window remains the real UI created by that session.}

\begin{mrmacextension}
An MMP timer callback starts as its own runtime execution session. If its source
is open in an observing debugger Bento and the callback has an enabled
breakpoint, it stops in that Bento like any other macro. The source pane,
Debugger Output, Variables and Watches remain the controls for that paused
callback; \texttt{F5} resumes it. Without an observing debugger, timer
callbacks retain their normal scheduling behaviour.
\end{mrmacextension}

\debuggerscreenshot{DBG_callbacl_BP.png}{fig:debugger-timer-callback}{An enabled breakpoint in the \texttt{MmpCanvasDemoTick} timer callback stops the independently scheduled callback at its source line. The modeless MMP windows remain visible while the debugger exposes the callback session and its local \texttt{Ticks} value.}

\begin{mrmacextension}
\section{Stepping}

Stepping requires a paused live session. \texttt{F10} Step Into advances to the
next source span and enters a call when that is the next operation.
\texttt{F11} Step Over executes a called operation without stopping inside it,
then pauses at the next source span in the current frame. \texttt{Shift-F11}
Step Out continues the current frame to its return and then pauses in its
caller. Continue runs until a breakpoint, pause request, stop request, runtime
error or normal completion. Nested local calls and starts of loaded macros
contribute frames to the displayed debug location when source mapping is
available.

At every pause, the source pane follows the mapped instruction, marks its
source span and centres the line. Debugger Output reports the macro, session,
state, stop reason, instruction offset, frame count, location, variable and
breakpoint counts, call stack and the latest log lines.

Debugger Output lists the call stack below its summary. Frame \texttt{\#0} is
the paused macro. Later frames are callers: \texttt{[CALL]} identifies a
nested call in that macro, while \texttt{[RUN\_MACRO]} identifies the macro
that started a loaded child macro. Each frame includes its mapped source
location when available.
\end{mrmacextension}

\debuggerscreenshot{DBG_MMP_timer.png}{fig:debugger-step}{A step stop advances the current source location and records \emph{step} as the stop reason in Debugger Output.}

\debuggerscreenshot{DBG_callmacro.png}{fig:debugger-call-stack}{A breakpoint in a child macro started by \texttt{RUN\_MACRO} leaves the child at frame \texttt{\#0} and its parent at frame \texttt{\#1} in the Call stack.}

\chapter{Inspecting State and Diagnosing Failures}

\begin{mrmacextension}
\section{Variables and collections}

The debugger panes show locals, file globals, application globals, closures and
session values. While a session is paused, select an existing integer, real,
string or character scalar in Variables to open its inline value field. Enter
validates and writes the new value, then refreshes Variables and Watches; Esc
cancels the edit. Names, types, hashes and arrays remain read-only in the
debugger.
\end{mrmacextension}

\debuggerscreenshot{DBG_var_editing.png}{fig:debugger-variable}{The source span and Variables pane identify the same scalar. Existing scalar values may be changed while the session is paused.}

\begin{mrmacextension}
Hashes and arrays appear as typed collection summaries. They may be inspected at
a pause, but the debugger does not permit collection mutation.
\end{mrmacextension}

\debuggerscreenshot{DBG_multilevel_hashes.png}{fig:debugger-hash}{A hash is shown as a typed key-count summary in Variables. The corresponding source assignments and collection are visible at the paused location.}

\begin{mrmacextension}
\section{Watches and expression evaluation}

\texttt{F7} opens Watch expression and stores the entered expression for the
current debug macro; it does not require a live session. \texttt{Shift-F7}
opens Remove watch and removes the watch whose expression exactly matches the
entered text. A stored watch is evaluated only for a paused live session;
otherwise its pane reports that no live debug session is available. Watches are
compiled through the canonical compiler in a restricted pure-expression mode.
\texttt{F4} evaluates an expression against the paused VM. A watch therefore
reports a value or a compiler/runtime diagnostic; it does not execute an
arbitrary procedure or create a second parser.
\end{mrmacextension}

\debuggerscreenshot{DBG_add_watch.png}{fig:debugger-watch}{F7 opens the Watch expression dialog for the selected debug macro. The expression is compiled in restricted pure-expression mode and is evaluated when the macro has a paused live session.}

\begin{mrmacextension}
F4 evaluates a pure expression immediately against the paused state. Debugger
Output records the expression and result without invoking an arbitrary macro
procedure.
\end{mrmacextension}

\debuggerscreenshot{DBG_eval.png}{fig:debugger-evaluation}{F4 evaluates an expression against paused values. Debugger Output records its result and the Variables pane shows the participating state.}

\begin{mrmacextension}
\section{Diagnostic model}

Compiler diagnostics identify a source line and a concrete syntax, type or
declaration problem. Runtime diagnostics identify the failed operation while
the debugger retains the current source location and snapshots. MRMAC has no
stable MEMAC-style numeric error catalogue; treat the displayed diagnostic text
and the paused state as the authoritative failure report.
\end{mrmacextension}

\chapter{Example Macros}

\section{Insert the current date}

\begin{verbatim}
$MACRO INSERT_DATE FROM EDIT;
    TEXT(DATE);
END_MACRO;
\end{verbatim}

\section{A compact modeless status panel}

\begin{mrmacextension}
\begin{verbatim}
$MACRO SHOW_STATUS FROM EDIT;
    UI_DIALOG(2, 2, 42, 7, 'MRMAC status');
    MMP_STATUS_FIELD(2, 2, 36, 'state', 'ready');
    MMP_ACTION_BUTTON(2, 4, 14, 1, 'Refresh', 'status^REFRESH_STATUS');
    UI_MODELESS_SHOW('StatusPanel');
END_MACRO;

$MACRO REFRESH_STATUS FROM EDIT;
    MMP_STATUS_SET('StatusPanel', 'state', 'updated');
END_MACRO;
\end{verbatim}
\end{mrmacextension}

\appendix
\chapter{MEMAC Compatibility Notes}

MRMAC retains a broad, practical subset of MEMAC syntax and editor commands,
but it does not emulate the DOS shell, EMS/video-card probing or a separate
word-processing runtime. In particular, historical DOS directory-shell commands
and device-specific video operations are not macro interfaces in MR. The nearby
footnotes in this manual identify comparable names where migration is likely.

\chapter{Command Index}

\printindex

\end{document}
